% Options for packages loaded elsewhere
% Options for packages loaded elsewhere
\PassOptionsToPackage{unicode}{hyperref}
\PassOptionsToPackage{hyphens}{url}
%
\documentclass[
  10pt,
  letterpaper,
]{scrbook}
\usepackage{xcolor}
\usepackage[paperwidth=6.25in,paperheight=9.25in,textwidth=4.5in,textheight=7.25in,left=0.75in,right=0.75in,top=0.75in,bottom=1in]{geometry}
\usepackage{amsmath,amssymb}
\setcounter{secnumdepth}{1}
\usepackage{iftex}
\ifPDFTeX
  \usepackage[T1]{fontenc}
  \usepackage[utf8]{inputenc}
  \usepackage{textcomp} % provide euro and other symbols
\else % if luatex or xetex
  \usepackage{unicode-math} % this also loads fontspec
  \defaultfontfeatures{Scale=MatchLowercase}
  \defaultfontfeatures[\rmfamily]{Ligatures=TeX,Scale=1}
\fi
\usepackage[]{libertinus}
\ifPDFTeX\else
  % xetex/luatex font selection
\fi
% Use upquote if available, for straight quotes in verbatim environments
\IfFileExists{upquote.sty}{\usepackage{upquote}}{}
\IfFileExists{microtype.sty}{% use microtype if available
  \usepackage[]{microtype}
  \UseMicrotypeSet[protrusion]{basicmath} % disable protrusion for tt fonts
}{}
\makeatletter
\@ifundefined{KOMAClassName}{% if non-KOMA class
  \IfFileExists{parskip.sty}{%
    \usepackage{parskip}
  }{% else
    \setlength{\parindent}{0pt}
    \setlength{\parskip}{6pt plus 2pt minus 1pt}}
}{% if KOMA class
  \KOMAoptions{parskip=half}}
\makeatother
% Make \paragraph and \subparagraph free-standing
\makeatletter
\ifx\paragraph\undefined\else
  \let\oldparagraph\paragraph
  \renewcommand{\paragraph}{
    \@ifstar
      \xxxParagraphStar
      \xxxParagraphNoStar
  }
  \newcommand{\xxxParagraphStar}[1]{\oldparagraph*{#1}\mbox{}}
  \newcommand{\xxxParagraphNoStar}[1]{\oldparagraph{#1}\mbox{}}
\fi
\ifx\subparagraph\undefined\else
  \let\oldsubparagraph\subparagraph
  \renewcommand{\subparagraph}{
    \@ifstar
      \xxxSubParagraphStar
      \xxxSubParagraphNoStar
  }
  \newcommand{\xxxSubParagraphStar}[1]{\oldsubparagraph*{#1}\mbox{}}
  \newcommand{\xxxSubParagraphNoStar}[1]{\oldsubparagraph{#1}\mbox{}}
\fi
\makeatother


\usepackage{longtable,booktabs,array}
\usepackage{calc} % for calculating minipage widths
% Correct order of tables after \paragraph or \subparagraph
\usepackage{etoolbox}
\makeatletter
\patchcmd\longtable{\par}{\if@noskipsec\mbox{}\fi\par}{}{}
\makeatother
% Allow footnotes in longtable head/foot
\IfFileExists{footnotehyper.sty}{\usepackage{footnotehyper}}{\usepackage{footnote}}
\makesavenoteenv{longtable}
\usepackage{graphicx}
\makeatletter
\newsavebox\pandoc@box
\newcommand*\pandocbounded[1]{% scales image to fit in text height/width
  \sbox\pandoc@box{#1}%
  \Gscale@div\@tempa{\textheight}{\dimexpr\ht\pandoc@box+\dp\pandoc@box\relax}%
  \Gscale@div\@tempb{\linewidth}{\wd\pandoc@box}%
  \ifdim\@tempb\p@<\@tempa\p@\let\@tempa\@tempb\fi% select the smaller of both
  \ifdim\@tempa\p@<\p@\scalebox{\@tempa}{\usebox\pandoc@box}%
  \else\usebox{\pandoc@box}%
  \fi%
}
% Set default figure placement to htbp
\def\fps@figure{htbp}
\makeatother





\setlength{\emergencystretch}{3em} % prevent overfull lines

\providecommand{\tightlist}{%
  \setlength{\itemsep}{0pt}\setlength{\parskip}{0pt}}



 


% Course Notes - PDF Preamble
% Fonts and title styling


\usepackage{anyfontsize}
\newcommand\titleFontSize{\fontsize{50pt}{55pt}\selectfont}
\usepackage{marginnote}
\renewcommand*{\marginfont}{\scriptsize\itshape}
\setkomafont{part}{\titleFontSize\bfseries} 
\setkomafont{title}{\titleFontSize\bfseries\sffamily\scshape}

% tikz-cd support
\usepackage{tikz-cd}
\usetikzlibrary{calc}
\usetikzlibrary{decorations.pathmorphing}

% Quiver curved arrow support
\tikzset{curve/.style={settings={#1},to path={(\tikztostart)
    .. controls ($(\tikztostart)!\pv{pos}!(\tikztotarget)!\pv{height}!270:(\tikztotarget)$)
    and ($(\tikztostart)!1-\pv{pos}!(\tikztotarget)!\pv{height}!270:(\tikztotarget)$)
    .. (\tikztotarget)\tikztonodes}},
    settings/.code={\tikzset{quiver/.cd,#1}
        \def\pv##1{\pgfkeysvalueof{/tikz/quiver/##1}}},
    quiver/.cd,pos/.initial=0.35,height/.initial=0}
\tikzset{tail reversed/.code={\pgfsetarrowsstart{tikzcd to}}}
\tikzset{2tail/.code={\pgfsetarrowsstart{Implies[reversed]}}}
\tikzset{2tail reversed/.code={\pgfsetarrowsstart{Implies}}}
\tikzset{no body/.style={/tikz/dash pattern=on 0 off 1mm}}% Chapter and section titles
% Sized for a professional math textbook feel
\setkomafont{chapter}{\LARGE\bfseries}
\setkomafont{section}{\Large\bfseries}
\setkomafont{subsection}{\large\bfseries}

% ============================================
% SPACING CONFIGURATION
% ============================================

% --------------------------------------------
% Paragraph Style
% --------------------------------------------
\KOMAoptions{parskip=never}
\setlength{\parindent}{1.5em}

% --------------------------------------------
% Section Headings
% Format: beforeskip = space above title
%         afterskip  = space below title
% --------------------------------------------
\RedeclareSectionCommand[
  beforeskip=1.5\baselineskip,
  afterskip=0.5\baselineskip
]{section}

\RedeclareSectionCommand[
  beforeskip=1.0\baselineskip,
  afterskip=0.3\baselineskip
]{subsection}

\RedeclareSectionCommand[
  beforeskip=0.75\baselineskip,
  afterskip=0.2\baselineskip
]{subsubsection}

% --------------------------------------------
% Theorem Environments
% Negative vspace compensates for amsthm's 
% built-in spacing around theorem blocks
% --------------------------------------------
\newcommand{\thmbeforeskip}{-0.15\baselineskip}
\newcommand{\thmafterskip}{-0.15\baselineskip}

\BeforeBeginEnvironment{theorem}{\vspace{\thmbeforeskip}}
\AfterEndEnvironment{theorem}{\vspace{\thmafterskip}}

\BeforeBeginEnvironment{definition}{\vspace{\thmbeforeskip}}
\AfterEndEnvironment{definition}{\vspace{\thmafterskip}}

\BeforeBeginEnvironment{proposition}{\vspace{\thmbeforeskip}}
\AfterEndEnvironment{proposition}{\vspace{\thmafterskip}}

\BeforeBeginEnvironment{example}{\vspace{\thmbeforeskip}}
\AfterEndEnvironment{example}{\vspace{\thmafterskip}}

\BeforeBeginEnvironment{corollary}{\vspace{\thmbeforeskip}}
\AfterEndEnvironment{corollary}{\vspace{\thmafterskip}}

\BeforeBeginEnvironment{remark}{\vspace{\thmbeforeskip}}
\AfterEndEnvironment{remark}{\vspace{\thmafterskip}}

\BeforeBeginEnvironment{proof}{\vspace{\thmbeforeskip}}
\AfterEndEnvironment{proof}{\vspace{\thmafterskip}}
\newcommand{\ZZ}{\mathbb{Z}}
\newcommand{\NN}{\mathbb{N}}
\newcommand{\QQ}{\mathbb{Q}}
\newcommand{\RR}{\mathbb{R}}
\newcommand{\CC}{\mathbb{C}}
\newcommand{\FF}{\mathbb{F}}

\newcommand{\ep}{\varepsilon}
\usepackage{tikz-cd}
\usetikzlibrary{calc}
\usetikzlibrary{decorations.pathmorphing}
% Quiver curved arrow support
\tikzset{curve/.style={settings={#1},to path={(\tikztostart)
    .. controls ($(\tikztostart)!\pv{pos}!(\tikztotarget)!\pv{height}!270:(\tikztotarget)$)
    and ($(\tikztostart)!1-\pv{pos}!(\tikztotarget)!\pv{height}!270:(\tikztotarget)$)
    .. (\tikztotarget)\tikztonodes}},
    settings/.code={\tikzset{quiver/.cd,#1}
        \def\pv##1{\pgfkeysvalueof{/tikz/quiver/##1}}},
    quiver/.cd,pos/.initial=0.35,height/.initial=0}
\tikzset{tail reversed/.code={\pgfsetarrowsstart{tikzcd to}}}
\tikzset{2tail/.code={\pgfsetarrowsstart{Implies[reversed]}}}
\tikzset{2tail reversed/.code={\pgfsetarrowsstart{Implies}}}
\tikzset{no body/.style={/tikz/dash pattern=on 0 off 1mm}}
\makeatletter
\@ifpackageloaded{bookmark}{}{\usepackage{bookmark}}
\makeatother
\makeatletter
\@ifpackageloaded{caption}{}{\usepackage{caption}}
\AtBeginDocument{%
\ifdefined\contentsname
  \renewcommand*\contentsname{Table of contents}
\else
  \newcommand\contentsname{Table of contents}
\fi
\ifdefined\listfigurename
  \renewcommand*\listfigurename{List of Figures}
\else
  \newcommand\listfigurename{List of Figures}
\fi
\ifdefined\listtablename
  \renewcommand*\listtablename{List of Tables}
\else
  \newcommand\listtablename{List of Tables}
\fi
\ifdefined\figurename
  \renewcommand*\figurename{Figure}
\else
  \newcommand\figurename{Figure}
\fi
\ifdefined\tablename
  \renewcommand*\tablename{Table}
\else
  \newcommand\tablename{Table}
\fi
}
\@ifpackageloaded{float}{}{\usepackage{float}}
\floatstyle{ruled}
\@ifundefined{c@chapter}{\newfloat{codelisting}{h}{lop}}{\newfloat{codelisting}{h}{lop}[chapter]}
\floatname{codelisting}{Listing}
\newcommand*\listoflistings{\listof{codelisting}{List of Listings}}
\makeatother
\makeatletter
\makeatother
\makeatletter
\@ifpackageloaded{caption}{}{\usepackage{caption}}
\@ifpackageloaded{subcaption}{}{\usepackage{subcaption}}
\makeatother
\usepackage{bookmark}
\IfFileExists{xurl.sty}{\usepackage{xurl}}{} % add URL line breaks if available
\urlstyle{same}
\hypersetup{
  pdftitle={Reaching for Infinity},
  pdfauthor={Steve Trettel},
  hidelinks,
  pdfcreator={LaTeX via pandoc}}


\title{Reaching for Infinity}
\author{Steve Trettel}
\date{2026-04-16}
\begin{document}
\frontmatter
\maketitle

\renewcommand*\contentsname{Table of contents}
{
\setcounter{tocdepth}{1}
\tableofcontents
}

\mainmatter
\bookmarksetup{startatroot}

\chapter*{Preface}\label{preface}
\addcontentsline{toc}{chapter}{Preface}

\markboth{Preface}{Preface}

Analysis Text, Draft III, opening motivation.

\part{Numbers}

\section*{Part I Introduction: Taming
Infinity}\label{part-i-introduction-taming-infinity}
\addcontentsline{toc}{section}{Part I Introduction: Taming Infinity}

\markright{Part I Introduction: Taming Infinity}

\subsection*{We need infinite processes (and always
have)}\label{we-need-infinite-processes-and-always-have}
\addcontentsline{toc}{subsection}{We need infinite processes (and always
have)}

\begin{itemize}
\tightlist
\item
  The diagonal of a unit square has length \(\sqrt{2}\), but
  \(\sqrt{2}\) is irrational --- no fraction represents it exactly
\item
  The circumference of a circle involves \(\pi\), which is not only
  irrational but transcendental --- no polynomial equation with integer
  coefficients has \(\pi\) as a root
\item
  These aren't exotic curiosities; they're the most basic measurements
  in geometry
\item
  The only access we have to these numbers is through \emph{processes
  that converge to them}: sequences of rational approximations that get
  closer and closer to a target we can't write down in finite terms
\item
  The Babylonians (\textasciitilde1800 BCE) computed \(\sqrt{2}\) by
  iterating \(x_{n+1} = \frac{1}{2}(x_n + 2/x_n)\), producing the
  approximations
  \(1, \frac{3}{2}, \frac{17}{12}, \frac{577}{408}, \ldots\) --- each
  one rational, each one closer to \(\sqrt{2}\) than the last
\item
  Archimedes (\textasciitilde250 BCE) trapped \(\pi\) between the
  perimeters of inscribed and circumscribed regular polygons: a 96-gon
  gave him \(3\frac{10}{71} < \pi < 3\frac{1}{7}\)
\item
  In both cases the \emph{method} is an infinite process --- a sequence
  of improving approximations --- and the \emph{target} is a number that
  can only be defined as the thing these approximations converge to
\item
  Fast-forward two millennia: Newton and Leibniz build calculus entirely
  on infinite processes. The derivative is a limit of difference
  quotients. The integral is a limit of sums. Power series represent
  functions as infinite sums of polynomials. By the 18th century,
  infinite processes aren't a niche tool for approximating \(\pi\) ---
  they are the language of mathematics itself
\end{itemize}

\subsection*{Infinite processes are
dangerous}\label{infinite-processes-are-dangerous}
\addcontentsline{toc}{subsection}{Infinite processes are dangerous}

\begin{itemize}
\tightlist
\item
  In the hands of a master like Euler, formal manipulation of infinite
  series produces extraordinary truths
\item
  Consider the \emph{Basel problem} (1735): find the exact value of
  \(1 + \frac{1}{4} + \frac{1}{9} + \frac{1}{16} + \cdots\)
\item
  Euler's approach: the function \(\sin(x)/x\) vanishes at
  \(x = \pm\pi, \pm 2\pi, \pm 3\pi, \ldots\) --- so treat it like a
  polynomial with these roots and factor it:
  \[\frac{\sin x}{x} = \left(1 - \frac{x^2}{\pi^2}\right)\left(1 - \frac{x^2}{4\pi^2}\right)\left(1 - \frac{x^2}{9\pi^2}\right)\cdots\]
\item
  Expand the right side and compare the coefficient of \(x^2\) with the
  known series \(\sin(x)/x = 1 - x^2/6 + \cdots\)
\item
  This gives
  \(\frac{1}{\pi^2} + \frac{1}{4\pi^2} + \frac{1}{9\pi^2} + \cdots = \frac{1}{6}\),
  hence \(\sum_{n=1}^{\infty} \frac{1}{n^2} = \frac{\pi^2}{6}\)
\item
  The answer is correct. The reasoning --- treating a power series as a
  polynomial with infinitely many roots --- has no logical foundation
  whatsoever. Euler knew this; he spent years trying to find a rigorous
  proof. But at the time, the formal algebra was all anyone had
\item
  The trouble is that the \emph{same} style of formal reasoning also
  produces nonsense
\item
  The geometric series gives
  \(\frac{1}{1-x} = 1 + x + x^2 + x^3 + \cdots\)
\item
  Substitute \(x = -1\): we get
  \(\frac{1}{2} = 1 - 1 + 1 - 1 + \cdots\), already dubious
\item
  Substitute \(x = 2\): we get \(-1 = 1 + 2 + 4 + 8 + \cdots\), plainly
  absurd
\item
  The formal algebra doesn't know which substitutions are legitimate; it
  treats them all the same. Something beyond algebra is needed to
  separate the valid from the invalid
\end{itemize}

\subsubsection*{The rearrangement shock}\label{the-rearrangement-shock}
\addcontentsline{toc}{subsubsection}{The rearrangement shock}

\begin{itemize}
\tightlist
\item
  Here is a more subtle danger --- one that shows infinite sums don't
  even obey the most basic rules of arithmetic
\item
  The \emph{alternating harmonic series} converges:
  \[1 - \frac{1}{2} + \frac{1}{3} - \frac{1}{4} + \frac{1}{5} - \frac{1}{6} + \cdots = \ln 2 \approx 0.693\]
\item
  For a finite sum, the order of the terms doesn't matter: addition is
  commutative
\item
  Now rearrange the same series --- take two positive terms, then one
  negative, repeating:
  \[\left(1 + \frac{1}{3}\right) - \frac{1}{2} + \left(\frac{1}{5} + \frac{1}{7}\right) - \frac{1}{4} + \left(\frac{1}{9} + \frac{1}{11}\right) - \frac{1}{6} + \cdots\]
\item
  Every term of the original series appears exactly once. But the
  partial sums of this rearranged series converge to
  \(\frac{3}{2}\ln 2 \approx 1.039\), not \(\ln 2\)
\item
  Same terms, different order, different sum --- impossible for finite
  addition, routine for infinite series
\item
  It gets worse: Riemann proved that a \emph{conditionally convergent}
  series (one that converges, but not absolutely) can be rearranged to
  converge to \emph{any real number whatsoever}, or to diverge to
  \(+\infty\) or \(-\infty\)
\item
  The commutative law of addition, which we take for granted, simply
  \emph{fails} for infinite sums in general
\end{itemize}

\subsubsection*{Limits don't always
commute}\label{limits-dont-always-commute}
\addcontentsline{toc}{subsubsection}{Limits don't always commute}

\begin{itemize}
\tightlist
\item
  Beyond rearrangement, there is the problem of \emph{interchanging}
  infinite processes
\item
  Cauchy (\textasciitilde1821) claimed to prove: if a sequence of
  continuous functions \(f_1, f_2, f_3, \ldots\) converges at every
  point, the limit function is also continuous
\item
  The theorem is false. Take a sequence of smooth bump functions that
  converge pointwise to a step function --- each \(f_n\) is continuous,
  their pointwise limit is not
\item
  Cauchy's error was not carelessness; it was a failure to distinguish
  between \emph{pointwise convergence} (for each fixed \(x\), the
  sequence \(f_n(x)\) converges) and \emph{uniform convergence} (the
  convergence happens at the same rate for all \(x\) simultaneously)
\item
  The distinction is subtle, invisible to informal reasoning, and
  crucial --- it's exactly the kind of thing a rigorous theory of
  convergence must make precise
\end{itemize}

\subsubsection*{The lesson}\label{the-lesson}
\addcontentsline{toc}{subsubsection}{The lesson}

\begin{itemize}
\tightlist
\item
  Infinite processes don't obey the rules of finite arithmetic
\item
  Rearrangement can change the sum; limits don't always commute with
  other operations; properties of approximations don't automatically
  pass to the limit
\item
  We need a \emph{theory} that tells us precisely when these operations
  are safe and when they aren't
\item
  Building that theory is the project of Part I
\end{itemize}

\subsection*{The foundations weren't
there}\label{the-foundations-werent-there}
\addcontentsline{toc}{subsection}{The foundations weren't there}

\begin{itemize}
\tightlist
\item
  In the early 19th century, mathematicians began trying to \emph{prove}
  the basic theorems of calculus from first principles
\item
  They discovered they couldn't --- not because the theorems were wrong,
  but because the number system hadn't been specified precisely enough
\item
  The pattern repeated across decades: someone would identify a property
  the real line ``obviously'' has, try to prove it, and find a gap that
  no amount of cleverness could fill
\end{itemize}

\subsubsection*{Bolzano and the Intermediate Value Theorem
(1817)}\label{bolzano-and-the-intermediate-value-theorem-1817}
\addcontentsline{toc}{subsubsection}{Bolzano and the Intermediate Value
Theorem (1817)}

\begin{itemize}
\tightlist
\item
  If a continuous function satisfies \(f(a) < 0\) and \(f(b) > 0\),
  there must be some \(c\) between \(a\) and \(b\) where \(f(c) = 0\)
\item
  Geometrically obvious: a continuous curve that starts below the axis
  and ends above it must cross somewhere
\item
  Bolzano recognized --- remarkably, for his time --- that
  ``geometrically obvious'' is not a proof
\item
  His argument essentially constructs a supremum: the least upper bound
  of the set \(\{x \in [a,b] : f(x) < 0\}\)
\item
  But he can't justify that this supremum \emph{exists} --- that step
  requires a property of the number system he doesn't have as an axiom
\end{itemize}

\subsubsection*{Cauchy and the convergence criterion
(\textasciitilde1821)}\label{cauchy-and-the-convergence-criterion-1821}
\addcontentsline{toc}{subsubsection}{Cauchy and the convergence
criterion (\textasciitilde1821)}

\begin{itemize}
\tightlist
\item
  Cauchy formulates a beautiful idea: a sequence should converge if its
  terms get arbitrarily close to each other
\item
  Formally: for every tolerance \(\varepsilon > 0\), there is a stage
  \(N\) such that \(|a_m - a_n| < \varepsilon\) for all \(m, n > N\)
\item
  This captures the \emph{idea} of convergence without reference to the
  limit value --- you can test whether a sequence converges without
  knowing \emph{what} it converges to
\item
  But the criterion doesn't work over the rational numbers: the
  Babylonian sequence
  \(\frac{3}{2}, \frac{17}{12}, \frac{577}{408}, \ldots\) for
  \(\sqrt{2}\) satisfies Cauchy's criterion perfectly --- the terms get
  closer and closer together --- yet converges to nothing in \(\QQ\).
  The sequence is reaching for a number that doesn't exist in the system
\item
  The criterion is correct in \(\RR\), where completeness fills the
  holes; it fails in \(\QQ\), where the holes are exactly the problem
\end{itemize}

\subsubsection*{The Archimedean
property}\label{the-archimedean-property}
\addcontentsline{toc}{subsubsection}{The Archimedean property}

\begin{itemize}
\tightlist
\item
  Since antiquity, mathematicians had used a more basic property: given
  any two positive quantities, some multiple of the smaller exceeds the
  larger
\item
  This rules out infinitely large and infinitely small elements, and
  it's what makes Archimedes' polygon method work --- the inscribed and
  circumscribed perimeters actually close in on their target, rather
  than hovering infinitely far away
\item
  Archimedes used this explicitly; Newton and Leibniz needed it
  implicitly whenever they argued that a sequence of approximations
  could be made ``as close as desired''
\item
  The property is clearly \emph{necessary} for the number line to behave
  as we expect
\item
  But it is not \emph{sufficient}: the rational numbers are Archimedean,
  and they are full of holes
\item
  The Archimedean property says the number line has no \emph{infinitely
  wide} gaps, but it says nothing about the \emph{infinitely thin} gaps
  --- the missing irrationals, the missing limits of Cauchy sequences
\end{itemize}

\subsubsection*{Nested intervals}\label{nested-intervals}
\addcontentsline{toc}{subsubsection}{Nested intervals}

\begin{itemize}
\tightlist
\item
  If you have a sequence of closed intervals
  \([a_1, b_1] \supseteq [a_2, b_2] \supseteq [a_3, b_3] \supseteq \cdots\)
  with lengths shrinking to zero, their intersection should be a single
  point --- the unique number that all the intervals close in on
\item
  This is the implicit logic behind \emph{every} approximation scheme:
  the Babylonian algorithm traps \(\sqrt{2}\) in nested intervals;
  Archimedes traps \(\pi\); the bisection method for the IVT traps the
  root
\item
  Yet in \(\QQ\), nested intervals can have \emph{empty intersection}
  --- the point they close in on might be irrational, so it simply isn't
  there
\end{itemize}

\subsubsection*{A common pattern}\label{a-common-pattern}
\addcontentsline{toc}{subsubsection}{A common pattern}

\begin{itemize}
\tightlist
\item
  The Archimedean property, Cauchy's criterion, Bolzano's supremum
  argument, nested intervals --- each identifies something the real line
  should satisfy
\item
  Each fails over the rational numbers
\item
  And the remarkable fact, which only became clear later: \emph{these
  are all manifestations of the same missing property}
\end{itemize}

\subsection*{The resolution, and a
surprise}\label{the-resolution-and-a-surprise}
\addcontentsline{toc}{subsection}{The resolution, and a surprise}

\subsubsection*{Completeness}\label{completeness}
\addcontentsline{toc}{subsubsection}{Completeness}

\begin{itemize}
\tightlist
\item
  In 1872, Dedekind and Cantor independently gave rigorous constructions
  of the real numbers from the rationals --- Dedekind via ``cuts''
  (partitions of \(\QQ\)), Cantor via equivalence classes of Cauchy
  sequences
\item
  The constructions look different, but they produce the same object: a
  \emph{complete ordered field}
\item
  We take the axiomatic route: rather than building \(\RR\) from
  scratch, we \emph{characterize} it by its properties
\item
  \(\RR\) is an ordered field (it has addition, multiplication, and an
  order, all compatible) --- this much \(\QQ\) already satisfies
\item
  The one additional axiom: \emph{every nonempty set of real numbers
  that is bounded above has a least upper bound (supremum)}
\item
  This single property --- \emph{completeness} --- is the one thing
  \(\QQ\) lacks
\item
  It is an axiom that goes beyond arithmetic and algebra; it is a
  statement about the \emph{topological} structure of the number line,
  asserting that it has no holes
\end{itemize}

\subsubsection*{Five faces of one
property}\label{five-faces-of-one-property}
\addcontentsline{toc}{subsubsection}{Five faces of one property}

\begin{itemize}
\tightlist
\item
  With completeness in hand, everything works
\item
  \(\sqrt{2}\) exists: it is the supremum of \(\{x > 0 : x^2 < 2\}\)
\item
  \(\pi\) exists: it is the common limit of Archimedes' inscribed and
  circumscribed polygon perimeters
\item
  The Intermediate Value Theorem holds; Bolzano's argument goes through
\item
  The Monotone Convergence Theorem holds: a bounded increasing sequence
  converges to its supremum
\item
  The nested interval property holds: completeness provides the point
  that the intervals close in on
\item
  Every bounded sequence has a convergent subsequence
  (Bolzano-Weierstrass)
\item
  Cauchy's convergence criterion works: every Cauchy sequence converges
\item
  And the deepest revelation: \emph{all five of these properties are
  logically equivalent.} Completeness, monotone convergence, nested
  intervals, Bolzano-Weierstrass, and the Cauchy criterion are different
  faces of the same axiom. Assume any one of them, and the other four
  follow. This is what Archimedes, Bolzano, and Cauchy were each
  grasping at from different directions: a single structural property of
  the continuum
\end{itemize}

\subsubsection*{The uncountability of the
continuum}\label{the-uncountability-of-the-continuum}
\addcontentsline{toc}{subsubsection}{The uncountability of the
continuum}

\begin{itemize}
\tightlist
\item
  Having constructed \(\RR\), Cantor made an unexpected discovery: the
  real numbers are \emph{uncountable}
\item
  The rationals, the algebraic numbers, even the computable numbers are
  all \emph{countable} --- they can be listed as
  \(a_1, a_2, a_3, \ldots\)
\item
  But \(\RR\) cannot be listed this way; any attempted list necessarily
  misses real numbers
\item
  We built \(\RR\) to fill the holes in \(\QQ\) --- missing roots,
  missing limits --- but we got far more than we asked for: an
  uncountably infinite collection, vastly larger than any number system
  humans had ever explicitly constructed
\item
  Most real numbers cannot be named, described, or computed by any
  finite procedure
\item
  Infinite processes are not just a convenience for reaching numbers
  like \(\sqrt{2}\) or \(\pi\) --- they are the \emph{only possible
  language} for the continuum. Even writing down every infinite series,
  every continued fraction, every algorithm we could ever devise, we
  could only name countably many real numbers. The rest exist by the
  completeness axiom and nothing else
\end{itemize}

\subsection*{What's ahead}\label{whats-ahead}
\addcontentsline{toc}{subsection}{What's ahead}

\begin{itemize}
\tightlist
\item
  Chapter 1 builds the stage: the real number system, characterized by
  the completeness axiom, and the geometry it unlocks --- lengths,
  areas, and angles defined rigorously for the first time
\item
  Chapter 2 develops the fundamental tool: the \emph{limit}. Every real
  number is the limit of a sequence of rationals, so limits are the
  bridge between the countable world of explicit computation and the
  uncountable continuum. The chapter builds from the
  \(\varepsilon\)-\(N\) definition through the limit laws and the
  Monotone Convergence Theorem to decimal representations and
  Archimedes' quadrature of the parabola
\item
  Chapter 3 asks what completeness \emph{forces}: even chaotic bounded
  sequences must contain convergent subsequences (Bolzano-Weierstrass),
  and the five properties we identified above --- least upper bounds,
  monotone convergence, nested intervals, Bolzano-Weierstrass, the
  Cauchy criterion --- are proved equivalent. Five faces, one axiom
\item
  Chapter 4 puts these tools to work on the infinite processes that
  motivated the whole enterprise: alternating series, comparison tests,
  contraction mappings, continued fractions, and the rigorous
  computation of \(\pi\) via Archimedes' polygon method --- finally
  completing the argument he started twenty-two centuries earlier
\item
  Chapter 5 returns to the dangers we saw above and gives precise
  answers. Absolute convergence is the condition that makes
  rearrangement safe; domination is the condition that makes limit--sum
  interchange safe. The Cauchy product of series is developed, the
  exponential function gets its functional equation, and dominated
  convergence tells us exactly when an infinite sum commutes with a
  limit. By the end, we know \emph{when} infinite processes behave like
  finite ones, and \emph{when} they don't
\end{itemize}

\chapter{The Real Line}\label{the-real-line}

\textbf{The first crisis: √2} - Pythagoreans (\textasciitilde500 BCE):
discovered the diagonal of a unit square is incommensurable with the
side - Proof that √2 is irrational --- no ratio of integers gives the
diagonal - Crisis: their philosophy held ``all is number'' (meaning
ratios of integers) - The response: separate number from magnitude;
geometry proceeds without arithmetic foundation

\textbf{The Greek resolution (or avoidance)} - Eudoxus
(\textasciitilde370 BCE): theory of proportion for magnitudes, avoiding
irrational ``numbers'' - Euclid's Elements: ratios of magnitudes, not
numbers - ``Continuous magnitude'' used but never defined - Arithmetic
and geometry remain separate for two millennia - Curved objects posed
further challenges: what is the circumference of a circle? its area? -
Archimedes trapped π between inscribed and circumscribed polygons ---
but π itself was not a ``number''

\textbf{Reunification without foundations} - Descartes
(\textasciitilde1637): coordinate geometry reunites number and space -
But what \emph{are} the numbers on the number line? - Newton, Leibniz
(\textasciitilde1680s): calculus uses ``infinitely small'' quantities,
limits, continuity - It works --- but what is ``the continuum'' they're
computing with? - 18th century: calculus develops rapidly, foundations
remain vague

\textbf{The 19th century reckoning} - Cauchy, Weierstrass: rigorize
limits and continuity (the subject of Chapter 2) - But limits of
\emph{what}? Need the real numbers themselves to be rigorous - Dedekind
(1872): defines reals via ``cuts'' --- partitions of ℚ - Cantor (1872):
defines reals via Cauchy sequences of rationals - Finally: rigorous
constructions of ℝ from ℚ - Surprise: Cantor discovers ℝ is uncountable
--- a new kind of infinity

\textbf{Our approach} - We inherit their work but take the axiomatic
path - Define ℝ by its properties: a complete ordered field -
Completeness is the ONE property ℚ lacks - We don't construct ℝ; we
characterize it and use it - Existence and uniqueness: there is
essentially one complete ordered field

\textbf{What we'll see in this chapter} - §1.1: Foundations --- the
language, algebraic rules, and order structure we need - §1.2: Why ℚ
isn't enough --- gaps that break calculus - §1.3: The completeness
property --- the key axiom - §1.4: The real line --- roots exist,
decimals make sense, ℝ is uncountable - §1.5: Measuring the world ---
with ℝ in hand, we can rigorously define lengths, areas, and angles -
The payoff: √2 finally exists; π finally exists; curved geometry is
finally rigorous

\section{§1.1 Foundations}\label{foundations}

\subsection{Narrative}\label{narrative}

This section establishes the language and rules for the course. Students
should know most of this material from a proofs course; we collect it
here for reference and to fix notation. The section is intentionally
brisk --- definitions, notation, key facts. We build up to ordered
fields, which provide the algebraic and order structure for
\(\mathbb{Q}\) and (soon) \(\mathbb{R}\).

\subsection{Content}\label{content}

\subsection{Logic and Sets}\label{logic-and-sets}

\begin{itemize}
\tightlist
\item
  Notation: \(\forall\) (for all), \(\exists\) (there exists),
  \(\Rightarrow\) (implies), \(\Leftrightarrow\) (iff), \(\neg\) (not)
\item
  Negation: \(\neg(\forall x: P(x))\) is \(\exists x: \neg P(x)\);
  \(\neg(\exists x: P(x))\) is \(\forall x: \neg P(x)\)
\item
  Remark: comfort with quantifiers is essential for
  \(\varepsilon\)-\(N\) arguments
\item
  Sets: \(x \in A\), \(A \subseteq B\), \(A = B\), \(\{x : P(x)\}\)
\item
  Operations: \(A \cup B\), \(A \cap B\), \(A \setminus B\),
  \(A \times B\)
\item
  Power set: \(\mathcal{P}(A) = \{S : S \subseteq A\}\)
\item
  Standard sets: \(\NN = \{1, 2, 3, \ldots\}\), \(\ZZ\),
  \(\QQ = \{p/q : p, q \in \ZZ, q \neq 0\}\)
\item
  Indexed unions/intersections: \(\bigcup_{i \in I} A_i\),
  \(\bigcap_{i \in I} A_i\)
\item
  Induction: if \(P(1)\) and \(P(n) \Rightarrow P(n+1)\) for all \(n\),
  then \(P(n)\) for all \(n \in \NN\)
\item
  Strong induction: if \(P(1), \ldots, P(n)\) imply \(P(n+1)\), then
  \(P(n)\) for all \(n\)
\item
  Well-ordering: every nonempty subset of \(\NN\) has a least element
\item
  Remark: induction, strong induction, and well-ordering are equivalent
\end{itemize}

\subsection{Functions and Cardinality}\label{functions-and-cardinality}

\begin{itemize}
\tightlist
\item
  Definition: \(f: A \to B\) assigns to each \(a \in A\) exactly one
  \(f(a) \in B\)
\item
  Remark: functions are assignments, not formulas
\item
  Injective: \(f(x_1) = f(x_2) \Rightarrow x_1 = x_2\)
\item
  Surjective: every \(y \in B\) is hit
\item
  Bijective: both; \(f^{-1}\) exists iff \(f\) is bijective
\item
  Composition: \((g \circ f)(x) = g(f(x))\)
\item
  Cardinality: \(|A| = |B|\) if there exists a bijection \(A \to B\)
\item
  Countable: finite or same cardinality as \(\NN\); uncountable: not
  countable
\item
  Theorem (Cantor): no surjection \(A \to \mathcal{P}(A)\) --- multiple
  sizes of infinity
\item
  Theorem: \(\NN \times \NN\) is countable; \(\QQ\) is countable
\item
  Theorem: finite product of countable sets is countable
\item
  Theorem: countable union of countable sets is countable
\end{itemize}

\subsection{Fields}\label{fields}

\begin{itemize}
\tightlist
\item
  Definition: a field \((F, +, \cdot)\) satisfies:

  \begin{itemize}
  \tightlist
  \item
    \((F, +)\): associative, commutative, identity \(0\), inverses
    \(-a\)
  \item
    \((F \setminus \{0\}, \cdot)\): associative, commutative, identity
    \(1\), inverses \(a^{-1}\)
  \item
    Distributivity: \(a(b + c) = ab + ac\)
  \end{itemize}
\item
  Example: \(\QQ\) is a field; \(\ZZ\) is not (no multiplicative
  inverses)
\item
  Notation: \(2 = 1 + 1\), \(3 = 2 + 1\), etc.
\item
  Theorem: \(0 \cdot a = 0\) (prove in text)
\item
  Theorem: \((-1) \cdot a = -a\) (prove in text)
\item
  Theorem: \(ab = 0 \Rightarrow a = 0\) or \(b = 0\)
\item
  Remark: all the algebra you learned in school follows from these
  axioms
\item
  ✎ Inline: Prove \((-a)(-b) = ab\)
\item
  ✎ Inline: Prove \(2 + 2 = 4\) from the axioms
\end{itemize}

\subsection{Order}\label{order}

\begin{itemize}
\tightlist
\item
  Definition: a total order \(<\) on a set \(S\) satisfies:

  \begin{itemize}
  \tightlist
  \item
    Transitivity: \(a < b\) and \(b < c \Rightarrow a < c\)
  \item
    Trichotomy: exactly one of \(a < b\), \(a = b\), \(a > b\) holds
  \end{itemize}
\item
  Notation: \(a \leq b\) means \(a < b\) or \(a = b\); similarly
  \(\geq\), \(>\)
\item
  Definition: an ordered field is a field with a total order \(<\)
  satisfying:

  \begin{itemize}
  \tightlist
  \item
    \(a < b \Rightarrow a + c < b + c\) (order compatible with addition)
  \item
    \(a < b\) and \(c > 0 \Rightarrow ac < bc\) (order compatible with
    multiplication)
  \end{itemize}
\item
  Example: \(\QQ\) is an ordered field
\item
  Theorem: \(a^2 \geq 0\) for all \(a\), with equality iff \(a = 0\)
  (prove in text)
\item
  Corollary: \(1 > 0\)
\item
  ✎ Inline: Prove \(a < b \Rightarrow a + c < b + c\)
\item
  ✎ Inline: Prove \(a < b\) and \(c < 0 \Rightarrow ac > bc\)
\item
  Remark (positivity): equivalently, define ordered field via a set
  \(P\) of positive elements satisfying
  \(a, b \in P \Rightarrow a + b, ab \in P\) and trichotomy
  (\(a \in P\), \(a = 0\), or \(-a \in P\)). Then
  \(a < b \Leftrightarrow b - a \in P\).
\end{itemize}

\subsection{Absolute Value and
Intervals}\label{absolute-value-and-intervals}

\begin{itemize}
\tightlist
\item
  Definition: \(|a| = a\) if \(a \geq 0\), \(|a| = -a\) if \(a < 0\)
\item
  Properties: \(|a| \geq 0\); \(|a| = 0\) iff \(a = 0\);
  \(|ab| = |a||b|\)
\item
  Theorem (Triangle Inequality): \(|a + b| \leq |a| + |b|\)
\item
  Corollary: \(||a| - |b|| \leq |a - b|\) (reverse triangle inequality)
\item
  Geometric: \(|a - b|\) is distance between \(a\) and \(b\)
\item
  Intervals: \((a, b)\), \([a, b]\), \([a, b)\), \((a, b]\); unbounded:
  \([a, \infty)\), \((-\infty, b]\), etc.
\item
  Notation: \(\infty\), \(-\infty\) are symbols, not numbers
\end{itemize}

\subsection{Algebraic Tools}\label{algebraic-tools}

\begin{itemize}
\tightlist
\item
  Powers: \(a^n = a \cdot a \cdots a\) (\(n\) times); \(a^0 = 1\);
  \(a^{-n} = 1/a^n\)
\item
  Laws: \(a^{m+n} = a^m a^n\), \((a^m)^n = a^{mn}\),
  \((ab)^n = a^n b^n\)
\item
  Notation: \(\sqrt[n]{a}\) denotes \(b\) with \(b^n = a\) (when such
  \(b\) exists)
\item
  Remark: existence of \(\sqrt[n]{a}\) for general \(a\) requires
  completeness (§1.3)
\item
  Theorem (Bernoulli): \((1 + x)^n \geq 1 + nx\) for \(x \geq -1\),
  \(n \in \NN\) (proof by induction)
\item
  ✎ Inline: Verify the induction step of Bernoulli's inequality
\item
  Theorem (Binomial):
  \((a + b)^n = \sum_{k=0}^{n} \binom{n}{k} a^{n-k} b^k\) where
  \(\binom{n}{k} = \frac{n!}{k!(n-k)!}\)
\item
  Pascal's identity: \(\binom{n}{k} + \binom{n}{k-1} = \binom{n+1}{k}\)
\item
  Remark: binomial theorem proved by induction using Pascal's identity
\item
  Theorem (AM-GM, two variables): \(\frac{a+b}{2} \geq \sqrt{ab}\) for
  \(a, b \geq 0\) (when \(\sqrt{ab}\) exists), equality iff \(a = b\)
\item
  Remark: AM-GM generalizes to \(n\) variables (exercise)
\end{itemize}

\subsection{Guided Exercise}\label{guided-exercise}

\subsection{\texorpdfstring{Base-\(b\) Representation of
Integers}{Base-b Representation of Integers}}\label{base-b-representation-of-integers}

\begin{enumerate}
\def\labelenumi{(\alph{enumi})}
\item
  Show: for any \(n \in \NN\) and \(b \geq 2\), there exist unique
  digits \(d_0, \ldots, d_k\) with \(0 \leq d_i < b\), \(d_k \neq 0\),
  such that \(n = \sum_{i=0}^{k} d_i b^i\).
\item
  Describe the greedy algorithm: repeatedly divide by \(b\), record
  remainders.
\item
  Convert \(42\) to base \(2\), base \(3\), base \(16\).
\item
  In base \(b\), what is the representation of \(b^m - 1\)?
\end{enumerate}

\subsection{Exercises}\label{exercises}

\emph{Logic and Sets} - Prove De Morgan: \((A \cup B)^c = A^c \cap B^c\)
and \((A \cap B)^c = A^c \cup B^c\) - Prove
\(A \setminus (B \cup C) = (A \setminus B) \cap (A \setminus C)\) -
Prove \((A \cup B) \setminus C = (A \setminus C) \cup (B \setminus C)\)

\emph{Functions and Cardinality} - Prove: \(f\) bijective iff \(f\) has
a two-sided inverse - Prove: composition of bijections is a bijection -
Prove Cantor's theorem: no surjection \(A \to \mathcal{P}(A)\) - Show
\(\NN \times \NN\) is countable - Show \(\QQ\) is countable - Show
finite product of countable sets is countable - Show countable union of
countable sets is countable - Prove: a set with \(n\) elements has
\(2^n\) subsets

\emph{Induction} - Prove \(1 + 2 + \cdots + n = n(n+1)/2\) - Prove
\(1^2 + 2^2 + \cdots + n^2 = n(n+1)(2n+1)/6\) - Prove
\(1^3 + 2^3 + \cdots + n^3 = (n(n+1)/2)^2\) (the sum of cubes equals the
square of the sum!) - Prove \(1 + 3 + 5 + \cdots + (2n-1) = n^2\) -
Prove \(n^3 - n\) is divisible by 6 for all \(n \geq 1\) - Prove
\(4^n - 1\) is divisible by 3 for all \(n \geq 1\) - Prove \(2^n > n\)
for all \(n \geq 1\) - Prove \(n! > 2^n\) for \(n \geq 4\) - Prove
\(2^n \geq n^2\) for \(n \geq 4\) - (Strong induction) Prove every
\(n \geq 2\) is a product of primes - (Strong induction) Prove
\(F_n < 2^n\) for all \(n \geq 1\), where \(F_n\) is the Fibonacci
number

\emph{Fields} - Prove uniqueness of additive inverse - Prove uniqueness
of multiplicative inverse - Prove \((-a)(-b) = ab\) - Prove
\((a+b)(a-b) = a^2 - b^2\)

\emph{Order} - Prove: \(0 < a < b \Rightarrow a^2 < b^2\) - Prove:
\(0 < a < b \Rightarrow 0 < 1/b < 1/a\) - Prove transitivity of \(<\)

\emph{Absolute Value} - Prove the reverse triangle inequality:
\(||a| - |b|| \leq |a - b|\) - Prove \(|a - b| \leq |a - c| + |c - b|\)

\emph{Algebraic Tools} - Prove Bernoulli's inequality by induction -
Prove strengthened Bernoulli:
\((1+x)^n \geq 1 + nx + \frac{n(n-1)}{2}x^2\) for \(x \geq 0\) - Prove
the binomial theorem by induction - Prove Pascal's identity - Prove
\(n! \geq 2^{n-1}\) for \(n \geq 1\) - Prove
\(\binom{n}{k} \leq n^k / k!\) - Prove AM-GM for two variables - ★ Prove
AM-GM for \(n\) variables

\subsection{Dependencies}\label{dependencies}

\textbf{Requires}: Nothing --- this is foundational.

\textbf{Used in}: All subsequent sections.

\section{\texorpdfstring{The Incompleteness of
\(\QQ\)}{The Incompleteness of \textbackslash QQ}}\label{the-incompleteness-of-qq}

\subsection{Narrative}\label{narrative-1}

\begin{itemize}
\tightlist
\item
  Greeks discover \(\sqrt{2}\) irrational: geometry demands a number
  \(\QQ\) doesn't have
\item
  First patch attempt: adjoin \(\sqrt{2}\) to get
  \(\QQ(\sqrt{2})\)---but now \(\sqrt{3}\) is missing
\item
  Adjoin more roots? Still miss \(\sqrt{5}\), \(\sqrt{7}\),
  \(\sqrt{2 + \sqrt{3}}\), \ldots{}
\item
  Go to infinity: constructible numbers (all nested square roots)---but
  miss \(\sqrt[3]{2}\)
\item
  Go to infinity again: algebraic numbers \(\AA\) (all polynomial
  roots)---but miss \(\pi\)
\item
  Hopelessness: every patch leaves holes, can't construct our way to
  completeness
\item
  Deeper problem: curved geometry requires \emph{definitions} we can't
  even state in \(\QQ\)
\item
  Length of a curve = sup of inscribed polygons --- but what if that sup
  doesn't exist?
\item
  We've seen sequences of rationals can converge to nothing
  (\(\sqrt{2}\))
\item
  The crisis isn't just missing numbers --- it's that basic definitions
  become meaningless
\item
  Finally: gaps also break calculus --- IVT, EVT, MVT all fail in
  \(\QQ\)
\end{itemize}

\subsection{Content}\label{content-1}

\subsection{\texorpdfstring{The First Crisis:
\(\sqrt{2}\)}{The First Crisis: \textbackslash sqrt\{2\}}}\label{the-first-crisis-sqrt2}

\begin{itemize}
\tightlist
\item
  Theorem: \(\sqrt{2}\) is irrational (prove in text)
\item
  The problem: geometry demands \(\sqrt{2}\) exist (diagonal of unit
  square), but \(\QQ\) doesn't have it
\item
  The Babylonian sequence \(a_1 = 1\),
  \(a_{n+1} = \frac{1}{2}(a_n + 2/a_n)\) gives better and better
  approximations
\item
  But approximations to \emph{what}? In \(\QQ\), there's no number being
  approximated
\end{itemize}

\subsection{\texorpdfstring{Attempts to Patch
\(\QQ\)}{Attempts to Patch \textbackslash QQ}}\label{attempts-to-patch-qq}

\begin{itemize}
\tightlist
\item
  Definition: \(\QQ(\sqrt{2}) = \{a + b\sqrt{2} : a, b \in \QQ\}\)
\item
  \(\QQ(\sqrt{2})\) is an ordered field, but
  \(\sqrt{3} \notin \QQ(\sqrt{2})\)
\item
  Definition: Constructible numbers --- iterated square roots from
  \(\QQ\)
\item
  Doubling the cube impossible: \(\sqrt[3]{2}\) not constructible
  (Wantzel 1837)
\item
  Definition: Algebraic numbers \(\AA\) --- roots of polynomials with
  integer coefficients
\item
  Definition: Transcendental --- not algebraic
\item
  \(\pi\) is transcendental (Lindemann 1882)
\item
  Remark: Constructibles satisfy polynomials of degree \(2^n\);
  \(\sqrt[3]{2}\) satisfies degree 3
\item
  The pattern: each fix reveals new holes --- can't construct our way
  out
\end{itemize}

\subsection{The Measurement Problem}\label{the-measurement-problem}

\begin{itemize}
\tightlist
\item
  The ancient challenge: what is the length of a curve? the area of a
  region?
\item
  Natural idea: approximate with inscribed polygons, take ``the limit''
\item
  Length of curve = sup of inscribed polygonal lengths
\item
  Area of region = sup of inscribed polygonal areas (or: when inner and
  outer agree)
\item
  But this definition requires the supremum to \emph{exist}
\item
  In \(\QQ\), suprema of bounded sets need not exist
\item
  Example: \(\{r \in \QQ : r^2 < 2\}\) is bounded above but has no
  supremum in \(\QQ\)
\item
  The Babylonian sequence shows: we can have a sequence of
  approximations converging to \emph{nothing}
\item
  This is a crisis at the level of \emph{definitions}, not theorems
\item
  We cannot even \emph{state} what ``length of a circle'' means without
  knowing the answer exists
\item
  Archimedes could trap \(\pi\) between \(3\frac{10}{71}\) and
  \(3\frac{1}{7}\) --- but what \emph{is} \(\pi\)?
\item
  Remark: once we have completeness (§1.3), these definitions become
  meaningful (§1.5)
\end{itemize}

\subsection{Gaps Break Calculus}\label{gaps-break-calculus}

\begin{itemize}
\tightlist
\item
  Even if we ignore definitional problems, the theorems of calculus fail
  in \(\QQ\)
\item
  IVT fails: \(f(x) = x^2 - 2\) on \([1,2] \cap \QQ\) has
  \(f(1) < 0 < f(2)\) but no root
\item
  EVT fails: \(f(x) = 1/(x^2 - 2)\) is continuous on \([1,2] \cap \QQ\)
  but unbounded
\item
  MVT fails (depends on EVT)
\item
  Bisection algorithm: produces nested intervals \([a_n, b_n]\) with
  \(\bigcap [a_n, b_n] \cap \QQ = \emptyset\)
\item
  Remark: these failures will be repaired by completeness (§1.3) and
  proved in Chapter 6
\end{itemize}

\subsection{Guided Exercises}\label{guided-exercises}

None for this section.

\subsection{Exercises}\label{exercises-1}

\emph{Irrationality} - Prove \(\sqrt{6}\) is irrational - Prove
\(\sqrt{p}\) is irrational for \(p\) prime - Prove
\(\sqrt{2} + \sqrt{3}\) is irrational

\emph{Field extensions} - Verify \(\QQ(\sqrt{2})\) is a field - ★ Prove
\(\sqrt{3} \notin \QQ(\sqrt{2})\) - ★ Prove \(\QQ(\sqrt{2}, \sqrt{3})\)
is a field

\emph{The measurement problem} - Show \(S = \{r \in \QQ : r^2 < 2\}\) is
bounded above but has no sup in \(\QQ\) - The bisection algorithm for
\(\sqrt{2}\): show \(a_n^2 < 2 < b_n^2\) and \(b_n - a_n = 2^{1-n}\) -
Show the set \(\{3, 3.1, 3.14, 3.141, 3.1415, \ldots\}\) (decimal
approximations to \(\pi\)) is bounded in \(\QQ\) but has no sup in
\(\QQ\)

\subsection{Dependencies}\label{dependencies-1}

\textbf{Requires}: §1.1 (ordered field structure, inequalities)

\textbf{Used in}: §1.3 (completeness resolves these issues), §1.4
(\(\AA\) countable), §1.5 (definitions now possible)

\section{Completeness}\label{completeness-1}

\subsection{Narrative}\label{narrative-2}

\begin{itemize}
\tightlist
\item
  Many infinite processes want a ``target'' --- we need a way to name it
\item
  Some are sequences: Babylonian √2, inscribed n-gon perimeters, nested
  intervals
\item
  Some are sets with no natural ordering: all rationals with r²
  \textless{} 2, all inscribed polygons, all rectangles under a curve
\item
  In each case: a collection of approximations, and we want the smallest
  number ≥ all of them
\item
  This is the supremum --- the least upper bound
\item
  Completeness: every nonempty bounded set has a supremum
\item
  This ONE property is what ℚ lacks and ℝ has
\item
  The ε-characterization is the workhorse for arguments
\item
  Nested intervals: the ancient method of trapping between bounds is now
  validated
\end{itemize}

\subsection{Content}\label{content-2}

\emph{(Opening motivation --- unnumbered)}

We've seen that ℚ has holes --- numbers that ``should'' exist but don't.
Now we ask: what property would fill those holes?

Consider the many situations where we want to ``complete'' an infinite
process:

\emph{Sequences approaching a target:} - Babylonian iteration for √2:
\(a_1 = 1\), \(a_{n+1} = \frac{1}{2}(a_n + 2/a_n)\) --- each \(a_n\) is
rational, approaching something - Inscribed regular n-gons: perimeters
\(P_6, P_{12}, P_{24}, \ldots\) approaching the circumference -
Bisection for √2: nested intervals
\([1, 2] \supset [1.4, 1.5] \supset [1.41, 1.42] \supset \cdots\) -
Archimedes' bounds: \(3\frac{10}{71} < \pi < 3\frac{1}{7}\), refining
with each polygon doubling

\emph{Sets with no natural ordering:} - \(\{r \in \QQ : r^2 < 2\}\) ---
all rationals whose square is less than 2 - All inscribed polygons in a
circle (not just regular n-gons --- any polygon) - All polygonal paths
inscribed in a curve - All rectangles contained under a curve
\(y = f(x)\)

In each case, we have a collection of ``approximations from below,'' and
we want the \emph{target} --- the number they're all pointing toward.

For a finite set, we'd take the maximum. But these are infinite sets,
and the maximum may not exist (there may be no largest element).

What we want: the \emph{smallest number that's ≥ all elements of the
set}.

This is the \textbf{supremum}.

\subsection{Bounds and Supremum}\label{bounds-and-supremum}

\begin{itemize}
\tightlist
\item
  Definition: \(M\) is an \emph{upper bound} for \(S\) if \(x \leq M\)
  for all \(x \in S\)
\item
  Definition: \(S\) is \emph{bounded above} if it has an upper bound;
  \emph{bounded below}, \emph{bounded} similarly
\item
  Definition: \(M = \sup S\) (supremum, or least upper bound) if:

  \begin{itemize}
  \tightlist
  \item
    \(M\) is an upper bound for \(S\), and
  \item
    No number smaller than \(M\) is an upper bound
  \end{itemize}
\item
  Definition: \(m = \inf S\) (infimum, or greatest lower bound) ---
  analogous
\item
  Theorem: Supremum is unique (if it exists)
\item
  Convention for unbounded sets:

  \begin{itemize}
  \tightlist
  \item
    If \(S\) is nonempty and unbounded above, write \(\sup S = +\infty\)
  \item
    If \(S\) is nonempty and unbounded below, write \(\inf S = -\infty\)
  \item
    With this convention, every nonempty set has a sup and inf (possibly
    \(\pm\infty\))
  \end{itemize}
\item
  Examples: \(\sup\{1 - 1/n : n \in \NN\} = 1 \notin S\);
  \(\sup[0,1] = 1 \in S\); \(\sup(0,1) = 1 \notin S\)
\item
  Examples: \(\sup \NN = +\infty\); \(\inf \NN = 1\);
  \(\sup \ZZ = +\infty\); \(\inf \ZZ = -\infty\)
\item
  Key example: \(S = \{r \in \QQ : r^2 < 2\}\) is bounded above in
  \(\QQ\) (e.g., by 2) but has no sup in \(\QQ\)
\item
  This is how holes manifest: bounded sets without least upper bounds
\item
  Definition: An ordered field is \emph{complete} if every nonempty set
  bounded above has a supremum
\item
  The Real Numbers: We define \(\RR\) to be a complete ordered field
\item
  Remark: \(\QQ\) satisfies all field and order axioms but not
  completeness --- this ONE property is what it lacks
\item
  Theorem: In a complete ordered field, infima also exist
\item
  Proof: \(\inf S = -\sup(-S)\) where \(-S = \{-x : x \in S\}\)
\end{itemize}

\subsection{The ε-Characterization}\label{the-ux3b5-characterization}

\begin{itemize}
\tightlist
\item
  Theorem: \(M = \sup S\) if and only if:

  \begin{itemize}
  \tightlist
  \item
    \begin{enumerate}
    \def\labelenumi{(\arabic{enumi})}
    \tightlist
    \item
      \(M\) is an upper bound for \(S\), and
    \end{enumerate}
  \item
    \begin{enumerate}
    \def\labelenumi{(\arabic{enumi})}
    \setcounter{enumi}{1}
    \tightlist
    \item
      For all \(\varepsilon > 0\), there exists \(x \in S\) with
      \(x > M - \varepsilon\)
    \end{enumerate}
  \end{itemize}
\item
  Proof in text
\item
  Interpretation: you can get arbitrarily close to sup from below using
  elements of \(S\)
\item
  This is the workhorse for supremum arguments throughout the course
\item
  ✎ Inline: Verify \(\sup\{1 - 1/n : n \in \NN\} = 1\) using the
  ε-characterization
\item
  Theorem: ε-characterization of inf (analogous)
\end{itemize}

\subsection{Working with Suprema}\label{working-with-suprema}

\begin{itemize}
\tightlist
\item
  Theorem: \(\sup(S + c) = \sup S + c\) where
  \(S + c = \{x + c : x \in S\}\)
\item
  Theorem: \(\sup(cS) = c \cdot \sup S\) for \(c > 0\)
\item
  ✎ Inline: Prove \(\sup(cS) = c \cdot \sup S\) for \(c > 0\)
\item
  Theorem: \(A \subseteq B \Rightarrow \sup A \leq \sup B\)
\item
  Theorem: \(\sup(A \cup B) = \max\{\sup A, \sup B\}\)
\item
  Analogous results for inf
\end{itemize}

\subsection{The Nested Interval
Property}\label{the-nested-interval-property}

\begin{itemize}
\tightlist
\item
  The ancients used a powerful technique: trap the target between upper
  and lower bounds, then squeeze
\item
  Babylonians: rectangle sides approaching √2 from above and below
\item
  Archimedes: inscribed polygon \textless{} circumference \textless{}
  circumscribed polygon
\item
  Bisection: \(a_n < \sqrt{2} < b_n\) with \(b_n - a_n \to 0\)
\item
  They trusted this method worked --- now we can prove it
\item
  Definition: Intervals
  \(I_1 \supseteq I_2 \supseteq I_3 \supseteq \cdots\) are \emph{nested}
\item
  Theorem (Nested Interval Property): If
  \([a_1, b_1] \supseteq [a_2, b_2] \supseteq \cdots\) are nested closed
  intervals in a complete ordered field, then
  \(\bigcap_{n=1}^{\infty} [a_n, b_n] \neq \emptyset\)
\item
  Proof: Let \(L = \sup\{a_n\}\). Show \(L \leq b_n\) for all \(n\).
  Then \(L \in [a_n, b_n]\) for all \(n\).
\item
  Theorem: If also the lengths \(b_n - a_n\) become arbitrarily small,
  the intersection is exactly one point
\item
  Remark: This is what completeness guarantees --- the trapping method
  always works
\item
  Remark: Open intervals can have empty intersection even when nested
  (example: \((0, 1/n)\))
\item
  Remark: In \(\QQ\), nested closed intervals can have empty
  intersection --- completeness fails
\end{itemize}

\subsection{Looking Ahead}\label{looking-ahead}

\begin{itemize}
\tightlist
\item
  With completeness, we can finally make rigorous definitions
\item
  §1.4: √n exists for all n; decimal expansions make sense; ℝ is
  uncountable
\item
  §1.5: Lengths of curves, areas of regions, angles --- all defined via
  suprema, now guaranteed to exist
\end{itemize}

\subsection{Guided Exercise}\label{guided-exercise-1}

\subsection{The Babylonian Square
Root}\label{the-babylonian-square-root}

The Babylonians (\textasciitilde1800 BCE) discovered an algorithm for
approximating \(\sqrt{2}\). Starting with a \(2 \times 1\) rectangle
(area 2), repeatedly replace it with a more ``square-like'' rectangle of
the same area: the new width is the average of the old width and height.

Let \(w_1 = 2\), \(h_1 = 1\), and for \(n \geq 1\):
\[w_{n+1} = \frac{w_n + h_n}{2}, \qquad h_{n+1} = \frac{2}{w_{n+1}}\]

We prove that the intervals \([h_n, w_n]\) are nested and shrink to
\(\sqrt{2}\).

\emph{Part I: The bounds}

\begin{enumerate}
\def\labelenumi{(\alph{enumi})}
\item
  Verify that \(w_n h_n = 2\) for all \(n\) (the area is always 2).
\item
  Using AM-GM, show that
  \(w_{n+1} = \frac{w_n + 2/w_n}{2} \geq \sqrt{2}\) for any \(w_n > 0\).
\item
  Conclude that \(w_n \geq \sqrt{2}\) for all \(n \geq 2\), and
  therefore \(h_n = 2/w_n \leq \sqrt{2}\).
\item
  Explain why this means \(\sqrt{2} \in [h_n, w_n]\) for all
  \(n \geq 2\).
\end{enumerate}

\emph{Part II: The nesting}

\begin{enumerate}
\def\labelenumi{(\alph{enumi})}
\setcounter{enumi}{4}
\item
  Show that \((w_n)\) is decreasing for \(n \geq 2\): use
  \(w_{n+1} = \frac{w_n + h_n}{2}\) and \(h_n \leq w_n\).
\item
  Show that \((h_n)\) is increasing for \(n \geq 2\): use
  \(h_n = 2/w_n\) and \((w_n)\) decreasing.
\item
  Conclude that \([h_{n+1}, w_{n+1}] \subseteq [h_n, w_n]\) for all
  \(n \geq 2\).
\end{enumerate}

\emph{Part III: The shrinking gap}

Let \(e_n = w_n - h_n\) denote the gap.

\begin{enumerate}
\def\labelenumi{(\alph{enumi})}
\setcounter{enumi}{7}
\item
  Show that \(w_{n+1}^2 - 2 = \frac{(w_n - h_n)^2}{4}\). \emph{Hint:}
  Expand \(w_{n+1}^2 = \left(\frac{w_n + h_n}{2}\right)^2\) and use
  \(w_n h_n = 2\).
\item
  Using \(w_{n+1} - h_{n+1} = \frac{w_{n+1}^2 - 2}{w_{n+1}}\), show that
  \(e_{n+1} = \frac{e_n^2}{4w_{n+1}}\).
\item
  Since \(w_{n+1} \geq 1\), conclude that
  \(e_{n+1} \leq \frac{e_n^2}{4}\).
\item
  Verify that \(e_2 < 1\). Then show by induction that
  \(e_n \leq \frac{1}{4^{n-2}}\) for \(n \geq 2\).
\item
  Using the Archimedean property, explain why for any
  \(\varepsilon > 0\), there exists \(n\) with \(e_n < \varepsilon\).
\end{enumerate}

\emph{Part IV: Conclusion}

\begin{enumerate}
\def\labelenumi{(\alph{enumi})}
\setcounter{enumi}{12}
\item
  Apply the Nested Interval Property to conclude that
  \(\bigcap_{n=2}^{\infty} [h_n, w_n]\) contains at least one point
  \(c\).
\item
  Explain why the intersection contains \emph{exactly} one point.
\item
  Show that \(w_n^2 - 2 = w_n \cdot e_n\) and
  \(2 - h_n^2 = h_n \cdot e_n\).
\item
  Using the bounds \(w_n \leq w_2\) and \(h_n \leq w_2\), show that both
  \(w_n^2 - 2\) and \(2 - h_n^2\) become arbitrarily small.
\item
  Since \(h_n^2 \leq c^2 \leq w_n^2\) and both approach 2, conclude that
  \(c^2 = 2\). \emph{Hint:} If \(c^2 \neq 2\), then \(|c^2 - 2| > 0\).
  Derive a contradiction.
\item
  Conclude: \(c = \sqrt{2}\), and the Babylonian algorithm traps
  \(\sqrt{2}\) in nested intervals.
\end{enumerate}

\subsection{Exercises}\label{exercises-2}

\emph{Infimum} - Prove uniqueness of inf - State and prove the
ε-characterization of inf - Prove: in a complete ordered field, every
nonempty set bounded below has an infimum

\emph{Working with sup and inf} - Prove:
\(A \subseteq B \Rightarrow \sup A \leq \sup B\) - Prove:
\(\sup(cS) = c \cdot \inf S\) for \(c < 0\) - Prove:
\(\sup(A \cup B) = \max\{\sup A, \sup B\}\) - Prove: if
\(\sup A < \sup B\), then some \(b \in B\) is an upper bound for \(A\)

\emph{Finding sup and inf} - Find sup and inf: \([1,3]\), \([1,3)\),
\((1,3)\), \(\{1/n : n \in \NN\}\) - Find sup and inf:
\(\{(-1)^n / n : n \in \NN\}\) - Find sup and inf:
\(\{n/(n+1) : n \in \NN\}\) - Find sup and inf:
\(\{r \in \QQ : r^2 < 3\}\) (in \(\RR\))

\emph{Nested intervals} - Give an example of nested open intervals with
empty intersection - Prove: if \(\sup\{a_n\} = \inf\{b_n\}\), then
\(\bigcap [a_n, b_n]\) is exactly one point - Construct nested closed
intervals in \(\QQ\) with empty intersection

\emph{Equivalences} - ★ Prove: sup-completeness is equivalent to
inf-completeness

\subsection{Dependencies}\label{dependencies-2}

\textbf{Requires}: §1.1 (ordered fields), §1.2 (motivation from
\(\QQ\)'s incompleteness)

\textbf{Used in}: §1.4 (roots, decimals), §1.5 (lengths, areas, angles
defined via sup)

\section{The Real Numbers}\label{the-real-numbers}

\subsection{Narrative}\label{narrative-3}

\begin{itemize}
\tightlist
\item
  \textbf{The bold move}: define \(\RR\) as a complete ordered field
\item
  Axiomatic method: say what reals \emph{do}, not what they \emph{are}
\item
  Existence and uniqueness theorem (stated)
\item
  \textbf{No monsters}: \(\NN\) unbounded, no infinitesimals
\item
  \textbf{Archimedean property}: \(na > b\) for some \(n\) (completeness
  implies this)
\item
  \textbf{Density of \(\QQ\)}: between any two reals lies a rational
\item
  \textbf{Density of irrationals}: between any two reals lies an
  irrational
\item
  \textbf{The payoff}: \(\sqrt{2}\) exists! \(\sup\{x : x^2 < 2\}\)
  works
\item
  \(n\)th roots exist; rational exponents; real exponents via sup
\item
  Logarithms via sup
\item
  \textbf{Surprise}: \(\RR\) is uncountable (nested intervals proof)
\item
  \(\QQ\) countable, \(\RR\) uncountable \(\Rightarrow\) ``most'' reals
  irrational
\item
  \(\AA\) countable \(\Rightarrow\) transcendentals exist, ``most''
  reals transcendental
\item
  Computables countable \(\Rightarrow\) ``most'' reals indescribable
\item
  \textbf{Closing}: We built \(\RR\) to fill holes in \(\QQ\); it's an
  uncountable ocean where rationals, algebraics, computables are
  countable dust
\end{itemize}

\subsection{Content}\label{content-3}

\subsection{\texorpdfstring{The Definition of
\(\RR\)}{The Definition of \textbackslash RR}}\label{the-definition-of-rr}

\begin{itemize}
\tightlist
\item
  Definition: \(\RR\) is a complete ordered field
\item
  Axiomatic method: we say what reals \emph{do}, not what they
  \emph{are}
\item
  Theorem (stated): A complete ordered field exists and is unique up to
  isomorphism
\item
  Remark: constructions (Dedekind cuts, Cauchy sequences) prove
  existence; we focus on using completeness
\end{itemize}

\subsection{No Monsters (Archimedean
Property)}\label{no-monsters-archimedean-property}

\begin{itemize}
\tightlist
\item
  Theorem: For any \(x \in \RR\), there exists \(n \in \NN\) with
  \(n > x\) (prove in text)
\item
  Proof idea: if \(\NN\) were bounded above, \(\sup \NN\) would exist;
  derive contradiction
\item
  Corollary: For any \(\varepsilon > 0\), there exists \(n \in \NN\)
  with \(1/n < \varepsilon\)
\item
  Interpretation: no infinitely large or infinitely small reals; every
  real lives between two integers
\end{itemize}

\subsection{Density}\label{density}

\begin{itemize}
\tightlist
\item
  Theorem: Between any two reals lies a rational (prove in text)
\item
  Proof uses Archimedean property to find suitable denominator
\item
  Theorem: Between any two reals lies an irrational (quick corollary
  using \(\sqrt{2}\))
\item
  Interpretation: rationals are everywhere dense --- we've filled gaps,
  not added alien regions
\end{itemize}

\subsection{Roots Exist}\label{roots-exist}

\begin{itemize}
\tightlist
\item
  Theorem: \(\sqrt{2}\) exists in \(\RR\) (prove in text --- THIS IS THE
  PAYOFF)
\item
  Let \(S = \{x \in \RR : x > 0, x^2 < 2\}\)
\item
  ✎ Inline: Verify \(1 \in S\) and \(2\) is an upper bound for \(S\)
\item
  By completeness, \(s = \sup S\) exists; prove \(s^2 = 2\) by ruling
  out \(s^2 < 2\) and \(s^2 > 2\)
\item
  Corollary: \(\sqrt{a}\) exists for any \(a \geq 0\) (same proof)
\item
  Definition: \(\sqrt[n]{a} = \sup\{x \geq 0 : x^n < a\}\) for \(a > 0\)
\item
  Theorem: \((\sqrt[n]{a})^n = a\) (exercise)
\item
  Remark: AM-GM (§1.2 exercise) now holds unconditionally in \(\RR\):
  for \(a, b \geq 0\), \(\frac{a+b}{2} \geq \sqrt{ab}\)
\end{itemize}

\subsection{Infinite Riches
(Uncountability)}\label{infinite-riches-uncountability}

\begin{itemize}
\tightlist
\item
  Theorem: \(\RR\) is uncountable (prove in text --- nested intervals
  argument)
\item
  Proof: given any list \(x_1, x_2, \ldots\), construct nested intervals
  avoiding each \(x_n\); intersection contains a real not on the list
\item
  Corollary: The irrationals are uncountable (if irrationals were
  countable, \(\RR = \QQ \cup (\RR \setminus \QQ)\) would be countable)
\item
  ``Most'' reals are irrational (\(\QQ\) countable, \(\RR\) uncountable)
\item
  \(\AA\) countable \(\Rightarrow\) transcendentals exist, ``most''
  reals transcendental (exercise)
\item
  Computables countable \(\Rightarrow\) ``most'' reals not computable
\item
  Definables countable \(\Rightarrow\) ``most'' reals not describable
\item
  Interpretation: we built \(\RR\) to fill holes in \(\QQ\); we got an
  uncountable ocean where rationals, algebraics, computables are
  countable dust
\end{itemize}

\subsection{Guided Exercises}\label{guided-exercises-1}

\subsection{Exponents and Logarithms}\label{exponents-and-logarithms}

We extend exponentiation from integer exponents to rational exponents to
real exponents, using completeness.

\begin{enumerate}
\def\labelenumi{(\alph{enumi})}
\item
  For \(a > 0\) and \(p/q \in \QQ\) with \(q > 0\), define
  \(a^{p/q} = (\sqrt[q]{a})^p\). Prove this is well-defined: if
  \(p/q = r/s\), then \((\sqrt[q]{a})^p = (\sqrt[s]{a})^r\).
\item
  Prove \(a^{r+s} = a^r \cdot a^s\) for \(a > 0\) and \(r, s \in \QQ\).
\item
  Prove \(a^{rs} = (a^r)^s\) for \(a > 0\) and \(r, s \in \QQ\).
\item
  Prove: if \(a > 1\) and \(r < s\) are rational, then \(a^r < a^s\).
\item
  For \(a > 1\) and \(x \in \RR\), define
  \(a^x = \sup\{a^r : r \in \QQ, r < x\}\). Prove this agrees with the
  previous definition when \(x \in \QQ\).
\item
  Prove \(a^{x+y} = a^x \cdot a^y\) for \(a > 1\) and \(x, y \in \RR\).
\item
  For \(a > 1\) and \(y > 0\), define
  \(\log_a(y) = \sup\{r \in \QQ : a^r < y\}\). Prove
  \(a^{\log_a(y)} = y\).
\item
  Prove \(\log_a(xy) = \log_a(x) + \log_a(y)\) for \(x, y > 0\).
\end{enumerate}

\subsection{Exercises}\label{exercises-3}

\emph{Archimedean Property} - Prove: for any \(x, y > 0\), there exists
\(n \in \NN\) with \(nx > y\) - Prove: for any \(x \in \RR\), there
exist integers \(m, n\) with \(m < x < n\)

\emph{Roots} - Prove \((\sqrt[n]{a})^n = a\) for \(a > 0\) and
\(n \in \NN\) - Prove \(\sqrt[n]{ab} = \sqrt[n]{a} \cdot \sqrt[n]{b}\) -
Prove \(\sqrt[m]{\sqrt[n]{a}} = \sqrt[mn]{a}\)

\emph{Exponents} - Prove rational exponents are well-defined:
\(a^{p/q}\) doesn't depend on the representation of \(p/q\) - Prove
\(a^{r+s} = a^r \cdot a^s\) for rational \(r, s\) - Prove
\((a^r)^s = a^{rs}\) for rational \(r, s\) - Prove: \(a^x < a^y\) for
\(a > 1\) and \(x < y\) (real exponents) - ★ Extend the definition of
\(a^x\) to \(0 < a < 1\) and verify the expected properties

\emph{Countability} - Prove \(\AA\) is countable (hint: countably many
polynomials with integer coefficients, each has finitely many roots) -
Prove the set of computable reals is countable (hint: countably many
algorithms) - Prove: if \(A\) is countable and \(B\) is uncountable,
then \(B \setminus A\) is uncountable

\subsection{Dependencies}\label{dependencies-3}

\textbf{Requires}: §1.1 (countability, Cantor), §1.2 (ordered field),
§1.3 (motivation), §1.4 (completeness, nested intervals)

\textbf{Used in}: §2.2 (Archimedean in \(\varepsilon\)-\(N\) proofs),
§2.3 (density of \(\QQ\)), §2.4 (roots in Babylonian, AM-GM)

\section{Geometry Reborn}\label{geometry-reborn}

\subsection{Narrative}\label{narrative-4}

\begin{itemize}
\tightlist
\item
  We've built \(\mathbb{R}\) with completeness --- now what can we do
  with it?
\item
  The Greeks measured straight lines easily, but curves and regions were
  hard
\item
  With sup/inf we can rigorously define: lengths of curves, areas of
  regions, even angle
\item
  This chapter's payoff: making curved geometry rigorous
\end{itemize}

\subsection{Content}\label{content-4}

\emph{(Opening --- unnumbered)} - Completeness lets us define
measurements the Greeks struggled with - Distance, length, area --- all
via sup/inf - We define trigonometry as a consequence of arc length

\subsection{The Plane}\label{the-plane}

\begin{itemize}
\tightlist
\item
  Definition: \(\mathbb{R}^2 = \{(x, y) : x, y \in \mathbb{R}\}\)
  (ordered pairs)
\item
  Definition:
  \(\mathbb{R}^n = \{(x_1, \ldots, x_n) : x_i \in \mathbb{R}\}\)
  (n-tuples)
\item
  The distance formula:
  \(d(p, q) = \sqrt{(x_2 - x_1)^2 + (y_2 - y_1)^2}\) for
  \(p = (x_1, y_1)\), \(q = (x_2, y_2)\)
\item
  Remark: The Greeks proved this as a theorem (Pythagoras). We take it
  as \emph{definition} of distance in \(\mathbb{R}^2\)
\item
  In \(\mathbb{R}^n\): \(d(p, q) = \sqrt{\sum_{i=1}^{n}(q_i - p_i)^2}\)
\item
  ✎ Inline: Verify \(d(p, q) = d(q, p)\) and \(d(p, p) = 0\)
\end{itemize}

\subsection{Measuring Sets}\label{measuring-sets}

\begin{itemize}
\tightlist
\item
  Definition: The \emph{diameter} of a set \(S \subseteq \mathbb{R}^2\)
  is \(\text{diam}(S) = \sup\{d(p, q) : p, q \in S\}\)
\item
  Exists (possibly \(= +\infty\)) by completeness
\item
  Definition: A set is \emph{bounded} iff \(\text{diam}(S) < \infty\)
\item
  Definition: \emph{Distance from a point to a set}:
  \(d(p, S) = \inf\{d(p, q) : q \in S\}\)
\item
  Exists by completeness (infimum of nonempty set bounded below by 0)
\item
  Definition: \emph{Distance between sets}:
  \(d(S, T) = \inf\{d(p, q) : p \in S, q \in T\}\)
\item
  Example: \(d(p, S) = 0\) does not imply \(p \in S\) (consider
  \(S = (0, 1)\), \(p = 0\))
\item
  ✎ Inline: Find the diameter of the unit disk
  \(\{(x,y) : x^2 + y^2 \leq 1\}\)
\end{itemize}

\subsection{Measuring Curves}\label{measuring-curves}

\begin{itemize}
\tightlist
\item
  Informally: a ``curve'' is a path in the plane
\item
  We don't yet have tools to make ``curve'' precise (need continuity,
  Chapter 5)
\item
  But we can define \emph{length} for well-behaved curves
\item
  Setup: Given points \(P_0, P_1, \ldots, P_n\) along a curve in order,
  the \emph{inscribed polygonal length} is
  \(\sum_{i=1}^{n} d(P_{i-1}, P_i)\)
\item
  Definition: A curve is \emph{rectifiable} if the supremum of inscribed
  polygonal lengths (over all finite inscribed polygons) is finite
\item
  Definition: The \emph{length} of a rectifiable curve is this supremum
\item
  Exists by completeness
\item
  Remark: Not all curves are rectifiable --- we'll see examples in
  Chapter 10
\item
  Example: A line segment from \(p\) to \(q\) has length \(d(p, q)\)
\item
  ✎ Inline: Why is the length of a segment equal to \(d(p, q)\)? (The
  segment itself is an inscribed polygon!)
\end{itemize}

\subsection{Measuring Regions}\label{measuring-regions}

\begin{itemize}
\tightlist
\item
  How to define area of a region \(R \subseteq \mathbb{R}^2\)?
\item
  \emph{Inner approximation}: inscribed polygons whose area we can
  compute
\item
  \emph{Outer approximation}: circumscribed polygons
\item
  Definition:
  \(A_{\text{inner}}(R) = \sup\{\text{area}(P) : P \subseteq R, P \text{ polygon}\}\)
\item
  Definition:
  \(A_{\text{outer}}(R) = \inf\{\text{area}(P) : R \subseteq P, P \text{ polygon}\}\)
\item
  Always \(A_{\text{inner}}(R) \leq A_{\text{outer}}(R)\)
\item
  Definition: \(R\) has \emph{area} if
  \(A_{\text{inner}}(R) = A_{\text{outer}}(R)\); the common value is
  \(\text{area}(R)\)
\item
  Remark: Not every region has area --- this leads to measure theory
  (much later)
\item
  Remark: The Koch snowflake has finite area but infinite perimeter
  (Chapter 10)
\item
  ✎ Inline: Does a line segment have area?
\end{itemize}

\subsection{Trigonometry}\label{trigonometry}

\begin{itemize}
\tightlist
\item
  The unit circle: \(\{(x, y) : x^2 + y^2 = 1\}\)
\item
  For \(\theta \geq 0\): there exists a unique point \(P(\theta)\) on
  the unit circle such that the arc length from \((1, 0)\) to
  \(P(\theta)\) (counterclockwise) equals \(\theta\)
\item
  Existence: completeness! Arc length is sup of inscribed polygons; for
  any \(\theta\) we can find the point where this sup equals \(\theta\)
\item
  Definition: \(\cos\theta\) = x-coordinate of \(P(\theta)\);
  \(\sin\theta\) = y-coordinate of \(P(\theta)\)
\item
  For \(\theta < 0\): measure clockwise; equivalently \(P(-\theta)\) is
  reflection of \(P(\theta)\) across x-axis
\item
  Immediate consequences:

  \begin{itemize}
  \tightlist
  \item
    \(\cos^2\theta + \sin^2\theta = 1\) (point is on unit circle)
  \item
    \(-1 \leq \cos\theta \leq 1\), \(-1 \leq \sin\theta \leq 1\)
  \item
    \(\cos 0 = 1\), \(\sin 0 = 0\)
  \item
    \(\cos(-\theta) = \cos\theta\), \(\sin(-\theta) = -\sin\theta\)
    (reflection symmetry)
  \end{itemize}
\item
  Key principle: \emph{equal arcs have equal chords} --- arcs of the
  same length subtend chords of the same length
\item
  ⚠️ FLAG: This principle needs justification. For now, state it as
  evident from the symmetry of the circle.
\item
  From this principle: the angle addition formulas (see guided exercise)
\item
  Remark: We have the algebraic laws of trigonometry, but we cannot yet
  compute \(\cos(1)\) or \(\sin(2)\)
\item
  Remark: To find specific values, we need to approximate --- this leads
  to Chapter 4 (Archimedes' method)
\end{itemize}

\subsection{Guided Exercise}\label{guided-exercise-2}

\subsection{The Metric Origins of
Trigonometry}\label{the-metric-origins-of-trigonometry}

We derive the angle addition formulas from the distance formula and the
principle that equal arcs have equal chords.

\emph{Part I: Symmetry}

\begin{enumerate}
\def\labelenumi{(\alph{enumi})}
\item
  Explain why \(P(-\theta)\) is the reflection of \(P(\theta)\) across
  the x-axis.
\item
  Conclude: \(\cos(-\theta) = \cos\theta\) and
  \(\sin(-\theta) = -\sin\theta\).
\item
  Verify \(\cos 0 = 1\), \(\sin 0 = 0\) directly from the definition.
\end{enumerate}

\emph{Part II: The Chord-Arc Principle}

Consider four points on the unit circle defined by arc length from
\((1, 0)\): - \(A = P(\alpha + \beta)\) - \(B = P(0) = (1, 0)\) -
\(C = P(\alpha)\) - \(D = P(-\beta)\)

\begin{enumerate}
\def\labelenumi{(\alph{enumi})}
\setcounter{enumi}{3}
\item
  What is the arc length from \(B\) to \(A\)? From \(D\) to \(C\)?
\item
  Using the principle ``equal arcs have equal chords,'' explain why
  \(d(A, B) = d(C, D)\).
\end{enumerate}

\emph{Part III: The Cosine Addition Formula}

\begin{enumerate}
\def\labelenumi{(\alph{enumi})}
\setcounter{enumi}{5}
\item
  Write \(A\), \(B\), \(C\), \(D\) in coordinates using sin and cos.
\item
  Compute \(d(A, B)^2\) using the distance formula.
\item
  Compute \(d(C, D)^2\) using the distance formula.
\item
  Set \(d(A, B)^2 = d(C, D)^2\), expand, and use \(\sin^2 + \cos^2 = 1\)
  to simplify.
\item
  Derive:
  \(\cos(\alpha + \beta) = \cos\alpha\cos\beta - \sin\alpha\sin\beta\).
\end{enumerate}

\emph{Part IV: The Sine Addition Formula}

\begin{enumerate}
\def\labelenumi{(\alph{enumi})}
\setcounter{enumi}{10}
\tightlist
\item
  Using a similar chord argument or the classic ``stacked triangles''
  construction, derive:
  \(\sin(\alpha + \beta) = \sin\alpha\cos\beta + \cos\alpha\sin\beta\).
\end{enumerate}

\emph{Part V: Double and Half Angle}

\begin{enumerate}
\def\labelenumi{(\alph{enumi})}
\setcounter{enumi}{11}
\item
  Set \(\beta = \alpha\) in the cosine formula to derive
  \(\cos(2\theta) = 2\cos^2\theta - 1\).
\item
  Derive \(\sin(2\theta) = 2\sin\theta\cos\theta\).
\item
  Solve the double angle formula for \(\cos(\theta/2)\) in terms of
  \(\cos\theta\).
\item
  Explain how the half-angle formula will connect to Archimedes' polygon
  doubling method (§4.5).
\end{enumerate}

\subsection{Exercises}\label{exercises-4}

\emph{The plane} - Prove the triangle inequality:
\(d(p, r) \leq d(p, q) + d(q, r)\) - Prove \(d(p, q) \geq 0\) with
equality iff \(p = q\)

\emph{Measuring sets} - Find the diameter of the square
\([0, 1] \times [0, 1]\) - Find \(d(p, S)\) where \(p = (0, 0)\) and
\(S = \{(x, y) : x^2 + y^2 = 1\}\) (unit circle) - Give an example where
\(d(S, T) = 0\) but \(S \cap T = \emptyset\)

\emph{Measuring curves} - Prove: the perimeter of a rectangle with sides
\(a\) and \(b\) is \(2a + 2b\) - Show that a circle is rectifiable
(hint: inscribed n-gons have perimeters bounded above)

\emph{Measuring regions} - Prove: a rectangle has area equal to base
times height - Show \(A_{\text{inner}} \leq A_{\text{outer}}\) for any
region

\emph{Trigonometry} - Verify:
\(\cos(\alpha - \beta) = \cos\alpha\cos\beta + \sin\alpha\sin\beta\) -
Derive \(\sin(2\theta) = 2\sin\theta\cos\theta\) from the sine addition
formula - Derive \(\cos^2\theta = (1 + \cos(2\theta))/2\)

\subsection{Dependencies}\label{dependencies-4}

\textbf{Requires}: §1.3 (completeness, sup/inf)

\textbf{Used in}: §4.5 (Archimedes' \(\pi\), polygon doubling), Chapter
10 (computing trig values)

\chapter{Limits}\label{limits}

\textbf{Limits before the definition} - Mathematicians across cultures
computed limits long before anyone defined them - Babylonians
(\textasciitilde1800 BCE): iterative algorithm for \(\sqrt{2}\),
accurate to six decimals - Greeks: Eudoxus' method of exhaustion
(\textasciitilde370 BCE); Archimedes' parabola area, π approximations
via 96-gons (\textasciitilde250 BCE) - China: Liu Hui (263 CE) and Zu
Chongzhi (\textasciitilde480 CE) pushed π to seven decimal places -
Islamic world: Al-Kashi (\textasciitilde1420) reached sixteen decimal
places for π - India: Madhava and the Kerala school (\textasciitilde1400
CE) discovered infinite series for π, arctan, sine, cosine --- Taylor
series 250 years before Taylor

\textbf{What they had in common} - All understood: infinite processes
can pin down finite answers - All could make errors arbitrarily small -
All knew their methods worked

\textbf{What they lacked} - No general \emph{definition} of limit - Each
convergence proof was bespoke, tailored to specific problem -
Archimedes' parabola argument doesn't help with the circle - Babylonian
√2 algorithm requires its own verification - Every new limit demanded
fresh ingenuity

\textbf{The resolution} - Cauchy and Weierstrass (19th century): replace
``approaching'' with precise quantification - The ε-N definition:
sequence converges to \(L\) if, for any tolerance, terms eventually stay
within that tolerance

\textbf{The payoff: general theorems} - Rigorous definition enables
\emph{general theorems}: limit laws, MCT, squeeze - Theorems proved
once, applied everywhere - No new argument for each limit --- combine
general tools - Computational power, not just philosophical clarity

\textbf{What we'll do} - Limit laws resolve Archimedes' parabola in one
line - Prove Babylonian algorithm converges to \(\sqrt{2}\) --- two
different methods - Show every real is limit of rationals --- bridge
between countable and uncountable - The ancients were right; now we can
prove it

\section{Sequences and Series}\label{sequences-and-series}

\subsection{Narrative}\label{narrative-5}

\begin{itemize}
\tightlist
\item
  Chapter 1: \(\RR\) is uncountable, finite descriptions are countable
\item
  \textbf{Key insight}: Infinite processes bridge the gap
\item
  Historical examples show sequences as tools in action --- before we
  even define them
\item
  Formal definition: what these examples have in common
\item
  Series and products are secretly sequences --- find closed form for
  \(S_n\) or \(P_n\)
\item
  Geometric and telescoping examples show how to collapse infinite
  sums/products into sequences
\item
  The question: we have \(S_n\) as a sequence, but what happens as
  \(n \to \infty\)?
\end{itemize}

\subsection{Content}\label{content-5}

\begin{itemize}
\tightlist
\item
  Chapter 1 showed \(\RR\) uncountable, finite descriptions countable
\item
  Infinite processes bridge the gap
\item
  Sequences let us \emph{name} numbers we can't write down finitely
\item
  Preview: every real is limit of a sequence of rationals (proved in
  §2.3)
\end{itemize}

\subsection{Historical Examples: Sequences in
Action}\label{historical-examples-sequences-in-action}

\begin{itemize}
\tightlist
\item
  \emph{Archimedes and the parabola}: area under parabola via inscribed
  triangles

  \begin{itemize}
  \tightlist
  \item
    Total area: \(T + T/4 + T/16 + \cdots\)
  \item
    Partial sums \(S_n = T(1 + 1/4 + \cdots + 1/4^n)\) --- Archimedes
    claimed \(S = 4T/3\)
  \end{itemize}
\item
  \emph{Babylonians and \(\sqrt{2}\)} (\textasciitilde1800 BCE):

  \begin{itemize}
  \tightlist
  \item
    Clay tablet shows approximation \(1;24,51,10\) in base 60
  \item
    Algorithm: \(x_1 = 2\), \(x_{n+1} = (x_n + 2/x_n)/2\)
  \item
    Compute: \(x_2 = 1.5\), \(x_3 = 1.4167\ldots\),
    \(x_4 = 1.41422\ldots\)
  \item
    Clearly ``approaches'' \(\sqrt{2}\) --- but what does that mean?
  \end{itemize}
\item
  Both produce lists of numbers that seem to approach a value
\item
  What do these examples have in common?
\end{itemize}

\subsection{What Is a Sequence?}\label{what-is-a-sequence}

\begin{itemize}
\tightlist
\item
  Definition: a sequence is a function \(a: \NN \to \RR\); write \(a_n\)
  for \(a(n)\)
\item
  Notation: \((a_n)\) or \((a_1, a_2, a_3, \ldots)\)
\item
  Two ways to specify: closed form (e.g., \(a_n = 1/n\)) vs recursive
  (e.g., \(a_{n+1} = f(a_n)\))
\item
  The Babylonian example is recursive; Archimedes' partial sums have a
  closed form (we'll find it)
\end{itemize}

\subsection{Series and Products as
Sequences}\label{series-and-products-as-sequences}

\begin{itemize}
\tightlist
\item
  A series \(\sum a_n\) is really the sequence of partial sums
  \(S_n = a_1 + \cdots + a_n\)
\item
  A product \(\prod a_n\) is really the sequence of partial products
  \(P_n = a_1 \cdots a_n\)
\item
  To understand a series or product, find a closed form for \(S_n\) or
  \(P_n\)
\item
  Then ask: what happens to this sequence as \(n \to \infty\)?
\end{itemize}

\emph{Geometric partial sums} - Goal: closed form for
\(S_n = 1 + r + r^2 + \cdots + r^n\) - Trick:
\(S_n - rS_n = 1 - r^{n+1}\) - Theorem: For \(r \neq 1\):
\(S_n = \frac{1 - r^{n+1}}{1 - r}\) (prove in text) - This resolves
Archimedes: \(S_n = T \cdot \frac{1 - (1/4)^{n+1}}{1 - 1/4}\)

\emph{Telescoping sums and products} - Idea: consecutive terms cancel,
leaving only endpoints - Sums: if \(a_k = t_k - t_{k-1}\), then
\(S_n = t_n - t_0\) - Products: if \(a_k = t_k / t_{k-1}\), then
\(P_n = t_n / t_0\) - Example:
\(\frac{1}{k(k+1)} = \frac{1}{k} - \frac{1}{k+1}\), so
\(\sum_{k=1}^{n} \frac{1}{k(k+1)} = 1 - \frac{1}{n+1}\) - Example:
\(1 - \frac{1}{k^2} = \frac{(k-1)(k+1)}{k^2}\) telescopes to give
\(\prod_{k=2}^{n}(1 - 1/k^2) = \frac{n+1}{2n}\) - ✎ Inline: Verify the
partial fractions \(\frac{1}{k(k+1)} = \frac{1}{k} - \frac{1}{k+1}\) - ✎
Inline: Verify that \(\prod_{k=1}^{n}(1 + 1/k) = n + 1\)

\subsection{The Question}\label{the-question}

\begin{itemize}
\tightlist
\item
  We have closed forms: \(S_n = \frac{1 - r^{n+1}}{1-r}\),
  \(S_n = 1 - \frac{1}{n+1}\), \(P_n = \frac{n+1}{2n}\)
\item
  We have the Babylonian sequence: \(1.5, 1.4167, 1.41422, \ldots\)
\item
  What happens as \(n \to \infty\)? What does ``\(\to\)'' mean?
\item
  This is \textbf{convergence} --- next section
\end{itemize}

\subsection{Guided Exercises}\label{guided-exercises-2}

None for this section.

\subsection{Exercises}\label{exercises-5}

\begin{itemize}
\tightlist
\item
  Prove the geometric partial sum formula by induction
\item
  Compute the first 5 terms of the Babylonian sequence for \(\sqrt{3}\)
\item
  (Pell structure) Write \(x_n = p_n/q_n\) for the Babylonian sequence
  with \(x_1 = 2/1\). Show that \(p_{n+1} = p_n + 2q_n\) and
  \(q_{n+1} = p_n + q_n\). Compute \((p_n, q_n)\) for
  \(n = 1, 2, 3, 4, 5\). Show that \(p_n^2 - 2q_n^2 = (-1)^n \cdot 2\)
  for all \(n \geq 1\).
\end{itemize}

\subsection{Dependencies}\label{dependencies-5}

\textbf{Requires}: §1.5 (uncountability of \(\RR\), existence of
\(\sqrt{2}\))

\textbf{Used in}: §2.2 (convergence definition), §2.3 (geometric series
sum), §2.4 (Babylonian convergence)

\section{Convergence}\label{convergence}

\subsection{Narrative}\label{narrative-6}

\begin{itemize}
\tightlist
\item
  What should ``approaches'' mean? Eventually stays arbitrarily close.
\item
  The ε-N game: adversary picks tolerance, you find threshold
\item
  Full example showing estimation technique --- this is what ε-N proofs
  look like
\item
  ``That was tedious. Let's identify fundamental sequences we can build
  from.''
\item
  Fundamental examples: \(1/n \to 0\), \(r^n \to 0\), \((-1)^n\)
  diverges
\item
  Basic properties: uniqueness, convergent ⟹ bounded, tails
\item
  Divergence to infinity: a different kind of behavior
\item
  Preview: limit laws (§2.3) will let us avoid estimation work
\end{itemize}

\subsection{Content}\label{content-6}

\subsection{The Definition}\label{the-definition}

\begin{itemize}
\tightlist
\item
  Motivation: ``approaches'' means ``eventually stays arbitrarily
  close''
\item
  The ε-N game: adversary picks any tolerance \(\varepsilon > 0\), you
  respond with threshold \(N\); you win if all terms past \(N\) are
  within \(\varepsilon\) of \(L\); converges iff you have a winning
  strategy for every \(\varepsilon\)
\item
  Definition: \((a_n)\) converges to \(L\), written \(a_n \to L\) or
  \(\lim a_n = L\), if
  \[\forall \varepsilon > 0 \; \exists N \in \NN \; \forall n > N: |a_n - L| < \varepsilon\]
\item
  Definition: a sequence converges if it converges to some \(L\);
  otherwise it diverges
\end{itemize}

\emph{A first example: the estimation technique} - Example:
\(\frac{n^2 + n}{2n^2 + 1} \to \frac{1}{2}\) - Compute:
\(\left| \frac{n^2 + n}{2n^2 + 1} - \frac{1}{2} \right| = \frac{2n - 1}{2(2n^2 + 1)}\)
- Estimate:
\(\frac{2n - 1}{2(2n^2 + 1)} < \frac{2n}{4n^2} = \frac{1}{2n}\) - Choose
\(N > \frac{1}{2\varepsilon}\); for \(n > N\): expression
\(< \frac{1}{2n} < \varepsilon\) - Remark: this works, but it's tedious
--- let's identify fundamental sequences we can combine

\subsection{Fundamental Examples}\label{fundamental-examples}

\begin{itemize}
\tightlist
\item
  \(1/n \to 0\): the building block (prove in text, uses Archimedean
  property)
\item
  Constant sequence \(c \to c\): trivial but worth stating
\item
  \(r^n \to 0\) for \(|r| < 1\): important result (prove in text using
  Bernoulli from §1.2)

  \begin{itemize}
  \tightlist
  \item
    Write \(|r| = 1/(1+h)\) with \(h > 0\); by Bernoulli
    \((1+h)^n \geq 1 + nh\)
  \item
    So \(|r|^n \leq 1/(1+nh) < 1/(nh)\); choose
    \(N > 1/(\varepsilon h)\)
  \end{itemize}
\item
  \(r^n\) diverges for \(|r| > 1\); diverges to \(+\infty\) for
  \(r > 1\)
\item
  Remark: \(r = 1\) is constant; \(r = -1\) is \((-1)^n\)
\item
  \((-1)^n\) diverges: prove directly --- consecutive terms can't both
  be within \(1/2\) of any \(L\)
\item
  \(r^{1/n} \to 1\) for \(r > 0\): (prove in text)

  \begin{itemize}
  \tightlist
  \item
    Case \(r \geq 1\): write \(r^{1/n} = 1 + h_n\); then
    \(r = (1+h_n)^n \geq 1 + nh_n\) (Bernoulli), so
    \(h_n \leq (r-1)/n \to 0\)
  \item
    Case \(0 < r < 1\): need \(1 - r^{1/n} < \varepsilon\), i.e.,
    \(r > (1-\varepsilon)^n\); since \((1-\varepsilon)^n \to 0\), this
    holds for large \(n\)
  \end{itemize}
\end{itemize}

\subsection{Basic Properties}\label{basic-properties}

\begin{itemize}
\tightlist
\item
  Theorem (Uniqueness): if \(a_n \to L\) and \(a_n \to M\), then
  \(L = M\) (proof as exercise)
\item
  Definition: \((a_n)\) is bounded if \(\exists M > 0\) with
  \(|a_n| \leq M\) for all \(n\)
\item
  Theorem: convergent sequences are bounded (prove in text --- bound the
  tail, bound the beginning)
\item
  Remark: converse is false --- \((-1)^n\) is bounded but divergent
\item
  Theorem (Tails): changing finitely many terms doesn't affect
  convergence or the limit
\end{itemize}

\subsection{Divergence to Infinity}\label{divergence-to-infinity}

\begin{itemize}
\tightlist
\item
  Definition: \(a_n \to +\infty\) if
  \(\forall M > 0 \; \exists N: n > N \Rightarrow a_n > M\)
\item
  Definition: \(a_n \to -\infty\) similarly
\item
  Example: \(n^2 \to +\infty\) (prove in text)
\item
  Remark: divergence to infinity is a specific behavior, not just
  ``doesn't converge''
\end{itemize}

\subsection{Guided Exercises}\label{guided-exercises-3}

None for this section.

\subsection{Exercises}\label{exercises-6}

\begin{itemize}
\tightlist
\item
  Prove uniqueness of limits
\item
  Prove \(n/(n+1) \to 1\) from the definition
\item
  Prove \((2n+1)/(3n+5) \to 2/3\) from the definition
\item
  Prove: if \(a_n \to L\), then \(|a_n| \to |L|\)
\item
  Prove: tails determine convergence (if \((a_n)\) and \((b_n)\) agree
  for \(n > N_0\), they either both converge to the same limit or both
  diverge)
\item
  Show \(n \to +\infty\)
\item
  Show \((-1)^n n\) diverges (but not to \(\pm\infty\))
\item
  ★ Find a sequence that is unbounded but does not diverge to
  \(+\infty\) or \(-\infty\)
\end{itemize}

\subsection{Dependencies}\label{dependencies-6}

\textbf{Requires}: §1.2 (Bernoulli's inequality), §1.5 (Archimedean
property)

\textbf{Used in}: §2.3 (limit laws use these fundamental limits), §2.4
(MCT)

\section{§2.3 Limit Theorems}\label{limit-theorems}

\subsection{Narrative}\label{narrative-7}

\begin{itemize}
\tightlist
\item
  Limit laws let us compute limits from known limits
\item
  Two types: limits respect order (inequalities), limits respect algebra
  (arithmetic)
\item
  These follow from ordered field axioms --- no completeness needed yet
\item
  Applications: finish the computations from §2.1 (geometric,
  telescoping, Archimedes)
\item
  Density payoff: every real is a limit of rationals --- sequences reach
  all of \(\RR\)
\item
  But limit laws have limits: they tell us \emph{what} the limit is, not
  \emph{that} it exists
\item
  Babylonian example shows we need a new tool\ldots{}
\end{itemize}

\subsection{Content}\label{content-7}

\subsection{Limits Respect Order}\label{limits-respect-order}

\begin{itemize}
\tightlist
\item
  Theorem: if \(a_n \to L\) and \(a_n \geq 0\) for all \(n\), then
  \(L \geq 0\) (prove in text)
\item
  Corollary: if \(a_n \to L\), \(b_n \to M\), and \(a_n \leq b_n\) for
  all \(n\), then \(L \leq M\)
\item
  Warning: strict inequality is NOT preserved --- \(1/n > 0\) for all
  \(n\), but \(\lim 1/n = 0\)
\item
  Theorem (Squeeze): if \(a_n \leq b_n \leq c_n\) and \(a_n \to L\) and
  \(c_n \to L\), then \(b_n \to L\) (prove in text)
\end{itemize}

\subsection{Limits Respect Algebra}\label{limits-respect-algebra}

\begin{itemize}
\tightlist
\item
  ✎ Inline: Prove directly that \(ca_n \to cL\) when \(a_n \to L\)
\item
  Remark: this also follows from the product law below
\item
  Theorem (Sum): if \(a_n \to L\) and \(b_n \to M\), then
  \(a_n + b_n \to L + M\) (prove in text --- triangle inequality,
  \(\varepsilon/2\))
\item
  ✎ Inline: Prove the difference law from the sum law
\item
  Theorem (Product): if \(a_n \to L\) and \(b_n \to M\), then
  \(a_n b_n \to LM\) (prove in text --- uses boundedness from §2.2)
\item
  Theorem (Quotient): if \(a_n \to L\) and \(b_n \to M\) with
  \(M \neq 0\) and \(b_n \neq 0\), then \(a_n/b_n \to L/M\) (exercise
  --- hint: first prove \(1/b_n \to 1/M\))
\item
  Theorem (Square Root): if \(a_n \to L\) with \(a_n \geq 0\) and
  \(L \geq 0\), then \(\sqrt{a_n} \to \sqrt{L}\) (prove in text ---
  conjugate trick)
\end{itemize}

\subsection{Applications: Finishing What We
Started}\label{applications-finishing-what-we-started}

\begin{itemize}
\tightlist
\item
  Discussion: using limit laws, work inward to outward --- we can't
  assume the full limit exists, but prove subparts converge, then use
  limit laws to conclude more until we get the whole thing
\item
  ✎ Inline: Compute \(\lim \frac{3n^2 + n}{2n^2 + 5}\) using limit laws
  (divide by \(n^2\), apply laws)
\end{itemize}

\emph{Geometric series} - From §2.1: \(S_n = \frac{1 - r^{n+1}}{1-r}\) -
From §2.2: \(r^n \to 0\) for \(|r| < 1\) - By limit laws:
\(S_n \to \frac{1-0}{1-r} = \frac{1}{1-r}\) - Theorem: for \(|r| < 1\):
\(\sum_{n=0}^{\infty} r^n = \frac{1}{1-r}\)

\emph{Archimedes resolved} - The area under the parabola:
\(T(1 + 1/4 + 1/16 + \cdots) = T \cdot \frac{1}{1-1/4} = \frac{4T}{3}\)
- Our tools have reached Archimedes' answer!

\emph{Telescoping series} - From §2.1: \(S_n = 1 - \frac{1}{n+1}\) - By
limit laws: \(S_n \to 1 - 0 = 1\) - Theorem:
\(\sum_{k=1}^{\infty} \frac{1}{k(k+1)} = 1\)

\emph{Telescoping product} - From §2.1:
\(P_n = \frac{n+1}{2n} = \frac{1}{2}(1 + \frac{1}{n})\) - By limit laws:
\(P_n \to \frac{1}{2}(1 + 0) = \frac{1}{2}\) - Theorem:
\(\prod_{k=2}^{\infty}(1 - \frac{1}{k^2}) = \frac{1}{2}\)

\emph{More applications} - Conjugate trick:
\(\sqrt{n+1} - \sqrt{n} = \frac{1}{\sqrt{n+1} + \sqrt{n}} \to 0\) -
\(n^{1/n} \to 1\): write \(n^{1/n} = 1 + h_n\); then
\(n = (1+h_n)^n \geq \binom{n}{2}h_n^2 = \frac{n(n-1)}{2}h_n^2\); so
\(h_n \leq \sqrt{\frac{2}{n-1}} \to 0\) by squeeze

\emph{Babylonian sequence converges} - From §2.1 exercise:
\(x_n = p_n/q_n\) where \(p_n^2 - 2q_n^2 = (-1)^n \cdot 2\) - So
\(x_n^2 - 2 = \frac{p_n^2 - 2q_n^2}{q_n^2} = \frac{(-1)^n \cdot 2}{q_n^2}\)
- Since \(q_n \to \infty\) (verify!), we have \(x_n^2 - 2 \to 0\), i.e.,
\(x_n^2 \to 2\) - By the square root law:
\(x_n = \sqrt{x_n^2} \to \sqrt{2}\) - \textbf{Payoff}: The Babylonian
sequence converges to \(\sqrt{2}\) --- no MCT needed! - Remark: In §2.4
we'll give an alternative proof using monotonicity and completeness

\subsection{Density and Approximation}\label{density-and-approximation}

\begin{itemize}
\tightlist
\item
  Recall from §1.5: between any two reals lies a rational (density of
  \(\QQ\))
\item
  Theorem: every real is the limit of a sequence of rationals

  \begin{itemize}
  \tightlist
  \item
    Proof: given \(x \in \RR\), use density to find \(r_n \in \QQ\) with
    \(|r_n - x| < 1/n\); by squeeze, \(r_n \to x\)
  \end{itemize}
\item
  \textbf{Payoff}: sequences of rationals reach all of \(\RR\)!
\item
  Remark: more generally, any dense subset can describe all reals as
  limits
\end{itemize}

\subsection{The Limits of Limit Laws}\label{the-limits-of-limit-laws}

\begin{itemize}
\tightlist
\item
  Key observation: limit laws tell us \emph{what} the limit is, not
  \emph{that} it exists
\item
  Example: Babylonian sequence \(x_1 = 2\),
  \(x_{n+1} = (x_n + 2/x_n)/2\)

  \begin{itemize}
  \tightlist
  \item
    \emph{If} \(x_n \to L\), then by limit laws:
    \(L = \frac{L + 2/L}{2}\), so \(L^2 = 2\), so \(L = \sqrt{2}\)
  \item
    But this doesn't prove convergence! We've only shown: \emph{if} it
    converges, \emph{then} the limit is \(\sqrt{2}\)
  \end{itemize}
\item
  We need a tool that uses completeness to guarantee existence --- next
  section
\end{itemize}

\subsection{Guided Exercise}\label{guided-exercise-3}

\subsection{The Koch Snowflake: Finite Area, Infinite
Perimeter}\label{the-koch-snowflake-finite-area-infinite-perimeter}

The Koch snowflake is a fractal curve constructed by repeatedly adding
``tents'' to each edge of an equilateral triangle.

\emph{Part I: The Construction}

Start with an equilateral triangle of side length 1.

\begin{enumerate}
\def\labelenumi{(\alph{enumi})}
\item
  Describe the iteration: each edge is divided into thirds, and the
  middle third is replaced by two sides of an equilateral ``tent''
  pointing outward.
\item
  Let \(E_n\) = number of edges and \(L_n\) = length of each edge at
  stage \(n\). Verify \(E_0 = 3\), \(L_0 = 1\).
\item
  Explain why \(E_{n+1} = 4E_n\) and \(L_{n+1} = L_n/3\).
\item
  Derive formulas: \(E_n = 3 \cdot 4^n\) and \(L_n = (1/3)^n\).
\end{enumerate}

\emph{Part II: Perimeter Diverges}

\begin{enumerate}
\def\labelenumi{(\alph{enumi})}
\setcounter{enumi}{4}
\item
  Show the perimeter at stage \(n\) is
  \(P_n = E_n \cdot L_n = 3 \cdot (4/3)^n\).
\item
  Since \(4/3 > 1\), conclude that \(P_n \to \infty\). The Koch
  snowflake has infinite perimeter.
\end{enumerate}

\emph{Part III: Area Converges}

\begin{enumerate}
\def\labelenumi{(\alph{enumi})}
\setcounter{enumi}{6}
\item
  Let \(A_0 = \frac{\sqrt{3}}{4}\) be the area of the original
  equilateral triangle (side 1). At stage 1, we add 3 new triangles,
  each with side \(1/3\). What is the area of each? What is the total
  added area at stage 1?
\item
  At stage \(n \geq 1\), we add \(E_{n-1} = 3 \cdot 4^{n-1}\) new
  triangles, each with side \(L_n = (1/3)^n\). Show the added area at
  stage \(n\) is:
  \[\Delta A_n = 3 \cdot 4^{n-1} \cdot \frac{\sqrt{3}}{4} \cdot \frac{1}{9^n} = A_0 \cdot \frac{1}{3} \cdot \left(\frac{4}{9}\right)^{n-1}\]
\item
  The total area is \(A_0 + \sum_{n=1}^{\infty} \Delta A_n\). Factor out
  and sum the geometric series:
  \[A = A_0 + A_0 \cdot \frac{1}{3} \sum_{k=0}^{\infty} \left(\frac{4}{9}\right)^k = A_0 \left(1 + \frac{1}{3} \cdot \frac{1}{1 - 4/9}\right) = A_0 \cdot \frac{8}{5}\]
\item
  Conclude: the Koch snowflake has finite area
  \(\frac{8}{5} \cdot \frac{\sqrt{3}}{4} = \frac{2\sqrt{3}}{5}\).
\end{enumerate}

\emph{Part IV: Reflection}

\begin{enumerate}
\def\labelenumi{(\alph{enumi})}
\setcounter{enumi}{10}
\item
  We have a bounded region (finite area) with an infinitely long
  boundary. This is impossible for polygons --- why?
\item
  What property of the Koch snowflake allows it to have infinite
  perimeter while remaining bounded? (Hint: think about how ``crinkled''
  the boundary becomes.)
\end{enumerate}

\subsection{Exercises}\label{exercises-7}

\begin{itemize}
\tightlist
\item
  Prove uniqueness of limits using squeeze (alternative to §2.2
  approach)
\item
  Prove the quotient law (hint: first show \(1/b_n \to 1/M\) when
  \(b_n \to M \neq 0\))
\item
  Prove: if \(a_n \to L\) and \(a_n \leq M\) for all \(n\), then
  \(L \leq M\)
\item
  Prove: if \(a_n \to +\infty\), then \(1/a_n \to 0\)
\item
  Prove the converse is false: if \(a_n \to 0\), it is not necessarily
  true that \(1/a_n \to +\infty\) or \(-\infty\) (but
  \(|1/a_n| \to +\infty\))
\item
  Compute \(\lim \frac{n^2 - 3n + 1}{4n^2 + 2}\) using limit laws
\item
  Compute \(\lim (\sqrt{n^2 + n} - n)\) using conjugate trick
\item
  Prove: \(|a_n - L| \leq b_n\) and \(b_n \to 0\) implies \(a_n \to L\)
\item
  Use the telescoping product to prove:
  \(\sum_{k=2}^{\infty} \ln(1 - 1/k^2) = -\ln 2\)
\item
  Prove: every real is the limit of a sequence of irrationals
\item
  ★ Prove: if \(a_n \to L\) and \(b_n \to M\), then
  \(\max\{a_n, b_n\} \to \max\{L, M\}\)
\end{itemize}

\subsection{Dependencies}\label{dependencies-7}

\textbf{Requires}: §2.1 (geometric and telescoping formulas), §2.2
(fundamental limits, boundedness), §1.5 (density of \(\QQ\))

\textbf{Used in}: §2.4 (MCT combines with limit laws), §2.5 (decimal
convergence)

\section{Monotone Convergence}\label{monotone-convergence}

\subsection{Narrative}\label{narrative-8}

\begin{itemize}
\tightlist
\item
  Limit laws derive new limits from old using ordered field axioms
\item
  But they can't prove existence from nothing --- can't prove rationals
  converge to irrational
\item
  Completeness is the missing ingredient
\item
  MCT: monotone + bounded ⟹ convergent
\item
  Babylonian sequence: the payoff --- our tool reaches \(\sqrt{2}\)
\item
  The number \(e\) defined via MCT
\item
  More examples: nested radicals, series
\end{itemize}

\subsection{Content}\label{content-8}

\subsection{The Key Insight}\label{the-key-insight}

\begin{itemize}
\tightlist
\item
  Limit laws use ordered field axioms to derive new limits from known
  limits
\item
  But they can't conjure convergence from nothing
\item
  Example: we showed \emph{if} Babylonian converges, \emph{then} limit
  is \(\sqrt{2}\) --- but didn't prove convergence
\item
  To prove a sequence of rationals converges to an irrational, we need
  completeness
\item
  The Monotone Convergence Theorem uses completeness to guarantee
  existence
\end{itemize}

\subsection{The Monotone Convergence
Theorem}\label{the-monotone-convergence-theorem}

\begin{itemize}
\tightlist
\item
  Definition: \((a_n)\) is increasing if \(a_n \leq a_{n+1}\) for all
  \(n\)
\item
  Definition: \((a_n)\) is strictly increasing if \(a_n < a_{n+1}\) for
  all \(n\)
\item
  Definition: \((a_n)\) is decreasing if \(a_n \geq a_{n+1}\) for all
  \(n\)
\item
  Definition: \((a_n)\) is monotone if it is increasing or decreasing
\item
  Theorem (MCT): A monotone bounded sequence converges. Moreover:

  \begin{itemize}
  \tightlist
  \item
    If \((a_n)\) is increasing and bounded above, then
    \(a_n \to \sup\{a_n : n \in \NN\}\)
  \item
    If \((a_n)\) is decreasing and bounded below, then
    \(a_n \to \inf\{a_n : n \in \NN\}\)
  \end{itemize}
\item
  Proof of increasing case (in text): Let \(L = \sup\{a_n\}\). By
  \(\varepsilon\)-characterization of sup (§1.4), for any
  \(\varepsilon > 0\), some \(a_N > L - \varepsilon\). Since \((a_n)\)
  increasing, \(a_n > L - \varepsilon\) for all \(n > N\). Since \(L\)
  is upper bound, \(a_n \leq L < L + \varepsilon\). So
  \(|a_n - L| < \varepsilon\) for \(n > N\).
\item
  ✎ Inline: Prove the decreasing case
\item
  Remark: MCT fails in \(\QQ\) --- the sequence
  \(1, 1.4, 1.41, 1.414, \ldots\) is increasing and bounded in \(\QQ\)
  but has no rational limit. This is where completeness enters the
  theory of limits.
\end{itemize}

\subsection{The Babylonian Sequence}\label{the-babylonian-sequence}

\begin{itemize}
\tightlist
\item
  Recall from §2.1: \(x_1 = 2\), \(x_{n+1} = (x_n + 2/x_n)/2\)
\item
  We proved convergence in §2.3 using the Pell structure and square root
  law
\item
  Here's an alternative proof using MCT --- this approach generalizes to
  sequences without closed forms
\item
  Theorem: \(x_n \to \sqrt{2}\)
\item
  Step 1: \(x_n^2 \geq 2\) for all \(n \geq 2\)

  \begin{itemize}
  \tightlist
  \item
    Compute:
    \(x_{n+1}^2 - 2 = \frac{(x_n + 2/x_n)^2}{4} - 2 = \frac{(x_n - 2/x_n)^2}{4} \geq 0\)
  \item
    Remark: alternatively, use AM-GM (§1.2 exercise):
    \(x_{n+1} = \frac{x_n + 2/x_n}{2} \geq \sqrt{x_n \cdot \frac{2}{x_n}} = \sqrt{2}\)
  \end{itemize}
\item
  Step 2: \((x_n)\) is decreasing for \(n \geq 2\)

  \begin{itemize}
  \tightlist
  \item
    Compute:
    \(x_{n+1} - x_n = \frac{2/x_n - x_n}{2} = \frac{2 - x_n^2}{2x_n} \leq 0\)
    since \(x_n^2 \geq 2\)
  \end{itemize}
\item
  Step 3: By MCT, \(x_n \to L\) for some \(L \geq \sqrt{2}\)
\item
  Step 4: By limit laws (§2.3): \(L = \frac{L + 2/L}{2}\), so
  \(L^2 = 2\), so \(L = \sqrt{2}\)
\item
  \textbf{Lesson}: MCT proves convergence for recursive sequences even
  when no closed form exists
\item
  Historical note: known to Babylonians \textasciitilde1800 BCE;
  converges remarkably fast (roughly doubles correct digits each
  iteration)
\end{itemize}

\subsection{\texorpdfstring{The Number
\(e\)}{The Number e}}\label{the-number-e}

\begin{itemize}
\tightlist
\item
  Definition: \(e = \lim_{n \to \infty}(1 + 1/n)^n\), provided this
  limit exists
\item
  This subsection is a guided exercise proving \(e\) exists and
  establishing its basic properties
\end{itemize}

\subsection{Guided Exercises}\label{guided-exercises-4}

\subsection{\texorpdfstring{The Number
\(e\)}{The Number e}}\label{the-number-e-1}

Let \(a_n = (1 + 1/n)^n\) and \(b_n = (1 + 1/n)^{n+1}\).

\begin{enumerate}
\def\labelenumi{(\alph{enumi})}
\item
  Use the binomial theorem to expand
  \(a_n = \sum_{k=0}^{n} \binom{n}{k} \frac{1}{n^k}\).
\item
  Show that each term in the expansion increases with \(n\), and
  conclude \((a_n)\) is increasing.
\item
  Show
  \(a_n < 1 + 1 + \frac{1}{2!} + \frac{1}{3!} + \cdots + \frac{1}{n!}\).
\item
  Use \(k! \geq 2^{k-1}\) for \(k \geq 1\) to show \(a_n < 3\) for all
  \(n\).
\item
  Conclude by MCT that \((a_n)\) converges. Define \(e = \lim a_n\).
\item
  Show that \((b_n)\) is decreasing and bounded below by \(a_1 = 2\).
\item
  Conclude \((b_n)\) converges and show \(\lim b_n = e\) as well.
\item
  The intervals \([a_n, b_n]\) are nested and contain \(e\). Since
  \(b_n - a_n = a_n/n \to 0\), these intervals pin down \(e\) uniquely.
\end{enumerate}

\subsection{More MCT Examples}\label{more-mct-examples}

\begin{itemize}
\tightlist
\item
  One example worked in text: nested radicals

  \begin{itemize}
  \tightlist
  \item
    Let \(a_1 = 1\), \(a_{n+1} = \sqrt{2 + a_n}\)
  \item
    Bounded above by 2: if \(a_n \leq 2\), then
    \(a_{n+1} = \sqrt{2 + a_n} \leq \sqrt{4} = 2\) (induction)
  \item
    Increasing:
    \(a_{n+1}^2 - a_n^2 = (2 + a_n) - a_n^2 = (2 - a_n)(1 + a_n) > 0\)
    when \(a_n < 2\)
  \item
    MCT ⟹ converges to some \(L\)
  \item
    Limit laws: \(L = \sqrt{2 + L}\), so \(L^2 = 2 + L\), so
    \(L^2 - L - 2 = 0\), so \(L = 2\) (taking positive root)
  \item
    Interpretation: \(\sqrt{2 + \sqrt{2 + \sqrt{2 + \cdots}}} = 2\)
  \end{itemize}
\end{itemize}

\emph{Series via MCT} - Theorem (Nonnegative Series Criterion): If
\(a_n \geq 0\), then \(\sum a_n\) converges iff partial sums \(S_n\) are
bounded - Proof: partial sums increasing (since \(a_n \geq 0\)); apply
MCT - Example: \(\sum 1/2^n\) converges --- partial sums
\(S_n = 2 - 1/2^n < 2\) - Example: \(\sum 1/n^2\) converges - For
\(n \geq 2\):
\(\frac{1}{n^2} < \frac{1}{n(n-1)} = \frac{1}{n-1} - \frac{1}{n}\)
(telescoping!) - So
\(S_n < 1 + \sum_{k=2}^{n}\left(\frac{1}{k-1} - \frac{1}{k}\right) = 1 + (1 - 1/n) < 2\)
- Bounded, so converges - Example: The harmonic series \(\sum 1/n\)
diverges (prove in text) - Partial sums
\(H_n = 1 + 1/2 + 1/3 + \cdots + 1/n\) are unbounded - Proof:
\(H_{2n} - H_n = \frac{1}{n+1} + \frac{1}{n+2} + \cdots + \frac{1}{2n} > n \cdot \frac{1}{2n} = \frac{1}{2}\)
- So \(H_{2^k} > k/2\), which is unbounded - Remark: this is the
``boundary case'' --- \(\sum 1/n^p\) converges for \(p > 1\), diverges
for \(p \leq 1\) - Remark: we proved convergence/divergence
\emph{without computing limits}. The sum \(\sum 1/n^2 = \pi^2/6\), but
proving this requires Fourier analysis!

\subsection{Exercises}\label{exercises-8}

\begin{itemize}
\tightlist
\item
  Prove: if \((a_n)\) is increasing and unbounded, then
  \(a_n \to +\infty\)
\item
  Let \(a_1 = 1\), \(a_{n+1} = (a_n + 3)/(a_n + 1)\). Prove convergence
  and find the limit.
\item
  Power tower: let \(a_1 = 1\), \(a_{n+1} = (\sqrt{2})^{a_n}\). Show
  \((a_n)\) is increasing, bounded above by 2, and find the limit.
\item
  Prove \(\sum 1/n!\) converges
\item
  Prove \(\sum 1/n^3\) converges
\item
  ★ Prove: if \(a_{n+1}/a_n \to L\) with \(0 \leq L < 1\) and
  \(a_n > 0\), then \(a_n \to 0\)
\item
  ★ For which values of \(c > 0\) does the sequence \(a_1 = c\),
  \(a_{n+1} = (\sqrt{2})^{a_n}\) converge?
\end{itemize}

\subsection{Dependencies}\label{dependencies-8}

\textbf{Requires}: §1.4 (\(\varepsilon\)-characterization of sup), §1.5
(completeness, AM-GM), §2.2 (convergence definition), §2.3 (limit laws)

\textbf{Used in}: §2.5 (decimal convergence), Chapter 3
(Bolzano-Weierstrass), Chapter 4 (series tests)

\section{\texorpdfstring{Representing Real Numbers
\emph{(Terminal)}}{Representing Real Numbers (Terminal)}}\label{representing-real-numbers-terminal}

\subsection{Narrative}\label{narrative-9}

\begin{itemize}
\tightlist
\item
  The naming problem: finite descriptions are countable, \(\RR\) is
  uncountable
\item
  Infinite processes are \emph{necessary} to name reals --- not just
  convenient
\item
  Decimals give a systematic way to represent any real
\item
  Every decimal converges (MCT); every real has a decimal (greedy
  algorithm)
\item
  This completes the story: we can now describe all real numbers
\item
  Uniqueness up to trailing 9s; \(0.999\ldots = 1\)
\item
  Rationals characterized by eventually periodic decimals
\end{itemize}

\subsection{Content}\label{content-9}

\subsection{The Naming Problem}\label{the-naming-problem}

\begin{itemize}
\tightlist
\item
  How do we describe a specific real number?
\item
  Rationals: finite data (two integers \(p, q\))
\item
  But \(\RR\) is uncountable, finite descriptions are countable
\item
  Conclusion: representing arbitrary reals \emph{requires} infinite
  processes
\item
  Limits aren't just convenient --- they're necessary
\item
  Decimals give a systematic infinite representation
\end{itemize}

\subsection{Decimal Expansions}\label{decimal-expansions}

\begin{itemize}
\tightlist
\item
  Review: integers in base 10, finite decimals
  \(0.d_1d_2\ldots d_n = \sum_{k=1}^{n} d_k/10^k\)
\item
  Definition: infinite decimal \(0.d_1d_2d_3\ldots\) represents the
  series \(\sum_{n=1}^{\infty} d_n/10^n\)
\item
  Theorem: Every infinite decimal converges to a real number in
  \([0,1]\)

  \begin{itemize}
  \tightlist
  \item
    Proof: partial sums \(S_n = \sum_{k=1}^{n} d_k/10^k\) are increasing
    (all \(d_k \geq 0\))
  \item
    Bounded:
    \(S_n \leq \sum_{k=1}^{n} 9/10^k < \sum_{k=1}^{\infty} 9/10^k = 1\)
  \item
    By MCT, the series converges
  \end{itemize}
\item
  Theorem: Every \(x \in [0,1)\) has a decimal expansion

  \begin{itemize}
  \tightlist
  \item
    Construction (greedy algorithm):
    \(d_n = \lfloor 10^n x \rfloor - 10\lfloor 10^{n-1} x \rfloor\)
  \item
    Equivalently: \(d_1 = \lfloor 10x \rfloor\), then repeat with
    \(10x - d_1\)
  \item
    Proof: by construction, \(S_n \leq x < S_n + 1/10^n\); by squeeze,
    \(S_n \to x\)
  \item
    Full narrative in text; parts of proof may be exercises
  \end{itemize}
\item
  ✎ Inline: Use greedy algorithm to find first 4 decimal digits of
  \(1/3\)
\item
  Remark: extend to all reals by handling integer part and sign
  separately
\end{itemize}

\subsection{Uniqueness}\label{uniqueness}

\begin{itemize}
\tightlist
\item
  Question: is the decimal representation unique?
\item
  Theorem: Decimal representations are unique except for trailing 9s:
  \[0.d_1 d_2 \cdots d_k 999\ldots = 0.d_1 d_2 \cdots (d_k + 1) 000\ldots\]
\item
  Theorem: \(0.999\ldots = 1\)

  \begin{itemize}
  \tightlist
  \item
    Proof:
    \(0.999\ldots = \sum_{n=1}^{\infty} 9/10^n = 9 \cdot \frac{1/10}{1 - 1/10} = 9 \cdot \frac{1}{9} = 1\)
  \end{itemize}
\item
  Convention: avoid trailing 9s; then every real has a unique decimal
  expansion
\end{itemize}

\subsection{Characterizing Rationals}\label{characterizing-rationals}

\begin{itemize}
\tightlist
\item
  Theorem: \(x \in \QQ\) if and only if its decimal expansion is
  eventually periodic
\item
  (⟸) A repeating block is a geometric series, hence rational
\item
  ✎ Inline: Convert \(0.\overline{142857}\) to a fraction
\item
  (⟹) Long division of \(p\) by \(q\) cycles through at most \(q\)
  remainders; once a remainder repeats, the digits repeat (discuss in
  text)
\item
  Corollary: \(x\) is irrational iff its decimal never becomes periodic
\item
  Examples: \(\sqrt{2} = 1.41421356\ldots\) (non-repeating),
  \(\pi = 3.14159265\ldots\) (non-repeating)
\end{itemize}

\subsection{The Big Picture (conclusion to above
section)}\label{the-big-picture-conclusion-to-above-section}

\begin{itemize}
\tightlist
\item
  Rationals: countable, eventually periodic decimals
\item
  Algebraic irrationals (like \(\sqrt{2}\)): countable, non-repeating
\item
  Transcendentals: uncountable, ``most'' reals
\item
  Most real numbers cannot be described by any finite means
\item
  Decimals give a systematic way to approximate any real, but the full
  decimal of a ``generic'' real cannot be computed by any algorithm
\end{itemize}

\subsection{Other Bases}\label{other-bases}

\begin{itemize}
\tightlist
\item
  Everything above works in any integer base \(b \geq 2\)
\item
  Base \(b\): digits \(d_k \in \{0, 1, \ldots, b-1\}\), expansion
  \(\sum d_k / b^k\)
\item
  Uniqueness: trailing \((b-1)\)s play the role of trailing \(9\)s
\item
  Rationals ↔ eventually periodic in \emph{every} integer base (same
  proof: long division cycles)
\item
  Remark: the characterization of rationals is base-independent --- a
  beautiful fact
\end{itemize}

\subsection{Guided Exercises}\label{guided-exercises-5}

\subsection{Constructing
Transcendentals}\label{constructing-transcendentals}

We know algebraic numbers are countable, so transcendentals exist ---
but can we exhibit one explicitly?

Define Liouville's number:
\(L = \sum_{n=1}^{\infty} 10^{-n!} = 0.110001000000000000000001\ldots\)

\begin{enumerate}
\def\labelenumi{(\alph{enumi})}
\item
  Verify that \(L\) is irrational by showing its decimal expansion is
  not eventually periodic.
\item
  For each \(n\), let \(p_n/q_n = \sum_{k=1}^{n} 10^{-k!}\) with
  \(q_n = 10^{n!}\). Show that
  \(|L - p_n/q_n| = \sum_{k=n+1}^{\infty} 10^{-k!} < 2 \cdot 10^{-(n+1)!}\).
\item
  Conclude that \(|L - p_n/q_n| < 1/q_n^n\) for all \(n\).
\item
  Liouville's theorem (stated): if \(\alpha\) is algebraic of degree
  \(d\), then there exists \(c > 0\) such that
  \(|\alpha - p/q| > c/q^d\) for all rationals \(p/q\). Using this,
  conclude \(L\) is transcendental.
\item
  Construct your own Liouville number and verify it is transcendental.
\end{enumerate}

\textbf{Irrational Bases (Exploratory)} \emph{(Optional --- may be
removed)}

\emph{{[}Note: This is exploratory material. Need to verify: what
numbers terminate? what do integers look like? Is this worth
including?{]}}

For \(\beta > 1\) irrational, we can attempt base-\(\beta\) expansions:
represent \(x\) as \(\sum d_k \beta^{-k}\) with digits
\(d_k \in \{0, 1, \ldots, \lfloor \beta \rfloor\}\).

\begin{enumerate}
\def\labelenumi{(\alph{enumi})}
\item
  In base \(\sqrt{2}\), we use digits \(\{0, 1\}\). Show that
  \(\sqrt{2} = (10)_{\sqrt{2}}\) and \(2 = (100)_{\sqrt{2}}\).
\item
  Find the base-\(\sqrt{2}\) representation of \(3\) and \(4\).
\item
  What is \((0.1)_{\sqrt{2}}\) as a real number?
\item
  ★ Investigate: which numbers have terminating base-\(\sqrt{2}\)
  expansions?
\end{enumerate}

\subsection{Exercises}\label{exercises-9}

\begin{itemize}
\tightlist
\item
  Prove: the greedy algorithm produces digits
  \(d_n \in \{0, 1, \ldots, 9\}\)
\item
  Prove: if \(0.d_1d_2d_3\ldots = 0.e_1e_2e_3\ldots\) and neither has
  trailing 9s, then \(d_n = e_n\) for all \(n\)
\item
  Find the decimal expansion of \(1/7\). What is the period?
\item
  Which rationals have terminating (finite) decimal expansions? (Hint:
  consider the denominator's prime factors)
\item
  Prove:
  \(0.\overline{d_1 d_2 \cdots d_k} = \frac{d_1 d_2 \cdots d_k}{10^k - 1}\)
  where the bar denotes repetition
\item
  Prove: in base \(b\), \(0.\overline{(b-1)} = 1\) (generalize
  \(0.999\ldots = 1\))
\item
  Prove: in base \(b\), a rational has eventually periodic expansion
  (same argument as base 10)
\item
  Which rationals have terminating expansions in base \(b\)? (Generalize
  the base-10 result)
\item
  Construct a decimal that is clearly irrational (no eventual period)
\item
  ★ Prove: every real is the limit of a sequence of terminating decimals
\item
  ★ Binary expansions: find the base-2 expansion of \(1/3\), \(1/5\),
  \(1/10\)
\end{itemize}

\subsection{Dependencies}\label{dependencies-9}

\textbf{Requires}: §2.3 (geometric series), §2.4 (MCT)

\textbf{Used in}: None --- this is a capstone section

\section{Quadrature of the Parabola}\label{quadrature-of-the-parabola}

\subsection{Narrative}\label{narrative-10}

\begin{itemize}
\tightlist
\item
  Archimedes (\textasciitilde250 BCE) found the area of a parabolic
  segment
\item
  We prove it rigorously using §1.5 (area via inner/outer
  approximations) and geometric series
\item
  The geometric facts are taken for granted in the main text; proofs are
  guided exercises
\item
  This connects Chapter 1 (rigorous definition of area) with Chapter 2
  (geometric series)
\end{itemize}

\subsection{Content}\label{content-10}

\emph{(Opening --- unnumbered)} - The parabolic segment: region bounded
by parabola \(y = x^2\) and a chord - One of Archimedes' greatest
achievements: finding this area - He used the ``method of exhaustion''
--- we'll use our rigorous framework from §1.5 plus geometric series

\subsection{The Setup}\label{the-setup}

\begin{itemize}
\tightlist
\item
  Consider \(y = x^2\) with chord from \((a, a^2)\) to \((b, b^2)\)
\item
  The inscribed triangle: third vertex at \(x = (a+b)/2\), the point on
  the parabola above the midpoint
\item
  Let \(T\) = area of this inscribed triangle
\item
  Goal: show the parabolic segment has area \(\frac{4T}{3}\)
\end{itemize}

\subsection{Geometric Facts}\label{geometric-facts}

\begin{itemize}
\tightlist
\item
  We need three facts (proofs as guided exercise):
\item
  Fact 1 (Inscribed triangle area): Using \(m = (a+b)/2\),
  \(h = (b-a)/2\), the inscribed triangle has area \(h^3\)
\item
  Fact 2 (Circumscribed parallelogram): The parallelogram circumscribing
  the segment (sides parallel to chord and tangent at vertex) has area
  \(2h^3 = 2T\)
\item
  Fact 3 (Child triangles): The two child triangles (inscribed in the
  two remaining segments) have combined area \(T/4\)
\end{itemize}

\subsection{Inner Approximations}\label{inner-approximations}

\begin{itemize}
\tightlist
\item
  Stage 0: One inscribed triangle, area \(T\)
\item
  Stage 1: Add 2 child triangles, total area \(T + T/4\)
\item
  Stage 2: Add 4 grandchild triangles, total area \(T + T/4 + T/16\)
\item
  Stage \(n\): Total inscribed area
  \(= T(1 + 1/4 + 1/16 + \cdots + 1/4^n) = T \cdot \frac{1 - (1/4)^{n+1}}{1 - 1/4}\)
\item
  The inscribed regions are polygons contained in the parabolic segment
\item
  By §1.5: inner area
  \(= \sup\{\text{inscribed polygon areas}\} \geq T \cdot \frac{1 - (1/4)^{n+1}}{1 - 1/4}\)
  for all \(n\)
\item
  As \(n \to \infty\): inner area
  \(\geq T \cdot \frac{1}{1 - 1/4} = \frac{4T}{3}\)
\end{itemize}

\subsection{Outer Approximations}\label{outer-approximations}

\begin{itemize}
\tightlist
\item
  At stage \(n\): replace the last batch of triangles with their
  circumscribed parallelograms
\item
  Each parallelogram has area twice its inscribed triangle
\item
  Outer figure at stage \(n\) = inner figure at stage \(n\) + (extra
  area from parallelograms instead of triangles)
\item
  Extra area = sum of last batch of triangles = \(T/4^n\)
\item
  So: outer area at stage \(n\) = inner area at stage \(n\) + \(T/4^n\)
\item
  These outer figures contain the parabolic segment
\item
  By §1.5: outer area
  \(= \inf\{\text{circumscribed polygon areas}\} \leq T \cdot \frac{1 - (1/4)^{n+1}}{1 - 1/4} + T/4^n\)
  for all \(n\)
\item
  As \(n \to \infty\): outer area \(\leq \frac{4T}{3}\)
\end{itemize}

\subsection{Conclusion}\label{conclusion}

\begin{itemize}
\tightlist
\item
  Inner area \(\geq \frac{4T}{3}\) and outer area \(\leq \frac{4T}{3}\)
\item
  Since inner area \(\leq\) outer area always, we have inner area \(=\)
  outer area \(= \frac{4T}{3}\)
\item
  By §1.5 definition: the parabolic segment has area, and it equals
  \(\frac{4T}{3}\)
\item
  Theorem (Archimedes): The area of a parabolic segment is
  \(\frac{4}{3}\) times the area of its inscribed triangle
\end{itemize}

\subsection{Historical Remark}\label{historical-remark}

\begin{itemize}
\tightlist
\item
  Archimedes proved this using the method of exhaustion --- essentially
  our inner/outer argument
\item
  He also had a mechanical method using levers and centers of gravity
\item
  Our proof makes rigorous what Archimedes knew intuitively: the inner
  and outer bounds converge
\end{itemize}

\subsection{Guided Exercise}\label{guided-exercise-4}

\subsection{The Geometric Facts}\label{the-geometric-facts}

We establish the three geometric facts needed for the quadrature.

\emph{Setup:} Consider the parabola \(y = x^2\) with chord from
\((a, a^2)\) to \((b, b^2)\). Let \(m = (a+b)/2\) and \(h = (b-a)/2\),
so the endpoints are \((m-h, (m-h)^2)\) and \((m+h, (m+h)^2)\).

\begin{enumerate}
\def\labelenumi{(\alph{enumi})}
\item
  The inscribed triangle has vertices at \((m-h, (m-h)^2)\),
  \((m+h, (m+h)^2)\), and \((m, m^2)\). Using the formula
  \[\text{Area} = \frac{1}{2}|x_1(y_2 - y_3) + x_2(y_3 - y_1) + x_3(y_1 - y_2)|\]
  show that the area equals \(h^3\).
\item
  The circumscribed parallelogram has two sides along the lines through
  the chord endpoints parallel to the tangent at \((m, m^2)\), and two
  sides along lines parallel to the chord. Show its area is \(2h^3\).
\item
  The right child segment has chord from \((m, m^2)\) to
  \((m+h, (m+h)^2)\). Its midpoint x-coordinate is \(m + h/2\), so its
  half-width is \(h' = h/2\). By part (a), its inscribed triangle has
  area \((h')^3 = h^3/8\).
\item
  Similarly, the left child triangle has area \(h^3/8\). Conclude: the
  two child triangles have combined area \(h^3/4 = T/4\).
\item
  Explain why, at each stage, the new batch of triangles has total area
  \(1/4\) of the previous batch.
\end{enumerate}

\subsection{Exercises}\label{exercises-10}

\emph{Geometric facts} - Verify part (a) of the guided exercise by
direct computation - ★ Prove part (b) without using calculus (hint: use
similar triangles or direct coordinate calculation)

\emph{The series} - Verify \(1 + 1/4 + 1/16 + \cdots = 4/3\) using the
geometric series formula - At stage \(n\), how many triangles have been
inscribed? What is the total area?

\emph{Variations} - For the parabola \(y = cx^2\) (any \(c > 0\)), show
the quadrature result still holds: area = \(\frac{4}{3}T\) - ★ Set up
the quadrature for \(y = x^2\) on \([0, 1]\) and compute the area
directly. Verify it equals \(\frac{4}{3}T\) where \(T\) is the inscribed
triangle area.

\emph{Historical} - Archimedes' original inscribed triangle for
\(y = x^2\) on \([-1, 1]\) has vertices at \((-1, 1)\), \((1, 1)\),
\((0, 0)\). Compute its area and hence the area of the parabolic
segment.

\emph{The rectangle method (preview of integration)}

These exercises develop the ``modern'' approach to area: approximating
with rectangles rather than triangles. This method generalizes to
arbitrary curves and becomes the definition of the integral (Chapter 8).

\begin{itemize}
\item
  Divide \([0, a]\) into \(k\) equal subintervals of width \(a/k\). The
  \emph{lower sum} uses rectangles with heights from the left endpoints:
  \(L_k = \sum_{j=0}^{k-1} \left(\frac{ja}{k}\right)^2 \cdot \frac{a}{k}\).
  The \emph{upper sum} uses right endpoints:
  \(U_k = \sum_{j=1}^{k} \left(\frac{ja}{k}\right)^2 \cdot \frac{a}{k}\).
  Draw a picture showing \(L_k \leq \text{area} \leq U_k\).
\item
  Show that \(L_k = \frac{a^3}{k^3} \sum_{j=0}^{k-1} j^2\) and
  \(U_k = \frac{a^3}{k^3} \sum_{j=1}^{k} j^2\).
\item
  Using the formula \(\sum_{j=1}^{k} j^2 = \frac{k(k+1)(2k+1)}{6}\) from
  §1.1, show that:
  \[U_k = \frac{a^3}{k^3} \cdot \frac{k(k+1)(2k+1)}{6} = \frac{a^3(k+1)(2k+1)}{6k^2}\]
\item
  Compute \(\lim_{k \to \infty} U_k\) using limit laws. (Hint: divide
  numerator and denominator by \(k^2\).)
\item
  Similarly compute \(\lim_{k \to \infty} L_k\) and verify both limits
  equal \(\frac{a^3}{3}\).
\item
  Conclude: the area under \(y = x^2\) from \(0\) to \(a\) equals
  \(\frac{a^3}{3}\).
\item
  Verify this matches Archimedes: for the segment from \((0,0)\) to
  \((a, a^2)\), the inscribed triangle has area \(\frac{a^3}{4}\) (check
  this!), so Archimedes predicts area
  \(= \frac{4}{3} \cdot \frac{a^3}{4} = \frac{a^3}{3}\). ✓
\item
  ★ Extend to \(y = x^3\): using
  \(\sum_{j=1}^{k} j^3 = \left(\frac{k(k+1)}{2}\right)^2\), show the
  area under \(y = x^3\) from \(0\) to \(a\) equals \(\frac{a^4}{4}\).
\end{itemize}

\subsection{Dependencies}\label{dependencies-10}

\textbf{Requires}: §1.5 (area via inner/outer approximations), §2.1
(geometric series)

\textbf{Used in}: Historical payoff; connects Chapters 1 and 2

\emph{(Terminal section --- may be deferred or omitted)}

\chapter{Studying Sequences}\label{studying-sequences}

\textbf{The question} - Chapter 2 gave us the Monotone Convergence
Theorem: monotone + bounded ⟹ convergent - But many bounded sequences
aren't monotone - \((-1)^n\) bounces between \(-1\) and \(1\) forever;
it's bounded but doesn't converge - What can we say about bounded
sequences in general? - What does completeness force on sequences that
lack monotone structure?

\textbf{The answer, in brief} - Bolzano-Weierstrass: every bounded
sequence has a convergent subsequence - Even in apparent chaos, we can
find order hiding inside - Limsup and liminf capture the ``spread'' of
possible subsequential limits - The Cauchy criterion: an intrinsic test
for convergence that doesn't require knowing the limit

\textbf{The deeper theme} - These turn out to be equivalent faces of
completeness - LUB, MCT, Nested Intervals, Bolzano-Weierstrass, Cauchy
--- any one implies all the others - Completeness is a single deep
property wearing different masks - In this chapter, we see more of those
masks

\textbf{Historical note} - Bolzano proved his theorem in 1817, decades
before completeness was properly understood - Cauchy's criterion emerged
around the same time - Only later was it recognized that these are
equivalent characterizations of the same underlying property - The full
picture --- five equivalent faces of completeness --- crystallized by
the late 19th century

\textbf{What's ahead} - §3.1: Subsequences and the Bolzano-Weierstrass
theorem - §3.2: Subsequential limits, limsup, and liminf - §3.3: The
Cauchy criterion --- convergence without knowing the limit - §3.4: The
many faces of completeness --- tying it all together

\section{Subsequences}\label{subsequences}

\subsection{Narrative}\label{narrative-11}

\begin{itemize}
\tightlist
\item
  Chapter 2 handled monotone sequences --- MCT guarantees convergence
\item
  But bounded sequences without monotone structure can behave wildly:
  \((-1)^n\) stays in \([-1, 1]\) but doesn't converge
\item
  Subsequences let us isolate and study these behaviors
\item
  Key results: subsequences inherit limits; union theorem; divergence
  test
\item
  The big theorem: Bolzano-Weierstrass --- every bounded sequence has a
  convergent subsequence
\item
  This uses completeness (nested intervals) and fails in ℚ
\item
  Ends with the question: what limits can be achieved by subsequences?
\end{itemize}

\subsection{Content}\label{content-11}

\subsection{Subsequences}\label{subsequences-1}

\begin{itemize}
\tightlist
\item
  Motivation: \((-1)^n\) is bounded but doesn't converge --- it
  oscillates between regions
\item
  Idea: examine parts of a sequence by selecting terms along an
  increasing index sequence
\item
  Definition: a subsequence of \((a_n)\) is a sequence \((a_{n_k})\)
  where \(n_1 < n_2 < n_3 < \cdots\)
\item
  Notation: \((a_{n_k})_{k=1}^{\infty}\) or just \((a_{n_k})\)
\item
  Examples:

  \begin{itemize}
  \tightlist
  \item
    Even-indexed terms: \(a_2, a_4, a_6, \ldots\) (here \(n_k = 2k\))
  \item
    Odd-indexed terms: \(a_1, a_3, a_5, \ldots\) (here \(n_k = 2k-1\))
  \item
    Every third term: \(a_3, a_6, a_9, \ldots\)
  \item
    Terms along squares: \(a_1, a_4, a_9, a_{16}, \ldots\)
  \end{itemize}
\item
  Remark: a subsequence ``samples'' the original sequence --- it can
  skip terms but must preserve order
\item
  Lemma: if \((n_k)\) is strictly increasing in \(\NN\), then
  \(n_k \geq k\) for all \(k\) (quick induction)
\item
  ✎ Inline: For \(a_n = (-1)^n\), write out the first few terms of the
  even-indexed and odd-indexed subsequences
\end{itemize}

\subsection{Subsequences and
Convergence}\label{subsequences-and-convergence}

\begin{itemize}
\tightlist
\item
  Theorem: If \(a_n \to L\), then every subsequence \((a_{n_k}) \to L\)
  (prove in text)
\item
  Proof idea: given \(\varepsilon\), get \(N\) for original sequence;
  since \(n_k \geq k\), we have \(n_k > N\) for \(k > N\)
\item
  Theorem: If \((a_n)\) is the union of finitely many subsequences, each
  converging to \(L\), then \(a_n \to L\) (prove in text)
\item
  Proof idea: given \(\varepsilon\), each subsequence eventually stays
  within \(\varepsilon\) of \(L\); take max of the thresholds
\item
  Corollary (Divergence Test): If two subsequences converge to different
  limits, the original sequence diverges
\item
  Example: \((-1)^n\) diverges because even terms → 1 and odd terms → -1
\item
  ✎ Inline: Prove \(a_n = \sin(n\pi/2)\) diverges by finding
  subsequences with different limits
\item
  Remark: this is cleaner than the direct \(\varepsilon\)-\(N\) proof
  from §2.2
\end{itemize}

\subsection{Bolzano-Weierstrass}\label{bolzano-weierstrass}

\begin{itemize}
\tightlist
\item
  Question: for a bounded sequence, can we always find \emph{some}
  convergent subsequence?
\item
  Theorem (Bolzano-Weierstrass): Every bounded sequence has a convergent
  subsequence
\item
  Proof via bisection (prove in text):

  \begin{itemize}
  \tightlist
  \item
    Let \((a_n)\) be bounded, say \(a_n \in [A, B]\) for all \(n\)
  \item
    Bisect \([A, B]\); at least one half contains infinitely many terms
    --- pick that half, call it \(I_1\)
  \item
    Bisect \(I_1\); again pick the half with infinitely many terms ---
    call it \(I_2\)
  \item
    Continue: nested intervals
    \(I_1 \supseteq I_2 \supseteq I_3 \supseteq \cdots\) with lengths
    \(\to 0\)
  \item
    By nested interval property, \(\bigcap I_n = \{L\}\) for some \(L\)
  \item
    Construct subsequence: pick \(a_{n_1} \in I_1\), then
    \(a_{n_2} \in I_2\) with \(n_2 > n_1\), etc.
  \item
    Since \(a_{n_k} \in I_k\) and \(I_k\) shrinks to \(L\), we have
    \(a_{n_k} \to L\)
  \end{itemize}
\item
  Remark: this proof uses completeness (nested intervals); BW fails in
  \(\QQ\)
\item
  Example: \((a_n)\) in \(\QQ\) with \(a_n \to \sqrt{2}\); the only
  candidate limit is \(\sqrt{2} \notin \QQ\)
\item
  Question: we now know convergent subsequences exist --- but what
  limits can they achieve? (→ §3.2)
\end{itemize}

\subsection{Guided Exercises}\label{guided-exercises-6}

None for this section.

\subsection{Exercises}\label{exercises-11}

\emph{Basic subsequences} - Write out the first 5 terms of the
subsequence of \(a_n = 1/n\) along indices \(n_k = 2^k\) - Write out the
first 5 terms of the subsequence of \(a_n = (-1)^n / n\) along even
indices and along odd indices - Prove: if \((n_k)\) is a strictly
increasing sequence of natural numbers, then \(n_k \geq k\) for all
\(k\)

\emph{Subsequences and limits} - Prove: if \(a_n \to L\) and
\((a_{n_k})\) is a subsequence, then \(a_{n_k} \to L\) - Prove: if
\((a_n)\) is covered by two subsequences both converging to \(L\), then
\(a_n \to L\) - Extend to finitely many subsequences - Give an example
where infinitely many subsequences all converge to \(L\), but
\(a_n \not\to L\)

\emph{Divergence via subsequences} - Prove \(a_n = (-1)^n (1 + 1/n)\)
diverges - Prove \(a_n = \cos(n\pi/3)\) diverges by finding two
subsequences with different limits - Prove \(a_n = (-1)^n + 1/n\)
diverges

\emph{Bolzano-Weierstrass} - Use BW to prove: if \((a_n)\) is bounded
and monotone, it converges (alternate proof of MCT) - Construct
explicitly a convergent subsequence of \(a_n = \sin(n)\) - Construct
explicitly a convergent subsequence of \(a_n = n \mod 7\) (remainder
when dividing by 7)

\emph{Harder} - ★ Prove: every sequence has a monotone subsequence
(either increasing or decreasing) - ★ Use the previous result and MCT to
give an alternate proof of Bolzano-Weierstrass

\subsection{Dependencies}\label{dependencies-11}

\textbf{Requires}: §1.4 (nested intervals), §2.2 (convergence
definition), §2.3 (limit laws)

\textbf{Used in}: §3.2 (subsequential limits, limsup/liminf), §3.3
(Cauchy criterion)

\section{Subsequential Limits}\label{subsequential-limits}

\subsection{Narrative}\label{narrative-12}

\begin{itemize}
\tightlist
\item
  BW guarantees bounded sequences have convergent subsequences --- but
  what limits can they achieve?
\item
  Definition: subsequential limit --- a limit achieved by some
  subsequence
\item
  For bounded sequences, the set of subsequential limits is nonempty
  (BW) and bounded
\item
  Key corollary: if there's only one subsequential limit, the sequence
  converges
\item
  Limsup and liminf: the largest and smallest subsequential limits
\item
  Alternative characterization via tails: limsup = lim(sup of tail)
\item
  These definitions agree --- proof is the heart of the section
\item
  Limsup and liminf are always achieved by some subsequence
\item
  Convergence criterion: \(a_n \to L\) iff limsup = liminf = \(L\)
\end{itemize}

\subsection{Content}\label{content-12}

\subsection{Subsequential Limits}\label{subsequential-limits-1}

\begin{itemize}
\tightlist
\item
  Definition: \(L\) is a subsequential limit of \((a_n)\) if some
  subsequence converges to \(L\)
\item
  Alternative terminology: \(L\) is a limit point or accumulation point
  of the sequence
\item
  Examples:

  \begin{itemize}
  \tightlist
  \item
    If \(a_n \to L\), then \(L\) is the only subsequential limit
  \item
    \((-1)^n\) has subsequential limits \(\{-1, 1\}\)
  \item
    \(a_n = (-1)^n(1 + 1/n)\) has subsequential limits \(\{-1, 1\}\)
  \item
    \(a_n = n\) has no subsequential limits (unbounded, no convergent
    subsequences)
  \end{itemize}
\item
  Theorem: If \((a_n)\) is bounded, the set of subsequential limits is
  nonempty and bounded
\item
  Proof: nonempty by BW; bounded because all subsequential limits lie
  within the bounds of \((a_n)\)
\item
  Theorem: If \((a_n)\) is bounded and has exactly one subsequential
  limit \(L\), then \(a_n \to L\) (prove in text)
\item
  Proof idea: if \(a_n \not\to L\), then infinitely many terms stay away
  from \(L\); apply BW to get another subsequential limit \(\neq L\)
\item
  Reframing: divergence of bounded sequences means \emph{multiple}
  subsequential limits
\item
  ✎ Inline: Find all subsequential limits of \(a_n = \sin(n\pi/2)\)
\end{itemize}

\subsection{Limsup and Liminf}\label{limsup-and-liminf}

\begin{itemize}
\tightlist
\item
  For bounded sequences, the set of subsequential limits is nonempty and
  bounded --- so it has a sup and inf
\item
  Definition: \(\limsup a_n\) = sup of the set of subsequential limits
\item
  Definition: \(\liminf a_n\) = inf of the set of subsequential limits
\item
  Convention for unbounded sequences:

  \begin{itemize}
  \tightlist
  \item
    If \((a_n)\) has a subsequence \(\to +\infty\), we set
    \(\limsup a_n = +\infty\)
  \item
    If \((a_n)\) has a subsequence \(\to -\infty\), we set
    \(\liminf a_n = -\infty\)
  \item
    If \((a_n)\) is unbounded above but has no finite subsequential
    limits, \(\limsup a_n = +\infty\)
  \item
    Similarly for unbounded below
  \end{itemize}
\item
  Remark: with this convention, \emph{every} sequence has a limsup and
  liminf (possibly \(\pm\infty\))
\item
  For bounded sequences, \(\liminf a_n \leq \limsup a_n\), with equality
  iff exactly one subsequential limit
\item
  Examples:

  \begin{itemize}
  \tightlist
  \item
    \(a_n = (-1)^n\): limsup = 1, liminf = -1
  \item
    \(a_n = 1/n\): limsup = liminf = 0
  \item
    \(a_n = (-1)^n/n\): limsup = liminf = 0
  \item
    \(a_n = n\): limsup = \(+\infty\), liminf = \(+\infty\)
  \item
    \(a_n = (-1)^n n\): limsup = \(+\infty\), liminf = \(-\infty\)
  \end{itemize}
\item
  ✎ Inline: Find limsup and liminf for \(a_n = (1 + (-1)^n)/2\)
\end{itemize}

\subsection{The Tail Characterization}\label{the-tail-characterization}

\begin{itemize}
\tightlist
\item
  Definition: the \(n\)th tail of \((a_k)\) is the sequence
  \((a_n, a_{n+1}, a_{n+2}, \ldots)\)
\item
  For bounded \((a_n)\), define
  \(\overline{s}_n = \sup\{a_k : k \geq n\}\) and
  \(\underline{s}_n = \inf\{a_k : k \geq n\}\)
\item
  Key observation: \(\overline{s}_n\) is decreasing (sup over smaller
  set), \(\underline{s}_n\) is increasing
\item
  Both are bounded, so by MCT they converge
\item
  Alternative definition:
  \(\limsup a_n = \lim_{n \to \infty} \overline{s}_n = \lim_{n \to \infty} \sup\{a_k : k \geq n\}\)
\item
  Alternative definition:
  \(\liminf a_n = \lim_{n \to \infty} \underline{s}_n = \lim_{n \to \infty} \inf\{a_k : k \geq n\}\)
\item
  Convention for unbounded sequences:

  \begin{itemize}
  \tightlist
  \item
    If \((a_n)\) is unbounded above, then \(\overline{s}_n = +\infty\)
    for all \(n\)
  \item
    We set \(\lim_{n \to \infty} (+\infty) = +\infty\) by convention
  \item
    So \(\limsup a_n = +\infty\), consistent with the subsequential
    limit definition
  \item
    Similarly for unbounded below
  \end{itemize}
\item
  Theorem: The two definitions of limsup agree (prove in text)
\item
  Proof sketch:

  \begin{itemize}
  \tightlist
  \item
    Let \(L^* = \lim \overline{s}_n\) (tail definition) and
    \(S^* = \sup\{\text{subsequential limits}\}\)
  \item
    Show \(L^* \leq S^*\): construct subsequence converging to \(L^*\)
    by picking \(a_{n_k}\) close to \(\overline{s}_{n_k}\)
  \item
    Show \(S^* \leq L^*\): any subsequential limit
    \(\leq \overline{s}_n\) for all \(n\), so \(\leq L^*\)
  \end{itemize}
\item
  Theorem: \(\limsup a_n\) is achieved --- some subsequence converges to
  it (prove in text)
\item
  Proof: the construction above produces a subsequence
  \(\to L^* = \limsup a_n\)
\item
  Analogous results for liminf (exercises)
\item
  ✎ Inline: For \(a_n = (-1)^n(1 - 1/n)\), compute \(\overline{s}_n\)
  and \(\underline{s}_n\) explicitly
\end{itemize}

\subsection{Convergence Criterion}\label{convergence-criterion}

\begin{itemize}
\tightlist
\item
  Theorem: \(a_n \to L\) if and only if
  \(\limsup a_n = \liminf a_n = L\) (prove in text)
\item
  Proof (⟹): if \(a_n \to L\), then \(L\) is the only subsequential
  limit, so limsup = liminf = \(L\)
\item
  Proof (⟸): if limsup = liminf = \(L\), then \(\overline{s}_n \to L\)
  and \(\underline{s}_n \to L\); since
  \(\underline{s}_n \leq a_n \leq \overline{s}_n\), squeeze gives
  \(a_n \to L\)
\item
  Remark: this gives a way to test convergence without knowing the limit
  in advance
\item
  Example: \(a_n = \frac{n + (-1)^n}{n}\): compute limsup and liminf,
  conclude convergence
\item
  Remark: the gap \(\limsup a_n - \liminf a_n\) measures ``oscillation''
  of the sequence --- leads to Cauchy criterion (§3.3)
\end{itemize}

\subsection{Guided Exercises}\label{guided-exercises-7}

None for this section.

\subsection{Exercises}\label{exercises-12}

\emph{Subsequential limits} - Find all subsequential limits of
\(a_n = \cos(n\pi/3)\) - Find all subsequential limits of
\(a_n = n/(n+1) \cdot (-1)^n\) - Prove: if \(a_n \to L\), then \(L\) is
the only subsequential limit - Prove: if \((a_n)\) is bounded with
unique subsequential limit \(L\), then \(a_n \to L\) - Give an example
of an unbounded sequence with exactly one subsequential limit

\emph{Limsup and liminf} - Find limsup and liminf:
\(a_n = (-1)^n + 1/n\) - Find limsup and liminf: \(a_n = \sin(n\pi/4)\)
- Find limsup and liminf: \(a_n = (2 + (-1)^n)^{1/n}\) - Prove:
\(\liminf a_n \leq \limsup a_n\) for any bounded sequence - Prove:
\(\liminf(-a_n) = -\limsup a_n\)

\emph{Tail characterization} - Prove:
\(\overline{s}_n = \sup\{a_k : k \geq n\}\) is a decreasing sequence -
Prove: the two definitions of liminf agree - Prove: \(\liminf a_n\) is
achieved by some subsequence

\emph{Convergence criterion} - Use limsup/liminf to prove
\(a_n = \frac{n + (-1)^n \sqrt{n}}{n}\) converges - Use limsup/liminf to
prove \(a_n = \frac{(-1)^n}{n} + \frac{1}{n^2}\) converges - Prove: if
\(\limsup a_n < \infty\) and \(\liminf a_n > -\infty\), then \((a_n)\)
is bounded

\emph{Algebra of limsup/liminf} - Prove:
\(\limsup(a_n + b_n) \leq \limsup a_n + \limsup b_n\) - Give an example
where the inequality is strict - ★ Prove: if \(a_n \to L\), then
\(\limsup(a_n + b_n) = L + \limsup b_n\)

\subsection{Dependencies}\label{dependencies-12}

\textbf{Requires}: §3.1 (subsequences, BW), §2.4 (MCT for tail
sequences)

\textbf{Used in}: §3.3 (Cauchy criterion derives from limsup = liminf)

\section{The Cauchy Criterion}\label{the-cauchy-criterion}

\subsection{Narrative}\label{narrative-13}

\begin{itemize}
\tightlist
\item
  Start from §3.2: convergence ⟺ limsup = liminf
\item
  What does this mean? The gap between sup(tail) and inf(tail) shrinks
  to zero
\item
  ``Tail oscillation'' → 0 is equivalent to convergence
\item
  Unpack: terms in the tail are all getting close to \emph{each other}
\item
  This is a condition stated entirely in terms of the sequence --- no
  limit mentioned!
\item
  The Cauchy condition: for any ε, eventually all pairs of terms are
  within ε
\item
  Key theorem: Cauchy ⟺ convergent (in ℝ)
\item
  Significance: intrinsic test, generalizes to metric spaces
\end{itemize}

\subsection{Content}\label{content-13}

\emph{(Opening motivation --- not a numbered subsection)} - Recall from
§3.2: \(a_n \to L\) iff \(\limsup a_n = \liminf a_n = L\) -
Equivalently: the gap \(\overline{s}_n - \underline{s}_n \to 0\), where
\(\overline{s}_n = \sup\{a_k : k \geq n\}\),
\(\underline{s}_n = \inf\{a_k : k \geq n\}\) - Definition: tail
oscillation \(\omega_n = \overline{s}_n - \underline{s}_n\) - Theorem:
\((a_n)\) converges iff \(\omega_n \to 0\) - Question: what does
\(\omega_n < \varepsilon\) actually say? - Answer: all terms in the
\(n\)th tail lie within an interval of length \(< \varepsilon\) -
Equivalently: for all \(j, k \geq n\), we have
\(|a_j - a_k| < \varepsilon\) - Key insight: this condition involves
only the sequence itself --- no limit appears! - We can test whether a
sequence ``deserves to converge'' without knowing where it converges to

\subsection{The Cauchy Condition}\label{the-cauchy-condition}

\begin{itemize}
\tightlist
\item
  Definition: \((a_n)\) is a Cauchy sequence if for every
  \(\varepsilon > 0\), there exists \(N\) such that for all
  \(m, n \geq N\), \(|a_m - a_n| < \varepsilon\)
\item
  Remark: ``terms eventually stay close to each other''
\item
  Remark: this is exactly the condition \(\omega_n \to 0\) unpacked
\item
  Examples:

  \begin{itemize}
  \tightlist
  \item
    \(a_n = 1/n\) is Cauchy: for \(m, n \geq N\),
    \(|1/m - 1/n| \leq 1/N < \varepsilon\) if \(N > 1/\varepsilon\)
  \item
    \(a_n = (-1)^n\) is not Cauchy: \(|a_{2k} - a_{2k+1}| = 2\) for all
    \(k\)
  \item
    Partial sums \(s_n = \sum_{k=1}^n 1/2^k\) are Cauchy (preview of
    series convergence)
  \end{itemize}
\item
  ✎ Inline: Verify directly that \(a_n = n/(n+1)\) is Cauchy
\end{itemize}

\subsection{Cauchy Implies Convergent}\label{cauchy-implies-convergent}

\begin{itemize}
\tightlist
\item
  Theorem: In \(\RR\), a sequence converges if and only if it is Cauchy
  (prove in text)
\item
  Proof (⟹): If \(a_n \to L\), then for \(\varepsilon/2\), get \(N\) so
  \(|a_n - L| < \varepsilon/2\) for \(n \geq N\). Then for
  \(m, n \geq N\):
  \(|a_m - a_n| \leq |a_m - L| + |L - a_n| < \varepsilon\)
\item
  Proof (⟸): This is the substantial direction.

  \begin{itemize}
  \tightlist
  \item
    Step 1: Cauchy ⟹ bounded (prove: fix \(\varepsilon = 1\), get \(N\);
    all terms within 1 of \(a_N\); bound by
    \(\max\{|a_1|, \ldots, |a_{N-1}|, |a_N| + 1\}\))
  \item
    Step 2: Bounded ⟹ convergent subsequence (by BW)
  \item
    Step 3: Cauchy + convergent subsequence ⟹ whole sequence converges
  \item
    For step 3: let \(a_{n_k} \to L\). Given \(\varepsilon\), get \(K\)
    so \(|a_{n_k} - L| < \varepsilon/2\) for \(k \geq K\), and \(N\) so
    \(|a_m - a_n| < \varepsilon/2\) for \(m, n \geq N\). For
    \(n \geq \max\{N, n_K\}\), pick \(k\) with \(n_k \geq N\) and
    \(k \geq K\); then
    \(|a_n - L| \leq |a_n - a_{n_k}| + |a_{n_k} - L| < \varepsilon\)
  \end{itemize}
\item
  ✎ Inline: Why does the proof of (⟸) fail in \(\QQ\)? Which step uses
  completeness?
\item
  Remark: the proof uses BW (hence nested intervals, hence
  completeness); Cauchy criterion fails in \(\QQ\)
\end{itemize}

\subsection{Significance}\label{significance}

\begin{itemize}
\tightlist
\item
  The Cauchy criterion is \emph{intrinsic}: it tests convergence using
  only the sequence, not the limit
\item
  Practical: useful when we don't know what the limit should be
\item
  Theoretical: the condition is \emph{metric} --- it uses only
  \(|a_m - a_n|\), not the order structure of \(\RR\)
\item
  This means Cauchy completeness generalizes: in any metric space, we
  can define Cauchy sequences and ask if they converge
\item
  Definition: a metric space is \emph{complete} if every Cauchy sequence
  converges
\item
  \(\RR\) is complete; \(\QQ\) is not
\item
  Historical note: Cantor constructed \(\RR\) from \(\QQ\) by
  ``completing'' it --- adding limits of Cauchy sequences of rationals
\item
  Remark: all five faces of completeness (LUB, MCT, Nested Intervals,
  BW, Cauchy) are equivalent in \(\RR\) --- see §3.4
\end{itemize}

\subsection{Guided Exercises}\label{guided-exercises-8}

None for this section.

\subsection{Exercises}\label{exercises-13}

\emph{Cauchy sequences} - Prove directly that \(a_n = 1 - 1/n\) is
Cauchy - Prove directly that \(a_n = \sum_{k=1}^{n} 1/k!\) is Cauchy -
Prove that \(a_n = \sum_{k=1}^{n} 1/k\) is not Cauchy (hint: show
\(|a_{2n} - a_n| \geq 1/2\); cf.~§2.4) - Prove: if \((a_n)\) is Cauchy,
then so is \((|a_n|)\) - Prove: if \((a_n)\) and \((b_n)\) are Cauchy,
then so is \((a_n + b_n)\)

\emph{The equivalence} - Prove: every convergent sequence is Cauchy -
Prove: every Cauchy sequence is bounded - Prove: if \((a_n)\) is Cauchy
and has a subsequence converging to \(L\), then \(a_n \to L\)

\emph{Cauchy in ℚ} - Construct a Cauchy sequence in \(\QQ\) that does
not converge in \(\QQ\) (hint: approximate \(\sqrt{2}\)) - The sequence
\(a_1 = 1\), \(a_{n+1} = (a_n + 2/a_n)/2\) is Cauchy in \(\QQ\). What
does it converge to in \(\RR\)?

\emph{Applications} - Prove: if \(|a_{n+1} - a_n| \leq r^n\) for some
\(0 < r < 1\), then \((a_n)\) is Cauchy (hence convergent) - Use the
previous result to prove the Babylonian sequence converges - ★ Prove: if
\(|a_{n+1} - a_n| \leq c/n^2\) for some constant \(c\), then \((a_n)\)
is Cauchy

\subsection{Dependencies}\label{dependencies-13}

\textbf{Requires}: §3.1 (BW), §3.2 (limsup/liminf, tail
characterization)

\textbf{Used in}: §3.4 (equivalence of completeness notions), later
chapters (series convergence tests)

\section{\texorpdfstring{§3.4 The Many Faces of Completeness
\emph{(Terminal)}}{§3.4 The Many Faces of Completeness (Terminal)}}\label{the-many-faces-of-completeness-terminal}

\subsection{Narrative}\label{narrative-14}

\begin{itemize}
\tightlist
\item
  Completeness has appeared throughout our work in different guises
\item
  LUB was our axiom; MCT, Nested Intervals, BW, Cauchy all followed
\item
  These aren't five different properties --- they're five faces of one
  property
\item
  Theorem: for an ordered field, all five are equivalent
\item
  We've already proved most of the cycle; one direction remains
\item
  Significance: Cauchy completeness generalizes beyond ordered fields to
  metric spaces
\end{itemize}

\subsection{Content}\label{content-14}

\subsection{The Five
Characterizations}\label{the-five-characterizations}

\begin{itemize}
\tightlist
\item
  We've encountered completeness in five forms:

  \begin{enumerate}
  \def\labelenumi{\arabic{enumi}.}
  \tightlist
  \item
    \textbf{LUB (Least Upper Bound)}: Every nonempty set bounded above
    has a supremum (§1.4 --- our axiom)
  \item
    \textbf{MCT (Monotone Convergence)}: Every monotone bounded sequence
    converges (§2.4)
  \item
    \textbf{Nested Intervals}: Nested closed intervals have nonempty
    intersection (§1.4)
  \item
    \textbf{Bolzano-Weierstrass}: Every bounded sequence has a
    convergent subsequence (§3.1)
  \item
    \textbf{Cauchy Completeness}: Every Cauchy sequence converges (§3.3)
  \end{enumerate}
\item
  Theorem: For an ordered field, these five properties are equivalent
\item
  Remark: \(\QQ\) satisfies none of them; \(\RR\) satisfies all of them
\end{itemize}

\subsection{The Equivalence Cycle}\label{the-equivalence-cycle}

\begin{itemize}
\tightlist
\item
  We prove the equivalences by tracing a cycle: LUB ⟹ MCT ⟹ NI ⟹ BW ⟹
  Cauchy ⟹ LUB
\item
  Most implications we've already proved; we recall them and complete
  the cycle
\end{itemize}

\emph{LUB ⟹ MCT} - This was the content of §2.4 - Given increasing
bounded \((a_n)\), let \(L = \sup\{a_n\}\); by
\(\varepsilon\)-characterization, \(a_n \to L\)

\emph{MCT ⟹ Nested Intervals} - Let
\([a_1, b_1] \supseteq [a_2, b_2] \supseteq \cdots\) be nested closed
intervals - \((a_n)\) is increasing and bounded above by \(b_1\); by
MCT, \(a_n \to a\) for some \(a\) - \((b_n)\) is decreasing and bounded
below by \(a_1\); by MCT, \(b_n \to b\) for some \(b\) - Since
\(a_n \leq b_n\) for all \(n\), we have \(a \leq b\) - Any point in
\([a, b]\) lies in all \([a_n, b_n]\)

\emph{Nested Intervals ⟹ BW} - This was the content of §3.1 - Bisection
produces nested intervals; their intersection gives the subsequential
limit

\emph{BW ⟹ Cauchy Completeness} - This was the content of §3.3 - Cauchy
⟹ bounded ⟹ (BW) convergent subsequence ⟹ whole sequence converges

\emph{Cauchy ⟹ LUB} (prove in text) - This closes the cycle --- the one
direction we haven't yet proved - Let \(S\) be nonempty and bounded
above; we construct \(\sup S\) - Idea: build a Cauchy sequence
approaching the sup from below - Construction: - Let \(b_0\) be an upper
bound for \(S\), and pick \(a_0 \in S\) - Bisect \([a_0, b_0]\): if the
midpoint is an upper bound for \(S\), set \$b\_1 = \$ midpoint,
\(a_1 = a_0\); otherwise pick \(a_1 \in S\) with \$a\_1 \textgreater{}
\$ midpoint, set \(b_1 = b_0\) - Continue: get nested intervals
\([a_n, b_n]\) with \(b_n\) always an upper bound, \(a_n \in S\), and
\(b_n - a_n = (b_0 - a_0)/2^n\) - The sequences \((a_n)\) and \((b_n)\)
are Cauchy (distance halves each step) - By Cauchy completeness,
\(a_n \to L\) and \(b_n \to L\) for some \(L\) - Claim: \(L = \sup S\) -
\(L\) is an upper bound: if \(s \in S\) with \(s > L\), then \(s > b_n\)
for large \(n\), contradicting that \(b_n\) is an upper bound - \(L\) is
the least upper bound: any \(M < L\) satisfies \(M < a_n\) for large
\(n\), and \(a_n \in S\), so \(M\) is not an upper bound - ✎ Inline:
Where does this construction fail in \(\QQ\)?

\subsection{Significance}\label{significance-1}

\begin{itemize}
\tightlist
\item
  The five characterizations are logically equivalent --- any one
  implies all others
\item
  Completeness is a single deep property wearing different masks
\item
  Different faces are useful in different contexts:

  \begin{itemize}
  \tightlist
  \item
    LUB: foundational, good for existence proofs
  \item
    MCT: natural for recursive sequences
  \item
    Nested Intervals: geometric intuition, good for bisection arguments
  \item
    BW: essential for compactness arguments (Chapter 6+)
  \item
    Cauchy: intrinsic, generalizes to metric spaces
  \end{itemize}
\item
  The Cauchy criterion is special: it uses only distance
  \(|a_m - a_n|\), not order
\item
  Definition: A metric space is \emph{complete} if every Cauchy sequence
  converges
\item
  \(\RR\) is complete; \(\QQ\) is not; there are complete metric spaces
  with no order structure
\item
  Historical note: Cantor (1872) constructed \(\RR\) from \(\QQ\) by
  ``completing'' it --- defining reals as equivalence classes of Cauchy
  sequences of rationals
\item
  This construction proves existence of a complete ordered field; our
  uniqueness theorem (§1.5) says there's essentially only one
\end{itemize}

\subsection{Guided Exercises}\label{guided-exercises-9}

None for this section.

\subsection{Exercises}\label{exercises-14}

\emph{The cycle} - Fill in details: prove MCT ⟹ Nested Intervals
carefully - Fill in details: prove that the sequences \((a_n)\),
\((b_n)\) in the Cauchy ⟹ LUB proof are Cauchy - Prove directly: LUB ⟹
Nested Intervals (without going through MCT) - Prove directly: BW ⟹ MCT
(without going around the cycle)

\emph{Failures in \(\QQ\)} - Give an explicit example showing LUB fails
in \(\QQ\) - Give an explicit example showing MCT fails in \(\QQ\) -
Give an explicit example showing Nested Intervals fails in \(\QQ\) -
Give an explicit example showing BW fails in \(\QQ\) - Give an explicit
example showing Cauchy completeness fails in \(\QQ\)

\emph{Alternative equivalences} - ★ Prove directly: Nested Intervals ⟹
LUB (hint: bisection to approximate sup) - ★ Prove directly: Cauchy ⟹ BW
(hint: a Cauchy sequence is bounded; what else do you need?)

\subsection{Dependencies}\label{dependencies-14}

\textbf{Requires}: §1.4 (LUB, Nested Intervals), §2.4 (MCT), §3.1 (BW),
§3.3 (Cauchy)

\textbf{Used in}: None --- this is a capstone section

\chapter{Proving Convergence}\label{proving-convergence}

\textbf{The Message}

\begin{itemize}
\tightlist
\item
  We have the theory (MCT, Cauchy, completeness) --- now how do you
  \emph{use} it?
\item
  Showing a specific series converges or an iteration reaches its target
  requires craft
\item
  This chapter develops that craft, illustrated by historically
  important examples
\end{itemize}

\textbf{The Techniques}

\begin{itemize}
\tightlist
\item
  \emph{Finding hidden monotonicity}: oscillating sequences may
  decompose into monotone pieces converging to the same limit
\item
  \emph{Comparison}: bound your unknown process by one you understand
  (e.g., geometric series)
\item
  \emph{Contraction}: functions that squeeze points together force
  iteration to a fixed point, with explicit rates
\end{itemize}

\textbf{The Examples}

\begin{itemize}
\tightlist
\item
  Alternating series (Leibniz): values are \(\ln 2\) and \(\pi/4\)
\item
  Continued fractions (Lagrange): complete the \(\sqrt{2}\) story,
  reveal \(\phi\) as ``most irrational''
\item
  Products (Wallis, Viète): factors approaching 1, comparison gives
  convergence
\item
  Archimedes' \(\pi\): inscribed/circumscribed polygons squeeze the
  circle
\end{itemize}

\textbf{The Takeaway}

\begin{itemize}
\tightlist
\item
  The examples are beautiful in themselves
\item
  But the techniques matter more --- workhorses of analysis, appearing
  throughout the course
\end{itemize}

\section{Finding Monotonicity}\label{finding-monotonicity}

\subsection{Narrative}\label{narrative-15}

\begin{itemize}
\tightlist
\item
  Not all sequences are obviously monotone, but some decompose into
  monotone pieces
\item
  The strategy: find monotone subsequences → MCT gives convergence →
  compute limits → if equal, invoke §3.1 (union theorem)
\item
  This technique applies to alternating series, products, continued
  fractions, and more
\item
  Three examples: alternating series, the √2 continued fraction, the
  Wallis product
\end{itemize}

\subsection{Content}\label{content-15}

\emph{(Opening --- unnumbered)} - Many interesting sequences aren't
monotone: partial sums of alternating series zig-zag, convergents of
continued fractions alternate above and below - But often we can
decompose into monotone pieces - The strategy: 1. Identify subsequences
(e.g., even-indexed and odd-indexed terms) 2. Show each subsequence is
monotone and bounded → MCT gives convergence 3. Compute or compare their
limits 4. If the limits agree, invoke §3.1 (union theorem): the full
sequence converges - This chapter applies this technique to various
``infinite processes''

\subsection{Alternating Series}\label{alternating-series}

\begin{itemize}
\tightlist
\item
  Definition: an alternating series has the form
  \(\sum_{n=1}^{\infty} (-1)^{n+1} b_n\) or
  \(\sum_{n=1}^{\infty} (-1)^n b_n\) with \(b_n > 0\)
\item
  Theorem (Alternating Series Test): If \((b_n)\) is decreasing and
  \(b_n \to 0\), then \(\sum (-1)^{n+1} b_n\) converges
\item
  Proof (in text):

  \begin{itemize}
  \tightlist
  \item
    Let \(S_n = \sum_{k=1}^{n} (-1)^{k+1} b_k\) be partial sums
  \item
    Even partial sums:
    \(S_{2n} = (b_1 - b_2) + (b_3 - b_4) + \cdots + (b_{2n-1} - b_{2n})\)
  \item
    Each parenthesis \(\geq 0\) since \((b_n)\) decreasing, so
    \((S_{2n})\) is increasing
  \item
    Also
    \(S_{2n} = b_1 - (b_2 - b_3) - (b_4 - b_5) - \cdots - (b_{2n-2} - b_{2n-1}) - b_{2n} \leq b_1\)
  \item
    So \((S_{2n})\) is increasing and bounded above → MCT →
    \(S_{2n} \to L\) for some \(L\)
  \item
    Odd partial sums: \(S_{2n+1} = S_{2n} + b_{2n+1}\)
  \item
    Since \(b_n \to 0\), we have \(S_{2n+1} \to L\) as well
  \item
    \((S_{2n})\) increasing, \((S_{2n+1})\) decreasing (check!), and
    \(S_{2n+1} - S_{2n} = b_{2n+1} \to 0\)
  \item
    By §3.1 (union theorem), \(S_n \to L\)
  \end{itemize}
\item
  ✎ Inline: Verify \((S_{2n+1})\) is decreasing
\item
  Example: Alternating harmonic series
  \(\sum (-1)^{n+1}/n = 1 - 1/2 + 1/3 - 1/4 + \cdots\) converges

  \begin{itemize}
  \tightlist
  \item
    The value is \(\ln 2\) --- we'll prove this later when we develop
    logarithms and Abel's theorem
  \item
    Remark: the harmonic series \(\sum 1/n\) diverges, but the
    alternating harmonic series converges --- signs matter!
  \end{itemize}
\item
  Example: Leibniz series
  \(\sum (-1)^{n}/(2n+1) = 1 - 1/3 + 1/5 - 1/7 + \cdots\) converges

  \begin{itemize}
  \tightlist
  \item
    The value is \(\pi/4\) --- a remarkable connection between an
    infinite series and the circle constant
  \item
    We'll prove this later when we develop inverse trigonometric
    functions
  \item
    For now: AST applies since \(1/(2n+1)\) is decreasing and \(\to 0\)
  \end{itemize}
\end{itemize}

\subsection{Continued Fraction for √2}\label{continued-fraction-for-2}

\begin{itemize}
\tightlist
\item
  Goal: express \(\sqrt{2}\) as a ``continued fraction'' and prove
  convergence
\item
  Start with the identity: \(\sqrt{2} = 1 + (\sqrt{2} - 1)\)
\item
  Rewrite: \(\sqrt{2} - 1 = \frac{1}{\sqrt{2} + 1}\), so
  \(\sqrt{2} = 1 + \frac{1}{1 + \sqrt{2}}\)
\item
  Now \(1 + \sqrt{2} = 2 + (\sqrt{2} - 1) = 2 + \frac{1}{1 + \sqrt{2}}\)
\item
  Substituting repeatedly:
  \[\sqrt{2} = 1 + \cfrac{1}{2 + \cfrac{1}{2 + \cfrac{1}{2 + \cdots}}}\]
\item
  Notation: write this as \([1; 2, 2, 2, \ldots]\) or \([1; \bar{2}]\)
\item
  Define convergents by truncating:

  \begin{itemize}
  \tightlist
  \item
    \(c_1 = 1\)
  \item
    \(c_2 = 1 + \frac{1}{2} = \frac{3}{2}\)
  \item
    \(c_3 = 1 + \frac{1}{2 + \frac{1}{2}} = 1 + \frac{1}{5/2} = \frac{7}{5}\)
  \item
    \(c_4 = 1 + \frac{1}{2 + \frac{1}{2 + \frac{1}{2}}} = \frac{17}{12}\)
  \end{itemize}
\item
  Pattern:
  \(c_1 = 1, c_2 = 3/2, c_3 = 7/5, c_4 = 17/12, c_5 = 41/29, \ldots\)
\item
  ✎ Inline: Compute \(c_5\) and verify it equals \(41/29\)
\item
  Key observation (proved in guided exercise below):

  \begin{itemize}
  \tightlist
  \item
    Even convergents \(c_2, c_4, c_6, \ldots\) are increasing
  \item
    Odd convergents \(c_1, c_3, c_5, \ldots\) are decreasing
  \item
    Every even convergent \textless{} \(\sqrt{2}\) \textless{} every odd
    convergent
  \item
    The gap \(c_{2n+1} - c_{2n} \to 0\)
  \end{itemize}
\item
  By MCT: both subsequences converge; by §3.1: \(c_n \to \sqrt{2}\)
\item
  Remark: this is our first continued fraction --- §4.4 develops the
  full theory
\item
  Remark: the convergents \(p_n/q_n = 1/1, 3/2, 7/5, 17/12, \ldots\)
  satisfy \(p_n^2 - 2q_n^2 = \pm 1\) --- the Pell equation from §2.1!
\end{itemize}

\subsection{The Wallis Product}\label{the-wallis-product}

\begin{itemize}
\tightlist
\item
  The Wallis product:
  \[W = \frac{2}{1} \cdot \frac{2}{3} \cdot \frac{4}{3} \cdot \frac{4}{5} \cdot \frac{6}{5} \cdot \frac{6}{7} \cdots = \prod_{n=1}^{\infty} \frac{4n^2}{4n^2 - 1}\]
\item
  Claim: \(W = \pi/2\) (proved in §4.5; here we prove convergence)
\item
  Define partial products:
  \(W_n = \prod_{k=1}^{n} \frac{4k^2}{4k^2 - 1}\)
\item
  Rewrite:
  \(\frac{4k^2}{4k^2 - 1} = \frac{4k^2}{(2k-1)(2k+1)} = \frac{2k}{2k-1} \cdot \frac{2k}{2k+1}\)
\item
  So
  \(W_n = \frac{2}{1} \cdot \frac{2}{3} \cdot \frac{4}{3} \cdot \frac{4}{5} \cdots \frac{2n}{2n-1} \cdot \frac{2n}{2n+1}\)
\item
  Key insight: separate into even and odd partial products

  \begin{itemize}
  \tightlist
  \item
    After \(2n\) factors: involves fractions up to \(\frac{2n}{2n+1}\)
  \item
    After \(2n-1\) factors: involves fractions up to \(\frac{2n}{2n-1}\)
  \end{itemize}
\item
  Define:

  \begin{itemize}
  \tightlist
  \item
    \(P_n = \frac{2}{1} \cdot \frac{4}{3} \cdot \frac{6}{5} \cdots \frac{2n}{2n-1}\)
    (products of ``big'' fractions)
  \item
    \(Q_n = \frac{2}{3} \cdot \frac{4}{5} \cdot \frac{6}{7} \cdots \frac{2n}{2n+1}\)
    (products of ``small'' fractions)
  \end{itemize}
\item
  Then \(W_{2n-1} = P_n \cdot Q_{n-1}\) and \(W_{2n} = P_n \cdot Q_n\)
\item
  Since each factor of \(Q_n\) is \(< 1\): \(Q_n\) is decreasing
\item
  Since each factor of \(P_n\) is \(> 1\): \(P_n\) is increasing
\item
  One can show: \(W_{2n}\) is increasing, \(W_{2n-1}\) is decreasing
  (verify in exercises)
\item
  Bounds: \(W_{2n} < W_{2n+1}\) (even \textless{} odd) and both are
  bounded
\item
  By MCT: both subsequences converge; by §3.1: \(W_n \to W\) for some
  \(W > 0\)
\item
  The value \(W = \pi/2\) is proved in §4.5 using Archimedes' polygon
  method
\item
  ✎ Inline: Compute \(W_1, W_2, W_3, W_4\) and verify \(W_2 < W_4\) and
  \(W_3 > W_1\)
\end{itemize}

\subsection{Guided Exercises}\label{guided-exercises-10}

\subsection{Convergence of the √2 Continued
Fraction}\label{convergence-of-the-2-continued-fraction}

Let \(c_n\) denote the \(n\)th convergent of \([1; \bar{2}]\), so
\(c_1 = 1\), \(c_2 = 3/2\), \(c_3 = 7/5\), \(c_4 = 17/12\), etc.

\begin{enumerate}
\def\labelenumi{(\alph{enumi})}
\item
  Prove that \(c_{2n} < \sqrt{2}\) for all \(n \geq 1\). (Hint: compute
  \(c_{2n}^2\) and compare to 2.)
\item
  Prove that \(c_{2n-1} > \sqrt{2}\) for all \(n \geq 1\).
\item
  Prove that the even convergents are increasing: \(c_{2n} < c_{2n+2}\).
\item
  Prove that the odd convergents are decreasing:
  \(c_{2n+1} < c_{2n-1}\).
\item
  Conclude by MCT that \((c_{2n})\) and \((c_{2n-1})\) both converge.
\item
  Prove that the gap \(c_{2n-1} - c_{2n} \to 0\). (Hint: express the gap
  in terms of the denominators \(q_n\).)
\item
  Apply §3.1 to conclude \(c_n \to L\) for some \(L\).
\item
  Prove that \(L = \sqrt{2}\).
\end{enumerate}

\subsection{Exercises}\label{exercises-15}

\emph{Alternating series} - Prove: \(\sum (-1)^{n+1}/n^2\) converges -
Prove: \(\sum (-1)^{n+1}/\sqrt{n}\) converges - Give an example where
\(b_n \to 0\) but \((b_n)\) is not decreasing, and \(\sum (-1)^n b_n\)
still converges - Give an example where \(b_n \to 0\) but \((b_n)\) is
not decreasing, and \(\sum (-1)^n b_n\) diverges - Prove: if
\(\sum (-1)^{n+1} b_n\) converges by the alternating series test, then
\(|S - S_n| \leq b_{n+1}\) (error bound)

\emph{√2 continued fraction} - Verify that \((7/5)^2 = 49/25\) and
compare to 2; verify \((17/12)^2 = 289/144\) and compare to 2 - Prove:
\(p_n^2 - 2q_n^2 = (-1)^n\) where \(p_n/q_n\) is the \(n\)th convergent
(connects to Pell equation) - Compute \(c_6\) and \(c_7\) - Show that
\(|c_n - \sqrt{2}| < 1/q_n^2\) where \(c_n = p_n/q_n\)

\emph{Wallis product} - Compute \(W_5\) and \(W_6\) as fractions -
Prove: \(W_{2n} < W_{2n+2}\) (even partial products increasing) - Prove:
\(W_{2n+1} > W_{2n+3}\) (odd partial products decreasing) - Prove:
\(W_{2n} < W_{2n+1}\) for all \(n\) - ★ Prove: \(W_n\) is bounded above
(hint: compare to an integral, if you know calculus, or use a clever
grouping)

\subsection{Dependencies}\label{dependencies-15}

\textbf{Requires}: §2.4 (MCT), §3.1 (union theorem for subsequences)

\textbf{Used in}: §4.4 (continued fractions), §4.5 (Wallis product value
= π/2)

\section{Comparison}\label{comparison}

\subsection{Narrative}\label{narrative-16}

\begin{itemize}
\tightlist
\item
  Series and products often can't be evaluated directly, but we can
  determine convergence by comparison
\item
  Comparison test: bound terms by a known convergent series
\item
  Limit comparison: if terms are asymptotically proportional, series
  behave the same
\item
  Absolute convergence: when \(\sum |a_n|\) converges, so does
  \(\sum a_n\)
\item
  Root test: compare with geometric series via \(n\)th roots
\item
  Comparison for products: if factors approach 1 fast enough, product
  converges
\item
  Applications: Viète product, Euler's product for sine
\end{itemize}

\subsection{Content}\label{content-16}

\subsection{Comparison Tests}\label{comparison-tests}

\begin{itemize}
\tightlist
\item
  Theorem (Comparison Test): If \(0 \leq a_n \leq b_n\) for all \(n\)
  and \(\sum b_n\) converges, then \(\sum a_n\) converges
\item
  Proof: partial sums of \(\sum a_n\) are increasing and bounded above
  by \(\sum b_n\); apply MCT
\item
  Contrapositive: if \(\sum a_n\) diverges, so does \(\sum b_n\)
\item
  Example: \(\sum 1/(n^2 + n)\) converges (compare:
  \(1/(n^2+n) < 1/n^2\), and \(\sum 1/n^2\) converges)
\item
  ✎ Inline: Prove \(\sum 1/(n^3 + 1)\) converges
\item
  Theorem (Limit Comparison): If \(a_n, b_n > 0\) and \(a_n/b_n \to L\)
  with \(0 < L < \infty\), then \(\sum a_n\) and \(\sum b_n\) either
  both converge or both diverge
\item
  Proof sketch: for large \(n\), \((L/2) b_n < a_n < (2L) b_n\); apply
  comparison both ways
\item
  Example: \(\sum n/(n^3 + 1)\) converges (compare to \(\sum 1/n^2\):
  ratio \(\to 1\))
\item
  {[}NOTE: Would like more interesting/historically important examples
  here{]}
\end{itemize}

\subsection{Absolute Convergence}\label{absolute-convergence}

\begin{itemize}
\tightlist
\item
  Definition: \(\sum a_n\) converges absolutely if \(\sum |a_n|\)
  converges
\item
  Definition: \(\sum a_n\) converges conditionally if it converges but
  not absolutely
\item
  Theorem: Absolute convergence implies convergence (prove in text)
\item
  Proof via Cauchy criterion: if \(\sum |a_n|\) converges, partial sums
  of \(|a_n|\) are Cauchy; then
  \(|\sum_{k=m}^{n} a_k| \leq \sum_{k=m}^{n} |a_k|\), so partial sums of
  \(a_n\) are also Cauchy
\item
  Example: \(\sum (-1)^n/n^2\) converges absolutely (since
  \(\sum 1/n^2\) converges)
\item
  Example: \(\sum (-1)^n/n\) converges conditionally (converges by AST,
  but \(\sum 1/n\) diverges)
\item
  Remark: absolute convergence is ``stronger'' --- the series converges
  regardless of signs
\item
  Remark: conditionally convergent series can be rearranged to sum to
  any value (Riemann rearrangement theorem --- stated, not proved)
\end{itemize}

\subsection{Comparison with Geometric
Series}\label{comparison-with-geometric-series}

\begin{itemize}
\tightlist
\item
  Key fact: geometric series \(\sum r^n\) converges iff \(|r| < 1\)
\item
  Idea: if terms of a series are eventually bounded by \(Cr^n\),
  comparison with geometric series gives convergence
\item
  Theorem (Root Test): Let \(\alpha = \limsup |a_n|^{1/n}\).

  \begin{itemize}
  \tightlist
  \item
    If \(\alpha < 1\), then \(\sum a_n\) converges absolutely
  \item
    If \(\alpha > 1\), then \(\sum a_n\) diverges
  \item
    If \(\alpha = 1\), the test is inconclusive
  \end{itemize}
\item
  Proof of \(\alpha < 1\) case (in text):

  \begin{itemize}
  \tightlist
  \item
    Pick \(r\) with \(\alpha < r < 1\)
  \item
    By definition of limsup, \(|a_n|^{1/n} < r\) for all sufficiently
    large \(n\)
  \item
    So \(|a_n| < r^n\) for large \(n\)
  \item
    By comparison with \(\sum r^n\), \(\sum |a_n|\) converges
  \end{itemize}
\item
  Proof of \(\alpha > 1\) case: \(|a_n|^{1/n} > 1\) infinitely often, so
  \(|a_n| > 1\) infinitely often, so \(a_n \not\to 0\)
\item
  Example: \(\sum 1/2^n\) converges (\(|a_n|^{1/n} = 1/2 < 1\))
\item
  Example: \(\sum n/2^n\) converges
  (\((n/2^n)^{1/n} = n^{1/n}/2 \to 1/2 < 1\))
\item
  Example: \(\sum 1/n^2\) --- root test inconclusive
  (\((1/n^2)^{1/n} = 1/n^{2/n} \to 1\))
\item
  ✎ Inline: Use the root test on \(\sum n^{10}/2^n\)
\item
  Remark: the ratio test (see guided exercise) is often easier to apply
  in practice
\end{itemize}

\subsection{Comparison for Products}\label{comparison-for-products}

\begin{itemize}
\tightlist
\item
  Infinite products: when does \(\prod_{n=1}^{\infty} a_n\) converge?
\item
  Write \(a_n = 1 + u_n\) where \(u_n = a_n - 1\)
\item
  If factors approach 1 ``fast enough,'' the product converges
\item
  Lemma: For \(t_1, \ldots, t_m \geq 0\) with \(S = \sum t_k < 1\):
  \(\prod(1-t_k) \geq 1 - S\) (prove in text by induction)
\item
  Lemma: For \(t_k \geq 0\) with \(S = \sum t_k < 1\):
  \(\prod(1+t_k) \leq \frac{1}{1-S}\)

  \begin{itemize}
  \tightlist
  \item
    Proof: use \(1 + t \leq \frac{1}{1-t}\) and the previous lemma
  \end{itemize}
\item
  Corollary: For \(t_k \geq 0\) with \(S = \sum t_k < 1\):
  \(0 \leq \prod(1+t_k) - 1 \leq \frac{S}{1-S}\)
\item
  Lemma: \(|\prod(1+u_k) - 1| \leq \prod(1+|u_k|) - 1\) (prove in text
  by induction)
\item
  Theorem: If \(\sum |a_n - 1| < \infty\) and \(a_n > 0\), then
  \(\prod a_n\) converges to a positive limit
\item
  Proof (in text):

  \begin{itemize}
  \tightlist
  \item
    Let \(P_N = \prod_{n=1}^{N} a_n\) and \(u_n = a_n - 1\)
  \item
    For \(M > N\):
    \(\left|\frac{P_M}{P_N} - 1\right| \leq \prod_{n=N+1}^{M}(1+|u_n|) - 1 \leq \frac{S_{N,M}}{1-S_{N,M}}\)
  \item
    where \(S_{N,M} = \sum_{n=N+1}^{M} |u_n|\)
  \item
    Since \(\sum |u_n| < \infty\), tails \(\to 0\), so \((P_N)\) is
    Cauchy, hence convergent
  \item
    Nonzero: for large \(N\), \(|P_M/P_N - 1| < 1/2\), so
    \(|P_M| > |P_N|/2 > 0\)
  \end{itemize}
\item
  \textbf{Viète Product}: \(\prod_{n=1}^{\infty} \frac{r_n}{2}\) where
  \(r_1 = \sqrt{2}\), \(r_{n+1} = \sqrt{2 + r_n}\)

  \begin{itemize}
  \tightlist
  \item
    Need to show \(\sum |r_n/2 - 1| < \infty\)
  \item
    Let \(e_n = 2 - r_n\). Then \(|r_n/2 - 1| = |r_n - 2|/2 = e_n/2\)
  \item
    Claim: \(e_{n+1} \leq e_n/2\) (geometric decay)
  \item
    Proof: \(e_{n+1} = 2 - \sqrt{4 - e_n} = 2(1 - \sqrt{1 - e_n/4})\)
  \item
    Use \(\sqrt{1-x} \geq 1 - x\) for \(0 \leq x \leq 1\) (proof: square
    both sides)
  \item
    So \(\sqrt{1 - e_n/4} \geq 1 - e_n/4\), giving
    \(e_{n+1} \leq 2 \cdot e_n/4 = e_n/2\)
  \item
    Thus \(e_n \leq (2-\sqrt{2}) \cdot 2^{-(n-1)}\), so
    \(\sum e_n < \infty\)
  \item
    By the theorem, \(\prod r_n/2\) converges to a positive limit
  \item
    The value is \(2/\pi\) --- proved in §4.5 using Archimedes' polygon
    method
  \end{itemize}
\item
  \textbf{Euler's Product}: \(\prod_{n=1}^{\infty}(1 - x^2/n^2)\)
  converges for any fixed \(x\)

  \begin{itemize}
  \tightlist
  \item
    Here \(u_n = -x^2/n^2\), so \(|u_n| = x^2/n^2\)
  \item
    We know \(\sum 1/n^2\) converges (§2.4), so
    \(\sum x^2/n^2 = x^2 \sum 1/n^2 < \infty\)
  \item
    By the theorem, the product converges
  \item
    Euler showed this equals \(\sin(\pi x)/(\pi x)\) --- a profound
    connection we'll explore later
  \end{itemize}
\item
  Remark: Wallis product \(\prod \frac{4n^2}{4n^2-1}\) also follows
  (exercise: \(|a_n - 1| = 1/(4n^2-1) \sim 1/n^2\))
\end{itemize}

\subsection{Guided Exercises}\label{guided-exercises-11}

\subsection{The Ratio Test}\label{the-ratio-test}

\begin{enumerate}
\def\labelenumi{(\alph{enumi})}
\item
  Suppose \(|a_{n+1}/a_n| \leq r < 1\) for all \(n \geq N\). Prove that
  \(|a_n| \leq |a_N| \cdot r^{n-N}\) for \(n \geq N\).
\item
  Conclude that \(\sum |a_n|\) converges by comparison with a geometric
  series.
\item
  State the ratio test: if \(\limsup |a_{n+1}/a_n| < 1\), then
  \(\sum a_n\) converges absolutely; if \(\liminf |a_{n+1}/a_n| > 1\),
  then \(\sum a_n\) diverges.
\item
  Apply the ratio test to prove \(\sum n!/n^n\) converges.
\item
  Apply the ratio test to prove \(\sum n^{100}/2^n\) converges.
\item
  Show the ratio test is inconclusive for \(\sum 1/n^2\) (even though it
  converges).
\end{enumerate}

\subsection{Exercises}\label{exercises-16}

\emph{Comparison} - Prove \(\sum 1/(n^2 - 1)\) converges (for
\(n \geq 2\)) - Prove \(\sum 1/\sqrt{n}\) diverges - Prove
\(\sum 1/(n \ln n)\) diverges for \(n \geq 2\) (hint: compare to
integral or use Cauchy condensation) - Use limit comparison to determine
convergence of \(\sum (n+1)/(n^3 + n)\) - Use limit comparison to
determine convergence of \(\sum (\sqrt{n+1} - \sqrt{n})\)

\emph{Absolute and conditional convergence} - Determine whether
\(\sum (-1)^n / \sqrt{n}\) converges absolutely, conditionally, or
diverges - Determine whether \(\sum (-1)^n n / (n^2 + 1)\) converges
absolutely, conditionally, or diverges - Prove: if \(\sum a_n\)
converges absolutely, then \(|\sum a_n| \leq \sum |a_n|\)

\emph{Root and ratio tests} - Use the root test:
\(\sum (n/(n+1))^{n^2}\) - Use the root test: \(\sum 1/(\ln n)^n\) for
\(n \geq 2\) - Use the ratio test: \(\sum 2^n / n!\) - Use the ratio
test: \(\sum n! / n^n\) - Give an example where the ratio test is
inconclusive but the root test gives convergence

\emph{Products} - Verify that Wallis product satisfies
\(|a_n - 1| = 1/(4n^2-1)\) and conclude convergence - Prove
\(\prod_{n=2}^{\infty}(1 - 1/n^2)\) converges and find its value (hint:
partial fractions, telescope) - Prove
\(\prod_{n=1}^{\infty}(1 + 1/n^2)\) converges - ★ Prove: if
\(\sum |u_n|\) converges and \(u_n > -1\) for all \(n\), then
\(\prod(1+u_n) = 0\) iff some \(u_n = -1\)

\subsection{Dependencies}\label{dependencies-16}

\textbf{Requires}: §2.4 (MCT, series via bounded partial sums), §3.2
(limsup), §3.3 (Cauchy criterion)

\textbf{Used in}: §4.5 (Viète product value = 2/π), later chapters
(series manipulations)

\section{Contraction Mappings}\label{contraction-mappings}

\subsection{Narrative}\label{narrative-17}

\begin{itemize}
\tightlist
\item
  We've seen recursive sequences \(a_{n+1} = f(a_n)\) that converge:
  nested radicals, Babylonian method
\item
  What's the common structure? The function \(f\) ``squeezes'' points
  closer together
\item
  Contraction mappings: \(|f(x) - f(y)| \leq c|x-y|\) for some \(c < 1\)
\item
  Fixed point theorem: contractions have unique fixed points, and all
  orbits converge
\item
  Convergence is geometric --- we get explicit error bounds
\item
  Application: the golden ratio as a fixed point
\end{itemize}

\subsection{Content}\label{content-17}

\emph{(Opening --- unnumbered)} - Recall two recursive sequences from
Chapter 2: - Nested radicals: \(a_1 = 1\), \(a_{n+1} = \sqrt{2 + a_n}\)
converges to 2 - Babylonian: \(x_1 = 2\), \(x_{n+1} = (x_n + 2/x_n)/2\)
converges to \(\sqrt{2}\) - Both have the form \(a_{n+1} = f(a_n)\) for
some function \(f\) - Both converge to a fixed point: a value \(L\) with
\(f(L) = L\) - What property of \(f\) guarantees convergence? - Key
observation: both functions ``squeeze'' points closer together - If
\(f\) shrinks distances by a factor \(c < 1\), iterates must converge

\subsection{Contraction Mappings}\label{contraction-mappings-1}

\begin{itemize}
\tightlist
\item
  Definition: Let \(I = [a, b]\) be a closed interval. A function
  \(f: I \to I\) is a \emph{contraction} (on \(I\)) if there exists
  \(c \in (0, 1)\) such that for all \(x, y \in I\):
  \[|f(x) - f(y)| \leq c|x - y|\]
\item
  The constant \(c\) is called the \emph{contraction constant}
\item
  Remark: we require \(f: I \to I\) (the interval maps to itself) so
  that iterates stay in \(I\)
\item
  Example: \(f(x) = x/2\) on \([0, 1]\) is a contraction with
  \(c = 1/2\)
\item
  Example: \(f(x) = x^2\) on \([0, 1/2]\) is a contraction (check:
  \(|x^2 - y^2| = |x+y||x-y| \leq 1 \cdot |x-y|\)\ldots{} actually need
  \(|x+y| < 1\), so on \([0, 1/2]\) we get \(c = 1\), not a contraction.
  On \([0, 1/3]\): \(|x+y| \leq 2/3\), so \(c = 2/3\) works.)
\item
  ✎ Inline: Verify \(f(x) = x/3 + 1\) is a contraction on \([1, 2]\) and
  find \(c\)
\end{itemize}

\subsection{Fixed Point Theorem}\label{fixed-point-theorem}

\begin{itemize}
\tightlist
\item
  Definition: \(L\) is a \emph{fixed point} of \(f\) if \(f(L) = L\)
\item
  Theorem (Contraction Mapping Theorem): Let \(f: [a,b] \to [a,b]\) be a
  contraction with constant \(c < 1\). Then:

  \begin{enumerate}
  \def\labelenumi{\arabic{enumi}.}
  \tightlist
  \item
    \(f\) has exactly one fixed point \(L \in [a, b]\)
  \item
    For any starting point \(a_1 \in [a, b]\), the sequence
    \(a_{n+1} = f(a_n)\) converges to \(L\)
  \end{enumerate}
\item
  Proof (in text):

  \begin{itemize}
  \tightlist
  \item
    \emph{Step 1: consecutive terms get closer.}

    \begin{itemize}
    \tightlist
    \item
      \(|a_{n+1} - a_n| = |f(a_n) - f(a_{n-1})| \leq c|a_n - a_{n-1}|\)
    \item
      By induction: \(|a_{n+1} - a_n| \leq c^{n-1}|a_2 - a_1|\)
    \end{itemize}
  \item
    \emph{Step 2: \((a_n)\) is Cauchy.}

    \begin{itemize}
    \tightlist
    \item
      For \(m > n\):
      \(|a_m - a_n| \leq |a_m - a_{m-1}| + \cdots + |a_{n+1} - a_n|\)
    \item
      \(\leq (c^{m-2} + \cdots + c^{n-1})|a_2 - a_1|\)
    \item
      \(\leq c^{n-1} \cdot \frac{1}{1-c} |a_2 - a_1|\)
    \item
      This \(\to 0\) as \(n \to \infty\), so \((a_n)\) is Cauchy
    \end{itemize}
  \item
    \emph{Step 3: limit is a fixed point.}

    \begin{itemize}
    \tightlist
    \item
      By completeness, \(a_n \to L\) for some \(L \in [a, b]\)
    \item
      We have \(a_{n+1} = f(a_n)\)
    \item
      Taking limits: \(L = \lim a_{n+1} = \lim f(a_n)\)
    \item
      Claim: \(\lim f(a_n) = f(L)\)
    \item
      Proof: \(|f(a_n) - f(L)| \leq c|a_n - L| \to 0\)
    \item
      So \(L = f(L)\)
    \end{itemize}
  \item
    \emph{Step 4: uniqueness.}

    \begin{itemize}
    \tightlist
    \item
      Suppose \(L\) and \(M\) are both fixed points
    \item
      Then \(|L - M| = |f(L) - f(M)| \leq c|L - M|\)
    \item
      Since \(c < 1\), this forces \(|L - M| = 0\), so \(L = M\)
    \end{itemize}
  \end{itemize}
\item
  Remark: the proof uses completeness (Cauchy sequences converge) ---
  this fails in \(\QQ\)
\item
  ✎ Inline: Where exactly does the proof fail if \(c = 1\)?
\end{itemize}

\subsection{Convergence Rates}\label{convergence-rates}

\begin{itemize}
\tightlist
\item
  The proof gives an explicit error bound
\item
  Theorem: With notation as above,
  \(|a_n - L| \leq \frac{c^{n-1}}{1-c}|a_2 - a_1|\)
\item
  Proof:
  \(|a_n - L| = \lim_{m \to \infty} |a_n - a_m| \leq \frac{c^{n-1}}{1-c}|a_2 - a_1|\)
\item
  Interpretation: error decreases geometrically with ratio \(c\)
\item
  Smaller \(c\) means faster convergence
\item
  Revisit nested radicals: \(f(x) = \sqrt{2 + x}\) on \([1, 2]\)

  \begin{itemize}
  \tightlist
  \item
    Check contraction:
    \(|f(x) - f(y)| = |\sqrt{2+x} - \sqrt{2+y}| = \frac{|x - y|}{\sqrt{2+x} + \sqrt{2+y}}\)
  \item
    On \([1, 2]\): denominator \(\geq 2\sqrt{3} > 3\), so \(c < 1/3\)
  \item
    Geometric convergence with \(c \leq 1/3\): error decreases by factor
    of 3 each step
  \end{itemize}
\item
  Revisit Babylonian: \(f(x) = (x + 2/x)/2\) on \([1, 2]\)

  \begin{itemize}
  \tightlist
  \item
    Check contraction:
    \(|f(x) - f(y)| = \frac{1}{2}|1 - 2/(xy)||x - y|\)
  \item
    On \([\sqrt{2}, 2]\): \(xy \geq 2\), so
    \(|1 - 2/(xy)| \leq 1 - 1 = 0\)\ldots{}
  \item
    Actually:
    \(f(x) - f(y) = \frac{1}{2}(x - y) - \frac{1}{xy}(x - y) = \frac{1}{2}(1 - 2/(xy))(x-y)\)
  \item
    Hmm, need \(xy > 2\) for this to be positive. On \([\sqrt{2}, 2]\):
    \(xy \geq 2\)
  \item
    \(|f(x) - f(y)| = \frac{1}{2}|1 - 2/(xy)||x - y| \leq \frac{1}{2}|x - y|\)
    when \(xy \geq 2\)
  \item
    So \(c = 1/2\) on \([\sqrt{2}, 2]\)
  \item
    But actually Babylonian converges much faster --- it has
    ``quadratic'' convergence (error squares each step), which our
    linear contraction bound doesn't capture
  \end{itemize}
\item
  Remark: some iterations converge faster than the contraction bound
  suggests; the bound is a worst-case guarantee
\end{itemize}

\subsection{The Golden Ratio}\label{the-golden-ratio}

\begin{itemize}
\tightlist
\item
  Consider the iteration \(x_{n+1} = 1 + 1/x_n\)
\item
  Fixed point equation: \(x = 1 + 1/x\), i.e., \(x^2 - x - 1 = 0\)
\item
  Solutions: \(x = (1 \pm \sqrt{5})/2\)
\item
  The positive solution is the golden ratio
  \(\phi = (1 + \sqrt{5})/2 \approx 1.618\)
\item
  Claim: \(f(x) = 1 + 1/x\) is a contraction on \([1.5, 2]\)
\item
  Proof:

  \begin{itemize}
  \tightlist
  \item
    First check \(f: [1.5, 2] \to [1.5, 2]\):
    \(f(1.5) = 1 + 2/3 = 5/3 \approx 1.67 \in [1.5, 2]\) ✓;
    \(f(2) = 1.5 \in [1.5, 2]\) ✓
  \item
    Contraction: \(|f(x) - f(y)| = |1/x - 1/y| = \frac{|x - y|}{xy}\)
  \item
    On \([1.5, 2]\): \(xy \geq 1.5^2 = 2.25\), so
    \(|f(x) - f(y)| \leq \frac{1}{2.25}|x - y| < 0.45|x - y|\)
  \item
    So \(c = 4/9 < 1/2\)
  \end{itemize}
\item
  By the fixed point theorem: starting from any \(x_1 \in [1.5, 2]\),
  \(x_n \to \phi\)
\item
  Example: \(x_1 = 2\), \(x_2 = 1.5\), \(x_3 = 5/3 \approx 1.667\),
  \(x_4 = 8/5 = 1.6\), \(x_5 = 13/8 = 1.625\), \ldots{}
\item
  Pattern: \(x_n = F_{n+1}/F_n\) where \(F_n\) are Fibonacci numbers!
  (exercise)
\item
  ✎ Inline: Verify \(x_6 = 21/13\) and compare to \(\phi\)
\item
  Remark: this iteration is secretly computing continued fraction
  convergents for \(\phi = [1; 1, 1, 1, \ldots]\) --- see §4.4
\item
  Remark: the golden ratio appears throughout mathematics: Fibonacci
  sequence, regular pentagons, phyllotaxis in plants
\end{itemize}

\subsection{Guided Exercises}\label{guided-exercises-12}

\subsection{Newton's Method for Cube
Roots}\label{newtons-method-for-cube-roots}

The iteration
\(x_{n+1} = \frac{1}{3}\left(2x_n + \frac{A}{x_n^2}\right)\) converges
to \(\sqrt[3]{A}\).

\begin{enumerate}
\def\labelenumi{(\alph{enumi})}
\item
  Verify that \(L = \sqrt[3]{A}\) is a fixed point of
  \(f(x) = \frac{1}{3}(2x + A/x^2)\).
\item
  Let \(A = 2\). Show that \(f\) maps \([1, 2]\) to itself.
\item
  Compute \(|f(x) - f(y)|\) in terms of \(|x - y|\) and show \(f\) is a
  contraction on \([1, 2]\). Find an explicit contraction constant
  \(c\).
\item
  Starting from \(x_1 = 1\), compute \(x_2, x_3, x_4\) (as fractions or
  decimals).
\item
  Use the error bound to estimate how many iterations are needed to get
  within \(10^{-6}\) of \(\sqrt[3]{2}\).
\item
  Generalize: for which intervals \([a, b]\) and which values of \(A\)
  is this a contraction?
\end{enumerate}

\subsection{Exercises}\label{exercises-17}

\emph{Basic contractions} - Prove \(f(x) = x/2 + 1\) is a contraction on
\([0, 3]\) and find its fixed point - Prove \(f(x) = \cos(x)\) is a
contraction on \([0, 1]\) (for those who know calculus: use
\(|\cos'(x)| = |\sin(x)| \leq \sin(1) < 1\)) - Give an example of
\(f: [0,1] \to [0,1]\) that is NOT a contraction but still has a unique
fixed point

\emph{Verifying contractions} - Verify nested radicals
\(f(x) = \sqrt{2 + x}\) is a contraction on \([1, 2]\) and find \(c\) -
Verify \(f(x) = \sqrt{1 + x}\) is a contraction on \([1, 2]\) and find
its fixed point (the golden ratio!) - Show \(f(x) = 2 + 1/x\) is a
contraction on \([2, 3]\) and find its fixed point (relates to
\(\sqrt{2}\) CF)

\emph{Golden ratio and Fibonacci} - Prove: if \(x_1 = 1\) and
\(x_{n+1} = 1 + 1/x_n\), then \(x_n = F_{n+1}/F_n\) where \(F_n\) is the
\(n\)th Fibonacci number - Conclude: \(F_{n+1}/F_n \to \phi\) - Prove:
\(F_n = (\phi^n - \psi^n)/\sqrt{5}\) where \(\psi = (1 - \sqrt{5})/2\)
(Binet's formula)

\emph{Convergence rates} - For nested radicals with \(c = 1/3\), how
many iterations to get within \(10^{-10}\) of the limit? - The iteration
\(x_{n+1} = x_n/2 + 1/x_n\) converges to \(\sqrt{2}\). Verify it's a
contraction and compare its rate to Babylonian.

\emph{Non-examples} - Show \(f(x) = x^2\) on \([0, 1]\) is NOT a
contraction (find points where ratio \(= 1\)) - Show that
\(f(x) = 2x(1-x)\) on \([0, 1]\) has fixed points at \(0\) and \(1/2\),
but is not a contraction on \([0, 1]\)

\subsection{Dependencies}\label{dependencies-17}

\textbf{Requires}: §3.3 (Cauchy criterion), §2.4 (Babylonian, nested
radicals examples)

\textbf{Used in}: §4.4 (continued fractions as iterations), later
chapters (fixed point methods)

\section{§4.4 Continued Fractions}\label{continued-fractions}

\subsection{Narrative}\label{narrative-18}

\begin{itemize}
\tightlist
\item
  We've seen \(\sqrt{2} = [1; \overline{2}]\) in §4.1 and
  \(\phi = [\overline{1}]\) from contractions in §4.3
\item
  What if CFs had been developed as primary notation instead of
  decimals?
\item
  Every real has unique CF; every CF converges; rationals ↔ finite;
  quadratic irrationals ↔ periodic
\item
  \(\phi\) is ``most irrational'' --- hardest to approximate by
  rationals
\end{itemize}

\subsection{Content}\label{content-18}

\emph{(Opening --- unnumbered)} - Decimals depend on base 10; in base 2
or 7, digits change - CFs like \(\sqrt{2} = [1; \overline{2}]\) are
base-independent - This section: every real has unique CF, every CF
converges, CFs reveal structure decimals hide

\subsection{Continued Fraction
Notation}\label{continued-fraction-notation}

\begin{itemize}
\tightlist
\item
  Definition:
  \([a_0; a_1, a_2, \ldots] = a_0 + \cfrac{1}{a_1 + \cfrac{1}{a_2 + \cfrac{1}{a_3 + \cdots}}}\)
  with \(a_0 \in \ZZ\), \(a_1, a_2, \ldots \in \NN\)
\item
  The \(a_i\) are \emph{partial quotients}; finite CFs
  \([a_0; a_1, \ldots, a_n]\) are rationals
\item
  \emph{Convergents}: truncations \(c_n = [a_0; a_1, \ldots, a_n]\),
  written \(c_n = p_n/q_n\) in lowest terms
\item
  Recurrence: \(p_{-1} = 1, p_0 = a_0, p_n = a_n p_{n-1} + p_{n-2}\);
  same for \(q_n\) with \(q_{-1} = 0, q_0 = 1\)
\item
  Example (\(\sqrt{2} = [1; 2, 2, 2, \ldots]\)): \textbar{} n \textbar{}
  \(a_n\) \textbar{} \(p_n\) \textbar{} \(q_n\) \textbar{} \(p_n/q_n\)
  \textbar{}
  \textbar---\textbar-------\textbar-------\textbar-------\textbar-----------\textbar{}
  \textbar{} 0 \textbar{} 1 \textbar{} 1 \textbar{} 1 \textbar{} 1
  \textbar{} \textbar{} 1 \textbar{} 2 \textbar{} 3 \textbar{} 2
  \textbar{} 1.5 \textbar{} \textbar{} 2 \textbar{} 2 \textbar{} 7
  \textbar{} 5 \textbar{} 1.4 \textbar{} \textbar{} 3 \textbar{} 2
  \textbar{} 17 \textbar{} 12 \textbar{} 1.4167\ldots{} \textbar{}
  \textbar{} 4 \textbar{} 2 \textbar{} 41 \textbar{} 29 \textbar{}
  1.4138\ldots{} \textbar{}
\item
  ✎ Inline: Verify \(p_4 = 2 \cdot 17 + 7 = 41\),
  \(q_4 = 2 \cdot 12 + 5 = 29\)
\item
  Theorem (Key Identity): \(p_n q_{n-1} - p_{n-1} q_n = (-1)^{n-1}\)
  (prove by induction)
\item
  Corollary: \(\gcd(p_n, q_n) = 1\); also
  \(\frac{p_n}{q_n} - \frac{p_{n-1}}{q_{n-1}} = \frac{(-1)^{n-1}}{q_n q_{n-1}}\)
\end{itemize}

\subsection{Every Continued Fraction
Converges}\label{every-continued-fraction-converges}

\begin{itemize}
\tightlist
\item
  Theorem: Every infinite CF converges to a real number
\item
  Proof: From corollary, \(c_n - c_{n-1} = (-1)^{n-1}/(q_n q_{n-1})\)
  alternates in sign

  \begin{itemize}
  \tightlist
  \item
    Even convergents increasing, odd decreasing (check via the identity)
  \item
    Every even \textless{} every odd
  \item
    Gap: \(|c_n - c_{n-1}| = 1/(q_n q_{n-1}) \to 0\) since
    \(q_n \geq F_n \to \infty\)
  \item
    By §4.1 (even-odd technique): \((c_n)\) converges
  \end{itemize}
\item
  Remark: This generalizes the \(\sqrt{2}\) argument from §4.1
\item
  ✎ Inline: Prove \(q_n \geq F_n\) by induction
\end{itemize}

\subsection{Every Real Number Has a Continued
Fraction}\label{every-real-number-has-a-continued-fraction}

\begin{itemize}
\tightlist
\item
  Greedy algorithm: \(a_0 = \lfloor x \rfloor\); if \(x \neq a_0\), set
  \(x_1 = 1/(x - a_0)\), \(a_1 = \lfloor x_1 \rfloor\); repeat
\item
  Terminates iff \(x \in \QQ\); convergents converge to \(x\);
  representation unique (convention: don't end in 1)
\item
  Example: \(17/7\): \(a_0 = 2\), \(x_1 = 7/3\), \(a_1 = 2\),
  \(x_2 = 3\), stop. So \(17/7 = [2; 2, 3]\)
\item
  ✎ Inline: Verify \([2; 2, 3] = 2 + 1/(2 + 1/3) = 17/7\)
\item
  Remark: Rationals have finite CFs --- base-independent!
\end{itemize}

\subsection{Quadratic Irrationals}\label{quadratic-irrationals}

\begin{itemize}
\tightlist
\item
  Definition: irrational root of \(ax^2 + bx + c = 0\) with
  \(a, b, c \in \ZZ\)
\item
  Theorem (Lagrange): eventually periodic CF ↔ quadratic irrational
\item
  Eventually periodic:
  \([a_0; a_1, \ldots, a_k, \overline{b_1, \ldots, b_m}]\); purely
  periodic: \([\overline{a_1, \ldots, a_m}]\)
\item
  Examples: \(\sqrt{2} = [1; \overline{2}]\),
  \(\sqrt{3} = [1; \overline{1, 2}]\), \(\sqrt{5} = [2; \overline{4}]\),
  \(\phi = [\overline{1}]\)
\item
  Connection to contractions: periodic tail \([\overline{a}]\) satisfies
  \(t = a + 1/t\), fixed point of \(f(x) = a + 1/x\) (§4.3)
\item
  Proof sketch (periodic ⟹ quadratic): tail satisfies quadratic from
  fixed point equation
\item
  Theorem (Pell): for non-square \(N\), convergents of \(\sqrt{N}\)
  satisfy \(p_n^2 - Nq_n^2 = (-1)^{n+1}\) at complete periods
\item
  The \(\sqrt{2}\) story complete: Ch1 (sup), Ch2 (Babylonian, Pell),
  §4.1 (CF convergence), §4.3 (contraction), here (CF =
  \([1; \overline{2}]\), convergents = Pell solutions)
\end{itemize}

\subsection{Beyond Quadratics}\label{beyond-quadratics}

\begin{itemize}
\tightlist
\item
  Cubics: NO known patterns
  (\(\sqrt[3]{2} = [1; 3, 1, 5, 1, 1, 4, ...]\) --- appears random)
\item
  \(e = [2; 1, 2, 1, 1, 4, 1, 1, 6, 1, 1, 8, \ldots]\) --- beautiful
  pattern \([2; \overline{1, 2k, 1}]\); proof in Ch5
\item
  \(\pi = [3; 7, 15, 1, 292, ...]\) --- no pattern; large 292 explains
  why \(355/113\) is exceptional
\item
  Remark: CFs reveal structure decimals hide --- \(\sqrt{2}\) periodic,
  \(e\) patterned, \(\pi\) structureless
\end{itemize}

\subsection{Approximation and the Golden
Ratio}\label{approximation-and-the-golden-ratio}

\begin{itemize}
\tightlist
\item
  Theorem: \(|x - p_n/q_n| < 1/(q_n q_{n+1}) < 1/q_n^2\)
\item
  Theorem: convergents are best approximations among fractions with
  denominator \(\leq q_n\)
\item
  From recurrence \(q_n = a_n q_{n-1} + q_{n-2}\): larger partial
  quotients ⟹ faster growth ⟹ better approximations sooner
\item
  \(\phi = [\overline{1}]\): all 1s (smallest possible), so
  \(q_n = F_{n+1}\) grows slowest
\item
  Theorem (Hurwitz): for any irrational \(x\), infinitely many \(p/q\)
  with \(|x - p/q| < 1/(\sqrt{5}q^2)\); for \(\phi\), constant
  \(1/\sqrt{5}\) is sharp
\item
  \(\phi\) is the ``most irrational'' number!
\item
  ✎ Inline: Verify convergents of \(\phi\) are \(F_{n+1}/F_n\)
\end{itemize}

\subsection{Guided Exercise}\label{guided-exercise-5}

\subsection{\texorpdfstring{The Continued Fraction for
\(\sqrt{3}\)}{The Continued Fraction for \textbackslash sqrt\{3\}}}\label{the-continued-fraction-for-sqrt3}

\begin{enumerate}
\def\labelenumi{(\alph{enumi})}
\item
  \(a_0 = \lfloor \sqrt{3} \rfloor = 1\). Show
  \(x_1 = 1/(\sqrt{3} - 1) = (\sqrt{3} + 1)/2\) by rationalizing.
\item
  \(a_1 = 1\). Show \(x_2 = 1/(x_1 - 1) = \sqrt{3} + 1\).
\item
  \(a_2 = 2\). Show \(x_3 = 1/(x_2 - 2) = x_1\), so
  \(\sqrt{3} = [1; \overline{1, 2}]\) with period 2.
\item
  Compute convergents for \(n = 0, 1, 2, 3, 4\) using the recurrence.
\item
  Compute \(p_n^2 - 3q_n^2\) for each. What pattern?
\item
  At period ends (\(n = 2, 4, ...\)), verify \(p_n^2 - 3q_n^2 = 1\)
  (Pell solutions).
\item
  Fundamental solution: \((p_2, q_2) = (2, 1)\). Verify
  \(2^2 - 3 \cdot 1^2 = 1\).
\end{enumerate}

\subsection{Exercises}\label{exercises-18}

\emph{Computing CFs} - Find CF for \(7/5\) and \(22/7\) using greedy
algorithm - Find first 5 partial quotients of \(\sqrt{7}\) and
\((1 + \sqrt{13})/2\)

\emph{Convergents and recurrence} - Prove recurrence
\(p_n = a_n p_{n-1} + p_{n-2}\) by induction - Prove key identity
\(p_n q_{n-1} - p_{n-1} q_n = (-1)^{n-1}\) by induction - Prove
\(q_n \geq F_n\); compute first 6 convergents of
\(\sqrt{5} = [2; \overline{4}]\)

\emph{Quadratic irrationals} - Verify \(\sqrt{2} = [1; \overline{2}]\)
by showing \(x_2 = x_1\) - Find CF and period for \(\sqrt{7}\) - Prove:
\([\overline{a}] = (a + \sqrt{a^2 + 4})/2\) - Silver ratio
\(\delta = 1 + \sqrt{2}\) has CF \([2; \overline{2}]\). Verify.

\emph{Pell equation} - Verify \(p_n^2 - 2q_n^2 = (-1)^{n+1}\) for first
5 convergents of \(\sqrt{2}\) - Find fundamental solution to
\(x^2 - 5y^2 = 1\) from \(\sqrt{5}\) convergents - ★ Prove: for
\(\sqrt{N} = [a_0; \overline{a_1, \ldots, a_k}]\),
\(p_{k-1}^2 - Nq_{k-1}^2 = (-1)^k\)

\emph{Approximation} - From \(\pi = [3; 7, 15, 1, 292, ...]\), find
\(355/113\). Why so good? - Prove \(|x - p_n/q_n| < 1/q_n^2\) - ★ Prove
convergents are best approximations

\emph{Golden ratio} - Prove convergents of \(\phi = [\overline{1}]\) are
\(F_{n+1}/F_n\) - Prove Cassini's identity:
\(F_{n+1}^2 - F_n F_{n+2} = (-1)^n\) - ★ Prove Hurwitz for \(\phi\)

\subsection{Dependencies}\label{dependencies-18}

\textbf{Requires}: §4.1 (even-odd convergence technique), §4.3
(contractions, golden ratio)

\textbf{Used in}: §4.5 (possible connections), later chapters
(approximation theory)

\section{Archimedes' π}\label{archimedes-ux3c0}

\subsection{Narrative}\label{narrative-19}

\begin{itemize}
\tightlist
\item
  Archimedes trapped the circle between inscribed and circumscribed
  polygons
\item
  A beautiful recurrence for angle-doubling lets us compute successive
  approximations
\item
  This gives a rigorous definition of π as a limit
\item
  The same constant governs both perimeter and area: Archimedes' great
  theorem
\end{itemize}

\subsection{Content}\label{content-19}

\emph{(Opening --- unnumbered)} - How to define π? ``Ratio of
circumference to diameter'' assumes we know what circumference means -
Archimedes' idea: trap the circle between polygons and squeeze - This
gives rigorous definition and computable approximations - Historical:
Archimedes (96-gon), Liu Hui (263 CE, 3072-gon), Zu Chongzhi (5th c.,
got 355/113), al-Kashi (1424, 16 digits)

\subsection{Archimedes' Method}\label{archimedes-method}

\begin{itemize}
\tightlist
\item
  Setup: unit circle (radius 1)
\item
  \(a_n\) = half-perimeter of regular inscribed n-gon
\item
  \(b_n\) = half-perimeter of regular circumscribed n-gon
\item
  Archimedes' Axiom: For convex curves, inscribed perimeter \textless{}
  curve length \textless{} circumscribed perimeter
\item
  Therefore: \(a_n < \text{(half-circumference)} < b_n\)
\item
  Starting values (hexagon, n = 6):

  \begin{itemize}
  \tightlist
  \item
    Inscribed regular hexagon: 6 equilateral triangles from center, each
    side = radius = 1
  \item
    Perimeter = 6, so \(a_6 = 3\)
  \item
    Circumscribed hexagon: by similar triangles,
    \(b_6 = 2\sqrt{3} \approx 3.464\)
  \end{itemize}
\item
  Already: \(3 < \text{(half-circumference)} < 3.465\)
\item
  The doubling recurrence (Archimedes): from n-gon to 2n-gon:
  \[b_{2n} = \frac{2a_n b_n}{a_n + b_n}, \qquad a_{2n} = \sqrt{a_n \cdot b_{2n}}\]
\item
  First is harmonic mean \(H(a_n, b_n)\); second is geometric mean
  \(G(a_n, b_{2n})\)
\item
  Derivation via trigonometry in guided exercise; Archimedes used pure
  geometry
\end{itemize}

\subsection{Approximating π}\label{approximating-ux3c0}

\begin{itemize}
\tightlist
\item
  Computation requires approximating square roots
\item
  Archimedes used: \(\frac{265}{153} < \sqrt{3} < \frac{1351}{780}\)
\item
  After 4 doublings (96-gon), Archimedes obtained:
  \[3\frac{10}{71} < \pi < 3\frac{1}{7}\]
\item
  That is: \(3.1408... < \pi < 3.1429...\)
\end{itemize}

\subsection{Convergence}\label{convergence-1}

\begin{itemize}
\tightlist
\item
  Theorem: Along the doubling sequence: \((a_n)\) increasing, \((b_n)\)
  decreasing, \(a_n < b_n\)
\item
  Proof: For \(a < b\): \(a < H(a,b) < b\)

  \begin{itemize}
  \tightlist
  \item
    \(b_{2n} = H(a_n, b_n) < b_n\) ✓
  \item
    \(b_{2n} > a_n\), so \(a_{2n} = \sqrt{a_n b_{2n}} > a_n\) ✓
  \end{itemize}
\item
  By MCT: \(a_n \nearrow L_a\), \(b_n \searrow L_b\), with
  \(L_a \leq L_b\)
\item
  Ratio trick: \(r_n = a_n/b_n\) satisfies
  \(r_{2n} = \sqrt{(r_n + 1)/2}\)
\item
  \((r_n)\) increasing, bounded by 1; MCT gives \(r_n \to L\)
\item
  Limit equation: \(L = \sqrt{(L+1)/2}\) gives \(L = 1\)
\item
  Conclusion: \(L_a = L_b\)
\end{itemize}

\subsection{The Definition of π}\label{the-definition-of-ux3c0}

\begin{itemize}
\tightlist
\item
  Definition: \(\pi := \lim a_n = \lim b_n\)
\item
  Corollary: \(a_n < \pi < b_n\) --- rigorous bounds
\item
  Corollary: Circumference of radius \(r\) circle is \(2\pi r\)
\item
  Areas: \(A_n\) = inscribed area, \(B_n\) = circumscribed area
\item
  Circumscribed:
  \(B_n = \frac{1}{2} \cdot \text{perimeter} \cdot 1 = b_n\) (triangles
  with height = radius)
\item
  Inscribed:
  \(A_n = \frac{1}{2} \cdot \text{perimeter} \cdot h_n = a_n h_n\) where
  \(h_n\) = height to midpoint of side
\item
  Pythagoras: \(h_n^2 + (a_n/n)^2 = 1\), so \(h_n \to 1\) as
  \(n \to \infty\)
\item
  Therefore \(A_n \to \pi\), \(B_n \to \pi\)
\item
  Theorem (Archimedes): Area =
  \(\frac{1}{2} \times \text{circumference} \times \text{radius} = \pi r^2\)
\item
  Same constant \(\pi\) for perimeter and area!
\end{itemize}

\subsection{Viète's Formula}\label{viuxe8tes-formula}

\begin{itemize}
\tightlist
\item
  From §4.2: Viète product converges; now find value
\item
  Start from square: \(a_4 = 2\sqrt{2}\)
\item
  Telescope:
  \(\pi = a_4 \cdot \prod_{j=0}^{\infty} \frac{a_{4 \cdot 2^{j+1}}}{a_{4 \cdot 2^j}}\)
\item
  Result:
  \(\frac{2}{\pi} = \frac{\sqrt{2}}{2} \cdot \frac{\sqrt{2+\sqrt{2}}}{2} \cdot \frac{\sqrt{2+\sqrt{2+\sqrt{2}}}}{2} \cdots\)
\end{itemize}

\subsection{Guided Exercise}\label{guided-exercise-6}

\subsection{Deriving the Doubling
Recurrence}\label{deriving-the-doubling-recurrence}

Let \(\theta_n\) be the central half-angle of the regular n-gon.

\begin{enumerate}
\def\labelenumi{(\alph{enumi})}
\item
  Define \(\cos\theta\), \(\sin\theta\) as coordinates on unit circle.
  Why is \(\sin^2\theta + \cos^2\theta = 1\)?
\item
  Show \(a_n = n\sin\theta_n\) and \(b_n = n\tan\theta_n\).
\item
  From angle-doubling identities
  \(\sin(2\alpha) = 2\sin\alpha\cos\alpha\),
  \(\cos(2\alpha) = 2\cos^2\alpha - 1\), derive:
  \(\tan(\theta/2) = \frac{\sin\theta}{1 + \cos\theta}\).
\item
  Using \(\theta_{2n} = \theta_n/2\), show
  \(b_{2n} = \frac{2a_n b_n}{a_n + b_n}\).
\item
  Show \(a_{2n} = \sqrt{a_n \cdot b_{2n}}\).
\item
  Compute \(b_{12}\) and \(a_{12}\) from \(a_6 = 3\),
  \(b_6 = 2\sqrt{3}\).
\end{enumerate}

\subsection{Exercises}\label{exercises-19}

\emph{Starting values} - Verify \(a_6 = 3\) geometrically - Verify
\(b_6 = 2\sqrt{3}\) - Show \(a_4 = 2\sqrt{2}\)

\emph{Computation} - Compute \(a_{24}, b_{24}\) - Verify Archimedes'
bounds for \(\sqrt{3}\) - Derive \(3\frac{10}{71} < \pi < 3\frac{1}{7}\)
from 96-gon

\emph{Viète} - Express \(a_8/a_4\) in nested radicals - ★ Complete
derivation of Viète's formula

\subsection{Dependencies}\label{dependencies-19}

\textbf{Requires}: §2.4 (MCT), §4.2 (Viète convergence)

\textbf{Used in}: Chapter 9 (rigorous trig, equivalence with integral
definition)

\emph{Note: Section deferrable; π redefined via integration in Chapter
9.}

\chapter{When Order Matters}\label{when-order-matters}

\textbf{The question} - Finite sums commute: \(a + b + c = c + a + b\)
regardless of order - Surely infinite sums work the same way? - And if
two limits exist, surely we can take them in either order? - Early
analysts assumed so, and often got correct answers - But sometimes they
got nonsense --- when does order matter for infinite operations?

\textbf{The answer, in brief} - Rearranging a series can change its sum
--- or make it diverge - Riemann: a conditionally convergent series can
be rearranged to converge to \emph{any} value - Iterated limits can
differ: \(\lim_m \lim_n a_{mn} \neq \lim_n \lim_m a_{mn}\) -
Interchanging a limit and an infinite sum can fail - Absolute
convergence and domination conditions are what make these operations
safe

\textbf{The deeper theme} - Absolute convergence provides ``room to
spare'' --- errors are controlled by convergent tails - Domination by a
convergent series gives uniform control across a parameter - These
conditions recur throughout analysis: rearrangements, double sums,
products, limit interchange - The chapter is about discipline ---
knowing exactly when infinite operations commute

\textbf{Historical note} - Cauchy identified absolute convergence as the
safety condition for rearrangements - Riemann (1854) proved the shocking
rearrangement theorem for conditional convergence - Mertens developed
the theory of multiplying series via Cauchy products - Tannery (1886)
gave precise conditions for interchanging limits and sums - By century's
end, analysts knew exactly when to trust infinite manipulations

\textbf{What's ahead} - Four variations on a single question: when can
we safely rearrange, reorder, or interchange infinite operations?

\section{§5.1 Rearrangements}\label{rearrangements}

\subsection{Narrative}\label{narrative-20}

\begin{itemize}
\tightlist
\item
  For finite sums, order doesn't matter: commutativity
\item
  For infinite sums, order can matter dramatically
\item
  Absolutely convergent series: order doesn't matter (robust)
\item
  Conditionally convergent series: can rearrange to any value (fragile)
\item
  Riemann's theorem is a shocker --- the sum is not determined by the
  terms alone
\end{itemize}

\subsection{Content}\label{content-20}

\subsection{Rearrangements}\label{rearrangements-1}

\begin{itemize}
\tightlist
\item
  For finite sums, commutativity says order doesn't matter:
  \(a_1 + a_2 + a_3 = a_3 + a_1 + a_2\)
\item
  What about infinite sums?
\item
  Definition: A \emph{rearrangement} of a series
  \(\sum_{n=1}^{\infty} a_n\) is a series
  \(\sum_{n=1}^{\infty} a_{\pi(n)}\) where
  \(\pi: \mathbb{N} \to \mathbb{N}\) is a bijection
\item
  Same terms, different order
\item
  Question: Does rearranging change the sum?
\end{itemize}

\subsection{Absolute Convergence is
Robust}\label{absolute-convergence-is-robust}

\begin{itemize}
\tightlist
\item
  Theorem: If \(\sum a_n\) converges absolutely to \(S\), then every
  rearrangement \(\sum a_{\pi(n)}\) converges absolutely to \(S\)
\item
  Proof:

  \begin{itemize}
  \tightlist
  \item
    Let \(T_N = \sum_{n=1}^{N} a_{\pi(n)}\) be the partial sums of the
    rearrangement
  \item
    Given \(\varepsilon > 0\), choose \(M\) such that
    \(\sum_{n=M+1}^{\infty} |a_n| < \varepsilon\)
  \item
    Since \(\pi\) is a bijection, there exists \(N_0\) such that
    \(\{1, 2, \ldots, M\} \subseteq \{\pi(1), \pi(2), \ldots, \pi(N_0)\}\)
  \item
    For \(N \geq N_0\): the sum \(T_N\) contains all of
    \(a_1, \ldots, a_M\), plus some terms \(a_j\) with \(j > M\)
  \item
    Let \(S_M = \sum_{n=1}^{M} a_n\). Then:
    \[|T_N - S_M| = \left| \sum_{\substack{n \leq N \\ \pi(n) > M}} a_{\pi(n)} \right| \leq \sum_{n=M+1}^{\infty} |a_n| < \varepsilon\]
  \item
    Since \(\sum a_n\) converges absolutely to \(S\), we have
    \(|S - S_M| < \varepsilon\) for \(M\) large
  \item
    Therefore \(|T_N - S| \leq |T_N - S_M| + |S_M - S| < 2\varepsilon\)
  \item
    So \(T_N \to S\)
  \end{itemize}
\item
  Absolute convergence of the rearrangement: same argument applied to
  \(\sum |a_{\pi(n)}|\)
\item
  Interpretation: for absolutely convergent series, the sum depends only
  on the terms, not their order
\end{itemize}

\subsection{Conditional Convergence is
Fragile}\label{conditional-convergence-is-fragile}

\begin{itemize}
\tightlist
\item
  For conditionally convergent series, the situation is completely
  different
\item
  Lemma: If \(\sum a_n\) converges conditionally, then
  \(\sum a_n^+ = +\infty\) and \(\sum a_n^- = -\infty\)

  \begin{itemize}
  \tightlist
  \item
    Here \(a_n^+ = \max(a_n, 0)\) and \(a_n^- = \min(a_n, 0)\)
  \end{itemize}
\item
  Proof of Lemma:

  \begin{itemize}
  \tightlist
  \item
    Note \(a_n = a_n^+ + a_n^-\) and \(|a_n| = a_n^+ - a_n^-\)
  \item
    If both \(\sum a_n^+\) and \(\sum a_n^-\) converged, then
    \(\sum |a_n| = \sum a_n^+ - \sum a_n^-\) would converge,
    contradicting conditional convergence
  \item
    If one converged and the other diverged, then
    \(\sum a_n = \sum a_n^+ + \sum a_n^-\) would diverge, contradicting
    convergence
  \item
    So both must diverge
  \item
    Since \(a_n^+ \geq 0\), we have \(\sum a_n^+ = +\infty\); since
    \(a_n^- \leq 0\), we have \(\sum a_n^- = -\infty\)
  \end{itemize}
\item
  Theorem (Riemann Rearrangement Theorem): If \(\sum a_n\) converges
  conditionally, then for any
  \(L \in \mathbb{R} \cup \{-\infty, +\infty\}\), there exists a
  rearrangement that converges to \(L\)
\item
  Proof (for finite \(L\)):

  \begin{itemize}
  \tightlist
  \item
    List the positive terms in order: \(p_1, p_2, p_3, \ldots\) (these
    are the nonzero \(a_n^+\))
  \item
    List the negative terms in order: \(q_1, q_2, q_3, \ldots\) (these
    are the nonzero \(a_n^-\), so \(q_i < 0\))
  \item
    By the lemma, \(\sum p_i = +\infty\) and \(\sum q_i = -\infty\)
  \item
    Also \(p_i \to 0\) and \(q_i \to 0\) (since \(a_n \to 0\))
  \item
    Construct the rearrangement:

    \begin{itemize}
    \tightlist
    \item
      Take positive terms \(p_1, p_2, \ldots\) until the partial sum
      first exceeds \(L\)
    \item
      Then take negative terms \(q_1, q_2, \ldots\) until the partial
      sum first drops below \(L\)
    \item
      Repeat: take positive terms until exceeding \(L\), then negative
      until dropping below \(L\)
    \end{itemize}
  \item
    This process uses every term exactly once (since both \(\sum p_i\)
    and \(\sum |q_i|\) diverge, we never run out)
  \item
    Let \(T_N\) denote the partial sums. By construction, after each
    ``switch'' the partial sum is within one term of \(L\)
  \item
    Since \(p_i \to 0\) and \(q_i \to 0\), these overshoots shrink to
    zero
  \item
    Therefore \(T_N \to L\)
  \end{itemize}
\item
  Proof (for \(L = +\infty\)):

  \begin{itemize}
  \tightlist
  \item
    Take positive terms until sum exceeds \(1\), then one negative term
  \item
    Take positive terms until sum exceeds \(2\), then one negative term
  \item
    Continue: always enough positive terms (since
    \(\sum p_i = +\infty\)), and the sum grows without bound
  \end{itemize}
\item
  Proof for \(L = -\infty\): symmetric
\item
  Remark: The sum of a conditionally convergent series depends on the
  order of terms, not just the terms themselves
\end{itemize}

\subsection{Guided Exercise}\label{guided-exercise-7}

\subsection{Rearranging the Alternating Harmonic
Series}\label{rearranging-the-alternating-harmonic-series}

The alternating harmonic series
\(\sum_{n=1}^{\infty} \frac{(-1)^{n+1}}{n} = 1 - \frac{1}{2} + \frac{1}{3} - \frac{1}{4} + \cdots\)
converges conditionally.

By Riemann's theorem, we can rearrange it to converge to any value. In
this exercise, we find explicit rearrangements converging to specific
targets.

Let
\(S = \sum_{n=1}^{\infty} \frac{(-1)^{n+1}}{n} = 1 - \frac{1}{2} + \frac{1}{3} - \frac{1}{4} + \cdots\)
(which converges by the alternating series test).

\emph{Part I: Two positives, one negative}

Consider the rearrangement:
\(1 + \frac{1}{3} - \frac{1}{2} + \frac{1}{5} + \frac{1}{7} - \frac{1}{4} + \frac{1}{9} + \frac{1}{11} - \frac{1}{6} + \cdots\)

\begin{enumerate}
\def\labelenumi{(\alph{enumi})}
\item
  Describe the pattern: the \(k\)-th block consists of which terms?
\item
  Write the partial sum \(T_N\) through \(N\) complete blocks. Express
  it in terms of partial sums of the harmonic series
  \(H_n = \sum_{j=1}^{n} 1/j\).
\item
  Let \(S_n\) denote the partial sums of the original series, and recall
  \(S_{2M} = H_{2M} - H_M\). Show that
  \(T_N = S_{4N} + \frac{1}{2}S_{2N}\). (Hint: express both the odd
  reciprocals and even reciprocals in terms of harmonic numbers
  \(H_n = \sum_{j=1}^{n} 1/j\).)
\item
  Conclude that \(T_N \to S + \frac{1}{2}S = \frac{3}{2}S\). The
  rearrangement has a different sum!
\end{enumerate}

\emph{Part II: One positive, two negatives}

Consider the rearrangement:
\(1 - \frac{1}{2} - \frac{1}{4} + \frac{1}{3} - \frac{1}{6} - \frac{1}{8} + \frac{1}{5} - \frac{1}{10} - \frac{1}{12} + \cdots\)

\begin{enumerate}
\def\labelenumi{(\alph{enumi})}
\setcounter{enumi}{4}
\item
  Describe the pattern for this rearrangement.
\item
  Show that after \(N\) complete blocks, \(T_N = \frac{1}{2}S_{2N}\).
  Conclude the rearrangement converges to \(\frac{1}{2}S\).
\end{enumerate}

\emph{Part III: Reflection}

\begin{enumerate}
\def\labelenumi{(\alph{enumi})}
\setcounter{enumi}{6}
\tightlist
\item
  Taking more positive terms relative to negative terms increased the
  sum from \(S\) to \(\frac{3}{2}S\); taking more negatives decreased it
  to \(\frac{1}{2}S\). The Riemann rearrangement theorem guarantees
  \emph{every} real number is achievable. But finding explicit formulas
  for general \(p\)-positive-\(q\)-negative rearrangements requires
  tools (logarithms) we haven't developed yet --- we'll return to this
  in Chapter 9.
\end{enumerate}

\subsection{Exercises}\label{exercises-20}

\emph{Absolute vs conditional} - The series
\(\sum_{n=1}^{\infty} \frac{(-1)^{n+1}}{n^2}\) converges. Can it be
rearranged to a different value? Justify your answer. - The series
\(\sum_{n=1}^{\infty} \frac{(-1)^{n+1}}{\sqrt{n}}\) converges. Can it be
rearranged to a different value? Justify your answer. - Determine
whether
\(\sum_{n=0}^{\infty} \frac{(-1)^n}{2n+1} = 1 - \frac{1}{3} + \frac{1}{5} - \frac{1}{7} + \cdots\)
converges absolutely or conditionally.

\emph{Constructing rearrangements} - Describe a rearrangement of
\(\sum_{n=1}^{\infty} (-1)^{n+1}/n\) that diverges to \(+\infty\). -
Describe a rearrangement of \(\sum_{n=1}^{\infty} (-1)^{n+1}/n\) that
diverges to \(-\infty\). - Describe a rearrangement of
\(\sum_{n=1}^{\infty} (-1)^{n+1}/\sqrt{n}\) that converges to \(0\).

\emph{Theoretical} - Prove: if \(\sum a_n\) converges absolutely, then
\(\sum a_n^2\) converges - Describe how to rearrange
\(\sum (-1)^n/\sqrt{n}\) to converge to \(100\) - Suppose
\(\pi: \mathbb{N} \to \mathbb{N}\) is a bijection with
\(|\pi(n) - n| \leq K\) for some constant \(K\) and all \(n\). Prove
that if \(\sum a_n\) converges (not necessarily absolutely), then
\(\sum a_{\pi(n)}\) converges to the same sum. - Suppose \(\sum a_n\)
converges and every rearrangement of \(\sum a_n\) converges to the same
value. Prove that \(\sum a_n\) converges absolutely. - ★ Prove the
Riemann rearrangement theorem in detail - ★ Prove: if \(\sum a_n\)
converges conditionally, there exists a rearrangement whose sequence of
partial sums is unbounded above and unbounded below (i.e., diverges by
oscillation).

\subsection{Dependencies}\label{dependencies-20}

\textbf{Requires}: §2.4 (series), §3.3 (Cauchy criterion), §4.1
(alternating series test), §4.2 (absolute vs conditional convergence)

\textbf{Used in}: §5.3 (double sums, reordering)

\section{Iterated Limits}\label{iterated-limits}

\subsection{Narrative}\label{narrative-21}

\begin{itemize}
\tightlist
\item
  Double sequences: two indices, two ways to take limits
\item
  Iterated limits: fix one index, take limit in the other, then take the
  remaining limit
\item
  These can differ! Order matters.
\item
  The double limit (both indices at once) is stronger
\item
  When double limit exists, it forces iterated limits to agree
\end{itemize}

\subsection{Content}\label{content-21}

\subsection{Double Sequences}\label{double-sequences}

\begin{itemize}
\tightlist
\item
  Definition: A \emph{double sequence} is a function
  \(a: \mathbb{N} \times \mathbb{N} \to \mathbb{R}\), written \(a_{mn}\)
  or \((a_{mn})_{m,n=1}^{\infty}\)
\item
  Think of it as an infinite matrix of real numbers
\item
  Examples:

  \begin{itemize}
  \tightlist
  \item
    \(a_{mn} = \frac{1}{m+n}\)
  \item
    \(a_{mn} = \frac{m}{m+n}\)
  \item
    \(a_{mn} = \frac{mn}{m^2 + n^2}\)
  \item
    \(a_{mn} = (1 + 1/m)^n\)
  \end{itemize}
\end{itemize}

\subsection{Iterated Limits}\label{iterated-limits-1}

\begin{itemize}
\item
  For a double sequence, we can take limits in two different orders
\item
  Definition: The \emph{iterated limit}
  \(\lim_{m \to \infty} \lim_{n \to \infty} a_{mn}\) means: first take
  \(\lim_{n \to \infty} a_{mn}\) (for each fixed \(m\)), then take
  \(\lim_{m \to \infty}\) of the result
\item
  Similarly for \(\lim_{n \to \infty} \lim_{m \to \infty} a_{mn}\)
\item
  These can be different!
\item
  Example: \(a_{mn} = \frac{m}{m+n}\)

  \begin{itemize}
  \tightlist
  \item
    Fix \(m\), let \(n \to \infty\): \(\frac{m}{m+n} \to 0\)
  \item
    Then \(\lim_{m \to \infty} 0 = 0\)
  \item
    Fix \(n\), let \(m \to \infty\): \(\frac{m}{m+n} \to 1\)
  \item
    Then \(\lim_{n \to \infty} 1 = 1\)
  \item
    So \(\lim_m \lim_n a_{mn} = 0\) but \(\lim_n \lim_m a_{mn} = 1\)
  \end{itemize}
\item
  Example: \(a_{mn} = \frac{m - n}{m + n}\)

  \begin{itemize}
  \tightlist
  \item
    \(\lim_m \lim_n a_{mn} = \lim_m (-1) = -1\)
  \item
    \(\lim_n \lim_m a_{mn} = \lim_n (1) = 1\)
  \end{itemize}
\item
  Example: \(a_{mn} = (1 + 1/m)^n\)

  \begin{itemize}
  \tightlist
  \item
    Fix \(m\), let \(n \to \infty\): \((1 + 1/m)^n \to +\infty\) (since
    \(1 + 1/m > 1\))
  \item
    Fix \(n\), let \(m \to \infty\): \((1 + 1/m)^n \to 1\)
  \item
    One iterated limit is \(+\infty\), the other is \(1\)
  \end{itemize}
\item
  Example: \(a_{mn} = (1 - 1/m)^n\) (for \(m \geq 2\))

  \begin{itemize}
  \tightlist
  \item
    Fix \(m\), let \(n \to \infty\): \((1 - 1/m)^n \to 0\) (since
    \(0 < 1 - 1/m < 1\))
  \item
    Fix \(n\), let \(m \to \infty\): \((1 - 1/m)^n \to 1\)
  \item
    So \(\lim_m \lim_n a_{mn} = 0\) but \(\lim_n \lim_m a_{mn} = 1\)
  \item
    Along the diagonal \(m = n\): \((1 - 1/n)^n \to 1/e\) (Chapter 2)
  \item
    Three different values from three different approaches!
  \end{itemize}
\item
  Remark: When iterated limits differ, there is no hope for a ``double
  limit'' to exist
\end{itemize}

\subsection{The Double Limit}\label{the-double-limit}

\begin{itemize}
\tightlist
\item
  Definition: We say \(\lim_{m,n \to \infty} a_{mn} = L\) if: for every
  \(\varepsilon > 0\), there exists \(N\) such that
  \(m, n > N \Rightarrow |a_{mn} - L| < \varepsilon\)
\item
  This is stronger than iterated limits: both indices must be large
  simultaneously
\item
  Theorem: If \(\lim_{m,n \to \infty} a_{mn} = L\) exists, and both
  iterated limits exist, then:
  \[\lim_{m \to \infty} \lim_{n \to \infty} a_{mn} = \lim_{n \to \infty} \lim_{m \to \infty} a_{mn} = L\]
\item
  Proof:

  \begin{itemize}
  \tightlist
  \item
    Let \(b_m = \lim_{n \to \infty} a_{mn}\) (this exists by hypothesis)
  \item
    We show \(b_m \to L\)
  \item
    Given \(\varepsilon > 0\), choose \(N\) such that
    \(m, n > N \Rightarrow |a_{mn} - L| < \varepsilon\)
  \item
    For \(m > N\): since \(|a_{mn} - L| < \varepsilon\) for all
    \(n > N\), taking \(n \to \infty\) gives
    \(|b_m - L| \leq \varepsilon\)
  \item
    So \(b_m \to L\), i.e.,
    \(\lim_{m \to \infty} \lim_{n \to \infty} a_{mn} = L\)
  \item
    Symmetric argument for the other iterated limit
  \end{itemize}
\item
  Corollary: If iterated limits exist and are different, the double
  limit does not exist
\item
  Remark: The converse is false --- iterated limits can be equal without
  the double limit existing
\item
  Example: \(a_{mn} = \frac{mn}{m^2 + n^2}\)

  \begin{itemize}
  \tightlist
  \item
    \(\lim_m \lim_n a_{mn} = \lim_m 0 = 0\)
  \item
    \(\lim_n \lim_m a_{mn} = \lim_n 0 = 0\)
  \item
    Iterated limits agree! Both are \(0\).
  \item
    But along diagonal \(m = n\): \(\frac{n^2}{2n^2} = \frac{1}{2}\)
  \item
    So the double limit cannot exist (diagonal gives different value)
  \end{itemize}
\item
  Theorem: If \(\lim_{m,n \to \infty} a_{mn} = L\), then
  \(\lim_{n \to \infty} a_{nn} = L\)
\item
  More generally: if \(f, g: \mathbb{N} \to \mathbb{N}\) satisfy
  \(f(n) \to \infty\) and \(g(n) \to \infty\), then
  \(\lim_{n \to \infty} a_{f(n), g(n)} = L\)
\item
  Proof: Given \(\varepsilon > 0\), choose \(N\) for the double limit.
  Choose \(N'\) such that \(n > N' \Rightarrow f(n) > N\) and
  \(g(n) > N\). Then
  \(n > N' \Rightarrow |a_{f(n), g(n)} - L| < \varepsilon\).
\end{itemize}

\subsection{Guided Exercise}\label{guided-exercise-8}

\subsection{Paths to Infinity}\label{paths-to-infinity}

The double sequence \(a_{mn} = \frac{mn}{m^2 + n^2}\) appeared in the
text as an example where both iterated limits exist and agree, yet the
double limit fails to exist. This exercise explores what goes wrong ---
and how rich the pathological behavior is.

\begin{enumerate}
\def\labelenumi{(\alph{enumi})}
\item
  Verify both iterated limits: compute
  \(\lim_{m \to \infty} \lim_{n \to \infty} a_{mn}\) and
  \(\lim_{n \to \infty} \lim_{m \to \infty} a_{mn}\). (Both should equal
  \(0\).)
\item
  Compute \(\lim_{n \to \infty} a_{nn}\) (the diagonal). What value do
  you get?
\item
  Compute \(\lim_{n \to \infty} a_{2n, n}\) (the path \(m = 2n\)). What
  value do you get?
\item
  More generally, compute \(\lim_{n \to \infty} a_{cn, n}\) for an
  arbitrary constant \(c > 0\). Express the answer as a function of
  \(c\). What is the maximum value, and where is it achieved?
\item
  Compute \(\lim_{n \to \infty} a_{n^2, n}\) (the path \(m = n^2\)).
  Does this give the same answer as any of the above?
\item
  Conclude: the double limit \(\lim_{m,n \to \infty} a_{mn}\) does not
  exist. How many distinct ``limit values'' can be achieved by different
  paths to infinity?
\item
  \textbf{Contrast}: Now consider \(b_{mn} = \frac{1}{m+n}\). Prove that
  \(\lim_{m,n \to \infty} b_{mn} = 0\). Verify that every path
  \(m = f(n)\) with \(f(n) \to \infty\) gives the same limit. Why does
  the theorem from the text guarantee this?
\end{enumerate}

\subsection{Exercises}\label{exercises-21}

\emph{Computing iterated limits} - Compute both iterated limits for
\(a_{mn} = \frac{n}{m + n}\) - Compute both iterated limits for
\(a_{mn} = \frac{m^2}{m^2 + n}\) - Compute both iterated limits for
\(a_{mn} = \frac{m + n}{m^2 + n^2}\) - Compute both iterated limits for
\(a_{mn} = \frac{(-1)^{m+n} mn}{m^2 + n^2}\) - Compute both iterated
limits for \(a_{mn} = \frac{(-1)^m n}{m + n}\)

\emph{Iterated limits and diagonal} - For
\(a_{mn} = \frac{mn}{m^2 + n^2}\), compute
\(\lim_{n \to \infty} a_{n, 2n}\). Compare to the iterated limits. - For
\(a_{mn} = \frac{m - n}{m + n}\), compute
\(\lim_{n \to \infty} a_{2n, n}\). - Give an example of a double
sequence where both iterated limits equal \(0\), but
\(\lim_{n \to \infty} a_{n, 2n} = 1\).

\emph{Double limits} - Prove:
\(\lim_{m,n \to \infty} \frac{1}{m + n} = 0\) - Prove:
\(\lim_{m,n \to \infty} \frac{m + n}{mn} = 0\) - Does
\(\lim_{m,n \to \infty} \frac{m}{n}\) exist? Prove or disprove.

\emph{Theoretical} - Prove: if \(a_{mn} \to L\) as \(m + n \to \infty\),
then the double limit exists and equals \(L\) - Give an example where
\(\lim_m \lim_n a_{mn}\) exists but \(\lim_n \lim_m a_{mn}\) does not -
Prove: if \(\lim_{m,n \to \infty} a_{mn} = L\) and
\(\lim_{m,n \to \infty} b_{mn} = M\), then
\(\lim_{m,n \to \infty} (a_{mn} + b_{mn}) = L + M\) - Prove: if
\(a_{mn} \leq b_{mn}\) for all \(m, n\), and both double limits exist,
then \(\lim_{m,n} a_{mn} \leq \lim_{m,n} b_{mn}\) - Give an example
where \(\lim_{m \to \infty} \lim_{n \to \infty} a_{mn}\) exists but
\(\lim_{n \to \infty} a_{mn}\) does not exist for some \(m\). - ★
Suppose \(\lim_{n \to \infty} a_{mn} = b_m\) exists for each \(m\), and
\(\lim_{m \to \infty} a_{mn} = c_n\) exists for each \(n\), and the
convergence \(a_{mn} \to b_m\) is uniform in \(m\) (i.e., for all
\(\varepsilon > 0\), there exists \(N\) such that
\(n > N \Rightarrow |a_{mn} - b_m| < \varepsilon\) for all \(m\)). Prove
that the double limit exists and equals both iterated limits.

\subsection{Dependencies}\label{dependencies-21}

\textbf{Requires}: §2.2 (definition of limits), §2.4
(\(e = \lim(1+1/n)^n\), geometric sequences)

\textbf{Used in}: §5.3 (double sums)

\section{§5.3 Double Sums}\label{double-sums}

\subsection{Narrative}\label{narrative-22}

\begin{itemize}
\tightlist
\item
  Double series: summing over a grid of terms
\item
  Two ways to sum: rows first or columns first (iterated sums)
\item
  These can differ! Just like iterated limits.
\item
  Absolute convergence saves us: when \(\sum_{m,n} |a_{mn}| < \infty\),
  any summation path works
\item
  Rows, columns, anti-diagonals, expanding squares --- all give the same
  answer
\item
  The Cauchy product: multiplying series by summing along anti-diagonals
\item
  Absolute convergence makes products work; Mertens weakens the
  hypothesis
\item
  Applications: squaring the geometric series, a remarkable identity for
  a mysterious series
\end{itemize}

\subsection{Content}\label{content-22}

\subsection{Double Series}\label{double-series}

\begin{itemize}
\tightlist
\item
  Definition: A \emph{double series} is a formal sum
  \(\sum_{m=1}^{\infty} \sum_{n=1}^{\infty} a_{mn}\) where \((a_{mn})\)
  is a double sequence
\item
  Two natural interpretations (iterated sums):

  \begin{itemize}
  \tightlist
  \item
    Sum rows first:
    \(\sum_{m=1}^{\infty} \left( \sum_{n=1}^{\infty} a_{mn} \right)\)
  \item
    Sum columns first:
    \(\sum_{n=1}^{\infty} \left( \sum_{m=1}^{\infty} a_{mn} \right)\)
  \end{itemize}
\item
  These are iterated sums, analogous to iterated limits from §5.2
\item
  Question: Do they always agree?
\end{itemize}

\subsection{When Order Matters}\label{when-order-matters-1}

\begin{itemize}
\tightlist
\item
  Example: Define \(a_{mn}\) by:
  \[a_{mn} = \begin{cases} 1 & \text{if } n = m \\ -1 & \text{if } n = m + 1 \\ 0 & \text{otherwise} \end{cases}\]
\item
  Visualize as a grid:

  \begin{itemize}
  \tightlist
  \item
    Row 1: \(1, -1, 0, 0, 0, \ldots\)
  \item
    Row 2: \(0, 1, -1, 0, 0, \ldots\)
  \item
    Row 3: \(0, 0, 1, -1, 0, \ldots\)
  \item
    (1's on diagonal, \(-1\)'s just above)
  \end{itemize}
\item
  Sum by rows: each row sums to \(1 + (-1) = 0\), so
  \(\sum_m \sum_n a_{mn} = 0\)
\item
  Sum by columns:

  \begin{itemize}
  \tightlist
  \item
    Column 1: only \(a_{11} = 1\), so sum is \(1\)
  \item
    Column \(n \geq 2\): \(a_{n-1,n} = -1\) and \(a_{nn} = 1\), so sum
    is \(0\)
  \item
    Therefore \(\sum_n \sum_m a_{mn} = 1 + 0 + 0 + \cdots = 1\)
  \end{itemize}
\item
  The iterated sums differ: \(0 \neq 1\)
\item
  Moral: order of summation matters for double series, just as for
  rearrangements (§5.1)
\item
  ✎ Inline: Construct a double sequence where row sums give \(2\) and
  column sums give \(-2\)
\end{itemize}

\subsection{When Order Doesn't Matter}\label{when-order-doesnt-matter}

\begin{itemize}
\tightlist
\item
  Definition: A double series \(\sum_{m,n} a_{mn}\) \emph{converges
  absolutely} if:
  \[\sum_{m=1}^{\infty} \sum_{n=1}^{\infty} |a_{mn}| < \infty\]
\item
  Equivalently: the partial sums
  \(\sum_{m=1}^{M} \sum_{n=1}^{N} |a_{mn}|\) are bounded as
  \(M, N \to \infty\)
\item
  Theorem: If \(\sum_{m,n} |a_{mn}| < \infty\), then:

  \begin{enumerate}
  \def\labelenumi{\arabic{enumi}.}
  \tightlist
  \item
    Both iterated sums converge absolutely and are equal
  \item
    In fact, \emph{any} ordering of the terms gives the same sum
  \end{enumerate}
\item
  Proof (in text):

  \begin{itemize}
  \tightlist
  \item
    Let \(S_{MN} = \sum_{m=1}^{M} \sum_{n=1}^{N} a_{mn}\) (partial sums
    over rectangles)
  \item
    \emph{Step 1}: \((S_{MN})\) is Cauchy in the double-index sense

    \begin{itemize}
    \tightlist
    \item
      Given \(\varepsilon > 0\), by absolute convergence choose
      \(M_0, N_0\) such that
      \[\sum_{m > M_0 \text{ or } n > N_0} |a_{mn}| < \varepsilon\]
    \item
      For \(M, M' > M_0\) and \(N, N' > N_0\):
      \(|S_{MN} - S_{M'N'}| < 2\varepsilon\)
    \end{itemize}
  \item
    \emph{Step 2}: The double limit \(S = \lim_{M,N \to \infty} S_{MN}\)
    exists
  \item
    \emph{Step 3}: Iterated sums equal \(S\)

    \begin{itemize}
    \tightlist
    \item
      For each fixed \(m\),
      \(\sum_n |a_{mn}| \leq \sum_{m,n} |a_{mn}| < \infty\), so row sums
      converge absolutely
    \item
      Let \(R_m = \sum_{n=1}^{\infty} a_{mn}\) (the \(m\)-th row sum)
    \item
      By §5.2 theorem: since double limit exists and
      \(\lim_{N \to \infty} S_{MN} = \sum_{m=1}^{M} R_m\) exists for
      each \(M\), we have \(\sum_{m=1}^{\infty} R_m = S\)
    \item
      Symmetric argument for columns
    \end{itemize}
  \item
    \emph{Step 4}: Any ordering works

    \begin{itemize}
    \tightlist
    \item
      Same argument as §5.1: absolute convergence allows rearrangement
    \item
      Any bijection
      \(\sigma: \mathbb{N} \to \mathbb{N} \times \mathbb{N}\) gives
      \(\sum_{k=1}^{\infty} a_{\sigma(k)} = S\)
    \end{itemize}
  \end{itemize}
\item
  Corollary: Under absolute convergence, we can sum by:

  \begin{itemize}
  \tightlist
  \item
    Rows: \(\sum_m (\sum_n a_{mn})\)
  \item
    Columns: \(\sum_n (\sum_m a_{mn})\)
  \item
    Expanding squares: \(\lim_{N \to \infty} \sum_{m,n \leq N} a_{mn}\)
  \item
    Anti-diagonals: group by \(m + n = k\)
  \item
    Any bijection \(\mathbb{N} \to \mathbb{N} \times \mathbb{N}\)
  \end{itemize}
\item
  Example: \(\sum_{m,n \geq 1} \frac{1}{m^2 n^2}\) converges absolutely

  \begin{itemize}
  \tightlist
  \item
    \(\sum_{m,n} \frac{1}{m^2 n^2} = \left(\sum_m \frac{1}{m^2}\right)\left(\sum_n \frac{1}{n^2}\right) = \left(\frac{\pi^2}{6}\right)^2\)
  \item
    (We'll prove \(\sum 1/n^2 = \pi^2/6\) via Fourier series in Chapter
    13; for now this shows the method)
  \end{itemize}
\item
  ✎ Inline: Verify that \(\sum_{m,n \geq 1} \frac{1}{2^{m+n}}\)
  converges absolutely and compute its sum two ways
\end{itemize}

\subsection{Multiplying Series}\label{multiplying-series}

\begin{itemize}
\tightlist
\item
  Setup: Given two series \(\sum_{j=0}^{\infty} a_j\) and
  \(\sum_{k=0}^{\infty} b_k\), what is their ``product''?
\item
  Natural idea: form the grid \(a_j b_k\) and sum it somehow
\item
  If we sum along anti-diagonals (where \(j + k\) is constant), we get:
\item
  Definition: The \emph{Cauchy product} of \(\sum_{j=0}^{\infty} a_j\)
  and \(\sum_{k=0}^{\infty} b_k\) is \(\sum_{n=0}^{\infty} c_n\) where:
  \[c_n = \sum_{k=0}^{n} a_k b_{n-k} = a_0 b_n + a_1 b_{n-1} + \cdots + a_n b_0\]
\item
  Interpretation: \(c_n\) collects all terms \(a_j b_k\) with
  \(j + k = n\) (the \(n\)-th anti-diagonal)
\item
  Connection to polynomials: if \(p(x) = \sum a_j x^j\) and
  \(q(x) = \sum b_k x^k\), then \(p(x)q(x) = \sum c_n x^n\)
\item
  Theorem: If \(\sum a_j\) and \(\sum b_k\) both converge absolutely,
  then their Cauchy product converges absolutely to
  \((\sum a_j)(\sum b_k)\)
\item
  Proof (in text):

  \begin{itemize}
  \tightlist
  \item
    The double series \(\sum_{j,k} a_j b_k\) converges absolutely:
    \[\sum_{j,k} |a_j b_k| = \left(\sum_j |a_j|\right)\left(\sum_k |b_k|\right) < \infty\]
  \item
    By the theorem above, any ordering gives the same sum
  \item
    Sum by rows:
    \(\sum_j a_j \left(\sum_k b_k\right) = \left(\sum_j a_j\right)\left(\sum_k b_k\right)\)
  \item
    Sum by anti-diagonals: \(\sum_{n=0}^{\infty} c_n\) (the Cauchy
    product)
  \item
    Therefore \(\sum c_n = (\sum a_j)(\sum b_k)\)
  \item
    Absolute convergence of \(\sum c_n\):
    \[\sum_{n=0}^{\infty} |c_n| \leq \sum_{n=0}^{\infty} \sum_{k=0}^{n} |a_k||b_{n-k}| = \sum_{j,k} |a_j b_k| < \infty\]
  \end{itemize}
\item
  Theorem (Mertens): If \(\sum a_j\) converges absolutely to \(A\) and
  \(\sum b_k\) converges (possibly conditionally) to \(B\), then the
  Cauchy product converges to \(AB\)
\item
  Proof (in text):

  \begin{itemize}
  \tightlist
  \item
    Let \(A_n = \sum_{j=0}^{n} a_j\), \(B_n = \sum_{k=0}^{n} b_k\),
    \(C_n = \sum_{m=0}^{n} c_m\)
  \item
    Let \(\beta_n = B_n - B\) (so \(\beta_n \to 0\))
  \item
    Compute:
    \[C_n = \sum_{m=0}^{n} \sum_{k=0}^{m} a_k b_{m-k} = \sum_{k=0}^{n} a_k \sum_{m=k}^{n} b_{m-k} = \sum_{k=0}^{n} a_k B_{n-k}\]
  \item
    Rewrite:
    \(C_n = \sum_{k=0}^{n} a_k (B + \beta_{n-k}) = A_n B + \sum_{k=0}^{n} a_k \beta_{n-k}\)
  \item
    Need to show \(\sum_{k=0}^{n} a_k \beta_{n-k} \to 0\)
  \item
    Let \(\alpha = \sum |a_j| < \infty\) and \(M = \sup_n |\beta_n|\)
  \item
    Given \(\varepsilon > 0\), choose \(N\) such that
    \(|\beta_m| < \varepsilon\) for \(m > N\)
  \item
    Split:
    \(\sum_{k=0}^{n} a_k \beta_{n-k} = \sum_{k=0}^{n-N-1} a_k \beta_{n-k} + \sum_{k=n-N}^{n} a_k \beta_{n-k}\)
  \item
    First sum:
    \(|\sum_{k=0}^{n-N-1} a_k \beta_{n-k}| \leq \varepsilon \sum_{k=0}^{n-N-1} |a_k| \leq \varepsilon \alpha\)
  \item
    Second sum:
    \(|\sum_{k=n-N}^{n} a_k \beta_{n-k}| \leq M \sum_{k=n-N}^{n} |a_k| \to 0\)
    as \(n \to \infty\) (tail of convergent series)
  \item
    Therefore \(\sum_{k=0}^{n} a_k \beta_{n-k} \to 0\), so
    \(C_n \to AB\)
  \end{itemize}
\item
  Remark: Without absolute convergence of at least one factor, the
  Cauchy product can diverge (see exercises)
\end{itemize}

\emph{Application: Squaring the Geometric Series} - For \(|x| < 1\):
\(\sum_{n=0}^{\infty} x^n = \frac{1}{1-x}\) converges absolutely -
Squaring via Cauchy product: take \(a_n = b_n = x^n\) -
\(c_n = \sum_{k=0}^{n} x^k \cdot x^{n-k} = \sum_{k=0}^{n} x^n = (n+1)x^n\)
- By the theorem: \[\frac{1}{(1-x)^2} = \sum_{n=0}^{\infty} (n+1)x^n\] -
✎ Inline: Verify at \(x = 1/2\): the series gives \(\sum (n+1)/2^n = ?\)
and \((1 - 1/2)^{-2} = 4\)

\emph{Application: A Remarkable Identity} - Define
\(E(x) = \sum_{n=0}^{\infty} \frac{x^n}{n!}\) for \(x \in \mathbb{R}\) -
This converges absolutely for all \(x\) (ratio test:
\(|x^{n+1}/(n+1)!| / |x^n/n!| = |x|/(n+1) \to 0\)) - Claim:
\(E(x) \cdot E(y) = E(x + y)\) for all \(x, y \in \mathbb{R}\) - Proof:
Both series converge absolutely, so we can use the Cauchy product -
\(c_n = \sum_{k=0}^{n} \frac{x^k}{k!} \cdot \frac{y^{n-k}}{(n-k)!} = \frac{1}{n!} \sum_{k=0}^{n} \frac{n!}{k!(n-k)!} x^k y^{n-k} = \frac{1}{n!} \sum_{k=0}^{n} \binom{n}{k} x^k y^{n-k}\)
- By the binomial theorem: \(c_n = \frac{(x+y)^n}{n!}\) - Therefore
\(E(x) E(y) = \sum c_n = E(x+y)\) - Remark: This series has a remarkable
property --- it turns addition into multiplication! - Remark: Setting
\(y = 0\): \(E(x) E(0) = E(x)\), so \(E(0) = 1\) - Remark: Setting
\(y = -x\): \(E(x) E(-x) = E(0) = 1\), so \(E(-x) = 1/E(x)\) - What
\emph{is} this function? In §5.4 we'll show
\(E(x) = \lim_{n \to \infty} (1 + x/n)^n\), connecting it to Chapter 2's
definition of \(e\). In Chapter 6 we'll call it the exponential function
and develop its properties systematically.

\subsection{Guided Exercise}\label{guided-exercise-9}

\subsection{A Conditionally Convergent Double
Series}\label{a-conditionally-convergent-double-series}

Consider the double series \(\sum_{m,n \geq 1} \frac{(-1)^{m+n}}{m+n}\).

\begin{enumerate}
\def\labelenumi{(\alph{enumi})}
\item
  Show that for each fixed \(m\), the series
  \(\sum_{n=1}^{\infty} \frac{(-1)^{m+n}}{m+n}\) converges
  (conditionally) by the alternating series test.
\item
  Show that the double series does NOT converge absolutely. (Hint:
  compare \(\sum_{m,n} \frac{1}{m+n}\) to something you know diverges.)
\item
  Compute the row sums: for fixed \(m\), what is
  \(R_m = \sum_{n=1}^{\infty} \frac{(-1)^{m+n}}{m+n}\)? (Hint: this is a
  tail of the alternating harmonic series, with a sign.)
\item
  Does \(\sum_{m=1}^{\infty} R_m\) converge? If so, what is its value?
\item
  Now compute the column sums
  \(C_n = \sum_{m=1}^{\infty} \frac{(-1)^{m+n}}{m+n}\) and
  \(\sum_{n=1}^{\infty} C_n\).
\item
  Do the row-sum and column-sum approaches give the same answer? What
  does this tell you?
\end{enumerate}

\subsection{Exercises}\label{exercises-22}

\emph{When order matters} - Construct a double sequence where
\(\sum_m \sum_n a_{mn} = 0\) and \(\sum_n \sum_m a_{mn} = 2\) -
Construct a double sequence where one iterated sum converges and the
other diverges - For the example in the text (\(1\)'s on diagonal,
\(-1\)'s above), what happens if you sum along anti-diagonals?

\emph{Absolute convergence} - Prove:
\(\sum_{m,n \geq 1} \frac{1}{(m+n)^3}\) converges absolutely - Prove:
\(\sum_{m,n \geq 1} \frac{1}{m^2 + n^2}\) diverges (hint: count terms
with \(m^2 + n^2 \leq R^2\)) - Compute
\(\sum_{m,n \geq 1} \frac{1}{2^m 3^n}\) by summing rows, then verify by
summing columns

\emph{Cauchy products} - Compute the Cauchy product of
\(\sum_{n=0}^{\infty} x^n\) with \(\sum_{n=0}^{\infty} (-1)^n x^n\).
What function does this represent? - Compute the Cauchy product of
\(\sum_{n=0}^{\infty} x^n/n!\) with itself. Verify it equals
\(\sum_{n=0}^{\infty} (2x)^n/n!\). - Use Cauchy products to find a power
series for \(1/(1-x)^3\). What are the coefficients?

\emph{Mertens and its limits} - Verify \((\sum x^n)^2 = \sum (n+1)x^n\)
by differentiating \(1/(1-x)\) - ★ Prove Mertens' theorem - ★ Give an
example where both \(\sum a_n\) and \(\sum b_n\) converge conditionally
but the Cauchy product diverges. (Hint: try
\(a_n = b_n = (-1)^n/\sqrt{n+1}\).) - Let
\(S = \sum_{n=1}^{\infty} (-1)^{n+1}/n\) (the alternating harmonic
series). Take \(a_n = (-1)^{n+1}/n\) (converges conditionally to \(S\))
and \(b_n = 1/2^n\) (converges absolutely to \(2\)). Verify that Mertens
applies and conclude that the Cauchy product converges to \(2S\). -
Compute the Cauchy product of \(\sum_{n=0}^{\infty} (-1)^n/2^n\) with
\(\sum_{n=0}^{\infty} (-1)^n/3^n\). Both converge absolutely; find
explicit formulas for the coefficients \(c_n\) and verify the product
equals \((2/3)(3/4) = 1/2\).

\emph{The series \(E(x)\)} - Prove \(E(x) > 0\) for all
\(x \in \mathbb{R}\). (Hint: \(E(x) = E(x/2)^2\).) - Prove
\(E(1) = \sum_{n=0}^{\infty} 1/n!\) lies between \(2\) and \(3\). -
Prove \(E(n) = E(1)^n\) for all \(n \in \mathbb{N}\). - ★ Prove
\(E(r) = E(1)^r\) for all \(r \in \mathbb{Q}\). (Hint: first show
\(E(1/q)^q = E(1)\).)

\subsection{Dependencies}\label{dependencies-22}

\textbf{Requires}: §5.1 (rearrangements), §5.2 (iterated limits, double
limits), §4.2 (absolute convergence, ratio test)

\textbf{Used in}: §5.4 (dominated convergence), §6.3 (power series,
Cauchy product)

\section{§5.4 Dominated Convergence}\label{dominated-convergence}

\subsection{Narrative}\label{narrative-23}

\begin{itemize}
\tightlist
\item
  New question: when can we interchange a limit and an infinite sum?
\item
  The danger: example where they don't commute (many small terms add up)
\item
  Dominated convergence: if terms are bounded by a convergent series,
  interchange is safe
\item
  The dominating series provides uniform control
\item
  The big application: \((1 + x/n)^n = \sum x^k/k!\)
\item
  This connects Chapter 2's definition of \(e\) to the series from §5.3
\item
  Payoff: the exponential series is ready for Chapter 6
\end{itemize}

\subsection{Content}\label{content-23}

\subsection{The Danger}\label{the-danger}

\begin{itemize}
\tightlist
\item
  Question: When can we interchange a limit and an infinite sum?
  \[\lim_{n \to \infty} \sum_{k=0}^{\infty} a_k(n) \stackrel{?}{=} \sum_{k=0}^{\infty} \lim_{n \to \infty} a_k(n)\]
\item
  Example: Let \(a_k(n) = 1/n\) if \(k \leq n\), and \(0\) otherwise
\item
  RHS: For each fixed \(k\), \(\lim_{n \to \infty} a_k(n) = 0\)
  (eventually \(n > k\), so \(a_k(n) = 1/n \to 0\))

  \begin{itemize}
  \tightlist
  \item
    Therefore
    \(\sum_{k=0}^{\infty} \lim_{n \to \infty} a_k(n) = \sum_{k=0}^{\infty} 0 = 0\)
  \end{itemize}
\item
  LHS:
  \(\sum_{k=0}^{\infty} a_k(n) = \sum_{k=0}^{n} \frac{1}{n} = \frac{n+1}{n} \to 1\)

  \begin{itemize}
  \tightlist
  \item
    Therefore \(\lim_{n \to \infty} \sum_{k=0}^{\infty} a_k(n) = 1\)
  \end{itemize}
\item
  The limit and sum don't commute: \(1 \neq 0\)
\item
  What went wrong? The terms \(a_k(n) = 1/n\) are small (each
  \(\to 0\)), but there are many of them (\(n+1\) terms), and
  \((n+1) \cdot (1/n) \to 1\) doesn't vanish
\item
  The number of nonzero terms grows with \(n\) --- no uniform control
\end{itemize}

\subsection{Dominated Convergence}\label{dominated-convergence-1}

\begin{itemize}
\tightlist
\item
  Theorem (Dominated Convergence for Series): Suppose:

  \begin{enumerate}
  \def\labelenumi{\arabic{enumi}.}
  \tightlist
  \item
    For each \(k\), the limit \(L_k = \lim_{n \to \infty} a_k(n)\)
    exists
  \item
    There exist \(M_k \geq 0\) with \(|a_k(n)| \leq M_k\) for all \(n\)
  \item
    \(\sum_{k=0}^{\infty} M_k < \infty\)
  \end{enumerate}

  Then:
  \[\lim_{n \to \infty} \sum_{k=0}^{\infty} a_k(n) = \sum_{k=0}^{\infty} L_k\]
\item
  Remark: The bound \(M_k\) is independent of \(n\) --- this is the
  ``domination'' condition
\item
  Remark: Condition 3 ensures both sides converge absolutely
\item
  Proof (in text):

  \begin{itemize}
  \tightlist
  \item
    Let \(S(n) = \sum_{k=0}^{\infty} a_k(n)\) and
    \(S = \sum_{k=0}^{\infty} L_k\)
  \item
    Both series converge absolutely: \(|a_k(n)| \leq M_k\) and
    \(|L_k| \leq M_k\) (taking limits preserves inequalities)
  \item
    Given \(\varepsilon > 0\), choose \(K\) such that
    \(\sum_{k > K} M_k < \varepsilon/3\)
  \item
    Split:
    \[|S(n) - S| \leq \sum_{k=0}^{K} |a_k(n) - L_k| + \sum_{k > K} |a_k(n)| + \sum_{k > K} |L_k|\]
  \item
    Third term:
    \(\sum_{k > K} |L_k| \leq \sum_{k > K} M_k < \varepsilon/3\)
  \item
    Second term:
    \(\sum_{k > K} |a_k(n)| \leq \sum_{k > K} M_k < \varepsilon/3\)
  \item
    First term: finitely many limits. Choose \(N\) such that for all
    \(n > N\):
    \[|a_k(n) - L_k| < \frac{\varepsilon}{3(K+1)} \quad \text{for each } k = 0, 1, \ldots, K\]
  \item
    Then
    \(\sum_{k=0}^{K} |a_k(n) - L_k| < (K+1) \cdot \frac{\varepsilon}{3(K+1)} = \varepsilon/3\)
  \item
    Total: \(|S(n) - S| < \varepsilon\) for \(n > N\)
  \end{itemize}
\item
  Why the counterexample fails:

  \begin{itemize}
  \tightlist
  \item
    We would need \(|a_k(n)| \leq M_k\) for all \(n\)
  \item
    But \(a_k(n) = 1/n\) for \(k \leq n\), so we'd need \(M_k \geq 1/n\)
    for all \(n \geq k\)
  \item
    This means \(M_k \geq \sup_{n \geq k} 1/n\)\ldots{} but wait, the
    sup over all \(n\) is taken at \(n = k\), giving \(M_k \geq 1/k\)
  \item
    Then \(\sum M_k \geq \sum 1/k = \infty\) --- condition 3 fails
  \end{itemize}
\item
  ✎ Inline: Verify that if \(|a_k(n)| \leq M_k\) for all \(n\), then
  \(|L_k| \leq M_k\)
\end{itemize}

\subsection{The Big Application}\label{the-big-application}

\begin{itemize}
\tightlist
\item
  Theorem: For all \(x \in \mathbb{R}\):
  \[\lim_{n \to \infty} \left(1 + \frac{x}{n}\right)^n = \sum_{k=0}^{\infty} \frac{x^k}{k!}\]
\item
  Proof (in text):

  \begin{itemize}
  \tightlist
  \item
    \emph{Step 1: Expand by binomial theorem}

    \begin{itemize}
    \tightlist
    \item
      \((1 + x/n)^n = \sum_{k=0}^{n} \binom{n}{k} \frac{x^k}{n^k}\)
    \end{itemize}
  \item
    Define \(a_k(n)\) for all \(k \geq 0\):
    \[a_k(n) = \begin{cases} \binom{n}{k} \frac{x^k}{n^k} & k \leq n \\ 0 & k > n \end{cases}\]
  \item
    Then \((1 + x/n)^n = \sum_{k=0}^{\infty} a_k(n)\)
  \item
    \emph{Step 2: Compute limits}

    \begin{itemize}
    \tightlist
    \item
      For \(k \leq n\):
      \[a_k(n) = \frac{n!}{k!(n-k)!} \cdot \frac{x^k}{n^k} = \frac{x^k}{k!} \cdot \frac{n(n-1)\cdots(n-k+1)}{n^k}\]
      \[= \frac{x^k}{k!} \cdot 1 \cdot \left(1 - \frac{1}{n}\right) \cdot \left(1 - \frac{2}{n}\right) \cdots \left(1 - \frac{k-1}{n}\right)\]
    \item
      For fixed \(k\), as \(n \to \infty\): each factor
      \((1 - j/n) \to 1\)
    \item
      Therefore \(\lim_{n \to \infty} a_k(n) = \frac{x^k}{k!}\)
    \end{itemize}
  \item
    \emph{Step 3: Find a dominating sequence}

    \begin{itemize}
    \tightlist
    \item
      For \(k \leq n\): each factor \((1 - j/n) \leq 1\), so
      \[|a_k(n)| = \frac{|x|^k}{k!} \cdot \left(1 - \frac{1}{n}\right) \cdots \left(1 - \frac{k-1}{n}\right) \leq \frac{|x|^k}{k!}\]
    \item
      For \(k > n\): \(a_k(n) = 0 \leq \frac{|x|^k}{k!}\)
    \item
      Let \(M_k = \frac{|x|^k}{k!}\). Then \(|a_k(n)| \leq M_k\) for all
      \(n\)
    \item
      \(\sum_{k=0}^{\infty} M_k = \sum_{k=0}^{\infty} \frac{|x|^k}{k!}\)
      converges (ratio test, or it equals \(E(|x|)\) from §5.3)
    \end{itemize}
  \item
    \emph{Step 4: Apply dominated convergence}

    \begin{itemize}
    \tightlist
    \item
      All hypotheses verified
    \item
      Therefore
      \(\lim_{n \to \infty} (1 + x/n)^n = \sum_{k=0}^{\infty} \frac{x^k}{k!}\)
    \end{itemize}
  \end{itemize}
\item
  Corollary: Setting \(x = 1\):
  \[e = \lim_{n \to \infty} \left(1 + \frac{1}{n}\right)^n = \sum_{k=0}^{\infty} \frac{1}{k!}\]
\item
  This connects the compound interest limit from Chapter 2 to the
  series!
\item
  Corollary: The series \(E(x) = \sum_{k=0}^{\infty} \frac{x^k}{k!}\)
  satisfies \(E(x) = \lim_{n \to \infty} (1 + x/n)^n\)
\item
  Combined with §5.3: \(E(x) E(y) = E(x+y)\), \(E(0) = 1\),
  \(E(-x) = 1/E(x)\)
\item
  Remark: In Chapter 6, we define the exponential function
  \(\exp(x) = E(x)\) and derive all its properties
\end{itemize}

\subsection{Guided Exercise}\label{guided-exercise-10}

\subsection{Variations on the Exponential
Limit}\label{variations-on-the-exponential-limit}

The dominated convergence theorem is robust --- small perturbations
don't affect the limit.

\begin{enumerate}
\def\labelenumi{(\alph{enumi})}
\tightlist
\item
  Show that
  \(\lim_{n \to \infty} \left(1 + \frac{x}{n} + \frac{y}{n^2}\right)^n = e^x\)
  for any fixed \(x, y \in \mathbb{R}\).
\end{enumerate}

\emph{Hint}: Expand using the binomial theorem (or generalized
binomial). The \(y/n^2\) term contributes \(O(1/n)\) to each binomial
coefficient.

\begin{enumerate}
\def\labelenumi{(\alph{enumi})}
\setcounter{enumi}{1}
\item
  More generally, show that if \(b_n \to 0\), then
  \(\lim_{n \to \infty} (1 + x/n + b_n/n)^n = e^x\).
\item
  Show that
  \(\lim_{n \to \infty} \left(1 + \frac{x}{n}\right)^{n + c} = e^x\) for
  any fixed \(c \in \mathbb{R}\).
\item
  Show that
  \(\lim_{n \to \infty} \left(1 - \frac{x}{n}\right)^{-n} = e^x\).
\item
  Conclude: many natural variations on the compound interest limit all
  give \(e^x\).
\end{enumerate}

\subsection{Exercises}\label{exercises-23}

\emph{The danger} - Let \(a_k(n) = 1/n^2\) for \(k \leq n\), and \(0\)
otherwise. Show that limit and sum DO commute in this case. Why doesn't
this contradict the counterexample? - Let \(a_k(n) = k/n^2\) for
\(k \leq n\), and \(0\) otherwise. Do limit and sum commute? - Construct
an example where the RHS converges but the LHS diverges.

\emph{Applying dominated convergence} - Use dominated convergence to
prove: \(\lim_{n \to \infty} \sum_{k=1}^{\infty} \frac{1}{k(k+n)} = 0\)
- Use dominated convergence to prove:
\(\lim_{n \to \infty} \sum_{k=0}^{\infty} \frac{n}{(n+k)^2} = 1\) - Let
\(a_k(n) = \frac{1}{n} \sin(k/n)\) for \(k \leq n\), \(0\) otherwise.
Find \(\lim_{n \to \infty} \sum_k a_k(n)\). (This is a Riemann sum!)

\emph{The exponential series} - Verify the domination bound: show
\(\frac{n(n-1)\cdots(n-k+1)}{n^k} \leq 1\) for \(k \leq n\) - Prove
directly that \(e = \sum 1/k!\) lies between \(2\) and \(3\) - Prove
\(\sum_{k=0}^{\infty} \frac{(-1)^k}{k!} = 1/e\) - ★ Use dominated
convergence to show:
\(\lim_{n \to \infty} \left(1 + \frac{x}{n} + \frac{x^2}{2n^2}\right)^n = e^x\)

\emph{Theoretical} - Prove: if \(\sum_k M_k < \infty\) and
\(|a_k(n)| \leq M_k\) for all \(n\), then \(\sum_k a_k(n)\) converges
uniformly in \(n\) - Give an example where dominated convergence fails
because condition 1 fails (limits don't exist) - ★ (Monotone convergence
for series) Prove: if \(0 \leq a_k(n) \leq a_k(n+1)\) for all \(k, n\),
and \(\sum_k a_k(n)\) converges for all \(n\), then
\(\lim_n \sum_k a_k(n) = \sum_k \lim_n a_k(n)\) (possibly \(= +\infty\))

\subsection{Dependencies}\label{dependencies-23}

\textbf{Requires}: §5.2 (iterated limits), §5.3 (the series \(E(x)\)),
§2.4 (\(e = \lim(1+1/n)^n\)), §4.2 (ratio test)

\textbf{Used in}: §6.3 (continuity of power series via dominated
convergence)

\part{Functions}

\section*{Part II Introduction: What Is a
Function?}\label{part-ii-introduction-what-is-a-function}
\addcontentsline{toc}{section}{Part II Introduction: What Is a
Function?}

\markright{Part II Introduction: What Is a Function?}

\subsection*{The old view}\label{the-old-view}
\addcontentsline{toc}{subsection}{The old view}

\begin{itemize}
\tightlist
\item
  For most of history, ``function'' meant ``formula'' --- an expression
  you could write down and evaluate
\item
  Euler (1748): functions are polynomials, trig expressions,
  exponentials, compositions of these
\item
  Everything in Part I was like this --- \(a^x\), \(\log_a\),
  \(E(x) = \sum x^k/k!\), geometric series
\item
  This world is tidy and well-behaved
\end{itemize}

\subsection*{Fourier's challenge}\label{fouriers-challenge}
\addcontentsline{toc}{subsection}{Fourier's challenge}

\begin{itemize}
\tightlist
\item
  Fourier (\textasciitilde1807) solving the heat equation: initial
  temperature of a bar can be anything --- one temperature on the left,
  different on the right
\item
  His solutions required representing arbitrary temperature
  distributions as infinite sums of sines and cosines
\item
  These ``functions'' aren't given by any single formula
\item
  Forced the question: what \emph{is} a function?
\end{itemize}

\subsection*{Dirichlet's answer}\label{dirichlets-answer}
\addcontentsline{toc}{subsection}{Dirichlet's answer}

\begin{itemize}
\tightlist
\item
  Dirichlet (1829): a function is \emph{any} rule assigning to each
  input exactly one output
\item
  To show he meant it: \(f(x) = 1\) on rationals, \(0\) on irrationals
\item
  No formula, no graph you could draw --- between any two points it
  oscillates infinitely between 0 and 1
\item
  A perfectly valid function under the new definition
\end{itemize}

\subsection*{The zoo}\label{the-zoo}
\addcontentsline{toc}{subsection}{The zoo}

\begin{itemize}
\tightlist
\item
  Once you accept Dirichlet's definition, the world of functions is vast
  and wild
\item
  The ruler function: \(f(p/q) = 1/q\) in lowest terms, \(0\) on
  irrationals --- somehow calm at irrationals, agitated at rationals
\item
  \(\sin(1/x)\) near zero: oscillates faster and faster, no way to
  assign a value at \(x = 0\) that ``fits''
\item
  Indicator functions of bizarre sets
\item
  The ``nice'' functions from calculus --- polynomials, exponentials,
  power series --- are a tiny minority
\item
  How do we sort the nice from the wild?
\end{itemize}

\subsection*{First lens: continuity}\label{first-lens-continuity}
\addcontentsline{toc}{subsection}{First lens: continuity}

\begin{itemize}
\tightlist
\item
  The most basic demand: nearby inputs should give nearby outputs
\item
  Bolzano (1817) wanted to prove the Intermediate Value Theorem --- a
  continuous function crossing from negative to positive must hit zero
\item
  Seemed geometrically obvious, but he couldn't prove it; needed a
  precise definition of continuity \emph{and} completeness of \(\RR\)
\item
  Cauchy (1821) gave a verbal definition --- close, but still vague
\item
  Weierstrass (1860s): the \(\varepsilon\)-\(\delta\) definition ---
  same quantifier game as \(\varepsilon\)-\(N\) for sequences, but local
\item
  Most monsters fail: Dirichlet's function is discontinuous everywhere
\item
  But some sneak through: Weierstrass (1872) constructs a function
  continuous at every point, yet so jagged it has no tangent line
  anywhere
\item
  Continuity tames the zoo, but doesn't fully domesticate it
\end{itemize}

\subsection*{Second lens: tangent
lines}\label{second-lens-tangent-lines}
\addcontentsline{toc}{subsection}{Second lens: tangent lines}

\begin{itemize}
\tightlist
\item
  Fermat (1630s): to find where a curve reaches its peak, look for where
  the tangent line is horizontal
\item
  His method of ``adequality'' --- set \(f(x)\) and \(f(x+e)\) nearly
  equal, divide by \(e\), then discard \(e\) --- mysteriously works
\item
  The tangent problem more generally: given a curve, what line best
  approximates it at a point?
\item
  Newton and Leibniz (1680s): the tangent slope is the limit of secant
  slopes
\item
  But what does ``limit'' mean here? Berkeley (1734): ``ghosts of
  departed quantities''
\item
  We now have limits (Chapter 2) --- so we can make this precise
\item
  The zoo: \(|x|\) has a corner at zero, no tangent. Weierstrass's
  monster has no tangent at \emph{any} point
\item
  Functions with tangent lines everywhere are a smaller, better-behaved
  class than merely continuous functions
\end{itemize}

\subsection*{Third lens: area}\label{third-lens-area}
\addcontentsline{toc}{subsection}{Third lens: area}

\begin{itemize}
\tightlist
\item
  We computed the area under a parabola in Chapter 2 (Archimedes'
  quadrature)
\item
  That worked beautifully --- but which functions in the zoo even
  \emph{have} a well-defined area?
\item
  For a smooth curve, intuition says yes --- but intuition hasn't met
  Dirichlet's function
\item
  What is the ``area'' under Dirichlet's function between 0 and 1? Upper
  approximation always 1, lower always 0 --- they never agree
\item
  Riemann (1854): the integral exists precisely when upper and lower
  approximations converge to the same value
\item
  Some functions have area, some don't --- understanding the boundary is
  one of the deep stories of Part II
\end{itemize}

\subsection*{What's ahead}\label{whats-ahead-1}
\addcontentsline{toc}{subsection}{What's ahead}

\begin{itemize}
\tightlist
\item
  Chapter 6: Continuity --- the definition, and what it forces (IVT,
  EVT)
\item
  Chapter 7: Differentiation --- tangent lines made rigorous, the Mean
  Value Theorem
\item
  Chapter 8: Integration --- axioms for area, the Riemann integral
\item
  Chapter 9: Calculus --- the Fundamental Theorem unifies
  differentiation and integration
\item
  Chapter 10: Extending the integral --- beyond Riemann, toward Lebesgue
\end{itemize}

\chapter{Continuity}\label{continuity}

\textbf{The problem} - The Part II introduction opened the zoo: once
``function'' means \emph{any} rule assigning values, most functions are
monsters - But we have a specific theorem we want to be true - If
\(f(a) < 0\) and \(f(b) > 0\), there should be some \(c\) between \(a\)
and \(b\) with \(f(c) = 0\) - For Euler's functions (polynomials, trig,
exponentials) this is obviously true - For general functions it's
obviously \emph{false}: a step function jumps from \(-1\) to \(1\)
without ever being zero - So the IVT holds for some class of functions
but not all --- which class?

\textbf{Cauchy and Bolzano try to answer} - Bolzano (1817): first to
insist on a purely analytic proof of IVT, no geometric appeals - He had
the right strategy --- essentially the least upper bound argument ---
but couldn't finish; the foundations of \(\RR\) weren't ready - Cauchy
(1821): defined continuity as ``an infinitely small change in \(x\)
produces an infinitely small change in \(f(x)\)'' - The right intuition
--- the step function fails because a tiny change at \(x = 0\) produces
a jump of \(2\) - Cauchy attempted IVT via bisection, but implicitly
assumed completeness without stating it - Both had the right instincts;
neither had the right tools

\textbf{Weierstrass makes it precise} - Weierstrass (1860s): the
\(\varepsilon\)-\(\delta\) definition - For every \(\varepsilon > 0\)
there exists \(\delta > 0\) such that \(|x - a| < \delta\) implies
\(|f(x) - f(a)| < \varepsilon\) - The same quantifier game as
\(\varepsilon\)-\(N\) for sequences, now played locally at a point -
Combined with completeness from Chapter 1, this gives everything needed
for IVT --- and much more

\textbf{The good news} - The class of continuous functions turns out to
be a well-behaved world - Sums, products, quotients, compositions of
continuous functions are continuous - Polynomials, power series,
exponentials --- all the functions we know --- are continuous - The
definition doesn't just identify the right class for IVT; it identifies
a class closed under the operations we care about

\textbf{But also surprises} - Cauchy (1821) claimed the limit of
continuous functions is continuous --- this is false (\(f_n(x) = x^n\)
on \([0,1]\)) - Weierstrass (1872) constructed a function continuous
everywhere but differentiable nowhere --- infinitely jagged at every
scale - Hermite: ``I turn away with fright and horror from this
lamentable plague of functions which have no derivatives'' - Continuity
is subtler and richer than it first appears

\textbf{What's ahead} - §6.1: The \(\varepsilon\)-\(\delta\) definition,
the connection to sequences, how continuity fails - §6.2: Building
continuous functions --- the good news made rigorous - §6.3: Power
series --- our most powerful source of continuous functions - §6.4: The
payoff --- IVT, EVT, continuous inverses: what Bolzano couldn't prove,
we now can - §6.5: Uniform continuity --- a stronger property for when
pointwise control isn't enough

\section{§6.1 Continuity}\label{continuity-1}

\subsection{Narrative}\label{narrative-24}

\begin{itemize}
\tightlist
\item
  Which functions respect closeness? The step function doesn't --- a
  tiny change at \(x = 0\) produces a jump of 2
\item
  The \(\varepsilon\)-\(\delta\) definition: control the output by
  controlling the input
\item
  Same game as \(\varepsilon\)-\(N\) from Chapter 2, but now local ---
  played at a point
\item
  A first example from the definition, showing the estimation technique
\item
  The bridge theorem: continuous iff sequences are preserved
  (\(x_n \to a \Rightarrow f(x_n) \to f(a)\))
\item
  This connects the new definition to five chapters of existing tools
\item
  Limits of functions: a variant where we don't care about the value at
  \(a\) itself
\item
  Continuity = the limit equals the value
\item
  One-sided limits: capturing directional behavior near a point
\end{itemize}

\subsection{Content}\label{content-24}

\subsection{The Definition}\label{the-definition-1}

\begin{itemize}
\tightlist
\item
  Motivation: which functions respect closeness? If \(x\) is close to
  \(a\), is \(f(x)\) close to \(f(a)\)?
\item
  The step function \(\text{sgn}(x) = -1\) for \(x < 0\), \(1\) for
  \(x \geq 0\): inputs \(-0.001\) and \(0.001\) are close, but outputs
  are \(-1\) and \(1\) --- not close
\item
  What property rules this out?
\item
  The \(\varepsilon\)-\(\delta\) game: adversary picks an output
  tolerance \(\varepsilon > 0\); you respond with an input tolerance
  \(\delta > 0\); you win if \(|x - a| < \delta\) guarantees
  \(|f(x) - f(a)| < \varepsilon\); \(f\) is continuous at \(a\) iff you
  have a winning strategy for every \(\varepsilon\)
\item
  Parallel with \(\varepsilon\)-\(N\): in Chapter 2, the adversary
  challenged output precision and you responded with ``go far enough in
  the sequence.'' Now the adversary challenges output precision and you
  respond with ``stay close enough to the point.''
\item
  Definition: \(f\) is \emph{continuous at \(a\)} if
  \[\forall \varepsilon > 0 \; \exists \delta > 0 \; \forall x: |x - a| < \delta \Rightarrow |f(x) - f(a)| < \varepsilon\]
\item
  Definition: \(f\) is \emph{continuous on a set \(S\)} if it is
  continuous at every point of \(S\)
\item
  Remark: \(\delta\) may depend on both \(\varepsilon\) \emph{and} the
  point \(a\) --- this dependence becomes important in §6.5
\end{itemize}

\emph{A first example: the estimation technique} - Example:
\(f(x) = 3x + 1\) is continuous at every \(a\) -
\(|f(x) - f(a)| = |3x + 1 - 3a - 1| = 3|x - a|\) - Choose
\(\delta = \varepsilon/3\); then
\(|x - a| < \delta \Rightarrow |f(x) - f(a)| = 3|x - a| < 3\delta = \varepsilon\)
- Example: \(f(x) = x^2\) is continuous at every \(a\) (prove in text) -
\(|x^2 - a^2| = |x - a| \cdot |x + a|\) - Need to control \(|x + a|\):
if \(|x - a| < 1\), then \(|x| < |a| + 1\), so \(|x + a| < 2|a| + 1\) -
Choose \(\delta = \min(1, \varepsilon/(2|a|+1))\) - Then
\(|f(x) - f(a)| < (2|a|+1) \cdot \delta \leq \varepsilon\) - Remark: the
\(\min(1, \ldots)\) trick --- restrict to a neighborhood first, then
estimate - Remark: this is the same estimation game as §2.2, just in a
new setting - ✎ Inline: Prove \(f(x) = 1/x\) is continuous at \(a = 2\)
from the \(\varepsilon\)-\(\delta\) definition

\emph{The step function is not continuous} - \(\text{sgn}(x)\) at
\(a = 0\): take \(\varepsilon = 1\). For any \(\delta > 0\), there exist
\(x = -\delta/2\) and \(x = \delta/2\) with \(|x - 0| < \delta\) but
\(|f(x) - f(0)| = 2 > \varepsilon\) - The adversary wins with
\(\varepsilon = 1\) --- no \(\delta\) works - Remark: to prove
discontinuity, we only need one \(\varepsilon\) for which no \(\delta\)
works

\subsection{The Sequential
Characterization}\label{the-sequential-characterization}

\begin{itemize}
\tightlist
\item
  Theorem: \(f\) is continuous at \(a\) if and only if: for every
  sequence \((x_n)\) with \(x_n \to a\), we have \(f(x_n) \to f(a)\)
\item
  This is the bridge between the new definition and five chapters of
  existing tools
\item
  Proof (\(\Rightarrow\)): Suppose \(f\) continuous at \(a\) and
  \(x_n \to a\). Given \(\varepsilon > 0\), get \(\delta > 0\) from
  continuity. Since \(x_n \to a\), there exists \(N\) with
  \(|x_n - a| < \delta\) for \(n > N\). Then
  \(|f(x_n) - f(a)| < \varepsilon\) for \(n > N\). So
  \(f(x_n) \to f(a)\).
\item
  Proof (\(\Leftarrow\), contrapositive): Suppose \(f\) is NOT
  continuous at \(a\). Then there exists \(\varepsilon_0 > 0\) such that
  for every \(\delta > 0\), some \(x\) has \(|x - a| < \delta\) but
  \(|f(x) - f(a)| \geq \varepsilon_0\). Take \(\delta = 1/n\): get
  \(x_n\) with \(|x_n - a| < 1/n\) but
  \(|f(x_n) - f(a)| \geq \varepsilon_0\). Then \(x_n \to a\) but
  \(f(x_n) \not\to f(a)\).
\item
  Remark: the reverse direction \emph{constructs} a bad sequence from
  the failure of \(\varepsilon\)-\(\delta\). This is a powerful
  technique.
\item
  Remark: from now on, we use whichever characterization is more
  convenient. The sequential version is often easier for proofs (invoke
  limit laws from Chapter 2); the \(\varepsilon\)-\(\delta\) version is
  often better for understanding what's happening locally.
\end{itemize}

\emph{Dirichlet's function is discontinuous everywhere} - Define
\(D(x) = 1\) if \(x \in \QQ\), \(D(x) = 0\) if \(x \notin \QQ\) - Claim:
\(D\) is discontinuous at every \(a \in \RR\) - Proof via sequential
characterization: - If \(a \in \QQ\): by density of irrationals, find
\(x_n \notin \QQ\) with \(x_n \to a\). Then
\(D(x_n) = 0 \not\to 1 = D(a)\). - If \(a \notin \QQ\): by density of
rationals, find \(x_n \in \QQ\) with \(x_n \to a\). Then
\(D(x_n) = 1 \not\to 0 = D(a)\). - Remark: density of \(\QQ\) and of
\(\RR \setminus \QQ\) (from §1.5) does the work; the sequential
characterization makes it clean

\subsection{Limits of Functions}\label{limits-of-functions}

\begin{itemize}
\tightlist
\item
  Sometimes a function ``wants'' to have a value at a point even though
  it doesn't --- or has the wrong one
\item
  Example: \(f(x) = \sin(x)/x\) is not defined at \(x = 0\), but for
  \(x\) near \(0\), \(f(x)\) is near \(1\)
\item
  Example: define \(g(x) = x^2\) for \(x \neq 0\) and \(g(0) = 17\). The
  value at \(0\) is ``wrong'' --- the function ``should be'' \(0\) there
\item
  We want to capture ``what \(f(x)\) approaches as \(x \to a\), ignoring
  what happens at \(a\) itself''
\item
  Definition: \(\lim_{x \to a} f(x) = L\) if
  \[\forall \varepsilon > 0 \; \exists \delta > 0 \; \forall x: 0 < |x - a| < \delta \Rightarrow |f(x) - L| < \varepsilon\]
\item
  The only difference from continuity: the condition \(0 < |x - a|\)
  excludes the point \(a\) itself
\item
  \(f\) doesn't even need to be defined at \(a\) for
  \(\lim_{x \to a} f(x)\) to exist
\item
  Theorem (Sequential characterization): \(\lim_{x \to a} f(x) = L\) iff
  for every sequence \((x_n)\) with \(x_n \to a\) and \(x_n \neq a\), we
  have \(f(x_n) \to L\)
\item
  Key observation: \(f\) is continuous at \(a\) if and only if
  \(\lim_{x \to a} f(x) = f(a)\)

  \begin{itemize}
  \tightlist
  \item
    Continuity = the function equals its limit
  \end{itemize}
\item
  ✎ Inline: Prove that \(\lim_{x \to 0} x \sin(1/x) = 0\)
\end{itemize}

\subsection{One-Sided Limits}\label{one-sided-limits}

\begin{itemize}
\tightlist
\item
  Definition: \(\lim_{x \to a^+} f(x) = L\) if the limit condition holds
  with \(0 < x - a < \delta\) (approaching from the right)
\item
  Definition: \(\lim_{x \to a^-} f(x) = L\) similarly (approaching from
  the left)
\item
  Theorem: \(\lim_{x \to a} f(x) = L\) iff \(\lim_{x \to a^+} f(x) = L\)
  and \(\lim_{x \to a^-} f(x) = L\)
\item
  Example: \(\text{sgn}(x)\) at \(0\): \(\lim_{x \to 0^+} = 1\),
  \(\lim_{x \to 0^-} = -1\); since these disagree,
  \(\lim_{x \to 0} \text{sgn}(x)\) does not exist
\end{itemize}

\subsection{Guided Exercises}\label{guided-exercises-13}

\subsection{A. Classifying
Discontinuities}\label{a.-classifying-discontinuities}

We've seen that a function can fail to be continuous in different ways.
In this exercise, you'll organize the possibilities into a taxonomy.

\begin{enumerate}
\def\labelenumi{(\alph{enumi})}
\tightlist
\item
  Consider the following discontinuous functions at the indicated point.
  Compute \(\lim_{x \to a^-} f(x)\) and \(\lim_{x \to a^+} f(x)\) for
  each (if they exist):
\end{enumerate}

\begin{itemize}
\tightlist
\item
  \(f(x) = \sin(x)/x\) at \(a = 0\) (set \(f(0) = 17\))
\item
  \(\text{sgn}(x)\) at \(a = 0\)
\item
  \(\lfloor x \rfloor\) (the floor function) at \(a = 1\)
\item
  \(\sin(1/x)\) at \(a = 0\) (set \(f(0) = 0\))
\item
  Dirichlet's function \(D(x)\) at \(a = 0\)
\end{itemize}

\begin{enumerate}
\def\labelenumi{(\alph{enumi})}
\setcounter{enumi}{1}
\item
  Sort these into three groups based on what goes wrong. In the mildest
  case, both one-sided limits exist and agree --- the function just has
  the ``wrong'' value. In the middle case, both one-sided limits exist
  but disagree. In the worst case, at least one one-sided limit fails to
  exist. Propose names for the three types.
\item
  For each type, explain: can the discontinuity be ``fixed'' by
  redefining \(f(a)\)? If so, what should the new value be? If not, why
  not?
\item
  We say \(f\) has a \emph{removable discontinuity} at \(a\) if
  \(\lim_{x \to a} f(x)\) exists but \(\neq f(a)\) (or \(f\) is
  undefined); a \emph{jump discontinuity} if both one-sided limits exist
  but differ; and an \emph{essential discontinuity} if at least one
  one-sided limit fails to exist. Verify that your grouping matches
  these definitions.
\item
  The \emph{jump} at \(a\) is defined as
  \(\lim_{x \to a^+} f(x) - \lim_{x \to a^-} f(x)\). Compute the jump of
  \(\text{sgn}\) at \(0\) and of \(\lfloor x \rfloor\) at \(1\).
\item
  Classify the discontinuity of \(f(x) = x \cdot D(x)\) at \(x = 0\).
  (This is a trick question --- think carefully.)
\end{enumerate}

\subsection{B. The Ruler Function (Thomae's
function)}\label{b.-the-ruler-function-thomaes-function}

Define \(R: [0,1] \to \RR\) by:
\[R(x) = \begin{cases} 1/q & \text{if } x = p/q \text{ in lowest terms, } q \geq 1 \\ 0 & \text{if } x \notin \QQ \end{cases}\]
(with \(R(0) = 1\)). This function is \(1/2\) at \(1/2\); \(1/3\) at
\(1/3\) and \(2/3\); \(1/4\) at \(1/4\) and \(3/4\); and so on. It looks
like an upside-down ruler --- hence the name.

\emph{Part I: Discontinuous at rationals}

\begin{enumerate}
\def\labelenumi{(\alph{enumi})}
\item
  Let \(a = p/q \in \QQ\) with \(q \geq 1\). Find a sequence of
  irrationals \(x_n \to a\).
\item
  Show \(R(x_n) = 0 \not\to 1/q = R(a)\), so \(R\) is discontinuous at
  \(a\).
\item
  Using the taxonomy from Guided Exercise A, classify the discontinuity
  at \(a\).
\end{enumerate}

\emph{Part II: Continuous at irrationals}

\begin{enumerate}
\def\labelenumi{(\alph{enumi})}
\setcounter{enumi}{3}
\item
  Let \(a \notin \QQ\) and let \(\varepsilon > 0\). Explain why there
  are only finitely many rationals \(p/q\) in \([0,1]\) with
  \(q < 1/\varepsilon\).
\item
  Conclude that there exists \(\delta > 0\) such that every rational
  \(p/q\) in \((a - \delta, a + \delta)\) has \(q \geq 1/\varepsilon\).
\item
  Show: for \(|x - a| < \delta\), we have
  \(|R(x) - R(a)| = |R(x)| \leq \varepsilon\) (whether \(x\) is rational
  or irrational).
\item
  Conclude \(R\) is continuous at every irrational.
\end{enumerate}

\emph{Part III: Reflection}

\begin{enumerate}
\def\labelenumi{(\alph{enumi})}
\setcounter{enumi}{7}
\item
  The ruler function is continuous at every irrational and discontinuous
  at every rational. So continuity can be a \emph{pointwise} property
  --- a function can be continuous at some points and discontinuous at
  others.
\item
  Can a function be continuous at every rational and discontinuous at
  every irrational? (We'll revisit this question in later chapters.)
\end{enumerate}

\subsection{Exercises}\label{exercises-24}

\emph{From the \(\varepsilon\)-\(\delta\) definition} - Prove
\(f(x) = 3x + 1\) is continuous at \(a = 2\) - Prove \(f(x) = x^3\) is
continuous at \(a = 1\) - Prove \(f(x) = 1/x\) is continuous at any
\(a > 0\) - Prove \(f(x) = \sqrt{x}\) is continuous at any \(a > 0\).
What happens at \(a = 0\)?

\emph{Using the sequential characterization} - Prove: if \(f\) is
continuous at \(a\) and \(f(a) > 0\), then there exists \(\delta > 0\)
such that \(f(x) > 0\) for all \(x\) with \(|x - a| < \delta\) (sign
preservation) - Prove: if \(f\) is continuous at \(a\) and
\(f(a) \neq 0\), then \(1/f\) is continuous at \(a\) - Show Dirichlet's
function is discontinuous at every point (done in text, but prove it
independently) - ★ Define \(f(x) = x\) for \(x \in \QQ\), \(f(x) = 0\)
for \(x \notin \QQ\). At which points is \(f\) continuous?

\emph{Limits of functions} - Prove
\(\lim_{x \to 1} \frac{x^2 - 1}{x - 1} = 2\) - Prove
\(\lim_{x \to 0} x \sin(1/x) = 0\) - Prove \(\lim_{x \to 0} \sin(1/x)\)
does not exist - Prove: if \(\lim_{x \to a} f(x) = L\) and
\(\lim_{x \to a} f(x) = M\), then \(L = M\) (uniqueness) - Prove:
\(\lim_{x \to a} f(x) = L\) iff \(\lim_{x \to a^+} f(x) = L\) and
\(\lim_{x \to a^-} f(x) = L\)

\emph{Theoretical} - Prove: \(f\) continuous at \(a\) iff for every
\(\varepsilon > 0\), the set \(\{x : |f(x) - f(a)| < \varepsilon\}\)
contains an interval around \(a\) - Prove: if \(f\) and \(g\) are both
continuous at \(a\), then \(\max(f, g)\) is continuous at \(a\) (hint:
use \(\max(s, t) = (s + t + |s - t|)/2\); but you'll need §6.2 for the
full proof) - If \(f\) and \(g\) are continuous and \(f = g\) on a dense
set, prove \(f = g\) everywhere

\subsection{Dependencies}\label{dependencies-24}

\textbf{Requires}: §1.5 (density of \(\QQ\) and irrationals), §2.2
(\(\varepsilon\)-\(N\) definition of convergence), §2.3 (limit laws)

\textbf{Used in}: §6.2 (algebra of continuous functions uses sequential
characterization + limit laws), §6.3 (power series continuity), §6.4
(IVT, EVT), §6.5 (uniform continuity refines the
\(\varepsilon\)-\(\delta\) definition), Ch 7 (\(\lim_{x \to a}\) is the
basis for the derivative)

\section{§6.2 Building Continuous
Functions}\label{building-continuous-functions}

\subsection{Narrative}\label{narrative-25}

\begin{itemize}
\tightlist
\item
  §6.1 defined continuity and gave us two tools: ε-δ and the sequential
  characterization
\item
  The sequential characterization now pays off: proving algebraic
  combinations are continuous reduces to invoking limit laws from
  Chapter 2
\item
  Sums, products, quotients, compositions of continuous functions are
  continuous
\item
  This lets us build a large library of continuous functions from a
  small set of base cases
\item
  Constants and \(f(x) = x\) are continuous; everything else follows by
  algebra and composition
\item
  Polynomials, rational functions, \(|f|\), \(\max(f,g)\), \(\min(f,g)\)
  --- all continuous
\item
  The proof technique: reduce to sequences, apply known limit laws, done
\end{itemize}

\subsection{Content}\label{content-25}

\subsection{The Sequential Shortcut}\label{the-sequential-shortcut}

\begin{itemize}
\tightlist
\item
  Observation: proving continuity from ε-δ requires estimation (as we
  saw for \(x^2\) in §6.1)
\item
  The sequential characterization offers an alternative: to show \(f\)
  is continuous at \(a\), show that
  \(x_n \to a \Rightarrow f(x_n) \to f(a)\)
\item
  If we already know facts about \emph{sequence} limits (Chapter 2), we
  can use them for \emph{function} continuity
\item
  This is the payoff of the bridge theorem
\end{itemize}

\subsection{Algebra of Continuous
Functions}\label{algebra-of-continuous-functions}

\begin{itemize}
\tightlist
\item
  Theorem: if \(f\) and \(g\) are continuous at \(a\), then so are
  \(f + g\), \(f - g\), \(f \cdot g\), and (if \(g(a) \neq 0\)) \(f/g\)
\item
  Proof (for sum): Let \(x_n \to a\). Since \(f\) continuous at \(a\),
  \(f(x_n) \to f(a)\). Since \(g\) continuous at \(a\),
  \(g(x_n) \to g(a)\). By the sum law for sequences (§2.3),
  \((f+g)(x_n) = f(x_n) + g(x_n) \to f(a) + g(a) = (f+g)(a)\). Done.
\item
  Remark: the proofs for difference, product, and quotient are identical
  --- just invoke the corresponding sequence law from §2.3
\item
  ✎ Inline: Write the proof for the product case
\item
  Remark: compare this with what a direct ε-δ proof of ``sum of
  continuous functions is continuous'' would look like. The sequential
  approach is dramatically cleaner.
\end{itemize}

\subsection{Composition}\label{composition}

\begin{itemize}
\tightlist
\item
  Theorem: if \(g\) is continuous at \(a\) and \(f\) is continuous at
  \(g(a)\), then \(f \circ g\) is continuous at \(a\)
\item
  Proof: Let \(x_n \to a\). Since \(g\) is continuous at \(a\),
  \(g(x_n) \to g(a)\). Since \(f\) is continuous at \(g(a)\) and
  \(g(x_n) \to g(a)\), we have \(f(g(x_n)) \to f(g(a))\). So
  \(f \circ g\) is continuous at \(a\).
\item
  Remark: the sequential proof is almost embarrassingly simple. The ε-δ
  proof requires chaining two ε-δ arguments (exercise).
\item
  ✎ Inline: Write a direct ε-δ proof that \(f \circ g\) is continuous at
  \(a\) (given \(\varepsilon\), get \(\delta_1\) for \(f\) at \(g(a)\),
  then get \(\delta_2\) for \(g\) at \(a\) with tolerance \(\delta_1\);
  take \(\delta = \delta_2\))
\end{itemize}

\subsection{Corollaries}\label{corollaries}

\begin{itemize}
\tightlist
\item
  Proposition: if \(f\) is continuous at \(a\), then \(|f|\) is
  continuous at \(a\)

  \begin{itemize}
  \tightlist
  \item
    Proof: \(|f| = h \circ f\) where \(h(t) = |t|\); \(h\) is continuous
    (exercise or prove: \(||x| - |a|| \leq |x - a|\))
  \end{itemize}
\item
  Proposition: if \(f\) and \(g\) are continuous at \(a\), then so are
  \(\max(f, g)\) and \(\min(f, g)\)

  \begin{itemize}
  \tightlist
  \item
    Proof: \(\max(f, g) = \frac{(f + g) + |f - g|}{2}\) (sum,
    difference, absolute value, scalar multiple --- all continuous)
  \item
    Similarly \(\min(f, g) = \frac{(f + g) - |f - g|}{2}\)
  \end{itemize}
\item
  Remark: the max/min identities are worth remembering --- they reduce
  order operations to algebraic ones
\end{itemize}

\subsection{Building from the Ground
Up}\label{building-from-the-ground-up}

\begin{itemize}
\tightlist
\item
  Example: \(f(x) = c\) (constant) is continuous everywhere

  \begin{itemize}
  \tightlist
  \item
    \(|f(x) - f(a)| = |c - c| = 0 < \varepsilon\) for any \(\delta\)
  \end{itemize}
\item
  Example: \(f(x) = x\) (identity) is continuous everywhere

  \begin{itemize}
  \tightlist
  \item
    \(|f(x) - f(a)| = |x - a| < \varepsilon\) with
    \(\delta = \varepsilon\)
  \end{itemize}
\item
  Proposition: every polynomial
  \(p(x) = a_n x^n + \cdots + a_1 x + a_0\) is continuous on \(\RR\)

  \begin{itemize}
  \tightlist
  \item
    Proof: \(x\) is continuous; by the product law,
    \(x^k = x \cdot x \cdots x\) is continuous; by the scalar multiple
    and sum laws, \(p(x)\) is continuous
  \end{itemize}
\item
  Corollary: every rational function \(p(x)/q(x)\) is continuous
  wherever \(q(x) \neq 0\)

  \begin{itemize}
  \tightlist
  \item
    Proof: quotient of continuous functions
  \end{itemize}
\item
  Example: \(f(x) = \sqrt{x}\) is continuous on \([0, \infty)\)

  \begin{itemize}
  \tightlist
  \item
    At \(a > 0\):
    \(|\sqrt{x} - \sqrt{a}| = \frac{|x - a|}{\sqrt{x} + \sqrt{a}} \leq \frac{|x - a|}{\sqrt{a}}\);
    take \(\delta = \varepsilon \sqrt{a}\)
  \item
    At \(a = 0\): \(|\sqrt{x} - 0| = \sqrt{x} < \varepsilon\) when
    \(x < \varepsilon^2\); take \(\delta = \varepsilon^2\)
  \item
    Remark: the conjugate trick again --- same technique as
    \(n^{1/n} \to 1\) in §2.3
  \item
    Remark: this is one case where ε-δ is cleaner than sequences. The
    conjugate gives the estimate directly.
  \end{itemize}
\item
  Corollary: \(f(x) = \sqrt{p(x)}\) is continuous wherever \(p(x) > 0\)
  (composition of continuous functions)
\end{itemize}

\subsection{Guided Exercise}\label{guided-exercise-11}

\subsection{Piecewise Functions and
Continuity}\label{piecewise-functions-and-continuity}

Continuous functions on overlapping or adjacent domains can be ``glued''
to build new continuous functions. This exercise explores when and how.

\begin{enumerate}
\def\labelenumi{(\alph{enumi})}
\item
  Define \(f: \RR \to \RR\) by \(f(x) = x^2\) for \(x \leq 1\) and
  \(f(x) = 2x - 1\) for \(x > 1\). Is \(f\) continuous at \(x = 1\)?
  Prove your answer using the sequential characterization. (Hint: any
  sequence \(x_n \to 1\) can be split into subsequences with
  \(x_{n_k} \leq 1\) and \(x_{n_j} > 1\).)
\item
  Define \(g: \RR \to \RR\) by \(g(x) = x^2\) for \(x \leq 1\) and
  \(g(x) = 3x - 1\) for \(x > 1\). Is \(g\) continuous at \(x = 1\)?
  Prove your answer.
\item
  Formulate a general theorem: suppose \(h(x) = f(x)\) for \(x \leq c\)
  and \(h(x) = g(x)\) for \(x > c\), where \(f\) and \(g\) are each
  continuous. Under what condition on \(f\), \(g\), and \(c\) is \(h\)
  continuous at \(c\)?
\item
  (Gluing Lemma) Suppose \(f\) is continuous on \([a, b]\) and \(g\) is
  continuous on \([b, c]\), with \(f(b) = g(b)\). Define \(h\) to be
  \(f\) on \([a, b]\) and \(g\) on \((b, c]\). Prove \(h\) is continuous
  on \([a, c]\).
\item
  Generalize: if \(f_1, f_2, \ldots, f_n\) are continuous on adjacent
  intervals \([a_0, a_1], [a_1, a_2], \ldots, [a_{n-1}, a_n]\) and agree
  at shared endpoints, prove the glued function is continuous on
  \([a_0, a_n]\).
\end{enumerate}

\subsection{Exercises}\label{exercises-25}

\emph{Algebra of continuous functions} - Prove: the product of
continuous functions is continuous (fill in the proof from the inline
exercise) - Prove: the quotient of continuous functions is continuous
(where denominator \(\neq 0\)) - Prove: \(|f|\) continuous from the
reverse triangle inequality \(||x| - |a|| \leq |x - a|\), without using
composition - Prove \(\max(f, g) = (f + g + |f - g|)/2\) and conclude
\(\max(f,g)\) is continuous when \(f\) and \(g\) are - Give an example
where \(f + g\) is continuous but neither \(f\) nor \(g\) is continuous

\emph{ε-δ reproofs of the algebra} - Give a direct ε-δ proof that
\(f + g\) is continuous at \(a\). (Hint: use \(\varepsilon/2\) for
each.) - Give a direct ε-δ proof that \(fg\) is continuous at \(a\).
(Hint: write
\(f(x)g(x) - f(a)g(a) = f(x)[g(x) - g(a)] + g(a)[f(x) - f(a)]\). You'll
need to know \(f\) is bounded near \(a\) --- why is this true?) - Give a
direct ε-δ proof that \(1/g\) is continuous at \(a\) when
\(g(a) \neq 0\). (Hint: first show \(|g(x)| \geq |g(a)|/2\) for \(x\)
near \(a\), then estimate
\(|1/g(x) - 1/g(a)| = |g(x) - g(a)| / |g(x)||g(a)|\).)

\emph{Building continuous functions} - Prove \(f(x) = x^n\) is
continuous on \(\RR\) for every \(n \in \NN\) - Prove
\(f(x) = 1/(1 + x^2)\) is continuous on \(\RR\) - Where is
\(f(x) = \sqrt{x^2 - 4}\) continuous? Prove it. - Prove: if \(f\) is
continuous and never zero on an interval, then either \(f > 0\) on the
whole interval or \(f < 0\) on the whole interval (hint: sign
preservation from §6.1 exercises)

\emph{Composition} - Give a direct ε-δ proof of the composition theorem
- Prove: if \(f\) is continuous at \(a\) and \(g\) is continuous at
\(a\), then \(h(x) = f(g(x) \cdot x)\) is continuous at \(a\) (assuming
\(f\) is continuous at \(g(a) \cdot a\)) - ★ Find continuous functions
\(f\) and \(g\) on \(\RR\) such that \(f \circ g\) is continuous but
\(g \circ f\) is not (impossible --- explain why)

\emph{Theoretical} - Prove: the set of continuous functions on \(\RR\)
is closed under addition and multiplication (i.e., forms a ring) - ★
Prove: if \(f\) is continuous at \(a\) and \(f(a) \neq 0\), then \(f\)
is bounded away from zero near \(a\): there exist \(\delta > 0\) and
\(c > 0\) with \(|f(x)| \geq c\) for \(|x - a| < \delta\) - ★ Suppose
\(f\) is continuous on \(\RR\) and \(f(r) = 0\) for every \(r \in \QQ\).
Prove \(f(x) = 0\) for all \(x \in \RR\). (Hint: density + sequential
characterization)

\subsection{Dependencies}\label{dependencies-25}

\textbf{Requires}: §6.1 (continuity definition, sequential
characterization), §2.3 (limit laws for sequences)

\textbf{Used in}: §6.3 (continuity of power series uses algebra), §6.4
(IVT/EVT proofs use these closure properties), §7.1 (differentiation
rules parallel these)

\section{§6.3 Power Series and
Continuity}\label{power-series-and-continuity}

\subsection{Narrative}\label{narrative-26}

\begin{itemize}
\tightlist
\item
  We know \(\sum x^n = 1/(1-x)\) --- but the left side \emph{defines a
  function} from an infinite process
\item
  General power series \(\sum a_n x^n\): a new kind of object, defined
  as a limit of polynomials
\item
  For which \(x\) does it converge? The Cauchy-Hadamard formula gives a
  sharp answer; the ratio test gives a practical one
\item
  At the boundary, anything can happen
\item
  On the interior, each partial sum is a polynomial (continuous), so the
  limit is a limit of continuous functions --- is it continuous?
\item
  Warning: no, in general (zigzag example)
\item
  But dominated convergence saves power series on the interior
\item
  The boundary is still dangerous (\(x^n\) on \([0,1]\)), but Abel's
  theorem says: if the series converges at the endpoint, continuity
  extends there
\end{itemize}

\subsection{Content}\label{content-26}

\subsection{Series as Functions}\label{series-as-functions}

\begin{itemize}
\tightlist
\item
  We already know \(\sum_{n=0}^{\infty} x^n = \frac{1}{1-x}\) for
  \(|x| < 1\)
\item
  New perspective: the left side \emph{defines a function}. Plug in
  \(x\), compute partial sums \(S_N(x) = 1 + x + x^2 + \cdots + x^N\),
  take the limit.
\item
  The partial sums are polynomials --- finite, tame objects we
  understand from §6.2
\item
  The limit is something new: a function defined by an infinite process
\item
  Definition: A \emph{power series} (centered at \(0\)) is
  \(\sum_{n=0}^{\infty} a_n x^n\); centered at \(c\), it is
  \(\sum_{n=0}^{\infty} a_n (x - c)^n\)
\item
  Where it converges, it defines a function
  \(f(x) = \sum_{n=0}^{\infty} a_n x^n\)
\item
  We've seen another example: the exponential series
  \(E(x) = \sum x^n/n!\) from §5.3, which converges for all \(x\)
\item
  Question: for a general power series, for which \(x\) does it
  converge?
\end{itemize}

\subsection{Convergence}\label{convergence-2}

\begin{itemize}
\tightlist
\item
  Theorem (Cauchy-Hadamard): For \(\sum a_n x^n\), define
  \(R = 1 / \limsup_{n \to \infty} |a_n|^{1/n}\) (with \(R = \infty\) if
  the limsup is \(0\), and \(R = 0\) if the limsup is \(\infty\)). Then:

  \begin{itemize}
  \tightlist
  \item
    \(\sum a_n x^n\) converges absolutely for \(|x| < R\)
  \item
    \(\sum a_n x^n\) diverges for \(|x| > R\)
  \end{itemize}
\item
  This \(R\) is the \emph{radius of convergence} --- it's an
  if-and-only-if characterization
\item
  Proof sketch: apply the root test (§4.2) to \(\sum a_n x^n\). The
  \(n\)-th root of \(|a_n x^n|\) is \(|a_n|^{1/n} |x|\), and the limsup
  is \(|x|/R\). The root test gives convergence when this is \(< 1\)
  (i.e., \(|x| < R\)) and divergence when \(> 1\).
\item
  Remark: the ratio test gives a practical sufficient condition: if
  \(|a_{n+1}/a_n| \to L\), then \(R = 1/L\). This is often easier to
  compute, but the Cauchy-Hadamard formula is the sharp result.
\item
  At \(x = \pm R\): anything can happen (exercise: find examples of all
  four boundary behaviors)
\item
  Examples:

  \begin{itemize}
  \tightlist
  \item
    \(\sum x^n\): \(R = 1\) (geometric series). Diverges at both
    \(x = 1\) and \(x = -1\).
  \item
    \(\sum x^n/n!\): \(R = \infty\) (exponential series). No boundary at
    all.
  \item
    \(\sum n! \, x^n\): \(R = 0\). Converges only at \(x = 0\).
  \item
    \(\sum x^n/n\): \(R = 1\). Diverges at \(x = 1\) (harmonic),
    converges at \(x = -1\) (alternating harmonic).
  \item
    \(\sum x^n/n^2\): \(R = 1\). Converges at both \(x = 1\) and
    \(x = -1\).
  \end{itemize}
\end{itemize}

\subsection{Continuity}\label{continuity-2}

\begin{itemize}
\tightlist
\item
  On \((-R, R)\), the power series defines a function
  \(f(x) = \sum a_n x^n\)
\item
  Each partial sum \(S_N(x) = \sum_{n=0}^{N} a_n x^n\) is a polynomial,
  hence continuous by §6.2
\item
  So \(f\) is a pointwise limit of continuous functions. Is \(f\)
  continuous?
\item
  It would be natural to hope the answer is always yes. It is not.
\item
  Example (zigzag approximation of a step function):

  \begin{itemize}
  \tightlist
  \item
    Define \(f_n: [0, 1] \to \RR\): equals \(0\) on
    \([0, \tfrac{1}{2} - \tfrac{1}{n}]\), rises linearly to \(1\) on
    \([\tfrac{1}{2} - \tfrac{1}{n}, \tfrac{1}{2}]\), equals \(1\) on
    \([\tfrac{1}{2}, 1]\)
  \item
    Each \(f_n\) is continuous (piecewise linear, glued continuously ---
    §6.2 guided exercise)
  \item
    Pointwise limit: step function jumping at \(x = 1/2\).
    Discontinuous.
  \end{itemize}
\item
  Remark: smoothing the ramp (e.g., with a scaled arctangent) gives the
  same phenomenon. The problem is pointwise convergence, not corners.
\item
  So we need conditions that guarantee continuity of the limit. For
  power series, dominated convergence (§5.4) provides them.
\item
  Theorem: If \(\sum a_n x^n\) has radius of convergence \(R > 0\), then
  \(f(x) = \sum a_n x^n\) is continuous on \((-R, R)\).
\item
  Proof sketch:

  \begin{itemize}
  \tightlist
  \item
    Fix \(c\) with \(|c| < R\) and a sequence \(x_k \to c\). Need to
    show \(f(x_k) \to f(c)\).
  \item
    This requires interchanging a limit and a sum:
    \(\lim_k \sum_n a_n x_k^n = \sum_n a_n c^n\)
  \item
    Pick \(\rho\) with \(|c| < \rho < R\). For \(k\) large,
    \(|x_k| \leq \rho\).
  \item
    Domination: \(|a_n x_k^n| \leq |a_n| \rho^n\), and
    \(\sum |a_n| \rho^n\) converges (since \(\rho < R\))
  \item
    Apply dominated convergence (§5.4): interchange is valid. Done.
  \end{itemize}
\item
  Corollary: \(\exp(x) = \sum x^n/n!\) is continuous on \(\RR\)
\item
  ✎ Inline: Verify the three hypotheses of dominated convergence
  explicitly for \(\exp\) at a point \(c\)
\end{itemize}

\subsection{Boundary Behavior}\label{boundary-behavior}

\begin{itemize}
\tightlist
\item
  The dominated convergence argument handles the interior. What about
  the boundary \(x = \pm R\)?
\item
  Example: \(f_n(x) = x^n\) on \([0, 1]\)

  \begin{itemize}
  \tightlist
  \item
    Each \(f_n\) is continuous
  \item
    Pointwise limit: \(0\) on \([0, 1)\), and \(1\) at \(x = 1\) ---
    discontinuous
  \item
    This IS a power series situation: the geometric series converges on
    \((-1, 1)\) but diverges at \(x = 1\)
  \end{itemize}
\item
  So even for power series, the boundary can produce discontinuities.
  But only when the series diverges there.
\item
  Lemma (Summation by parts): For sequences \((a_k)\) and \((b_k)\) with
  partial sums \(B_n = \sum_{k=0}^{n} b_k\):
  \[\sum_{k=0}^{n} a_k b_k = a_n B_n - \sum_{k=0}^{n-1} (a_{k+1} - a_k) B_k\]
\item
  Proof: substitute \(b_k = B_k - B_{k-1}\) and rearrange (discrete
  analogue of \(\int u \, dv = uv - \int v \, du\))
\item
  ✎ Inline: Verify the formula for \(n = 2\) by expanding both sides
\item
  Theorem (Abel): If \(\sum_{n=0}^{\infty} a_n\) converges to \(S\),
  then \[\lim_{x \to 1^-} \sum_{n=0}^{\infty} a_n x^n = S\]
\item
  In words: if the series converges at \(x = 1\), the power series
  extends continuously to \(x = 1\) from the left
\item
  Proof sketch:

  \begin{itemize}
  \tightlist
  \item
    Write \(S_N = \sum_{n=0}^{N} a_n\) (partial sums, \(S_N \to S\))
  \item
    Summation by parts:
    \(\sum_{n=0}^{N} a_n x^n = S_N x^N + (1 - x) \sum_{n=0}^{N-1} S_n x^n\)
  \item
    Let \(N \to \infty\): since \(|x| < 1\), \(S_N x^N \to 0\), so
    \(f(x) = (1 - x) \sum_{n=0}^{\infty} S_n x^n\)
  \item
    Write \(S_n = S + \varepsilon_n\) where \(\varepsilon_n \to 0\)
  \item
    Then
    \(f(x) = S \cdot (1 - x) \sum x^n + (1 - x) \sum \varepsilon_n x^n = S + (1 - x) \sum \varepsilon_n x^n\)
  \item
    Show \((1 - x) \sum \varepsilon_n x^n \to 0\) as \(x \to 1^-\):
    split at a threshold \(N\) where \(|\varepsilon_n| < \varepsilon\);
    the tail is bounded by \(\varepsilon\), the head vanishes because of
    the \((1-x)\) factor
  \end{itemize}
\item
  Remark: the \(x^n\) counterexample fails Abel's hypothesis ---
  \(\sum 1^n\) diverges --- so the discontinuity is consistent
\item
  Remark: generalizes to any power series centered at \(c\) with radius
  \(R\): convergence at \(c + R\) gives left-continuity there,
  convergence at \(c - R\) gives right-continuity
\end{itemize}

\subsection{Guided Exercise}\label{guided-exercise-12}

\subsection{Verifying Dominated Convergence for the Exponential
Series}\label{verifying-dominated-convergence-for-the-exponential-series}

We proved that power series are continuous on their interior using
dominated convergence. In this exercise, you verify every hypothesis
explicitly for \(\exp(x) = \sum_{n=0}^{\infty} x^n / n!\).

\begin{enumerate}
\def\labelenumi{(\alph{enumi})}
\item
  Fix \(c \in \RR\) and a sequence \(x_k \to c\). Explain why there
  exists \(\rho > 0\) such that \(|x_k| \leq \rho\) for all \(k\).
\item
  Write down the terms \(a_n(k) = x_k^n / n!\) that play the role of the
  summands in dominated convergence. For each fixed \(n\), compute
  \(\lim_{k \to \infty} a_n(k)\).
\item
  Find a dominating sequence \(M_n\) independent of \(k\) such that
  \(|a_n(k)| \leq M_n\) for all \(k\) and \(\sum M_n < \infty\).
\item
  State the conclusion of dominated convergence and interpret: what does
  it say about \(\exp\)?
\item
  Now consider boundary behavior. The exponential series has
  \(R = \infty\), so there is no boundary. Give an example of a power
  series with \(R = 1\) where the dominated convergence argument works
  on \((-1, 1)\) but says nothing about \(x = 1\).
\item
  For \(\sum x^n / n\) at \(x = -1\): the series \(\sum (-1)^n / n\)
  converges (alternating series test). What does Abel's theorem tell you
  about \(\lim_{x \to (-1)^+} \sum x^n / n\)?
\end{enumerate}

\subsection{Exercises}\label{exercises-26}

\emph{Convergence} - Find the radius of convergence using the ratio
test: \(\sum n x^n\), \(\sum x^n / n^2\), \(\sum 2^n x^n / n!\),
\(\sum n! \, x^n / n^n\) - Prove: if \(\sum a_n x_0^n\) converges for
some \(x_0 \neq 0\), then \(\sum a_n x^n\) converges absolutely for all
\(|x| < |x_0|\) - (Boundary zoo) Find power series with radius \(R = 1\)
that: (i) diverge at both \(x = 1\) and \(x = -1\); (ii) converge at
both; (iii) converge at \(x = 1\) but diverge at \(x = -1\); (iv)
converge at \(x = -1\) but diverge at \(x = 1\) - ★ Prove the
Cauchy-Hadamard formula: \(1/R = \limsup |a_n|^{1/n}\)

\emph{Continuity of power series} - Verify explicitly that
\(f(x) = \sum x^n / n^2\) is continuous at \(x = 1/2\): identify
\(\rho\), \(M_n\), check \(\sum M_n < \infty\) - Prove:
\(f(x) = \sum x^n / n!\) is continuous at \(x = 0\) directly from
\(\varepsilon\)-\(\delta\) (without dominated convergence) - Prove: if
\(f(x) = \sum a_n x^n\) and \(g(x) = \sum b_n x^n\) both have radius
\(\geq R\), then \(f + g\) is continuous on \((-R, R)\). What about
\(f \cdot g\)?

\emph{Limits of continuous functions} - Verify directly that
\(f_n(x) = x^n\) on \([0, 1]\) converges pointwise to a discontinuous
limit - Construct a sequence of continuous functions on \([0,1]\)
converging pointwise to Dirichlet's function restricted to \([0,1]\)
(hint: this is impossible --- explain why using cardinality or
measurability, or simply argue from properties of pointwise limits of
continuous functions) - ★ Give an example of continuous
\(f_n: [0, 1] \to \RR\) such that \(\lim f_n(x)\) exists for all \(x\)
but is discontinuous at every point of some interval

\emph{Abel's theorem applications} - Use Abel's theorem to show
\(\ln(2) = 1 - 1/2 + 1/3 - 1/4 + \cdots\) (students may use
\(\ln(2) = \int_0^1 1/(1+x)\,dx\) from calculus, or treat as forward
reference to ch8) - Prove \(\sum x^n/n\) is continuous on \((-1, 1)\)
and left-continuous at \(x = 1\) - If \(\sum a_n x^n = 0\) for all
\(x \in (-r, r)\), prove \(a_n = 0\) for all \(n\)

\emph{Summation by parts and Abel's theorem} - Prove the summation by
parts formula by induction on \(n\) - (Abel's test) Prove: if
\(\sum b_n\) converges and \((a_n)\) is monotone and bounded, then
\(\sum a_n b_n\) converges - (Dirichlet's test) Prove: if the partial
sums \(\sum_{n=0}^{N} b_n\) are bounded and \(a_n \to 0\) monotonically,
then \(\sum a_n b_n\) converges - Use Dirichlet's test to give a new
proof that \(\sum (-1)^n / n\) converges - ★ (Tauber's converse) Prove:
if \(\sum a_n x^n \to S\) as \(x \to 1^-\) AND \(a_n = O(1/n)\), then
\(\sum a_n\) converges to \(S\). (This is a partial converse to Abel's
theorem --- Abel requires convergence to get continuity, Tauber recovers
convergence from continuity plus a growth condition.)

\subsection{Dependencies}\label{dependencies-26}

\textbf{Requires}: §3.2 (limsup), §4.2 (ratio/root tests), §5.4
(dominated convergence), §6.2 (polynomials are continuous)

\textbf{Used in}: Ch 7 (term-by-term differentiation --- Power Series
II), Ch 9 (evaluating power series at endpoints, continuity of
elementary functions)

\section{§6.4 Global Properties}\label{global-properties}

\subsection{Narrative}\label{narrative-27}

\begin{itemize}
\tightlist
\item
  The local theory is done: we know what continuity means at a point,
  how to build continuous functions, and that power series are
  continuous
\item
  Now: what happens when continuity interacts with the domain?
\item
  Continuity alone is local --- global conclusions require continuity
  \emph{plus} structure (intervals, closed and bounded sets,
  monotonicity)
\item
  IVT: continuity + interval → image is an interval (no gaps)
\item
  EVT: continuity + closed bounded interval → bounded, attains max and
  min
\item
  Unification: continuous image of \([a,b]\) is \([m, M]\) --- a closed
  bounded interval
\item
  Monotone continuous functions have continuous inverses --- the tool
  that will make \(\log\), \(\arcsin\), \(\arctan\) automatic in Chapter
  9
\end{itemize}

\subsection{Content}\label{content-27}

\subsection{The Intermediate Value
Theorem}\label{the-intermediate-value-theorem}

\begin{itemize}
\tightlist
\item
  Motivating question: if \(f\) is continuous on \([a,b]\) with
  \(f(a) < 0\) and \(f(b) > 0\), must \(f\) have a zero somewhere in
  between?
\item
  For polynomials this feels obvious --- the graph can't ``jump'' over
  the axis
\item
  But we've seen that general functions CAN jump (step functions,
  Dirichlet). The hypothesis doing the work is continuity.
\item
  Theorem (IVT): If \(f\) is continuous on \([a,b]\) and \(y\) is any
  value between \(f(a)\) and \(f(b)\), then \(f(c) = y\) for some
  \(c \in [a,b]\).
\item
  Proof sketch (via completeness):

  \begin{itemize}
  \tightlist
  \item
    WLOG \(f(a) < y < f(b)\)
  \item
    Let \(c = \sup\{x \in [a,b] : f(x) < y\}\)
  \item
    This set is nonempty (\(a\) is in it) and bounded above (by \(b\)),
    so the sup exists by completeness
  \item
    Show \(f(c) = y\): if \(f(c) < y\), continuity gives a neighborhood
    where \(f < y\), contradicting that \(c\) is the sup; if
    \(f(c) > y\), continuity gives a neighborhood where \(f > y\),
    contradicting that points just below \(c\) have \(f < y\)
  \end{itemize}
\item
  Remark: completeness is essential --- the function \(f(x) = x^2 - 2\)
  on \(\QQ\) is continuous, \(f(0) < 0\), \(f(2) > 0\), but there is no
  rational \(c\) with \(f(c) = 0\)
\item
  Corollary: the continuous image of an interval is an interval

  \begin{itemize}
  \tightlist
  \item
    If \(f\) is continuous on an interval \(I\) and \(f\) takes values
    \(p\) and \(q\), it takes every value between them
  \end{itemize}
\end{itemize}

\emph{Applications} - Application (nth roots exist): For \(c > 0\) and
\(n \in \NN\), the equation \(x^n = c\) has a positive solution -
\(f(x) = x^n\) is continuous, \(f(0) = 0 < c\), and \(f(M) > c\) for
\(M\) large enough - By IVT, \(f(x_0) = c\) for some \(x_0 \in (0, M)\)
- Application (Fixed point theorem): If \(f: [a,b] \to [a,b]\) is
continuous, then \(f(c) = c\) for some \(c\) - Consider
\(g(x) = f(x) - x\); then \(g(a) = f(a) - a \geq 0\) and
\(g(b) = f(b) - b \leq 0\) - By IVT, \(g(c) = 0\) for some \(c\) -
Application: Every odd-degree polynomial has a real root -
\(p(x) \to +\infty\) as \(x \to +\infty\) and \(p(x) \to -\infty\) as
\(x \to -\infty\) (or vice versa) - By IVT, \(p\) takes the value \(0\)
- ✎ Inline: Use IVT to prove \(x^3 + x = 1\) has exactly one solution in
\((0, 1)\) (existence via IVT; uniqueness because \(x^3 + x\) is
strictly increasing)

\subsection{Boundedness and Extreme
Values}\label{boundedness-and-extreme-values}

\begin{itemize}
\tightlist
\item
  Can a continuous function on \([a,b]\) be unbounded? On an open
  interval it can: \(f(x) = 1/x\) on \((0,1)\)
\item
  But on a \emph{closed bounded} interval, continuity forces boundedness
\item
  Theorem (Boundedness): If \(f\) is continuous on \([a,b]\), then \(f\)
  is bounded on \([a,b]\)
\item
  Proof sketch (by contradiction via Bolzano-Weierstrass):

  \begin{itemize}
  \tightlist
  \item
    If not, there exist \(x_n \in [a,b]\) with \(|f(x_n)| > n\)
  \item
    By Bolzano-Weierstrass (§3.1), \((x_n)\) has a convergent
    subsequence \(x_{n_k} \to c \in [a,b]\)
  \item
    By continuity, \(f(x_{n_k}) \to f(c)\), so \(f(x_{n_k})\) is bounded
    --- contradiction
  \end{itemize}
\item
  Theorem (Extreme Value Theorem): If \(f\) is continuous on \([a,b]\),
  then \(f\) attains its maximum and minimum. That is, there exist
  \(c, d \in [a,b]\) with \(f(c) \leq f(x) \leq f(d)\) for all
  \(x \in [a,b]\).
\item
  Proof sketch (for the maximum):

  \begin{itemize}
  \tightlist
  \item
    \(f\) is bounded, so \(M = \sup\{f(x) : x \in [a,b]\}\) exists
  \item
    Choose \(x_n\) with \(f(x_n) > M - 1/n\)
  \item
    By BW, extract \(x_{n_k} \to d \in [a,b]\)
  \item
    By continuity, \(f(x_{n_k}) \to f(d)\), and by construction
    \(f(x_{n_k}) \to M\)
  \item
    Therefore \(f(d) = M\)
  \end{itemize}
\item
  Remark: closed AND bounded are both needed. \(f(x) = x\) on
  \([0, \infty)\) is continuous, closed domain, unbounded.
  \(f(x) = 1/x\) on \((0,1)\) is continuous, bounded domain, but doesn't
  attain its sup.
\item
  Corollary (Continuous image of \([a,b]\)): If \(f\) is continuous on
  \([a,b]\), then \(f([a,b]) = [m, M]\) where \(m\) and \(M\) are the
  minimum and maximum values of \(f\).

  \begin{itemize}
  \tightlist
  \item
    EVT gives \(m\) and \(M\) are attained; IVT fills in everything
    between.
  \item
    This is the unification: the continuous image of a closed bounded
    interval is a closed bounded interval.
  \end{itemize}
\end{itemize}

\subsection{Monotone Functions and
Inverses}\label{monotone-functions-and-inverses}

\begin{itemize}
\tightlist
\item
  Setup: we want to define functions like \(\log\), \(\arcsin\),
  \(\arctan\) as inverses of continuous functions in Chapter 9. When is
  the inverse continuous?
\item
  Definition: \(f\) is \emph{increasing} on \(S\) if
  \(x < y \Rightarrow f(x) \leq f(y)\); \emph{strictly increasing} if
  \(x < y \Rightarrow f(x) < f(y)\). Similarly for decreasing.
  \emph{Strictly monotone} means strictly increasing or strictly
  decreasing.
\item
  Proposition: Strictly monotone \(\Rightarrow\) injective (one-to-one)

  \begin{itemize}
  \tightlist
  \item
    Proof: immediate from the definition
  \end{itemize}
\item
  Theorem: If \(f\) is strictly monotone and continuous on \([a,b]\),
  then \(f^{-1}\) exists and is continuous on \(f([a,b])\).
\item
  Proof sketch:

  \begin{itemize}
  \tightlist
  \item
    \(f\) is injective (strict monotonicity) and surjective onto its
    range
  \item
    By the corollary above, \(f([a,b]) = [m, M]\) (a closed bounded
    interval)
  \item
    So \(f^{-1}: [m, M] \to [a,b]\) exists
  \item
    \(f^{-1}\) is strictly monotone (same direction as \(f\))
  \item
    Continuity of \(f^{-1}\): suppose \(y_n \to y\) in \([m, M]\). Let
    \(x_n = f^{-1}(y_n)\). The sequence \((x_n)\) is bounded (in
    \([a,b]\)). Any convergent subsequence \(x_{n_k} \to c\) satisfies
    \(f(c) = \lim f(x_{n_k}) = \lim y_{n_k} = y\), so \(c = f^{-1}(y)\).
    Since every convergent subsequence has the same limit,
    \(x_n \to f^{-1}(y)\).
  \end{itemize}
\item
  Application (nth root function): \(f(x) = x^n\) is strictly increasing
  and continuous on \([0, \infty)\). By the theorem,
  \(f^{-1}(x) = x^{1/n}\) is continuous on \([0, \infty)\).

  \begin{itemize}
  \tightlist
  \item
    This completes the nth root story: IVT gave existence, the inverse
    theorem gives continuity
  \end{itemize}
\item
  Remark: in Chapter 9, we define \(\exp(x) = \sum x^n/n!\), show it's
  strictly increasing, and define \(\log = \exp^{-1}\). Continuity of
  \(\log\) is then automatic. Similarly for \(\arcsin\), \(\arctan\).
\item
  Theorem: If \(f\) is monotone (not necessarily continuous) on
  \((a,b)\), then at every point \(c \in (a,b)\), both one-sided limits
  \(\lim_{x \to c^+} f(x)\) and \(\lim_{x \to c^-} f(x)\) exist.
\item
  Proof sketch: for increasing \(f\),
  \(\lim_{x \to c^-} f(x) = \sup\{f(x) : x < c\}\) (bounded above by
  \(f(c)\))
\item
  Corollary: a monotone function can only have jump discontinuities (no
  removable or essential)
\end{itemize}

\subsection{Guided Exercise}\label{guided-exercise-13}

\subsection{The Bisection Method}\label{the-bisection-method}

The IVT proof shows a zero exists but doesn't tell us how to find it.
The bisection method is an algorithm that does.

\begin{enumerate}
\def\labelenumi{(\alph{enumi})}
\item
  Suppose \(f\) is continuous on \([a,b]\) with \(f(a) < 0 < f(b)\). Let
  \(m = (a+b)/2\). Explain why at least one of the intervals \([a, m]\)
  or \([m, b]\) has the property that \(f\) changes sign on it.
\item
  Define sequences \((a_n)\) and \((b_n)\) by repeatedly bisecting:
  \([a_0, b_0] = [a, b]\), and at each step, replace \([a_n, b_n]\) with
  the half where \(f\) changes sign. Show that
  \(b_n - a_n = (b-a)/2^n\).
\item
  Show \((a_n)\) is increasing and bounded above, and \((b_n)\) is
  decreasing and bounded below. Conclude both converge, and
  \(\lim a_n = \lim b_n = c\) for some \(c \in [a,b]\).
\item
  Show \(f(a_n) \leq 0\) and \(f(b_n) \geq 0\) for all \(n\) (or the
  algorithm terminated early at a zero). Use continuity to conclude
  \(f(c) = 0\).
\item
  Apply the bisection method to \(f(x) = x^2 - 2\) on \([1, 2]\).
  Compute \(a_n\) and \(b_n\) for \(n = 0, 1, 2, 3, 4\). How close is
  your estimate to \(\sqrt{2}\)?
\item
  How many steps are needed to approximate \(\sqrt{2}\) to within
  \(10^{-6}\)?
\end{enumerate}

\subsection{Exercises}\label{exercises-27}

\emph{Intermediate Value Theorem} - Use IVT to prove:
\(x^3 + x - 1 = 0\) has a solution in \((0, 1)\) - Use IVT to prove:
\(E(x) = 3x\) has a solution in \((0, 2)\), where \(E(x) = \sum x^n/n!\)
is the exponential function (continuous by §6.3; use \(E(0) = 1\) and
\(E(2) > 4\)) - Prove: if \(f\) is continuous on \([0,1]\) with
\(f(0) = f(1)\), then there exists \(c \in [0, 1/2]\) with
\(f(c) = f(c + 1/2)\). (Hint: consider \(g(x) = f(x) - f(x + 1/2)\).) -
Prove: the equation \(x^3 + 2x = 5\) has exactly one solution (use IVT
for existence, strict monotonicity of \(x^3 + 2x\) for uniqueness) -
Show \(f(x) = x + \sin(x)\) has a unique fixed point (students may use
known properties of \(\sin\) from calculus) - ★ Prove: if \(p(x)\) has
odd degree and real coefficients, then \(p\) has at least one real root

\emph{Extreme Value Theorem} - Give an example showing EVT fails on
\((0, 1)\) (open interval) - Give an example showing EVT fails on
\([0, \infty)\) (unbounded domain) - Give an example of a
\emph{discontinuous} bounded function on \([0,1]\) that does not attain
its supremum - Prove: if \(f\) is continuous on \([a,b]\) and
\(f(x) > 0\) for all \(x \in [a,b]\), then there exists \(c > 0\) with
\(f(x) \geq c\) for all \(x \in [a,b]\) (bounded away from zero --- use
EVT) - ★ Prove: if \(f\) is continuous on \(\RR\) and
\(\lim_{x \to \pm\infty} f(x) = 0\), then \(f\) is bounded and attains
its maximum or minimum (or both)

\emph{Monotone functions and inverses} - Prove: if \(f\) is continuous
and injective on an interval, then \(f\) is strictly monotone. (Hint:
use IVT --- if \(f\) is not monotone, find three points where the
function ``turns around'' and use IVT to find two points with the same
image.) - Verify that \(f(x) = x^3\) is strictly increasing on \(\RR\),
and describe \(f^{-1}\). - Prove: \(f(x) = x^3 + x\) is strictly
increasing on \(\RR\). (Hint: show
\(b^3 - a^3 = (b - a)(b^2 + ab + a^2)\) and verify
\(b^2 + ab + a^2 > 0\) when \(a \neq b\).) - ★ (Countability of monotone
discontinuities) Prove: if \(f\) is monotone on \((a,b)\), then \(f\)
has at most countably many discontinuities. (Hint: at each jump
discontinuity, the one-sided limits define a nonempty open interval of
values that \(f\) skips. These intervals are disjoint. Each contains a
rational.) - ★ Construct a strictly increasing function on \([0,1]\)
that is discontinuous at every rational. (Hint: enumerate the rationals
\(r_1, r_2, \ldots\) and define \(f(x) = \sum_{r_n < x} 1/2^n\).)

\emph{Theoretical} - Prove: the continuous image of a closed bounded
interval is a closed bounded interval (combine IVT + EVT) - ★ Suppose
\(f\) is continuous on \([0,1]\) and \(f(x) \in \QQ\) for all \(x\).
Prove \(f\) is constant. (Hint: if \(f(a) < f(b)\), use density of
irrationals to find \(r \notin \QQ\) with \(f(a) < r < f(b)\), then
apply IVT.) - Prove: if \(f\) is continuous on \([a,b]\) and injective,
then \(f^{-1}\) is continuous (without assuming monotonicity --- but you
may use the fact that continuous injective functions on intervals are
monotone) - ★ Prove: if \(f: \RR \to \RR\) is continuous and injective,
then \(f\) is strictly monotone. (Harder than the interval case --- you
need to rule out ``turning around at infinity.'')

\subsection{Dependencies}\label{dependencies-27}

\textbf{Requires}: §3.1 (Bolzano-Weierstrass), §6.1 (continuity,
one-sided limits), §6.2 (algebra of continuous functions), completeness
of \(\RR\)

\textbf{Used in}: §6.5 (uniform continuity on \([a,b]\)), Ch 7 (Rolle's
theorem, mean value theorem use IVT), Ch 9 (\(\log\), \(\arcsin\),
\(\arctan\) defined as inverses of continuous monotone functions)

\section{§6.5 Uniform Continuity}\label{uniform-continuity}

\subsection{Narrative}\label{narrative-28}

\begin{itemize}
\tightlist
\item
  In §6.1 we noted that \(\delta\) depends on both \(\varepsilon\) and
  the point \(a\) --- now we ask: can one \(\delta\) work for all
  points?
\item
  Uniform continuity: a stronger condition where \(\delta\) depends only
  on \(\varepsilon\)
\item
  Non-examples show this is genuinely stronger: \(1/x\) on \((0,1)\),
  \(x^2\) on \(\RR\)
\item
  The key theorem: on a closed bounded interval, continuity
  automatically upgrades to uniform continuity
\item
  This explains why the non-examples needed open or unbounded domains
\item
  Continuous extensions: uniform continuity on \((a,b)\) is exactly
  what's needed to extend continuously to \([a,b]\)
\end{itemize}

\subsection{Content}\label{content-28}

\subsection{The Definition}\label{the-definition-2}

\begin{itemize}
\tightlist
\item
  In §6.1, the definition of continuity at \(a\) says: for every
  \(\varepsilon > 0\), there exists \(\delta > 0\) such that
  \(|x - a| < \delta \Rightarrow |f(x) - f(a)| < \varepsilon\)
\item
  The \(\delta\) can depend on both \(\varepsilon\) and the point \(a\)
\item
  Example: \(f(x) = x^2\) at \(a\). We showed
  \(\delta = \min(1, \varepsilon/(2|a|+1))\) works. As \(|a|\) grows,
  \(\delta\) shrinks --- the function changes faster far from the
  origin.
\item
  Example: \(f(x) = 1/x\) on \((0,1)\). Near \(x = 0\), the function is
  steep --- \(\delta\) must shrink as \(a \to 0\).
\item
  Question: is this dependence on \(a\) essential, or is it an artifact
  of our proofs? Can we sometimes find one \(\delta\) that works
  uniformly across all points?
\item
  Definition: \(f\) is \emph{uniformly continuous on \(S\)} if:
  \[\forall \varepsilon > 0 \; \exists \delta > 0 \; \forall x, y \in S: |x - y| < \delta \Rightarrow |f(x) - f(y)| < \varepsilon\]
\item
  Remark: the quantifier order matters. In ordinary continuity,
  \(\delta\) comes after the point \(a\) is fixed. In uniform
  continuity, \(\delta\) comes first and must work for ALL pairs
  \(x, y\).
\item
  Proposition: uniformly continuous on \(S\) \(\Rightarrow\) continuous
  at every point of \(S\) (immediate from the definition)
\item
  The converse is false. The non-examples:
\item
  Example: \(f(x) = 1/x\) is NOT uniformly continuous on \((0,1)\)

  \begin{itemize}
  \tightlist
  \item
    Take \(x_n = 1/n\), \(y_n = 1/(n+1)\). Then
    \(|x_n - y_n| = 1/(n(n+1)) \to 0\), but
    \(|f(x_n) - f(y_n)| = |n - (n+1)| = 1 \not\to 0\)
  \item
    No single \(\delta\) works for \(\varepsilon = 1\): no matter how
    small \(\delta\) is, points near \(0\) are close together but their
    function values are far apart
  \end{itemize}
\item
  Example: \(f(x) = x^2\) is NOT uniformly continuous on \(\RR\)

  \begin{itemize}
  \tightlist
  \item
    Take \(x_n = n\), \(y_n = n + 1/n\). Then
    \(|x_n - y_n| = 1/n \to 0\), but
    \(|f(x_n) - f(y_n)| = |2 + 1/n^2| \to 2\)
  \item
    The function grows too fast --- far from the origin, small changes
    in \(x\) produce large changes in \(x^2\)
  \end{itemize}
\item
  Remark: these proofs use a common technique for showing
  non-uniform-continuity: find sequences \(x_n, y_n\) with
  \(|x_n - y_n| \to 0\) but \(|f(x_n) - f(y_n)| \geq \varepsilon_0\) for
  some fixed \(\varepsilon_0\)
\item
  ✎ Inline: Prove \(f(x) = \sin(1/x)\) is not uniformly continuous on
  \((0, 1)\)
\item
  Positive example: \(f(x) = \sqrt{x}\) IS uniformly continuous on
  \([0, \infty)\)

  \begin{itemize}
  \tightlist
  \item
    \(|\sqrt{x} - \sqrt{y}| = \frac{|x - y|}{\sqrt{x} + \sqrt{y}} \leq \frac{|x - y|}{\max(\sqrt{x}, \sqrt{y})}\)\ldots{}
    but actually the easier estimate:
    \(|\sqrt{x} - \sqrt{y}| \leq \sqrt{|x - y|}\) (exercise). Take
    \(\delta = \varepsilon^2\).
  \end{itemize}
\item
  Remark: Lipschitz functions are uniformly continuous. If
  \(|f(x) - f(y)| \leq L|x - y|\) for some constant \(L\), take
  \(\delta = \varepsilon / L\). But uniform continuity is strictly
  weaker --- \(\sqrt{x}\) is uniformly continuous but not Lipschitz
  (exercise).
\end{itemize}

\subsection{Closed Intervals}\label{closed-intervals}

\begin{itemize}
\tightlist
\item
  The non-examples both involved domains that were open or unbounded. Is
  this a coincidence?
\item
  Theorem: If \(f\) is continuous on \([a,b]\), then \(f\) is uniformly
  continuous on \([a,b]\).
\item
  This is a major theorem --- continuity, a pointwise condition,
  automatically upgrades to the uniform condition on closed bounded
  intervals.
\item
  Proof sketch (by contradiction, via Bolzano-Weierstrass):

  \begin{itemize}
  \tightlist
  \item
    Suppose \(f\) is continuous on \([a,b]\) but not uniformly
    continuous
  \item
    Then there exists \(\varepsilon_0 > 0\) and sequences
    \(x_n, y_n \in [a,b]\) with \(|x_n - y_n| < 1/n\) but
    \(|f(x_n) - f(y_n)| \geq \varepsilon_0\)
  \item
    By BW, \((x_n)\) has a convergent subsequence:
    \(x_{n_k} \to c \in [a,b]\)
  \item
    Since \(|x_{n_k} - y_{n_k}| < 1/n_k \to 0\), also \(y_{n_k} \to c\)
  \item
    By continuity at \(c\): \(f(x_{n_k}) \to f(c)\) and
    \(f(y_{n_k}) \to f(c)\)
  \item
    So \(|f(x_{n_k}) - f(y_{n_k})| \to 0\) --- contradicting
    \(\geq \varepsilon_0\)
  \end{itemize}
\item
  Remark: the proof uses BW, which requires the domain to be bounded.
  And \(c \in [a,b]\) uses that the domain is closed. Both hypotheses
  are essential.
\item
  Remark: this theorem will be important for integration (Chapter 8) ---
  the proof that continuous functions are Riemann integrable relies on
  uniform continuity
\item
  ✎ Inline: Where exactly does the proof fail if the domain is
  \((0, 1)\) instead of \([0, 1]\)? (Answer: the subsequential limit
  \(c\) might be \(0\), which is not in the domain, so we can't invoke
  continuity at \(c\).)
\end{itemize}

\subsection{Continuous Extensions}\label{continuous-extensions}

\begin{itemize}
\tightlist
\item
  Situation: \(f\) is continuous on \((a, b)\) --- can we extend it to
  \([a, b]\)?
\item
  Not always: \(f(x) = \sin(1/x)\) on \((0, 1)\) has no limit as
  \(x \to 0^+\), so there's no way to define \(f(0)\) to make it
  continuous
\item
  Not always: \(f(x) = 1/x\) on \((0, 1)\) is unbounded, so certainly
  can't extend continuously
\item
  Uniform continuity is exactly the condition that makes extension work
\item
  Theorem (Continuous Extension): If \(f\) is uniformly continuous on
  \((a, b)\), then \(f\) extends to a continuous function on \([a, b]\).
\item
  More precisely: \(\lim_{x \to a^+} f(x)\) and
  \(\lim_{x \to b^-} f(x)\) both exist, and defining \(f(a)\) and
  \(f(b)\) to be these limits gives a continuous (in fact uniformly
  continuous) function on \([a, b]\).
\item
  Proof sketch:

  \begin{itemize}
  \tightlist
  \item
    Show \(\lim_{x \to a^+} f(x)\) exists via Cauchy criterion for
    limits:
  \item
    If \(x_n, y_n \to a^+\), then \(|x_n - y_n| \to 0\), so by uniform
    continuity \(|f(x_n) - f(y_n)| \to 0\)
  \item
    This means \((f(x_n))\) is Cauchy whenever \(x_n \to a^+\), so the
    limit exists
  \item
    And any two sequences approaching \(a\) give the same limit (same
    argument)
  \item
    Define \(f(a)\) to be this limit; similarly for \(f(b)\)
  \item
    Continuity at the endpoints follows from the definition
  \end{itemize}
\item
  Remark: the converse also holds --- if \(f\) is continuous on
  \((a,b)\) and extends to a continuous function on \([a,b]\), then
  \(f\) is uniformly continuous on \((a,b)\) (because the extension is
  continuous on \([a,b]\), hence uniformly continuous by the theorem
  above)
\item
  Corollary: uniformly continuous on a bounded set \(\Rightarrow\)
  bounded

  \begin{itemize}
  \tightlist
  \item
    The function extends to the closure, which is a closed bounded
    interval, and continuous functions on \([a,b]\) are bounded (§6.4)
  \end{itemize}
\end{itemize}

\subsection{Guided Exercise}\label{guided-exercise-14}

\subsection{Uniform Continuity and
Sequences}\label{uniform-continuity-and-sequences}

This exercise develops a sequential characterization of uniform
continuity, paralleling the sequential characterization of continuity
from §6.1.

\begin{enumerate}
\def\labelenumi{(\alph{enumi})}
\item
  Prove: \(f\) is NOT uniformly continuous on \(S\) if and only if there
  exist sequences \((x_n)\) and \((y_n)\) in \(S\) with
  \(|x_n - y_n| \to 0\) but \(|f(x_n) - f(y_n)| \not\to 0\).
\item
  Use this to give a second proof that \(f(x) = x^2\) is not uniformly
  continuous on \(\RR\).
\item
  Use this to prove \(f(x) = x \sin(x)\) is not uniformly continuous on
  \(\RR\). (Hint: consider \(x_n = 2\pi n\) and \(y_n = 2\pi n + 1/n\).)
\item
  Prove: if \(f\) is uniformly continuous on \(S\) and \((x_n)\) is a
  Cauchy sequence in \(S\), then \((f(x_n))\) is a Cauchy sequence.
  (This is stronger than what ordinary continuity gives --- continuity
  only preserves convergent sequences, not Cauchy sequences.)
\item
  Give an example showing ordinary (non-uniform) continuity does NOT
  preserve Cauchy sequences. (Hint: \(f(x) = 1/x\) on \((0,1)\) and a
  Cauchy sequence converging to \(0\).)
\item
  Use part (d) to give another proof of the extension theorem: if \(f\)
  is uniformly continuous on \((a,b)\), show \(\lim_{x \to a^+} f(x)\)
  exists.
\end{enumerate}

\subsection{Exercises}\label{exercises-28}

\emph{The definition} - Prove directly from the definition:
\(f(x) = 3x + 1\) is uniformly continuous on \(\RR\) - Prove directly
from the definition: \(f(x) = x/(1 + x)\) is uniformly continuous on
\([0, \infty)\) - Prove: every Lipschitz function is uniformly
continuous - Prove: \(f(x) = \sqrt{x}\) is uniformly continuous on
\([0, \infty)\). (Hint: show
\(|\sqrt{x} - \sqrt{y}| \leq \sqrt{|x-y|}\).) - Show \(f(x) = \sqrt{x}\)
is NOT Lipschitz on \([0, \infty)\). So uniform continuity is strictly
weaker than Lipschitz.

\emph{Non-examples} - Prove \(f(x) = x^3\) is not uniformly continuous
on \(\RR\) - Prove \(f(x) = \sin(x^2)\) is not uniformly continuous on
\(\RR\) - Prove \(f(x) = 1/(1-x)\) is not uniformly continuous on
\((0, 1)\) - ★ Prove: if \(f\) is uniformly continuous on \(\RR\), then
there exist constants \(A, B\) with \(|f(x)| \leq A|x| + B\) for all
\(x\). (So uniformly continuous functions on \(\RR\) grow at most
linearly.)

\emph{The theorem for closed intervals} - Where exactly does the proof
of ``continuous on \([a,b]\) ⟹ uniformly continuous'' break down for:
(i) \((0, 1)\)? (ii) \([0, \infty)\)? - Prove: if \(f\) is continuous on
\([0, \infty)\) and \(\lim_{x \to \infty} f(x)\) exists, then \(f\) is
uniformly continuous on \([0, \infty)\) - ★ Prove: if \(f\) is
continuous on \([0, \infty)\) and \(\lim_{x \to \infty} f(x) = 0\), then
\(f\) is uniformly continuous on \([0, \infty)\). (Hint: \(f\) is
uniformly continuous on \([0, N]\) by the theorem; choose \(N\) large
enough that the tail is small.)

\emph{Extension and Cauchy} - Prove: \(f(x) = x \sin(1/x)\) is uniformly
continuous on \((0, 1)\). (Hint: what is \(\lim_{x \to 0^+} f(x)\)?) -
Prove: if \(f\) is uniformly continuous on \((a, b)\), then \(f\) is
bounded on \((a, b)\) - ★ Give an example of a function continuous on
\((0,1)\) that is bounded but NOT uniformly continuous. (Hint:
\(\sin(1/x)\).) - ★ Prove: \(f\) is uniformly continuous on \((a,b)\) if
and only if \(f\) extends to a continuous function on \([a,b]\)

\subsection{Dependencies}\label{dependencies-28}

\textbf{Requires}: §3.1 (Bolzano-Weierstrass), §3.3 (Cauchy sequences),
§6.1 (\(\varepsilon\)-\(\delta\) definition, sequential
characterization), §6.4 (continuous on \([a,b]\) ⟹ bounded)

\textbf{Used in}: Ch 8 (Riemann integrability of continuous functions),
Ch 11 (uniform convergence revisits ``one parameter works for all'')

\chapter{Differentiation}\label{differentiation}

\textbf{Ancient tangent lines} - Apollonius (\textasciitilde200 BCE):
tangent lines to conic sections, purely geometric constructions -
Archimedes: tangents to spirals --- arguably the first derivative
calculation - But no general method; each curve required its own
argument - The question ``what is the tangent line?'' remained
geometric, not algebraic

\textbf{Fermat's method of adequality (1630s)} - Fermat: to find the
maximum of f, set f(x+e) ``adequal'' to f(x), simplify, divide by e, set
e = 0 - Mysteriously correct: produces f'(x) = 0 at extrema - Also found
tangent lines to curves by the same method - The reasoning was unclear
--- what does it mean to divide by e and then set e = 0? - Descartes had
a competing method via normal lines and algebraic geometry

\textbf{Newton and Leibniz (1680s)} - Newton: quantities in motion,
``fluxions'' as rates of change - Leibniz: infinitesimal differences dy
and dx, the derivative as their ratio - Both developed the rules ---
product, quotient, chain --- and the connection to integration -
Calculus worked brilliantly, but what \emph{are} these infinitely small
quantities?

\textbf{The critique} - Berkeley (1734): ``ghosts of departed
quantities'' --- the foundations are logically incoherent - Working
mathematicians didn't stop --- calculus was too powerful - Euler, the
Bernoullis: 18th century calculus develops rapidly on shaky ground -
Correct results obtained by incorrect (or at least unjustified)
reasoning

\textbf{Rigorization} - Cauchy (\textasciitilde1820s): first serious
attempt to define limits and continuity - Weierstrass
(\textasciitilde1860s): the ε-δ definition makes everything precise -
The derivative is a limit --- no infinitesimals needed - Surprise:
Weierstrass (1872) constructs a function continuous everywhere but
differentiable nowhere - Continuity and differentiability are far apart
--- differentiability is a much stronger condition

\textbf{Our approach} - We define the derivative via Weierstrass's
framework - The rules are mechanical; the power comes from the Mean
Value Theorem - Taylor's theorem: the derivatives of a function
reconstruct it as a power series - The exponential: built entirely from
calculus, the crown jewel of the chapter

\section{§7.1 Differentiation Rules}\label{differentiation-rules}

\subsection{Narrative}\label{narrative-29}

\begin{itemize}
\tightlist
\item
  The derivative asks: how well does a linear function approximate \(f\)
  near a point?
\item
  This section builds the basic toolkit: definition, relationship to
  continuity, algebraic rules, chain rule, inverse function derivative
\item
  By the end we can differentiate anything built from elementary
  operations --- specific functions (\(\exp\), \(\sin\), \(\cos\))
  require §7.3--7.4
\end{itemize}

\subsection{Content}\label{content-29}

\emph{(Opening --- unnumbered)} - Brief motivation: Fermat's
optimization idea, Newton/Leibniz tangent slopes, now we have limits to
make it precise

\subsection{The Derivative}\label{the-derivative}

\begin{itemize}
\tightlist
\item
  Definition: \(f\) differentiable at \(a\) if
  \(f'(a) = \lim_{h \to 0} \frac{f(a+h) - f(a)}{h}\) exists
\item
  Equivalent form: \(\lim_{x \to a} \frac{f(x) - f(a)}{x - a}\)
\item
  Notation: \(f'(a)\), \(\frac{df}{dx}\big|_{x=a}\), \(Df(a)\)
\item
  The tangent line: \(y = f(a) + f'(a)(x - a)\)
\item
  Definition: \(f\) differentiable on an interval \(I\) if
  differentiable at every point of \(I\); \(f'\) is then a function on
  \(I\)
\item
  Definition: if \(f'\) is itself differentiable, \(f'' = (f')'\) is the
  \emph{second derivative}; iterate to get \(f^{(n)}\). Notation:
  \(f^{(0)} = f\)
\item
  One-sided derivatives: \(f'_+(a)\), \(f'_-(a)\); differentiable iff
  both exist and agree
\item
  Examples: \(f(x) = c\) (compute: \(f'(a) = 0\)), \(f(x) = x\)
  (\(f'(a) = 1\)), \(f(x) = x^2\) (\(f'(a) = 2a\), compute from
  definition), \(f(x) = 1/x\) (\(f'(a) = -1/a^2\))
\item
  ✎ Inline: Compute \(f'(a)\) for \(f(x) = x^3\) directly from the
  definition
\end{itemize}

\subsection{Differentiability Implies
Continuity}\label{differentiability-implies-continuity}

\begin{itemize}
\tightlist
\item
  Theorem: \(f\) differentiable at \(a\) \(\Longrightarrow\) \(f\)
  continuous at \(a\) (prove in text: write
  \(f(x) - f(a) = \frac{f(x)-f(a)}{x-a} \cdot (x-a) \to f'(a) \cdot 0 = 0\))
\item
  Warning: converse false --- \(|x|\) at \(0\) (continuous, corner,
  one-sided derivatives \(-1\) and \(+1\) differ)
\item
  ✎ Inline: Compute \(f'_+(0)\) and \(f'_-(0)\) for \(f(x) = |x|\)
\item
  Remark: The gap is vast; Takagi's function (§7.3) is continuous
  everywhere, differentiable nowhere
\end{itemize}

\subsection{The Algebraic Rules}\label{the-algebraic-rules}

\begin{itemize}
\tightlist
\item
  Theorem (Linearity): \((f + g)' = f' + g'\), \((cf)' = cf'\) (prove in
  text: direct from limit laws)
\item
  Theorem (Product Rule): \((fg)' = f'g + fg'\) (prove in text:
  add-subtract trick \(f(a)g(x)\); uses differentiability
  \(\Rightarrow\) continuity for \(f\))
\item
  Theorem (Quotient Rule): \((f/g)' = (f'g - fg')/g^2\) when
  \(g(a) \neq 0\)
\item
  ✎ Inline: Derive quotient rule from product rule and the derivative of
  \(1/g\)
\item
  Theorem (Power Rule): \((x^n)' = nx^{n-1}\) for \(n \in \NN\) (prove
  in text: induction via product rule, or factor \(x^n - a^n\)); extends
  to \(n \in \ZZ\) via quotient rule
\item
  Corollary: polynomials differentiable on \(\RR\); rational functions
  differentiable on their domain
\end{itemize}

\subsection{The Chain Rule}\label{the-chain-rule}

\begin{itemize}
\tightlist
\item
  Theorem: \((f \circ g)'(a) = f'(g(a)) \cdot g'(a)\) (prove in text)
\item
  Leibniz notation:
  \(\frac{dy}{dx} = \frac{dy}{du} \cdot \frac{du}{dx}\)
\item
  Note the subtlety: the naive proof divides by \(g(x) - g(a)\), which
  may be zero; the correct proof defines
  \(\varphi(t) = \frac{f(t) - f(g(a))}{t - g(a)}\) for \(t \neq g(a)\),
  \(\varphi(g(a)) = f'(g(a))\); then \(\varphi\) is continuous at
  \(g(a)\) and \(f(g(x)) - f(g(a)) = \varphi(g(x)) \cdot (g(x) - g(a))\)
  for all \(x\)
\item
  ✎ Inline: Carry out the division by \(x - a\) and take limits to
  finish the proof
\end{itemize}

\subsection{The Inverse Function
Derivative}\label{the-inverse-function-derivative}

\begin{itemize}
\tightlist
\item
  Setup: \(f\) continuous, strictly monotone on interval \(I\);
  \(f^{-1}\) exists and is continuous by §6.4
\item
  Theorem: If \(f'(a) \neq 0\), then \((f^{-1})'(f(a)) = 1/f'(a)\)
  (prove in text: with \(y = f(x)\), \(b = f(a)\), write
  \(\frac{f^{-1}(y) - f^{-1}(b)}{y - b} = \frac{1}{(f(x)-f(a))/(x-a)}\);
  continuity of \(f^{-1}\) gives \(x \to a\) as \(y \to b\))
\item
  Leibniz notation: \(\frac{dx}{dy} = \frac{1}{dy/dx}\)
\item
  Example: \((\sqrt{x})' = 1/(2\sqrt{x})\);
  \((x^{1/n})' = \frac{1}{n}x^{1/n - 1}\)
\item
  Corollary: power rule extends to rational exponents:
  \((x^{p/q})' = \frac{p}{q} x^{p/q - 1}\)
\item
  ✎ Inline: Derive \((\sqrt[3]{x})'\) from the inverse function theorem
\end{itemize}

\subsection{Guided Exercise}\label{guided-exercise-15}

\subsection{The Linear Approximation Perspective
(Carathéodory)}\label{the-linear-approximation-perspective-carathuxe9odory}

\begin{enumerate}
\def\labelenumi{(\alph{enumi})}
\item
  Show \(f\) differentiable at \(a\) iff there exists \(\varphi\)
  continuous at \(a\) with \(f(x) - f(a) = \varphi(x)(x - a)\); then
  \(f'(a) = \varphi(a)\).
\item
  Reprove the chain rule cleanly using (a): compose the factorizations
  for \(f\) and \(g\).
\item
  Reprove the product rule using (a).
\end{enumerate}

\subsection{Exercises}\label{exercises-29}

\emph{Direct computation} - \(f'(a)\) from the definition for
\(f(x) = x^4\), \(f(x) = 1/(x+1)\), \(f(x) = \sqrt{2x+1}\)

\emph{Differentiability and continuity} - Show \(f(x) = x^2 \sin(1/x)\)
(with \(f(0) = 0\)) is differentiable at \(0\) with \(f'(0) = 0\), but
\(f'\) is not continuous at \(0\). (Students may use known properties of
\(\sin\) and \(\cos\) from calculus; rigorous treatment in §7.3.) - Show
\(f(x) = x|x|\) is differentiable everywhere; find \(f'\) - Give an
example of \(f\) differentiable on \(\RR\) with \(f'\) discontinuous at
a point

\emph{Algebraic rules} - General Leibniz rule:
\((fg)^{(n)} = \sum_{k=0}^{n} \binom{n}{k} f^{(k)} g^{(n-k)}\) - Prove
power rule by induction using product rule - Product rule for three
functions: \((fgh)' = ?\) - Prove: \(f\) differentiable at \(a\) and
\(f(a) \neq 0\) \(\Longrightarrow\) \(1/f\) differentiable at \(a\)

\emph{Chain rule and inverse} - If \(f(f(x)) = x\) for all \(x\) and
\(f\) differentiable, what can you say about \(f'(x)\)? - Extend the
power rule to rational exponents via inverse function derivative

\emph{Harder} - ★ Show: \(f(x) = x + 2x^2 \sin(1/x)\) (with
\(f(0) = 0\)) has \(f'(0) = 1 > 0\), yet \(f\) is not increasing on any
interval containing \(0\). (May use known properties of \(\sin\),
\(\cos\).) - ★ If \(f\) differentiable on \((a,b)\) and \(f'\) bounded,
show \(f\) is Lipschitz (preview of §7.2)

\subsection{Dependencies}\label{dependencies-29}

\textbf{Requires}: §2.2--2.3 (limits, limit laws), §6.1 (continuity,
function limits), §6.2 (algebra of continuous functions), §6.4 (monotone
inverses)

\textbf{Used in}: All subsequent sections of this chapter

\section{§7.2 The Mean Value Theorem}\label{the-mean-value-theorem}

\subsection{Narrative}\label{narrative-30}

\begin{itemize}
\tightlist
\item
  The derivative is defined locally, but the MVT connects it to global
  behavior
\item
  The chain: Fermat (local extrema) → Rolle → MVT → Cauchy MVT
\item
  Payoff: \(f'\) controls \(f\) --- constant, increasing/decreasing,
  Lipschitz, L'Hôpital
\item
  Darboux's theorem (guided exercise) reveals that derivatives are an
  unusual class of functions: they always satisfy the IVP, even when
  discontinuous
\end{itemize}

\subsection{Content}\label{content-30}

\emph{(Opening --- unnumbered)} - The derivative tells us about \(f\) at
a single point; can it control \(f\) on an entire interval? - The Mean
Value Theorem says yes: somewhere on the interval, the instantaneous
rate of change equals the average rate of change - Historical note:
Rolle (1691) proved his theorem for polynomials; the general MVT is
attributed to Cauchy (1823) and Bonnet; Darboux (1875) discovered the
surprising intermediate value property of derivatives

\subsection{A First Step}\label{a-first-step}

\begin{itemize}
\tightlist
\item
  Theorem (Fermat): If \(f\) has a local extremum at \(c\) and \(f\) is
  differentiable at \(c\), then \(f'(c) = 0\) (prove in text: one-sided
  difference quotients have opposite signs, both must equal \(f'(c)\))
\item
  Remark: converse is false --- \(f(x) = x^3\) has \(f'(0) = 0\) but no
  extremum at \(0\)
\item
  ✎ Inline: Why does the proof require a \emph{local} extremum, not just
  \(f'(c) = 0\)?
\item
  Theorem (Rolle): \(f\) continuous on \([a,b]\), differentiable on
  \((a,b)\), \(f(a) = f(b)\) \(\Longrightarrow\) \(\exists c \in (a,b)\)
  with \(f'(c) = 0\) (prove in text: if \(f\) constant, any \(c\) works;
  otherwise EVT gives max or min in interior, then apply Fermat)
\item
  ✎ Inline: Where does the proof use each hypothesis --- continuity on
  \([a,b]\), differentiability on \((a,b)\), \(f(a) = f(b)\)?
\end{itemize}

\subsection{The Mean Value Theorem}\label{the-mean-value-theorem-1}

\begin{itemize}
\tightlist
\item
  Theorem (MVT): \(f\) continuous on \([a,b]\), differentiable on
  \((a,b)\) \(\Longrightarrow\) \(\exists c \in (a,b)\) with
  \(f'(c) = \frac{f(b) - f(a)}{b - a}\) (prove in text: apply Rolle to
  \(g(x) = f(x) - \frac{f(b)-f(a)}{b-a}(x - a)\))
\item
  Geometric interpretation: some tangent line is parallel to the secant
  through \((a, f(a))\) and \((b, f(b))\)
\item
  ✎ Inline: Verify the MVT for \(f(x) = x^3\) on \([0, 2]\): find \(c\)
  explicitly
\item
  Theorem (Cauchy's MVT): \(f, g\) continuous on \([a,b]\),
  differentiable on \((a,b)\) \(\Longrightarrow\)
  \(\exists c \in (a,b)\) with
  \((f(b) - f(a))g'(c) = (g(b) - g(a))f'(c)\) (prove in text: apply
  Rolle to \(h(x) = (f(b)-f(a))g(x) - (g(b)-g(a))f(x)\))
\item
  Remark: reduces to MVT when \(g(x) = x\)
\item
  Remark: the motivation for this generalization is L'Hôpital's rule
  below
\end{itemize}

\subsection{Consequences of the MVT}\label{consequences-of-the-mvt}

\begin{itemize}
\tightlist
\item
  Theorem: \(f\) differentiable on \((a,b)\) with \(f' = 0\)
  \(\Longrightarrow\) \(f\) constant (prove in text: for any \(x, y\),
  MVT gives \(f(x) - f(y) = f'(c)(x-y) = 0\))
\item
  Corollary: if \(f' = g'\) on an interval, then \(f = g + C\) for some
  constant \(C\)
\item
  Theorem: \(f' > 0\) on \((a,b)\) \(\Longrightarrow\) \(f\) strictly
  increasing (prove in text via MVT)
\item
  Theorem: \(f' \geq 0\) \(\Longleftrightarrow\) \(f\) increasing;
  \(f' \leq 0\) \(\Longleftrightarrow\) \(f\) decreasing
\item
  Remark: local information (sign of \(f'\)) determines global behavior
  (monotonicity) --- this is what makes the derivative useful
\item
  Definition: \(f\) is \emph{Lipschitz} with constant \(K\) if
  \(|f(x) - f(y)| \leq K|x - y|\) for all \(x, y\)
\item
  Theorem: If \(|f'(x)| \leq K\) on an interval, then \(f\) is Lipschitz
  with constant \(K\) (prove in text via MVT)
\item
  Corollary: Lipschitz \(\Longrightarrow\) uniformly continuous
  (callback to §6.5)
\item
  Remark: if \(|f'| < 1\) on an interval, then \(f\) is a contraction
  (callback to §4.3; used in §7.5)
\item
  ✎ Inline: Show \(\sin\) is Lipschitz with constant 1 (students may use
  \(|\cos x| \leq 1\) from calculus; rigorous treatment in §7.3)
\end{itemize}

\subsection{L'Hôpital's Rule}\label{lhuxf4pitals-rule}

\begin{itemize}
\tightlist
\item
  Theorem (0/0 form): If \(f(a) = g(a) = 0\), \(g'(x) \neq 0\) near
  \(a\), and \(\lim_{x \to a} f'(x)/g'(x) = L\), then
  \(\lim_{x \to a} f(x)/g(x) = L\) (prove in text via Cauchy MVT)
\item
  Warning: the hypothesis is that \(\lim f'/g'\) exists; L'Hôpital does
  not apply otherwise
\item
  Warning: L'Hôpital can cycle --- applying it to \(x/\sqrt{x^2+1}\)
  gives \(1/\sqrt{1 + 1/x^2} \cdot \ldots\) which is no simpler
\end{itemize}

\subsection{Guided Exercise}\label{guided-exercise-16}

\subsection{Darboux's Theorem and Discontinuous
Derivatives}\label{darbouxs-theorem-and-discontinuous-derivatives}

Derivatives are a peculiar class of functions: they always satisfy the
intermediate value property, even when they are discontinuous. This is
Darboux's theorem (1875).

\begin{enumerate}
\def\labelenumi{(\alph{enumi})}
\item
  Let \(f\) be differentiable on \([a,b]\) with \(f'(a) < k < f'(b)\).
  Define \(g(x) = f(x) - kx\). Show \(g'(a) < 0 < g'(b)\).
\item
  Show \(g\) attains its minimum on \([a,b]\) (which theorem?), and that
  the minimum cannot occur at \(a\) or \(b\) (use the signs of \(g'(a)\)
  and \(g'(b)\)).
\item
  Conclude: \(g'(c) = 0\) for some \(c \in (a,b)\), hence \(f'(c) = k\).
  This is Darboux's theorem.
\item
  Corollary: derivatives cannot have jump discontinuities. Explain why.
\item
  Now exhibit a discontinuous derivative that satisfies the IVP. Let
  \(f(x) = x^2 \sin(1/x)\) for \(x \neq 0\), \(f(0) = 0\). Compute
  \(f'(x)\) for \(x \neq 0\) (using standard rules; students may use
  known properties of \(\sin\) and \(\cos\)) and \(f'(0)\) from the
  definition.
\item
  Show \(f'\) is discontinuous at \(0\): the \(\cos(1/x)\) term
  oscillates. Yet by Darboux, \(f'\) must satisfy the IVP.
\item
  Reflection: what kinds of discontinuities \emph{can} a derivative
  have?
\end{enumerate}

\subsection{Exercises}\label{exercises-30}

\emph{Applying the MVT} - Find \(c\) explicitly for \(f(x) = \sqrt{x}\)
on \([1, 4]\) - Prove: \(|\sin x - \sin y| \leq |x - y|\) for all
\(x, y\) (students may use \(|\cos| \leq 1\) from calculus) - Prove: if
\(f'(x) \neq 0\) for all \(x \in (a,b)\), then \(f\) is injective on
\((a,b)\)

\emph{Derivative tests} - First derivative test: if \(f'\) changes from
positive to negative at \(c\), then \(f\) has a local maximum at \(c\)
(and vice versa) - Second derivative test: if \(f'(c) = 0\) and
\(f''(c) > 0\), then \(f\) has a local minimum at \(c\); if
\(f''(c) < 0\), local maximum - Give an example where \(f'(c) = 0\) and
\(f''(c) = 0\) and \(c\) is a local minimum; another where it is neither

\emph{Convexity} - Definition: \(f\) is \emph{convex} on an interval if
\(f(\lambda x + (1-\lambda)y) \leq \lambda f(x) + (1-\lambda)f(y)\) for
all \(x, y\) and \(\lambda \in [0,1]\) - Prove: if \(f\) twice
differentiable, then \(f\) convex \(\Longleftrightarrow\) \(f'' \geq 0\)
- Prove: \(f\) convex \(\Longrightarrow\)
\(f(y) \geq f(x) + f'(x)(y - x)\) (tangent lines lie below graph) -
Prove: \(f\) convex on open interval \(\Longrightarrow\) \(f\)
continuous - Prove Jensen's inequality: if \(f\) convex and
\(\lambda_1 + \cdots + \lambda_n = 1\) with \(\lambda_i \geq 0\), then
\(f(\lambda_1 x_1 + \cdots + \lambda_n x_n) \leq \lambda_1 f(x_1) + \cdots + \lambda_n f(x_n)\)
- (Deferred to Chapter 9) Derive AM-GM from Jensen applied to
\(f(x) = -\ln x\)

\emph{L'Hôpital extensions} - Prove L'Hôpital for the \(\infty/\infty\)
form - Prove L'Hôpital as \(x \to \infty\) - Evaluate specific limits
using L'Hôpital

\emph{Harder} - Construct \(f\) with \(f'\) discontinuous but bounded -
★ Prove: if \(f\) is differentiable on \([a,b]\) and \(f'\) is monotone,
then \(f'\) is continuous (hint: Darboux + monotone means no oscillatory
discontinuities)

\subsection{Dependencies}\label{dependencies-30}

\textbf{Requires}: §6.4 (EVT --- for Rolle/MVT), §6.5 (uniform
continuity --- for Lipschitz connection), §7.1 (differentiability, chain
rule)

\textbf{Used in}: §7.3 (Taylor's theorem uses MVT repeatedly), §7.4
(identifying exp), §7.5 (Newton's method uses Lipschitz/contraction)

\section{§7.3 Power Series \& Taylor's
Theorem}\label{power-series-taylors-theorem}

\subsection{Narrative}\label{narrative-31}

\begin{itemize}
\tightlist
\item
  We already know power series from Chapters 4--6; now we ask: can we
  differentiate them?
\item
  Start from what we know (power series), discover the coefficients
  encode the derivatives
\item
  This motivates Taylor polynomials: given any function, build a power
  series candidate from its derivatives
\item
  Taylor's theorem measures the error --- when does the candidate
  converge to the function?
\item
  Power series are special: they always equal their Taylor series.
  General smooth functions need not (example in §7.4)
\end{itemize}

\subsection{Content}\label{content-31}

\subsection{Differentiating Power
Series}\label{differentiating-power-series}

\begin{itemize}
\tightlist
\item
  Lemma: \(\sum a_n x^n\) and \(\sum n a_n x^{n-1}\) have the same
  radius of convergence
\item
  Theorem: If \(f(x) = \sum a_n x^n\) on \((-R, R)\), then
  \(f'(x) = \sum n a_n x^{n-1}\) on \((-R, R)\) (prove in text via
  Tannery/dominated convergence from §5.4)
\item
  Corollary: power series are infinitely differentiable on \((-R, R)\)
\item
  Corollary: \(a_n = f^{(n)}(0)/n!\) --- the coefficients encode the
  derivatives
\end{itemize}

\subsection{Taylor Polynomials}\label{taylor-polynomials}

\begin{itemize}
\tightlist
\item
  Definition: For \(f\) with \(n\) derivatives at \(a\), the \(n\)th
  Taylor polynomial is
  \(P_n(x) = \sum_{k=0}^{n} \frac{f^{(k)}(a)}{k!}(x - a)^k\)
\item
  Theorem: \(P_n\) is the unique polynomial of degree \(\leq n\)
  matching \(f, f', \ldots, f^{(n)}\) at \(a\)
\item
  The question: for a general function, does \(P_n \to f\) as
  \(n \to \infty\)?
\end{itemize}

\subsection{Taylor's Theorem}\label{taylors-theorem}

\begin{itemize}
\tightlist
\item
  Theorem (Generalized Rolle): If \(f, f', \ldots, f^{(n)}\) all vanish
  at \(a\) and \(f(b) = 0\), then \(f^{(n+1)}(c) = 0\) for some \(c\)
  between \(a\) and \(b\) (prove in text: induction on \(n\) using
  Rolle's theorem)
\item
  Theorem (Polynomial MVT): If \(f\) and \(h\) share their first \(n\)
  derivatives at \(a\) and \(f(b) = h(b)\), then
  \(f^{(n+1)}(c) = h^{(n+1)}(c)\) for some \(c\) between \(a\) and \(b\)
  (prove in text: apply generalized Rolle to \(f - h\))
\item
  Theorem (Lagrange Remainder):
  \(f(x) = P_n(x) + \frac{f^{(n+1)}(c)}{(n+1)!}(x - a)^{n+1}\) for some
  \(c\) between \(a\) and \(x\) (prove in text: apply polynomial MVT
  with \(h = P_n\))
\item
  Error bound: \(|f(x) - P_n(x)| \leq \frac{M}{(n+1)!}|x - a|^{n+1}\)
  where \(M\) bounds \(|f^{(n+1)}|\)
\item
  Remark: the integral remainder form (which does not require a mystery
  point \(c\)) will appear in Chapter 9 once we have integration
\item
  Remark: For power series, the Taylor series always converges to the
  function (by the corollary in subsection 1). For general smooth
  functions, this can fail --- a counterexample requires the exponential
  function, built in §7.4
\end{itemize}

\subsection{Guided Exercise}\label{guided-exercise-17}

\subsection{Takagi's Function: Continuous Everywhere, Differentiable
Nowhere}\label{takagis-function-continuous-everywhere-differentiable-nowhere}

Let \(s(x) = \mathrm{dist}(x, \ZZ)\) (distance to nearest integer) and
define \(T(x) = \sum_{n=0}^{\infty} \frac{s(2^n x)}{2^n}\).

\begin{enumerate}
\def\labelenumi{(\alph{enumi})}
\item
  Show \(|s(2^n x)/2^n| \leq 1/2^{n+1}\) for all \(x\), and conclude
  \(T\) is defined for all \(x\) (the series converges absolutely and
  uniformly).
\item
  Each partial sum \(T_N(x) = \sum_{n=0}^{N} s(2^n x)/2^n\) is
  continuous (why?). Use uniform convergence to conclude \(T\) is
  continuous everywhere.
\item
  Show that \(T\) is differentiable nowhere. (Proof strategy: for any
  \(x\), choose \(h = \pm 1/2^m\) so that the difference quotient of the
  \(m\)th term contributes \(\pm 1/2\) while higher terms contribute 0;
  the difference quotient doesn't converge as \(m \to \infty\).)
\item
  Reflect: each summand \(s(2^n x)/2^n\) is differentiable except at
  finitely many points. The uniform limit is continuous but
  differentiable nowhere. What does this say about limits preserving
  differentiability?
\end{enumerate}

\subsection{Exercises}\label{exercises-31}

\begin{itemize}
\tightlist
\item
  Compute the Taylor series of \(1/(1-x)^2\) and \(1/(1-x)^3\) by
  differentiating the geometric series
\end{itemize}

\subsection{Dependencies}\label{dependencies-31}

\textbf{Requires}: §5.4 (Tannery/dominated convergence --- for
term-by-term differentiation), §6.3 (power series, radius of
convergence), §7.1 (derivative definition, higher derivatives), §7.2
(Rolle's theorem, MVT)

\textbf{Used in}: §7.4 (Taylor series for exp), §7.5 (Taylor expansion
for quadratic convergence of Newton)

\section{§7.4 Exponentials and
Logarithms}\label{exponentials-and-logarithms}

\subsection{Narrative}\label{narrative-32}

\begin{itemize}
\tightlist
\item
  We have aˣ from ch1 (defined via supremum) and know it's continuous
  --- but is it differentiable?
\item
  Proving differentiability at 0 is the key: the functional equation
  does the rest
\item
  Every exponential has a proportionality constant; we define the
  natural exponential as the one where this constant is 1
\item
  This gives us exp' = exp, which determines the Taylor series --- and
  we recognize it as E(x) from ch5
\item
  Defining e = exp(1) gives the \emph{right} definition of e; we then
  recover e = lim(1+1/n)ⁿ and e = Σ1/n! as theorems
\item
  Then we define the natural logarithm as the inverse of exp: the
  functional equation, derivative, and Taylor series all follow from
  what we already know
\item
  The mystery constant a'(0) is identified as log(a), completing the
  story of exponential functions
\end{itemize}

\subsection{Content}\label{content-32}

\subsection{Differentiability of
Exponentials}\label{differentiability-of-exponentials}

\begin{itemize}
\tightlist
\item
  Theorem: aˣ is differentiable at 0 (prove in text)
\item
  Theorem: If aˣ is differentiable at 0, then aˣ is differentiable
  everywhere, and (aˣ)' = aˣ · a'(0) (prove in text: functional
  equation)
\item
  Remark: every exponential is determined by its derivative factor a'(0)
\end{itemize}

\subsection{The Natural Exponential}\label{the-natural-exponential}

\begin{itemize}
\tightlist
\item
  Definition: the \emph{natural exponential} exp is the exponential
  function with exp'(0) = 1, i.e.~exp' = exp
\item
  exp(x) = aˣ for some base a; define e = exp(1), so exp(x) = eˣ
\item
  Remark: this is the proper definition of e --- the base whose
  exponential is its own derivative
\end{itemize}

\subsection{The Taylor Series for exp}\label{the-taylor-series-for-exp}

\begin{itemize}
\tightlist
\item
  exp' = exp and exp(0) = 1 determine all derivatives: exp⁽ⁿ⁾(0) = 1 for
  all n
\item
  Taylor polynomials: Pₙ(x) = Σₖ₌₀ⁿ xᵏ/k!
\item
  Remark: the series Σ xⁿ/n! converges for all x (ratio test:
  \textbar aₙ₊₁/aₙ\textbar{} = \textbar x\textbar/(n+1) → 0)
\item
  Theorem: exp(x) = Σ xⁿ/n! for all x (prove in text: Lagrange remainder
  → 0, since \textbar exp⁽ⁿ⁺¹⁾(c)\textbar{} = exp(c) is bounded on any
  fixed interval)
\item
  Remark: this is the series E(x) from ch5! The functional equation
  exp(x+y) = exp(x)exp(y) was already proved there via Cauchy product
\item
  ✎ Inline: verify the Lagrange remainder estimate in detail for a
  specific x
\end{itemize}

\subsection{The Number e}\label{the-number-e-2}

\begin{itemize}
\tightlist
\item
  e = exp(1) = Σ 1/n!
\item
  Theorem: e = lim(1 + 1/n)ⁿ (callback to ch2; proved via ch5's
  dominated convergence result that (1+x/n)ⁿ → E(x))
\item
  Error bound: 0 \textless{} e - Sₙ \textless{} 1/(n! · n) (prove in
  text)
\item
  Example: S₁₀ gives e ≈ 2.7182818 to 7 decimal places
\item
  Theorem: e is irrational (prove in text: multiply through by q!, get
  integer + something in (0,1))
\end{itemize}

\subsection{The Natural Logarithm}\label{the-natural-logarithm}

\begin{itemize}
\tightlist
\item
  Definition: log = exp⁻¹ (the inverse function of exp)
\item
  exp is strictly increasing and continuous, so by §6.4 the inverse
  exists and is continuous
\item
  log: (0, ∞) → ℝ, with log(eˣ) = x and e\^{}\{log y\} = y
\item
  Theorem: log(xy) = log(x) + log(y) for all x, y \textgreater{} 0
\item
  Proof: exp(log(x) + log(y)) = exp(log x) · exp(log y) = xy, so log(x)
  + log(y) = log(xy)
\item
  Theorem: log'(y) = 1/y
\item
  Proof: inverse function derivative (§7.1): log'(y) = 1/exp'(log y) =
  1/exp(log y) = 1/y
\item
  Basic properties: log(1) = 0, log(e) = 1, strictly increasing, log(x)
  → ∞ as x → ∞, log(x) → −∞ as x → 0⁺
\item
  The mystery constant identified: aˣ = exp(x log a) (since exp(log(aˣ))
  = aˣ and log(aˣ) = x log a by the functional equation). By chain rule:
  (aˣ)' = exp(x log a) · log a = aˣ · log a. Therefore a'(0) = log(a).
\item
  Remark: all exponentials and their derivatives are now completely
  determined --- (aˣ)' = aˣ log a
\item
  Theorem: log(x)/xᵅ → 0 as x → ∞ for any α \textgreater{} 0 (logarithms
  grow slower than any positive power)
\item
  Proof: L'Hôpital (§7.2): lim log(x)/xᵅ = lim (1/x)/(αxᵅ⁻¹) = lim
  1/(αxᵅ) = 0
\item
  ✎ Inline: show log(1/x) = −log(x) directly from the functional
  equation
\end{itemize}

\subsection{Guided Exercises}\label{guided-exercises-14}

\subsection{The Continued Fraction for
e}\label{the-continued-fraction-for-e}

\emph{(Fulfilling the promise from §4.4)}

\emph{(Details to be worked out --- the key idea is using the series for
e to derive the regular continued fraction {[}2; 1, 2, 1, 1, 4, 1, 1, 6,
\ldots{]})}

\subsection{Smooth ≠ Analytic}\label{smooth-analytic}

Let f(x) = e\^{}\{-1/x²\} for x ≠ 0, f(0) = 0.

\begin{enumerate}
\def\labelenumi{(\alph{enumi})}
\item
  Show f is continuous at 0.
\item
  Show f'(0) = 0 from the definition. (Hint: e\^{}\{-1/h²\}/h → 0 as h →
  0.)
\item
  Show f is smooth on ℝ ~\{0\} (standard rules).
\item
  Show by induction that f⁽ⁿ⁾(0) = 0 for all n.~(Strategy: show each
  derivative at nonzero x has the form p(1/x) · e\^{}\{-1/x²\} for some
  polynomial p, and these all → 0 as x → 0.)
\item
  Conclude: the Taylor series of f at 0 is identically 0, but f ≠ 0. So
  f is smooth but not analytic.
\item
  Reflect: power series are always equal to their Taylor series (§7.3
  corollary). This example shows general smooth functions need not be.
  What makes power series special?
\end{enumerate}

\subsection{Exercises}\label{exercises-32}

\emph{Exponential} - Prove exp(x) ≥ 1 + x for all x - Prove e
\textless{} 3 using partial sums - Show eˣ/xⁿ → ∞ as x → ∞ (exponential
beats any polynomial) - Compute e to 10 decimal places; how many terms
of the series are needed?

\emph{Logarithm} - Prove log(xⁿ) = n log(x) for n ∈ ℤ; extend to
log(x\^{}\{p/q\}) = (p/q) log(x) for rational exponents - Prove log(x) ≤
x − 1 for x \textgreater{} 0, with equality iff x = 1 - Prove log(x) ≥ 1
− 1/x for x \textgreater{} 0, with equality iff x = 1 - Taylor series
for log: show log(1+x) = x − x²/2 + x³/3 − x⁴/4 + ⋯ for
\textbar x\textbar{} \textless{} 1 (hint: log'(1+x) = 1/(1+x) =
Σ(−1)ⁿxⁿ; match coefficients using §7.3 term-by-term differentiation) -
Abel's theorem gives log(2) = 1 − 1/2 + 1/3 − 1/4 + ⋯. How many terms
for 2 decimal places? - Faster computation: show log((1+x)/(1−x)) = 2(x
+ x³/3 + x⁵/5 + ⋯) for \textbar x\textbar{} \textless{} 1. Use x = 1/3
to compute log(2) to 6 decimal places. How many terms are needed? -
Compute log(3) using the fast series with appropriate x - Prove Σ 1/n
diverges by comparing to ∫₁ⁿ 1/t dt (preview --- uses the integral
formula log(x) = ∫₁ˣ 1/t dt from §9.2)

\subsection{Note}\label{note}

The proof that aˣ is differentiable at 0 could potentially move to §7.1
if this section feels overloaded.

\subsection{Dependencies}\label{dependencies-32}

\textbf{Requires}: §4.4 (continued fractions --- for guided exercise),
§5.3 (Cauchy product --- for recognizing functional equation), §5.4
(dominated convergence --- for (1+x/n)ⁿ → E(x)), §6.4 (continuous
monotone inverses --- for log), §7.1 (derivative definition, chain rule,
inverse function derivative), §7.2 (L'Hôpital --- for log growth), §7.3
(term-by-term differentiation, Taylor's theorem with Lagrange remainder)

\textbf{Used in}: §7.5 (exp used in Newton's method convergence
analysis), §8.4 (∫E = (E(b)−E(a))/log(a)), §9.2--9.6 (log available for
examples and computations)

\section{§7.5 Newton's Method}\label{newtons-method}

\subsection{Narrative}\label{narrative-33}

\begin{itemize}
\tightlist
\item
  Newton's method connects differentiation back to the iteration and
  contraction ideas of Chapter 4
\item
  The algorithm is simple: follow the tangent line to the x-axis, repeat
\item
  Convergence follows from contraction (§4.3); the rate is quadratic,
  explained by Taylor (§7.3)
\item
  The Babylonian method from Chapter 2 is a special case --- the story
  comes full circle
\end{itemize}

\subsection{Content}\label{content-33}

\subsection{The Algorithm}\label{the-algorithm}

\begin{itemize}
\tightlist
\item
  Problem: find a root of f(x) = 0
\item
  Newton iteration: x\_\{n+1\} = x\_n - f(x\_n)/f'(x\_n)
\item
  Geometric interpretation: follow the tangent line at (x\_n, f(x\_n))
  to the x-axis
\end{itemize}

\subsection{Convergence via
Contraction}\label{convergence-via-contraction}

\begin{itemize}
\tightlist
\item
  Define g(x) = x - f(x)/f'(x); fixed points of g are roots of f
\item
  Compute: g'(x) = f(x)f'\,`(x)/(f'(x))²
\item
  At a simple root r: g'(r) = 0
\item
  Theorem: for x₀ sufficiently close to r, Newton's method converges
  (prove in text: g'(r) = 0 implies \textbar g'\textbar{} \textless{} 1
  near r, so g is a contraction; callback to §4.3)
\end{itemize}

\subsection{Quadratic Convergence via
Taylor}\label{quadratic-convergence-via-taylor}

\begin{itemize}
\tightlist
\item
  Theorem: ε\_\{n+1\} ≈ {[}f'\,`(r)/(2f'(r)){]} εₙ² (prove in text:
  Taylor expand f(xₙ) around r)
\item
  Consequence: errors square at each step --- the number of correct
  digits roughly doubles each iteration
\end{itemize}

\subsection{Examples}\label{examples}

\begin{itemize}
\tightlist
\item
  √2: x\_\{n+1\} = (xₙ + 2/xₙ)/2 --- the Babylonian method! (callback to
  ch2)
\item
  ∛5: x\_\{n+1\} = (2xₙ + 5/xₙ²)/3
\end{itemize}

\subsection{Failure Modes}\label{failure-modes}

\begin{itemize}
\item
  Multiple roots: g'(r) ≠ 0, convergence only linear
\item
  Bad starting points: iteration may cycle, diverge, or find the wrong
  root
\item
  f'(xₙ) = 0: division by zero
\item
  Remark: Newton's method is fixed-point iteration x\_\{n+1\} = g(xₙ).
  The same idea in function space --- Picard iteration --- solves
  differential equations; the contraction mapping theorem (§4.3)
  guarantees convergence in both settings.
\end{itemize}

\subsection{Guided Exercise}\label{guided-exercise-18}

\subsection{Division by
Multiplication}\label{division-by-multiplication}

Apply Newton's method to f(x) = 1/x - a (for fixed a \textgreater{} 0).

\begin{enumerate}
\def\labelenumi{(\alph{enumi})}
\item
  Derive the iteration x\_\{n+1\} = xₙ(2 - axₙ).
\item
  Show this computes 1/a using only multiplication and subtraction ---
  no division.
\item
  Prove convergence for suitable starting x₀.
\item
  Demonstrate quadratic convergence numerically for a specific value of
  a.
\item
  Remark: this is how computers actually perform division --- they
  implement multiplication in hardware and use Newton's method for
  reciprocals.
\end{enumerate}

\subsection{Exercises}\label{exercises-33}

\begin{itemize}
\tightlist
\item
  Apply Newton's method to find \(\sqrt[4]{7}\)
\item
  What happens applying Newton to \(f(x) = x^3\) starting at
  \(x_0 = 1\)? (A failure mode)
\item
  Modified Newton for multiple roots: if \(r\) is a root of multiplicity
  \(m\), show the iteration \(x_{n+1} = x_n - m \cdot f(x_n)/f'(x_n)\)
  restores quadratic convergence
\end{itemize}

\subsection{Notes}\label{notes}

\begin{itemize}
\tightlist
\item
  We want interesting Newton's method calculations to include as
  exercises; candidates will open up once trig and log are available
  (ch8/9). Flag for later: Newton on ln(x) = 1 to compute e efficiently.
\item
  Terminal section: nothing later depends on §7.5.
\end{itemize}

\subsection{Dependencies}\label{dependencies-33}

\textbf{Requires}: §4.3 (contraction mappings), §7.1 (derivative, chain
rule), §7.2 (MVT --- for bounding g'), §7.3 (Taylor's theorem --- for
quadratic convergence)

\textbf{Used in}: None (terminal section)

\chapter{Integration}\label{integration}

\textbf{The method of exhaustion} - What is the area under a curve? The
length of a curved line? These questions are ancient - Archimedes
(\textasciitilde250 BCE): computed the area of a parabolic segment by
trapping it between inscribed and circumscribed figures - Result: area
under y = x² on {[}0,1{]} equals 1/3 --- two thousand years before
Riemann sums - Key idea: squeeze the unknown between known bounds, show
the gap shrinks to zero - He also studied arc length of convex curves,
but needed an axiom he couldn't prove (we will prove it in §8.3)

\textbf{Cavalieri's indivisibles} - Cavalieri (1630s): imagined regions
as made of infinitely many parallel line segments - Compare two regions
by comparing their cross-sections at every height - Cavalieri's
principle: two solids with equal cross-sectional areas at every height
have equal volume - Philosophically suspect --- what is an ``infinitely
thin'' slice? --- but computationally powerful - We will make this
rigorous in §8.3

\textbf{Fermat's integration of powers} - Fermat (1630s): computed ∫₀¹
xⁿ = 1/(n+1) without calculus - Insight: use geometric partitions
(points at \ldots, r³, r², r, 1) instead of uniform ones - Upper sums
become geometric series; as r → 1, the sum approaches 1/(n+1) - Forty
years before Newton and Leibniz

\textbf{The message} - Integration was a computational subject for two
millennia before the Fundamental Theorem - Archimedes, Cavalieri, and
Fermat computed areas, volumes, and arc lengths by direct methods - The
common thread: trapping quantities between bounds and showing the gap
vanishes

\textbf{Our approach} - We formalize what they were doing: axiomatize
the integral, see what the axioms force - For continuous functions, the
axioms determine the value uniquely - Then we verify a construction
(Darboux) that satisfies the axioms - Direct computation is possible
(§8.4) but laborious --- a better method comes in Chapter 9

\section{The Axiomatic Integral}\label{the-axiomatic-integral}

\subsection{Narrative}\label{narrative-34}

\begin{itemize}
\tightlist
\item
  We don't construct the integral --- we ask what properties it must
  have
\item
  Three axioms (rectangles, monotonicity, additivity) are
  extraordinarily constraining
\item
  The trapping theorem: any integral satisfying the axioms is squeezed
  between upper and lower sums
\item
  For continuous functions, uniform continuity forces the squeeze shut
  --- the value is uniquely determined, independent of any theory of
  integration
\item
  But existence remains open: does an integral satisfying the three
  axioms actually exist? (→ §8.2)
\end{itemize}

\subsection{Content}\label{content-34}

\subsection{The Axioms}\label{the-axioms}

\begin{itemize}
\tightlist
\item
  Motivation: what properties should an integral have? We don't ask what
  it \emph{is} --- we ask how it must \emph{behave}
\item
  Definition: An \emph{integral} on {[}a,b{]} is an assignment
  ∫\_\{{[}a,b{]}\} f ∈ ℝ for certain bounded functions f, satisfying:

  \begin{itemize}
  \tightlist
  \item
    (I1) Rectangles: ∫\_\{{[}a,b{]}\} c = c(b−a) for any constant c
  \item
    (I2) Monotonicity: f ≤ g on {[}a,b{]} ⟹ ∫\emph{\{{[}a,b{]}\} f ≤
    ∫}\{{[}a,b{]}\} g
  \item
    (I3) Additivity: ∫\emph{\{{[}a,b{]}\} f = ∫}\{{[}a,c{]}\} f +
    ∫\_\{{[}c,b{]}\} f for any c ∈ (a,b)
  \end{itemize}
\item
  Theorem (Bounds): If m ≤ f ≤ M on {[}a,b{]} and f is integrable, then
  m(b−a) ≤ ∫\_\{{[}a,b{]}\} f ≤ M(b−a)
\item
  Proof: Apply monotonicity to m ≤ f ≤ M, then the rectangle axiom to ∫m
  and ∫M
\item
  Remark: the integral of any bounded integrable function is finite ---
  this follows from the axioms alone
\item
  ✎ Inline: If f is integrable on {[}0,1{]} with 2 ≤ f ≤ 5, what bounds
  can you give on ∫\_\{{[}0,1{]}\} f?
\item
  Definition (Improper integrals): If f is integrable on {[}a,c{]} for
  every c ∈ (a,b), define ∫\emph{\{{[}a,b)\} f = lim\_\{c→b⁻\}
  ∫\emph{\{{[}a,c{]}\} f when the limit exists. Similarly ∫}\{(a,b{]}\}
  f = lim}\{c→a⁺\} ∫\emph{\{{[}c,b{]}\} f.~For unbounded domains:
  ∫}\{{[}a,∞)\} f = lim\_\{b→∞\} ∫\_\{{[}a,b{]}\} f.
\end{itemize}

\subsection{The Trapping Theorem}\label{the-trapping-theorem}

\begin{itemize}
\tightlist
\item
  Idea: the bounds theorem uses one rectangle --- what if we use many?
\item
  Definition: A \emph{partition} of {[}a,b{]} is a finite set P = \{x₀,
  x₁, \ldots, xₙ\} with a = x₀ \textless{} x₁ \textless{} ⋯ \textless{}
  xₙ = b
\item
  Definition: For a bounded function f on {[}a,b{]} and a partition P,
  define mᵢ = inf f on {[}xᵢ₋₁, xᵢ{]}, Mᵢ = sup f on {[}xᵢ₋₁, xᵢ{]},
  and:

  \begin{itemize}
  \tightlist
  \item
    L(f, P) = Σ mᵢ(xᵢ − xᵢ₋₁) (lower sum)
  \item
    U(f, P) = Σ Mᵢ(xᵢ − xᵢ₋₁) (upper sum)
  \end{itemize}
\item
  Theorem (Trapping): If f is integrable, then L(f, P) ≤
  ∫\_\{{[}a,b{]}\} f ≤ U(f, P) for every partition P
\item
  Proof: Apply the bounds theorem on each subinterval {[}xᵢ₋₁, xᵢ{]},
  then sum using additivity
\item
  Remark: the bounds theorem is the special case P = \{a, b\}
\item
  Corollary: sup\_P L(f, P) ≤ ∫\_\{{[}a,b{]}\} f ≤ inf\_P U(f, P)
\item
  Observation: if sup\_P L(f, P) = inf\_P U(f, P), then there is exactly
  one value the integral can take --- the value is forced
\item
  ✎ Inline: Compute L(f, P) and U(f, P) for f(x) = x on {[}0,1{]} with
  the partition P = \{0, 1/2, 1\}
\end{itemize}

\subsection{Integrating Continuous
Functions}\label{integrating-continuous-functions}

\begin{itemize}
\tightlist
\item
  Question: for which functions does sup L = inf U?
\item
  Theorem: If f is continuous on {[}a,b{]}, then sup\_P L(f, P) = inf\_P
  U(f, P). Consequently, any integral satisfying the axioms assigns the
  same value to f.
\item
  Proof: f continuous on {[}a,b{]} (compact) ⟹ uniformly continuous.
  Given ε \textgreater{} 0, choose δ so \textbar x − y\textbar{}
  \textless{} δ ⟹ \textbar f(x) − f(y)\textbar{} \textless{} ε/(b−a).
  Take any partition with subintervals of length \textless{} δ. On each
  subinterval, Mᵢ − mᵢ \textless{} ε/(b−a). Then U − L = Σ(Mᵢ − mᵢ)(xᵢ −
  xᵢ₋₁) \textless{} ε. Since ε arbitrary, sup L = inf U.
\item
  Example: ∫\_\{{[}0,1{]}\} x = 1/2. Uniform partition Pₙ gives L =
  (n−1)/(2n), U = (n+1)/(2n), both → 1/2.
\item
  Example: ∫\_\{{[}0,1{]}\} x² = 1/3. Uniform partition Pₙ and Σi² =
  n(n+1)(2n+1)/6 gives U → 1/3. Similarly L.
\item
  Remark: the axioms \emph{force} these values --- we haven't
  constructed anything, just deduced what must be true
\item
  The question: we know what the integral must be for continuous
  functions. But does an integral satisfying the three axioms actually
  \emph{exist}? → §8.2
\item
  ✎ Inline: Use the uniform partition to show that if x³ is integrable
  on {[}0,1{]}, then ∫\_\{{[}0,1{]}\} x³ = 1/4
\end{itemize}

\subsection{Guided Exercises}\label{guided-exercises-15}

None for this section.

\subsection{Exercises}\label{exercises-34}

TBD

\subsection{Dependencies}\label{dependencies-34}

\textbf{Requires}: §1.3 (supremum and infimum), §6.4 (uniform
continuity), §6.5 (compactness of {[}a,b{]})

\textbf{Used in}: §8.2 (Darboux integral verifies existence), §8.3
(geometric applications), §9.1 (FTC)

\section{The Darboux Integral}\label{the-darboux-integral}

\subsection{Narrative}\label{narrative-35}

\begin{itemize}
\tightlist
\item
  §8.1 showed the axioms can force the value of an integral --- we
  define integrability to mean exactly this
\item
  A function is integrable if and only if its value is forced: sup L =
  inf U
\item
  This is the minimal integral: we declare a function integrable
  precisely when the axioms leave no choice
\item
  We then verify this definition actually satisfies the three axioms ---
  so the integral exists
\item
  Payoff: continuous functions are integrable (already proved), and so
  are some discontinuous ones
\end{itemize}

\subsection{Content}\label{content-35}

\subsection{The Minimal Integral}\label{the-minimal-integral}

\begin{itemize}
\tightlist
\item
  Idea: §8.1 ended with a question --- does an integral satisfying the
  axioms exist? We build one.
\item
  Strategy: declare a function integrable when the axioms force its
  value, and define the integral to be that forced value
\item
  Definition: lower integral = sup\_P L(f, P), upper integral = inf\_P
  U(f, P)
\item
  Definition: f is \emph{integrable} on {[}a,b{]} if the lower and upper
  integrals are equal. The common value is ∫\_\{{[}a,b{]}\} f.
\item
  Remark: for continuous functions, §8.1 already showed sup L = inf U
  --- so continuous functions are integrable under this definition
\item
  Remark: this is the Darboux integral
\end{itemize}

\subsection{Refinements}\label{refinements}

\begin{itemize}
\tightlist
\item
  Definition: a partition Q \emph{refines} P if P ⊆ Q
\item
  Lemma (Refinements Improve Bounds): if P ⊆ Q, then L(f, P) ≤ L(f, Q) ≤
  U(f, Q) ≤ U(f, P)
\item
  Proof: add one point at a time; splitting a subinterval can only
  improve inf/sup estimates
\item
  Corollary: every lower sum ≤ every upper sum (take common refinement P
  ∪ Q)
\item
  Corollary: lower integral ≤ upper integral
\end{itemize}

\subsection{Verification of the
Axioms}\label{verification-of-the-axioms}

\begin{itemize}
\tightlist
\item
  Theorem: the Darboux integral satisfies all three axioms
\item
  Verify (I1) Rectangles: for constant c, every partition gives L = U =
  c(b−a)
\item
  Verify (I2) Monotonicity: f ≤ g ⟹ inf/sup on each subinterval respects
  this ⟹ L(f,P) ≤ L(g,P), so sup L(f) ≤ sup L(g)
\item
  Verify (I3) Additivity: partitions of {[}a,b{]} containing c split
  into partitions of {[}a,c{]} and {[}c,b{]}; sup over product = sum of
  sups
\item
  Remark: the integral exists --- the forced value actually defines an
  integral satisfying the axioms
\end{itemize}

\subsection{Beyond Continuous
Functions}\label{beyond-continuous-functions}

\begin{itemize}
\tightlist
\item
  The Darboux definition applies to any bounded function with sup L =
  inf U
\item
  Theorem: bounded functions with finitely many discontinuities are
  integrable
\item
  Proof sketch: enclose discontinuities in intervals of small total
  length; on the rest, uniform continuity controls U − L
\item
  Theorem (exercise): monotone functions on {[}a,b{]} are integrable
\item
  Remark: we don't belabor integrability criteria here --- a fuller
  picture comes with the Lebesgue integral
\item
  ✎ Inline: explain why the Dirichlet function (1 on rationals, 0 on
  irrationals) is not integrable
\end{itemize}

\subsection{Guided Exercises}\label{guided-exercises-16}

None for this section.

\subsection{Exercises}\label{exercises-35}

TBD

\subsection{Dependencies}\label{dependencies-35}

\textbf{Requires}: §8.1 (axioms, trapping theorem, continuous functions
forced)

\textbf{Used in}: §8.3 (geometric applications), §8.4 (computing
integrals), §9.1 (FTC)

\section{Geometry}\label{geometry}

\subsection{Narrative}\label{narrative-36}

\begin{itemize}
\tightlist
\item
  Integration meets the geometric definitions from §1.5
\item
  We prove: the integral \emph{computes} what the supremum
  \emph{defines}
\item
  Arc length: §1.5 definition = ∫√(1+(f')²)
\item
  Area: §1.5 definition = ∫f
\item
  Having validated the integral against geometric intuition, we use it
  to \emph{define} volume
\item
  Archimedes' inequality: he needed it as an axiom, we can prove it
\end{itemize}

\subsection{Content}\label{content-36}

\subsection{Arc Length}\label{arc-length}

\begin{itemize}
\tightlist
\item
  Recall §1.5: arc length = sup of inscribed polygon lengths
\item
  For a partition P, the inscribed polygon has length
  \(\sum \sqrt{(\Delta x_i)^2 + (\Delta f_i)^2}\)
\item
  By MVT: \(\Delta f_i = f'(c_i) \Delta x_i\) for some \(c_i\)
\item
  So polygon length \(= \sum \sqrt{1 + (f'(c_i))^2} \cdot \Delta x_i\)
  --- a Riemann sum!
\item
  This is trapped between \(L(\sqrt{1+(f')^2}, P)\) and
  \(U(\sqrt{1+(f')^2}, P)\)
\item
  Theorem: If \(f'\) is continuous on \([a,b]\), then arc length
  \(= \int_{[a,b]} \sqrt{1 + (f'(x))^2} \, dx\)
\item
  Proof: Polygon lengths are Riemann sums trapped between L and U;
  continuity forces sup L = inf U; polygon lengths converge to this
  common value
\item
  Remark: This confirms that π (defined in §4.5 via polygon perimeters)
  equals the arc length of the unit semicircle
\item
  Example: Compute arc length of \(y = x^{3/2}\) on \([0, 1]\)
\item
  ✎ Inline: Set up the arc length integral for \(y = x^2\) on \([0, 1]\)
\end{itemize}

\subsection{Area}\label{area}

\begin{itemize}
\tightlist
\item
  Recall §1.5: area exists when inner (inscribed) and outer
  (circumscribed) polygon areas agree
\item
  Lower sums use inscribed rectangles; upper sums use circumscribed
  rectangles
\item
  Rectangles are polygons, so:
  \(L(f, P) \leq \text{inner area} \leq \text{outer area} \leq U(f, P)\)
\item
  Theorem: If \(f \geq 0\) is continuous on \([a,b]\), the region under
  the curve has area \(= \int_{[a,b]} f(x) \, dx\)
\item
  Proof: sup L ≤ inner ≤ outer ≤ inf U; continuity gives sup L = inf U;
  all four are equal
\item
  Area between curves: if \(f \geq g\) on \([a,b]\), area
  \(= \int_{[a,b]} [f(x) - g(x)] \, dx\)
\item
  Example: Area of unit disk
  \(= \int_{[-1,1]} 2\sqrt{1-x^2} \, dx = \pi\) (same constant!)
\item
  ✎ Inline: Set up the integral for area between \(y = x\) and
  \(y = x^2\) on \([0,1]\)
\end{itemize}

\subsection{Archimedes' Inequality}\label{archimedes-inequality}

\begin{itemize}
\tightlist
\item
  Archimedes needed this as an axiom for his method; we can prove it
\item
  Theorem: Among convex curves with the same endpoints, the one closer
  to the chord is shorter
\item
  Proof: see guided exercise below (uses convexity from Chapter 7;
  details TBD)
\item
  Corollary: The straight line is the shortest path between two points
\item
  Corollary: A convex curve's length is trapped between inscribed and
  circumscribed polygon lengths
\item
  Corollary: The two definitions of π agree (polygon limits from §4.5 =
  arc length integral)
\item
  Remark: This validates Archimedes' method --- his axiom is now a
  theorem
\end{itemize}

\subsection{Volume}\label{volume}

\begin{itemize}
\tightlist
\item
  The integral agrees with geometric intuition for length and area
\item
  For volume, polygonal approximations become unwieldy (polyhedra in 3D)
\item
  We take the integral as our \emph{definition} of volume
\item
  Definition (Volume by slicing): A solid with cross-sectional area
  \(A(x)\) at position \(x\) has volume \(\int_{[a,b]} A(x) \, dx\)
\item
  Theorem (Cavalieri's Principle): Solids with equal cross-sectional
  areas at every height have equal volumes
\item
  Proof: immediate from the definition --- same integrand gives same
  integral
\item
  Example (Pyramid): Base \(s \times s\), height \(h\). Cross-section at
  height \(x\): square of side \(sx/h\). Volume
  \(= \int_{[0,h]} (sx/h)^2 \, dx = \frac{s^2 h}{3} = \frac{1}{3} \cdot \text{base} \cdot \text{height}\)
\item
  Example (Cone): Radius \(r\), height \(h\). Cross-section at height
  \(x\): disk of radius \(rx/h\). Volume
  \(= \int_{[0,h]} \pi(rx/h)^2 \, dx = \frac{\pi r^2 h}{3}\)
\item
  Example (Sphere): Radius \(R\). Cross-section at height \(x\): disk of
  radius \(\sqrt{R^2 - x^2}\). Volume
  \(= \int_{[-R,R]} \pi(R^2 - x^2) \, dx = \frac{4}{3}\pi R^3\)
\end{itemize}

\subsection{Surface Area of
Revolution}\label{surface-area-of-revolution}

\begin{itemize}
\tightlist
\item
  Rotate \(y = f(x)\) around the \(x\)-axis
\item
  Approximate by rotating line segments --- gives frustums of cones
\item
  Definition: Surface area
  \(= \int_{[a,b]} 2\pi f(x) \sqrt{1 + (f'(x))^2} \, dx\)
\item
  Remark: Combines arc length integrand with circumference \(2\pi f(x)\)
\item
  Example: Surface area of sphere --- rotate \(y = \sqrt{R^2 - x^2}\)
  around \(x\)-axis
\item
  ✎ Inline: Set up (but do not evaluate) the surface area integral for
  rotating \(y = x^2\) on \([0, 1]\)
\end{itemize}

\subsection{Guided Exercises}\label{guided-exercises-17}

\subsection{Archimedes' Inequality}\label{archimedes-inequality-1}

\emph{Setup:} Let f and g be convex functions on {[}a,b{]} with f(a) =
g(a), f(b) = g(b), and f ≤ g on (a,b).

\emph{Goal:} Prove that the arc length of f ≤ the arc length of g (the
curve closer to the chord is shorter).

\emph{Steps:} TBD --- requires careful argument using convexity from
Chapter 7.

\emph{Corollaries:} (a) The straight line is the shortest path between
two points (b) A convex curve's length is trapped between inscribed and
circumscribed polygon lengths (c) The two definitions of π agree
(polygon limits from §4.5 = arc length integral)

\subsection{Volumes in Higher
Dimensions}\label{volumes-in-higher-dimensions}

\emph{Setup:} - Define n-ball:
\(B^n(R) = \{x \in \mathbb{R}^n : |x| \leq R\}\) - Let \(V_n(R)\) =
n-dimensional volume of \(B^n(R)\)

\begin{enumerate}
\def\labelenumi{(\alph{enumi})}
\item
  Verify base cases: \(V_1(R) = 2R\) (length of interval \([-R, R]\))
  and \(V_2(R) = \pi R^2\) (area of disk).
\item
  Explain why slicing \(B^n(R)\) perpendicular to the \(x_1\)-axis at
  height \(t \in [-R, R]\) gives an \((n-1)\)-ball of radius
  \(\sqrt{R^2 - t^2}\).
\item
  Write the slicing formula:
  \(V_n(R) = \int_{-R}^{R} V_{n-1}(\sqrt{R^2 - t^2}) \, dt\)
\item
  Define \(C_n = V_n(1)\). Substitute \(t = Ru\) in the integral from
  (c) to show that \(V_n(R) = C_n R^n\).
\item
  Substitute \(V_{n-1}(r) = C_{n-1} r^{n-1}\) into (c) and use the
  substitution from (d) to derive:
  \[C_n = C_{n-1} \int_{-1}^{1} (1 - u^2)^{(n-1)/2} \, du\]
\item
  Define \(\omega_n = \int_{-1}^{1} (1 - u^2)^{(n-1)/2} \, du\), so
  \(C_n = \omega_n \cdot C_{n-1}\).
\item
  Verify:
  \(\omega_2 = \int_{-1}^{1} \sqrt{1-u^2} \, du = \frac{\pi}{2}\) (area
  of unit semicircle). Check this gives
  \(C_2 = \omega_2 \cdot C_1 = \frac{\pi}{2} \cdot 2 = \pi\). ✓
\item
  Compute \(\omega_3 = \int_{-1}^{1} (1-u^2) \, du\) directly and verify
  \(C_3 = \frac{4\pi}{3}\).
\end{enumerate}

\emph{In Chapter 10, we evaluate \(\omega_n\) for general \(n\) and
derive formulas for \(C_n\).}

\subsection{Exercises}\label{exercises-36}

\emph{Arc length} - Set up the arc length integral for \(y = x^2\) on
\([0, 1]\) - Prove: a curve with continuous \(f'\) has finite arc length
(is rectifiable)

\emph{Area} - Set up the integral for area between \(y = x\) and
\(y = x^2\) on \([0,1]\) - Set up the integral for area enclosed by the
ellipse \(\frac{x^2}{a^2} + \frac{y^2}{b^2} = 1\)

\emph{Volume} - Set up the slicing integral for a pyramid with square
base \(s \times s\) and height \(h\) - Set up the slicing integral for a
cone with radius \(r\) and height \(h\) - Set up the slicing integral
for a sphere of radius \(R\) - A solid has base the unit disk and
cross-sections perpendicular to the \(x\)-axis are squares. Set up the
volume integral. - ★ Use Cavalieri's principle (without computing
integrals) to show a hemisphere has the same volume as a cylinder minus
a cone

\emph{Surface area} - Set up the surface area integral for a sphere of
radius \(R\) - Set up the surface area integral for the cone formed by
rotating \(y = x\) on \([0, h]\)

\emph{Higher dimensions} - Set up (but do not evaluate)
\(\omega_4 = \int_{-1}^{1} (1-u^2)^{3/2} \, du\)

\subsection{Dependencies}\label{dependencies-36}

\textbf{Requires}: §1.5 (arc length and area definitions), §4.5
(definition of π), §7 (convexity), §8.1-8.2 (integration)

\textbf{Used in}: Chapter 10 (π computations, higher-dimensional
volumes)

\section{Computing Integrals}\label{computing-integrals}

\subsection{Narrative}\label{narrative-37}

\begin{itemize}
\tightlist
\item
  We compute integrals directly from the definition --- no shortcuts
\item
  Powers via sum formulas, exponentials via geometric series
\item
  These calculations are possible but laborious; a better method comes
  in Chapter 9
\end{itemize}

\subsection{Content}\label{content-37}

\subsection{The Direct Method}\label{the-direct-method}

\begin{itemize}
\tightlist
\item
  Lemma: to compute ∫\_\{{[}a,b{]}\} f, it suffices to find any sequence
  of partitions Pₙ with U(f, Pₙ) − L(f, Pₙ) → 0. The common limit is the
  integral.
\item
  Strategy: choose partitions, compute L and U explicitly, show both
  converge to the same limit
\end{itemize}

\subsection{Powers}\label{powers}

\begin{itemize}
\tightlist
\item
  Theorem: ∫\_\{{[}0,b{]}\} xⁿ = bⁿ⁺¹/(n+1) for any n ∈ ℕ
\item
  Proof: uniform partitions on {[}0,b{]}, monotonicity gives explicit L
  and U, compute using sum formulas
\item
  Extend to ∫\_\{{[}a,b{]}\} xⁿ = (bⁿ⁺¹ − aⁿ⁺¹)/(n+1) by subdivision
\item
  ✎ Inline: verify explicitly for n = 1 using Σi = n(n+1)/2 and for n =
  2 using Σi² = n(n+1)(2n+1)/6
\end{itemize}

\subsection{Exponentials}\label{exponentials}

\begin{itemize}
\tightlist
\item
  Theorem: if E is an exponential function, then E is integrable on
  {[}a,b{]} and ∫\_\{{[}a,b{]}\} E = (E(b) − E(a))/E'(0)
\item
  Proof: uniform partitions; law of exponents turns L into a geometric
  series; closed form via geometric sum formula; limit involves (E(Δₙ) −
  1)/Δₙ → E'(0); upper sum satisfies U(E, Pₙ) = E(Δₙ)·L(E, Pₙ), and
  E(Δₙ) → 1, so U and L have the same limit
\item
  Corollary: ∫\_\{{[}a,b{]}\} exp = exp(b) − exp(a)
\end{itemize}

\subsection{Guided Exercises}\label{guided-exercises-18}

None for this section.

\subsection{Exercises}\label{exercises-37}

TBD

\subsection{Dependencies}\label{dependencies-37}

\textbf{Requires}: §8.1--8.2 (integration theory), §2 (geometric
series), §7 (exponential functions, E'(0))

\textbf{Used in}: §9.1 (FTC)

\chapter{Calculus}\label{calculus}

\textbf{The missing link} - By the 1650s, mathematicians could
differentiate (Fermat's tangent lines, Chapter 7) and integrate
(Cavalieri's areas, Fermat's power sums, Chapter 8) --- but these were
separate subjects - Each problem required its own trick: a clever
algebraic manipulation for a tangent, a geometric partition for an area
- Nobody suspected the two were related

\textbf{Barrow's near miss} - Isaac Barrow (1660s), Newton's teacher at
Cambridge, proved a geometric theorem: if you take the area under a
curve and draw the tangent to the area function, you recover the
original curve - This \emph{is} the Fundamental Theorem --- but stated
geometrically, buried in a book of optics, and Barrow didn't see its
computational power - He resigned his professorship in favor of his
27-year-old student, who did

\textbf{Newton and Leibniz} - Newton (1660s, published much later):
discovered that finding areas reduces to ``reversing'' differentiation -
Leibniz (1670s): arrived at the same insight independently, with better
notation --- his ∫ and d/dx made the inverse relationship visible on the
page - The consequence was immediate and explosive: every derivative
formula in Chapter 7 became an integration formula. The hard
computations of Chapter 8 --- pages of Riemann sum estimates ---
collapsed to one-line calculations - The priority dispute between Newton
and Leibniz consumed decades and split European mathematics, but the
theorem itself was never in doubt

\textbf{What FTC unlocks} - Existence of antiderivatives: every
continuous function has one, guaranteed --- even when no formula can be
written down. Functions like the error function (probability), the sine
integral (signal processing), and the logarithmic integral (prime
numbers) are \emph{defined} this way - Differential equations: FTC turns
existence questions into integration problems. Wide classes of ODEs that
model physical phenomena can be solved explicitly, and where explicit
solutions fail, FTC still guarantees solutions exist - Rigorous
trigonometry: the integral ∫₀ʸ 1/√(1−t²) dt gives arc length a precise
definition, and inverting it constructs sin and cos from scratch --- two
millennia of geometric trigonometry, rebuilt from calculus - New tools
for computing π: Archimedes needed a 96-sided polygon for two decimal
places. With FTC and power series, Machin computed 100 digits by hand in
1706; the Gregory-Leibniz series reveals π as a sum of simple fractions;
the Wallis product builds it from integers alone

\textbf{Our approach} - We prove FTC, develop its computational
consequences, then push into the territory it opens up: differential
equations, trigonometry, and the computation of π

\section{§9.1 The Fundamental Theorem}\label{the-fundamental-theorem}

\subsection{Narrative}\label{narrative-38}

\begin{itemize}
\tightlist
\item
  In Chapter 8 we defined ∫\_\{{[}a,b{]}\} f --- the integral over a
  set. Now we let the upper bound vary and study F(x) = ∫ₐˣ f as a
  function
\item
  Signed notation lets x move freely, even before a
\item
  Two properties of F come quickly: continuity (from bounds) and the MVT
  for integrals (average value achieved)
\item
  Then the main event: FTC I says F' = f (differentiation undoes
  integration), FTC II says ∫f = F(b) − F(a) (the evaluation formula)
\item
  The chapter 8 Riemann sum calculations are instantly explained
\end{itemize}

\subsection{Content}\label{content-38}

\subsection{The Integral Function}\label{the-integral-function}

\begin{itemize}
\tightlist
\item
  In ch8 we defined ∫\_\{{[}a,b{]}\} f for a \textless{} b --- the
  integral over a set
\item
  Fix a basepoint a and let x vary: for x ≥ a, define F(x) =
  ∫\_\{{[}a,x{]}\} f
\item
  To let x move freely (including before a), introduce signed notation:

  \begin{itemize}
  \tightlist
  \item
    Definition: ∫ₐᵇ f = ∫\_\{{[}a,b{]}\} f for a \textless{} b
  \item
    Definition: ∫ᵇₐ f = −∫ₐᵇ f
  \item
    Definition: ∫ₐₐ f = 0
  \end{itemize}
\item
  Now F(x) = ∫ₐˣ f is defined for all x in the domain of f
\item
  Theorem (Additivity): ∫ₐᶜ f = ∫ₐᵇ f + ∫ᵇᶜ f for any ordering of a, b,
  c
\item
  Proof: reduces to the ch8 additivity axiom when a \textless{} b
  \textless{} c; the other cases follow from the sign convention
\item
  ✎ Inline: verify additivity when c \textless{} a \textless{} b
\end{itemize}

\subsection{Properties of the Integral
Function}\label{properties-of-the-integral-function}

\begin{itemize}
\tightlist
\item
  Theorem (Continuity): If f is integrable with m ≤ f ≤ M, then F(x) =
  ∫ₐˣ f is continuous
\item
  Proof: F(x+h) − F(x) = ∫ₓˣ⁺ʰ f.~By the bounds theorem (§8.1), m·h ≤
  F(x+h) − F(x) ≤ M·h. As h → 0, both bounds → 0, so F(x+h) → F(x).
\item
  Theorem (MVT for Integrals): If f is continuous on {[}a,b{]}, then
  there exists c ∈ (a,b) with ∫ₐᵇ f = f(c)(b−a)
\item
  Proof: By EVT, f attains its min m and max M on {[}a,b{]}. Bounds
  theorem gives m(b−a) ≤ ∫ₐᵇ f ≤ M(b−a). The average value (1/(b−a)) ∫ₐᵇ
  f lies between m and M. By IVT, f achieves this value at some c ∈
  (a,b).
\item
  Remark: the \emph{average value} of f on {[}a,b{]} is (1/(b−a)) ∫ₐᵇ f,
  and f always equals its average somewhere
\item
  ✎ Inline: find c explicitly for f(x) = x² on {[}0,1{]}
\end{itemize}

\subsection{The Fundamental Theorem, Part
I}\label{the-fundamental-theorem-part-i}

\begin{itemize}
\tightlist
\item
  Theorem (FTC I): If f is continuous on an interval containing a, then
  F(x) = ∫ₐˣ f is differentiable and F'(x) = f(x)
\item
  Proof: F'(x) = lim\_\{h→0\} (1/h)(F(x+h) − F(x)) = lim\_\{h→0\} (1/h)
  ∫ₓˣ⁺ʰ f.~By MVT for integrals, ∫ₓˣ⁺ʰ f = f(c)·h for some c between x
  and x+h. As h → 0, c → x, and continuity of f gives f(c) → f(x). So
  F'(x) = f(x).
\item
  Corollary: every continuous function has an antiderivative
\item
  Remark: this is an existence theorem --- we assert F exists without
  finding a formula
\item
  ✎ Inline: what is F'(x) when F(x) = ∫₀ˣ eᵗ² dt? (cannot find F
  explicitly, but FTC I answers immediately)
\end{itemize}

\subsection{The Fundamental Theorem, Part
II}\label{the-fundamental-theorem-part-ii}

\begin{itemize}
\tightlist
\item
  Theorem (FTC II): If f is continuous on {[}a,b{]} and F is any
  antiderivative of f (meaning F' = f), then ∫ₐᵇ f = F(b) − F(a)
\item
  Proof: Let G(x) = ∫ₐˣ f.~By FTC I, G' = f = F'. So (G − F)' = 0, hence
  G − F is constant by MVT (§7.2). Evaluating at x = a: G(a) − F(a) = 0
  − F(a) = −F(a), so G(x) = F(x) − F(a) for all x. At x = b: ∫ₐᵇ f =
  G(b) = F(b) − F(a).
\item
  Notation: F(x)\textbar ₐᵇ = F(b) − F(a), read ``F evaluated from a to
  b''
\item
  Callback: in ch8 we computed ∫\_\{{[}0,1{]}\} eˣ = e − 1 by Riemann
  sums and observed this equals e¹ − e⁰. FTC II explains why: (eˣ)' =
  eˣ, so ∫₀¹ eˣ dx = eˣ\textbar₀¹ = e − 1. The pattern was the
  Fundamental Theorem.
\item
  ✎ Inline: use FTC II to compute ∫₀¹ x³ dx and verify it matches the
  ch8 Riemann sum prediction
\end{itemize}

\subsection{Guided Exercises}\label{guided-exercises-19}

None for this section.

\subsection{Exercises}\label{exercises-38}

\emph{FTC I} - Compute d/dx ∫₀ˣ² eᵗ² dt (chain rule + FTC I) - If F(x) =
∫₀ˣ f and F is differentiable with F' continuous, what can you say about
f?

\emph{FTC II} - Use FTC II to compute ∫₀² x³ dx, ∫₁⁴ √x dx, ∫₀¹ (3x² +
2x) dx - Prove: if F' = G' = f on {[}a,b{]}, then F and G differ by a
constant

\emph{Integral function} - Prove: if f ≥ 0 is continuous on {[}a,b{]}
with ∫ₐᵇ f = 0, then f ≡ 0 - Prove: if f is continuous and ∫ₐᵇ f = 0 for
all b, then f ≡ 0

\subsection{Dependencies}\label{dependencies-38}

\textbf{Requires}: §7.2 (MVT --- for uniqueness of antiderivatives),
§8.1 (axiomatic integral, bounds theorem, additivity), §8.2 (Darboux
integral existence), §6.4 (EVT, IVT --- for MVT for integrals)

\textbf{Used in}: §9.2 (antidifferentiation techniques), §9.3 (ODE
existence), §9.4--9.5 (defining functions via integrals)

\section{§9.2 Antidifferentiation}\label{antidifferentiation}

\subsection{Narrative}\label{narrative-39}

\begin{itemize}
\tightlist
\item
  FTC II turns integration into a search for antiderivatives --- but how
  do we find them?
\item
  Each technique reverses a differentiation rule: linearity reverses
  linearity, substitution reverses the chain rule, parts reverses the
  product rule
\item
  For power series, term-by-term integration completes a trilogy running
  through ch6--ch7--ch9
\item
  Taylor's integral remainder connects back to ch7 and demystifies the
  Lagrange form
\item
  The section closes with a payoff: the integral formula for log, giving
  a concrete formula for a function we previously only knew abstractly
\end{itemize}

\subsection{Content}\label{content-39}

\subsection{Term-by-Term Integration}\label{term-by-term-integration}

\begin{itemize}
\tightlist
\item
  Corollary (Linearity): ∫ₐᵇ (αf + βg) = α ∫ₐᵇ f + β ∫ₐᵇ g
\item
  Proof: If F' = f and G' = g, then (αF + βG)' = αf + βg. Apply FTC II.
\item
  Theorem (Power Series III): If f(x) = Σ aₙxⁿ has radius of convergence
  R, then for \textbar x\textbar{} \textless{} R: ∫₀ˣ f(t) dt = Σ
  aₙxⁿ⁺¹/(n+1) and this series also has radius R.
\item
  Proof: Define G(x) = Σ aₙxⁿ⁺¹/(n+1). Same radius of convergence (ratio
  test). By Power Series II (§7.3), G'(x) = Σ aₙxⁿ = f(x). By FTC II,
  ∫₀ˣ f = G(x) − G(0) = G(x).
\item
  Remark: The Power Series Trilogy --- I: continuous (§6.3), II:
  differentiable term-by-term (§7.3), III: integrable term-by-term
  (§9.2). Power series behave perfectly under all basic operations of
  calculus.
\item
  ✎ Inline: integrate the geometric series 1/(1−x) = Σxⁿ term-by-term to
  obtain −log(1−x) = Σ xⁿ⁺¹/(n+1)
\end{itemize}

\subsection{Integration by
Substitution}\label{integration-by-substitution}

\begin{itemize}
\tightlist
\item
  Theorem: If g is differentiable with continuous derivative and f is
  continuous, then ∫ₐᵇ f(g(x)) g'(x) dx = ∫\_\{g(a)\}\^{}\{g(b)\} f(u)
  du
\item
  Proof: Let F be an antiderivative of f (exists by FTC I). By the chain
  rule, (F ∘ g)' = (f ∘ g) · g'. Apply FTC II to both sides.
\item
  Remark: substitution is the chain rule in reverse
\item
  Example: ∫₀¹ x · eˣ² dx. Let u = x², du = 2x dx. Then ∫₀¹ x · eˣ² dx =
  ½ ∫₀¹ eᵘ du = ½(e − 1).
\item
  Example: ∫₀¹ x²√(1 − x³) dx. Let u = 1 − x³, du = −3x² dx. Then ∫₀¹
  x²√(1−x³) dx = −⅓ ∫₁⁰ √u du = ⅓ ∫₀¹ u\^{}\{1/2\} du = ⅓ · ⅔ = 2/9.
\item
  Example: ∫₁ᵉ log(t)/t dt. Let u = log(t), du = dt/t. Then ∫₁ᵉ log(t)/t
  dt = ∫₀¹ u du = 1/2.
\item
  ✎ Inline: compute ∫₀¹ x/(1+x²)² dx by substitution
\end{itemize}

\subsection{Integration by Parts}\label{integration-by-parts}

\begin{itemize}
\tightlist
\item
  Theorem: ∫ₐᵇ f · g' = fg\textbar ₐᵇ − ∫ₐᵇ f' · g
\item
  Proof: By the product rule, (fg)' = f'g + fg'. Rearrange: fg' = (fg)'
  − f'g. Integrate both sides and apply FTC II.
\item
  Remark: integration by parts is the product rule in reverse
\item
  Example: ∫₀¹ x · eˣ dx. Let f(x) = x, g'(x) = eˣ, so f'(x) = 1, g(x) =
  eˣ. Then ∫₀¹ x eˣ dx = xeˣ\textbar₀¹ − ∫₀¹ eˣ dx = e − (e − 1) = 1.
\item
  Example: ∫₀¹ x² · eˣ dx. Apply IBP twice: let f = x², g' = eˣ first,
  then f = 2x, g' = eˣ. Result: e − 2.
\item
  ✎ Inline: derive a reduction formula for ∫ xⁿ eˣ dx
\end{itemize}

\subsection{Taylor's Theorem with Integral
Remainder}\label{taylors-theorem-with-integral-remainder}

\begin{itemize}
\tightlist
\item
  Theorem: If f has n+1 continuous derivatives on an interval containing
  a, then: f(x) = Σₖ₌₀ⁿ f⁽ᵏ⁾(a)(x−a)ᵏ/k! + (1/n!) ∫ₐˣ (x−t)ⁿ f⁽ⁿ⁺¹⁾(t)
  dt
\item
  Proof: Start from FTC II: f(x) = f(a) + ∫ₐˣ f'(t) dt. Apply IBP with u
  = f'(t), dv = dt written as d(t−x), to get f(a) + f'(a)(x−a) + ∫ₐˣ
  (x−t) f'\,'(t) dt. Continue inductively.
\item
  Connection to Lagrange form: applying MVT for integrals to the
  integral remainder gives (1/n!) ∫ₐˣ (x−t)ⁿ f⁽ⁿ⁺¹⁾(t) dt =
  f⁽ⁿ⁺¹⁾(c)/(n+1)! · (x−a)ⁿ⁺¹ for some c between a and x. This is the
  Lagrange remainder from §7.3 --- the mystery point c is demystified.
\item
  Remark: the integral form is often more useful for estimates since it
  avoids the unknown c
\item
  ✎ Inline: write out the integral remainder for eˣ at a = 0 and verify
  it matches the Lagrange remainder
\end{itemize}

\subsection{The Integral Formula for
log}\label{the-integral-formula-for-log}

\begin{itemize}
\tightlist
\item
  Theorem: log(x) = ∫₁ˣ 1/t dt for all x \textgreater{} 0
\item
  Proof: We know log'(x) = 1/x (§7.4) and log(1) = 0. By FTC II, ∫₁ˣ 1/t
  dt = log(t)\textbar₁ˣ = log(x) − log(1) = log(x).
\item
  Remark: this gives a concrete computational formula for log ---
  previously we could only evaluate it via the Taylor series (§7.4
  exercises). Now we also have an integral representation.
\item
  Example: ∫₁ᵉ 1/t dt = log(e) = 1
\item
  Example: ∫₁ᵉ log(x) dx. IBP with f = log(x), g' = 1: result is 1.
\item
  ✎ Inline: compute ∫ log(x) dx by parts (find the antiderivative)
\end{itemize}

\subsection{Guided Exercises}\label{guided-exercises-20}

\subsection{Integration Discovers the
Logarithm}\label{integration-discovers-the-logarithm}

Forget everything you know about log. Start from scratch with
integration alone.

\begin{enumerate}
\def\labelenumi{(\alph{enumi})}
\item
  Define L(x) = ∫₁ˣ 1/t dt for x \textgreater{} 0. Explain why L is
  well-defined for all x \textgreater{} 0 (including x \textless{} 1).
\item
  Compute L'(x) directly from FTC I. Conclude L is strictly increasing.
\item
  Prove L(xy) = L(x) + L(y) for all x, y \textgreater{} 0. (Hint: write
  L(xy) = ∫₁ˣ 1/t dt + ∫ₓˣʸ 1/t dt, then substitute u = t/x in the
  second integral.)
\item
  Prove L(1) = 0 and L(1/x) = −L(x).
\item
  Prove L(xⁿ) = nL(x) for n ∈ ℤ.
\item
  A function satisfying f(xy) = f(x) + f(y), f continuous and
  nonconstant, is called a \emph{logarithm}. You have just proved L is a
  logarithm, using nothing but FTC and substitution.
\item
  Identify L: its inverse E = L⁻¹ satisfies E'(y) = E(y) (by the inverse
  function derivative and L'(x) = 1/x). From §7.4, the unique such
  function with E(0) = 1 is exp. Since L(1) = 0, we have E(0) = 1.
  Therefore L = log.
\item
  Reflect: two completely independent paths --- one through
  differentiation (§7.4), one through integration (here) --- construct
  the same function. The Fundamental Theorem guarantees they must agree.
\end{enumerate}

\subsection{Exercises}\label{exercises-39}

\emph{Substitution} - Compute ∫₀¹ x³ eˣ⁴ dx - Compute ∫₀¹ x(1 + x²)⁵ dx
- Compute ∫₀¹ x² / (1 + x³)² dx - Compute ∫₁ᵉ² 1/(t√(log t)) dt

\emph{Integration by parts} - Compute ∫₀¹ x³ eˣ dx - Prove the reduction
formula: ∫₀¹ xⁿ eˣ dx = e − n ∫₀¹ xⁿ⁻¹ eˣ dx - Use the reduction formula
to evaluate ∫₀¹ x⁴ eˣ dx

\emph{Taylor remainder} - Use the integral remainder to bound
\textbar eˣ − Σₖ₌₀ⁿ xᵏ/k!\textbar{} on {[}0, 1{]} - Use the integral
remainder to prove eˣ = Σ xⁿ/n! for all x (reprove via integral form)

\emph{Logarithm} - Compute ∫₁ᵉ² 1/t dt - Compute ∫₁ᵉ (log x)² dx using
IBP - Prove Σ 1/n diverges by comparing to ∫₁ⁿ 1/t dt = log(n) (integral
test preview) - AM-GM from Jensen: derive the AM-GM inequality by
applying Jensen's inequality (§7.2 exercises) to f(x) = −log(x)

\emph{After §9.4}: additional substitution and IBP exercises involving
trig functions appear in §9.4 exercises.

\subsection{Dependencies}\label{dependencies-39}

\textbf{Requires}: §7.1 (chain rule, product rule), §7.3 (Power Series
II, Taylor's theorem with Lagrange remainder), §7.4 (log and its
derivative), §9.1 (FTC I and II)

\textbf{Used in}: §9.3 (ODE solutions via integration), §9.4 (trig via
integrals), §9.5 (π computations)

\section{\texorpdfstring{§9.3 Differential Equations
\emph{(Terminal)}}{§9.3 Differential Equations (Terminal)}}\label{differential-equations-terminal}

\subsection{Narrative}\label{narrative-40}

\begin{itemize}
\tightlist
\item
  FTC says every continuous function has an antiderivative --- this
  immediately solves certain differential equations
\item
  Three progressively complex classes where existence follows from our
  tools: quadrature, first-order linear, separable
\item
  Along the way we meet functions defined purely by integrals --- they
  exist by FTC even though no formula is possible
\item
  The section closes by showing where elementary methods fail,
  previewing the general existence theory in Chapter 12
\item
  Terminal section: nothing later depends on §9.3
\end{itemize}

\subsection{Content}\label{content-40}

\subsection{Quadrature}\label{quadrature}

\begin{itemize}
\tightlist
\item
  The simplest ODE: y' = f(x), y(a) = y₀
\item
  Theorem: If f is continuous, then y(x) = y₀ + ∫ₐˣ f(t) dt is the
  unique solution
\item
  Existence: FTC I gives y' = f
\item
  Uniqueness: if y₁, y₂ both solve it, (y₁ − y₂)' = 0, so y₁ − y₂
  constant by MVT (§7.2); equal at a, so identical
\item
  Remark: FTC guarantees existence of antiderivatives even when no
  elementary formula exists. This leads to important functions
  \emph{defined} as integrals:
\item
  The error function: erf(x) = (2/√π) ∫₀ˣ e⁻ᵗ² dt --- fundamental in
  probability (cumulative normal distribution)
\item
  The sine integral: Si(x) = ∫₀ˣ (sin t)/t dt --- appears in signal
  processing and Fourier analysis (sin t/t has a removable singularity
  at 0)
\item
  The Fresnel integrals: S(x) = ∫₀ˣ sin(t²) dt, C(x) = ∫₀ˣ cos(t²) dt
  --- describe diffraction patterns in optics
\item
  The logarithmic integral: li(x) = ∫₂ˣ 1/(log t) dt --- gives the best
  simple approximation to the prime counting function (the Prime Number
  Theorem says π(x) \textasciitilde{} li(x))
\item
  ✎ Inline: verify erf'(x) = (2/√π) e⁻ˣ² and erf(0) = 0 directly from
  the definition
\end{itemize}

\subsection{First-Order Linear
Equations}\label{first-order-linear-equations}

\begin{itemize}
\tightlist
\item
  The equation: y' + P(x)y = Q(x), y(a) = y₀, with P and Q continuous
\item
  Theorem: The unique solution is y(x) = (1/μ(x))(y₀ + ∫ₐˣ μ(t) Q(t) dt)
  where μ(x) = e\^{}\{∫ₐˣ P(t) dt\} is the integrating factor
\item
  Proof idea: multiply the equation by μ(x). Since μ' = Pμ, the left
  side becomes (μy)' = μQ. Integrate both sides using FTC.
\item
  Existence: all integrals exist by FTC (P, Q continuous)
\item
  Uniqueness: the difference of two solutions satisfies the homogeneous
  equation with zero initial condition
\item
  Example: y' + 2y = eˣ with y(0) = 1. Integrating factor μ = e²ˣ. Then
  (e²ˣ y)' = e³ˣ. Integrate: e²ˣ y = e³ˣ/3 + C. From y(0) = 1: C = 2/3.
  Solution: y = (1/3)eˣ + (2/3)e⁻²ˣ.
\item
  ✎ Inline: solve y' − y/x = x² with y(1) = 0 for x \textgreater{} 0
\end{itemize}

\subsection{Separable Equations}\label{separable-equations}

\begin{itemize}
\tightlist
\item
  The equation: y' = f(x)g(y), with g(y₀) ≠ 0
\item
  If g(y) ≠ 0, formally separate: dy/g(y) = f(x) dx
\item
  Integrate: G(y) = F(x) + C where G' = 1/g, F' = f
\item
  Theorem: If f is continuous near x₀ and g is continuous and nonzero
  near y₀, then y' = f(x)g(y) with y(x₀) = y₀ has a solution (at least
  locally)
\item
  Proof:

  \begin{enumerate}
  \def\labelenumi{\arabic{enumi}.}
  \tightlist
  \item
    F(x) = ∫\_\{x₀\}ˣ f(t) dt exists by FTC
  \item
    G(y) = ∫\_\{y₀\}ʸ 1/g(s) ds exists by FTC (since g ≠ 0 means 1/g
    continuous)
  \item
    G'(y) = 1/g(y) ≠ 0, so G is strictly monotone
  \item
    Strictly monotone continuous functions have continuous inverses
    (§6.4)
  \item
    Solution: y = G⁻¹(F(x) + C) with C chosen so y(x₀) = y₀
  \end{enumerate}
\item
  Remark: we proved existence without computing G⁻¹ explicitly --- FTC
  and the inverse function theorem do the work
\item
  Example: y' = y/x with y(1) = 1 for x \textgreater{} 0. Separate: dy/y
  = dx/x. Integrate: log\textbar y\textbar{} = log\textbar x\textbar{} +
  C. Solution: y = x. Remark: log arises naturally in separation
  whenever 1/g involves 1/y.
\item
  Example: y' = xy with y(0) = 1. Separate: dy/y = x dx. Integrate:
  log\textbar y\textbar{} = x²/2 + C. Solution: y = e\^{}\{x²/2\}..
  Verify: (e\^{}\{x²/2\})' = x · e\^{}\{x²/2\} = xy. ✓
\item
  Example: y' = y² with y(0) = 1. Separate: dy/y² = dx. Integrate: −1/y
  = x + C. Solution: y = 1/(1−x). This blows up at x = 1 --- the
  solution exists only on (−∞, 1). Even when our existence theorem
  applies, the solution may not exist globally.
\item
  ✎ Inline: solve y' = y(1−y) with y(0) = 1/2 (the logistic equation)
\end{itemize}

\subsection{The Limits of Elementary
Methods}\label{the-limits-of-elementary-methods}

\begin{itemize}
\tightlist
\item
  Consider y' = x + y² --- we cannot separate variables, there is no
  integrating factor, no algebraic trick works
\item
  Does a solution even exist? Yes --- but proving it requires
  fundamentally new tools:

  \begin{itemize}
  \tightlist
  \item
    Picard iteration: construct approximating functions that converge to
    a solution
  \item
    Contraction mappings on function spaces (callback to §4.3, but now
    in infinite dimensions)
  \item
    Completeness of the function space
  \end{itemize}
\item
  This is the content of Chapter 12
\end{itemize}

\subsection{Guided Exercises}\label{guided-exercises-21}

None for this section.

\subsection{Exercises}\label{exercises-40}

\emph{Quadrature} - Verify: erf is odd, and erf(x) → 1 as x → ∞ - Show
Si(x) is bounded (hint: integrate by parts)

\emph{Linear equations} - Solve y' + 3y = e²ˣ with y(0) = 1 - Solve y' +
y = x with y(0) = 0

\emph{Separable equations} - Solve y' = y/x with y(1) = 2 - Solve y' =
x/y with y(0) = 1 - Show that solutions to y' = 1 + y² blow up in finite
time

\emph{Extensions} - Bernoulli equations: show y' + P(x)y = Q(x)yⁿ
reduces to linear via v = y\^{}\{1−n\} - Homogeneous equations: show y'
= f(y/x) reduces to separable via v = y/x - Second-order constant
coefficient: solve y'\,' + ay' + by = 0 via the characteristic equation
(use exp from §7.4)

\subsection{Dependencies}\label{dependencies-40}

\textbf{Requires}: §6.4 (continuous monotone inverses), §7.2 (MVT ---
uniqueness), §7.4 (exp), §9.1 (FTC I and II), §9.2 (substitution, IBP)

\textbf{Used in}: None (terminal section). The general ODE existence
theory appears in Chapter 12.

\section{§9.4 Trigonometry}\label{trigonometry-1}

\subsection{Narrative}\label{narrative-41}

\begin{itemize}
\tightlist
\item
  The trig functions join exp and log as functions constructed
  rigorously from calculus
\item
  arcsin is defined as an integral (like log in the guided exercise of
  §9.2), then sin and cos as its inverse and complement
\item
  The key insight: sin and cos satisfy y'\,' = −y, and the ODE
  uniqueness lemma (energy argument) proves \emph{everything} ---
  addition formulas, periodicity, special values
\item
  Taylor series follow immediately from the derivative cycle and the
  Lagrange remainder
\item
  The section connects π (known from geometry, §4.5 and §8.3) to the
  integral world: arcsin(1) = π/2 is a theorem, not a definition
\end{itemize}

\subsection{Content}\label{content-41}

\subsection{Defining the Trigonometric
Functions}\label{defining-the-trigonometric-functions}

\begin{itemize}
\tightlist
\item
  Recall: on the unit circle, arcsin(y) should give the angle (= arc
  length) whose sine is y
\item
  Definition: arcsin(y) = ∫₀ʸ 1/√(1−t²) dt for y ∈ (−1, 1)
\item
  This is an improper integral at y = ±1 (§8.1); it converges (compare
  1/√(1−t²) to 1/√(1−t) near t = 1)
\item
  arcsin'(y) = 1/√(1−y²) by FTC I
\item
  arcsin is strictly increasing (derivative positive), so it has an
  inverse
\item
  Theorem: arcsin(1) = π/2. Proof: ∫₀¹ 1/√(1−t²) dt is the arc length of
  the quarter circle from (1,0) to (0,1), which equals π/2 by §8.3.
\item
  So arcsin maps {[}−1, 1{]} → {[}−π/2, π/2{]}
\item
  Definition: sin: {[}−π/2, π/2{]} → {[}−1, 1{]} is the inverse of
  arcsin
\item
  Definition: cos(θ) = sin(π/2 − θ), or equivalently cos(θ) = √(1 −
  sin²θ) for θ ∈ {[}−π/2, π/2{]}
\item
  Immediate: sin²θ + cos²θ = 1
\item
  Derivatives:

  \begin{itemize}
  \tightlist
  \item
    sin'(θ) = 1/arcsin'(sin θ) = √(1 − sin²θ) = cos(θ)
  \item
    cos'(θ) = −sin(θ) (differentiate sin² + cos² = 1, or use cos =
    sin(π/2 − θ) and chain rule)
  \end{itemize}
\item
  ✎ Inline: verify cos'(θ) = −sin(θ) both ways
\end{itemize}

\subsection{The Differential Equation}\label{the-differential-equation}

\begin{itemize}
\tightlist
\item
  Theorem: sin and cos both satisfy y'\,' = −y. Proof: sin'\,' = (cos)'
  = −sin, cos'\,' = (−sin)' = −cos.
\item
  Lemma (Uniqueness): if y'\,' = −y on an interval with y(0) = y'(0) =
  0, then y ≡ 0.
\item
  Proof: define E(t) = y(t)² + y'(t)² (the ``energy''). Then E'(t) =
  2yy' + 2y'y'\,' = 2yy' + 2y'(−y) = 0. So E is constant. E(0) = 0, so
  E(t) = 0 for all t, hence y(t) = 0 for all t.
\item
  Corollary: solutions to y'\,' = −y are uniquely determined by y(0) and
  y'(0)
\item
  Theorem (Addition formulas):

  \begin{itemize}
  \tightlist
  \item
    sin(α + β) = sin α cos β + cos α sin β
  \item
    cos(α + β) = cos α cos β − sin α sin β
  \end{itemize}
\item
  Proof: fix α. Define f(t) = sin(α + t) and g(t) = sin(α)cos(t) +
  cos(α)sin(t). Both satisfy y'\,' = −y. Check: f(0) = sin(α) = g(0),
  f'(0) = cos(α) = g'(0). By uniqueness, f = g. Similarly for cosine.
\item
  Remark: the addition formulas are proved purely from the ODE, not from
  geometry
\item
  Special values:

  \begin{itemize}
  \tightlist
  \item
    sin(0) = 0, cos(0) = 1 (from arcsin(0) = 0 and cos = √(1−sin²))
  \item
    sin(π/2) = 1, cos(π/2) = 0 (from arcsin(1) = π/2)
  \item
    sin(π) = 0, cos(π) = −1 (addition formulas with α = β = π/2)
  \item
    sin(2π) = 0, cos(2π) = 1
  \item
    sin(π/4) = cos(π/4) = 1/√2
  \end{itemize}
\item
  Theorem: sin and cos are periodic with period 2π
\item
  Proof: define f(t) = sin(t + 2π). Then f'\,' = −f, f(0) = sin(2π) = 0,
  f'(0) = cos(2π) = 1. But sin also satisfies y'\,' = −y with y(0) = 0,
  y'(0) = 1. By uniqueness, f = sin. Similarly for cos.
\item
  Theorem: 2π is the smallest period (equivalently, π is the smallest
  positive zero of sin)
\item
  ✎ Inline: derive sin(π/6) = 1/2 from sin(π/2) = 1 and the
  addition/half-angle formulas
\end{itemize}

\subsection{Taylor Series}\label{taylor-series}

\begin{itemize}
\tightlist
\item
  sin and cos are infinitely differentiable: they satisfy y'\,' = −y, so
  all derivatives are ±sin or ±cos
\item
  The derivative cycle: sin → cos → −sin → −cos → sin gives the Taylor
  coefficients at 0
\item
  Theorem: sin(x) = x − x³/3! + x⁵/5! − ⋯ = Σ (−1)ⁿ x²ⁿ⁺¹/(2n+1)! for
  all x
\item
  Theorem: cos(x) = 1 − x²/2! + x⁴/4! − ⋯ = Σ (−1)ⁿ x²ⁿ/(2n)! for all x
\item
  Convergence: ratio test gives \textbar aₙ₊₁/aₙ\textbar{} =
  \textbar x\textbar²/((2n+2)(2n+3)) → 0, so both series converge for
  all x
\item
  Proof of equality: Lagrange remainder. All derivatives of sin and cos
  are bounded by 1, so \textbar Rₙ(x)\textbar{} ≤
  \textbar x\textbar ⁿ⁺¹/(n+1)! → 0 for any fixed x.
\item
  Remark: these are the series S(x) and C(x) from ch5, if introduced
  there. The functional equations (Cauchy product) now have a geometric
  interpretation as the addition formulas for sin and cos.
\item
  ✎ Inline: verify the Taylor series for cos by differentiating the
  Taylor series for sin term-by-term (Power Series II)
\end{itemize}

\subsection{Archimedes' Theorem}\label{archimedes-theorem}

\begin{itemize}
\tightlist
\item
  In §8.3, Archimedes' inequality (proved geometrically) showed the
  perimeter and area constants of the circle are the same π
\item
  Here we give a modern, purely computational proof using trig
  substitution
\item
  Theorem: ∫₀¹ √(1−x²) dx = π/4 (so the area of the unit disk is π)
\item
  Proof: IBP on ∫₀¹ √(1−x²) dx connects the area integral to ∫₀¹
  1/√(1−x²) dx = arcsin(1) = π/2
\item
  Remark: the substitution x = sin(θ) transforms the area integral
  directly into an arc length integral --- calculus links the two
  constants without any geometric argument
\item
  Note: this subsection could alternatively be a guided exercise
\end{itemize}

\subsection{Guided Exercises}\label{guided-exercises-22}

\subsection{The Weierstrass Product for
Sine}\label{the-weierstrass-product-for-sine}

\emph{Goal:} Prove that sin(πx)/(πx) = ∏ₙ₌₁\^{}∞ (1 − x²/n²)

\emph{Note:} Elementary proof TBD. The standard proof uses complex
analysis (Hadamard factorization), but we seek an approach using only
real analysis tools from this chapter. The identity is used in §9.5 to
derive the Wallis product for π.

\emph{Consequences:} (a) Setting x = 1/2: the Wallis product π/2 = ∏
(2n)²/((2n−1)(2n+1)) (b) Expanding the product and matching x²
coefficients: Σ 1/n² = π²/6 (Basel problem --- proved independently via
Fourier in Chapter \_\_)

\subsection{Exercises}\label{exercises-41}

\emph{Definitions and derivatives} - Prove that arcsin is an odd
function: arcsin(−y) = −arcsin(y) - Define arccos(y) = ∫ᵧ¹ 1/√(1−t²) dt.
Show arccos(y) = π/2 − arcsin(y) and that cos is the inverse of arccos
on {[}0, π{]}. - Show that sec'(θ) = sec(θ) tan(θ) where sec = 1/cos

\emph{The ODE and its consequences} - Prove: if y'\,' = −y and y(0) =
y'(0) = 0, then y ≡ 0 using the Wronskian: if y₁, y₂ are solutions, show
W = y₁y₂' − y₁'y₂ is constant - Derive cos(π/3) = 1/2 and sin(π/3) =
√3/2 from the addition formulas - Prove the double-angle formulas:
sin(2θ) = 2 sin θ cos θ, cos(2θ) = cos²θ − sin²θ - Prove:
\textbar sin(θ)\textbar{} ≤ \textbar θ\textbar{} for all θ (hint:
energy, or integrate \textbar cos\textbar{} ≤ 1)

\emph{Taylor series} - Compute sin(1) to 5 decimal places. How many
terms are needed? - Prove sin(x)/x → 1 as x → 0 using the Taylor series
- Prove: \textbar sin(x) − x\textbar{} ≤ \textbar x\textbar³/6 and
\textbar cos(x) − 1\textbar{} ≤ x²/2 for all x

\emph{Integration with trig} (uses §9.2) - Compute ∫₀\^{}\{π/2\} sin²(x)
dx and ∫₀\^{}\{π/2\} cos²(x) dx using the identity sin²x = (1 − cos
2x)/2 - Compute ∫₀\^{}\{π/4\} tan(x) dx by writing tan = sin/cos and
substituting u = cos(x) - Compute ∫₀\^{}\{π/2\} sin(x) cos(x) dx two
ways: substitution and double-angle formula - Derive the reduction
formula: ∫₀\^{}\{π/2\} sinⁿ(x) dx = ((n−1)/n) ∫₀\^{}\{π/2\} sinⁿ⁻²(x) dx
(integration by parts) - Use the reduction formula to derive the Wallis
integrals: ∫₀\^{}\{π/2\} sin²ⁿ(x) dx = π/2 · (2n)!/(2²ⁿ(n!)²)

\subsection{Dependencies}\label{dependencies-41}

\textbf{Requires}: §4.5 (π from polygon limits), §6.4 (continuous
monotone inverses), §7.1 (inverse function derivative, chain rule), §7.3
(Power Series II --- term-by-term differentiation), §8.1 (improper
integrals), §8.3 (arc length --- for arcsin(1) = π/2), §9.1 (FTC I and
II), §9.2 (Power Series III, substitution, IBP)

\textbf{Used in}: §9.5 (arctan and π computations)

\section{\texorpdfstring{§9.5 Calculating π
\emph{(Terminal)}}{§9.5 Calculating π (Terminal)}}\label{calculating-ux3c0-terminal}

\subsection{Narrative}\label{narrative-42}

\begin{itemize}
\tightlist
\item
  Each subsection presents a different strategy for computing π,
  progressing in sophistication
\item
  Newton's method: fast convergence, but computationally circular (needs
  trig values to find π)
\item
  Inverse trig functions: π as an output rather than a zero --- arcsin
  has irrational summands, arctan has rational ones
\item
  Gregory-Leibniz: power series for arctan, beautiful but slow
\item
  Machin-type formulas: arctan identities that accelerate convergence
  --- the dominant method for 250 years
\item
  Wallis product: π from an infinite product, completely different from
  series
\item
  Historical thread woven throughout: Archimedes' polygons (§4.5/§8.3),
  Madhava/Gregory/Leibniz, Wallis, Machin, Euler, and remarks on what
  lies beyond (Ramanujan, Chudnovsky, AGM)
\end{itemize}

\subsection{Content}\label{content-42}

\subsection{Newton's Method}\label{newtons-method-1}

\begin{itemize}
\tightlist
\item
  We know cos(π/2) = 0, so π/2 is a root of cosine
\item
  Newton's method (§7.5): xₙ₊₁ = xₙ − cos(xₙ)/sin(xₙ), starting near π/2
\item
  Converges quadratically --- digits double each iteration
\item
  Work through a few iterations explicitly (using Taylor polynomial
  approximations for sin and cos)
\item
  Problem: each step requires evaluating cos and sin at xₙ, which means
  computing Taylor series to high precision
\item
  Computationally expensive per step, and somewhat circular --- we need
  trig values to find π
\item
  This motivates a different strategy: find a function that
  \emph{outputs} π as a value, not one that has π as a zero
\item
  Historical remark: Newton himself computed π to 15 digits (1665),
  though using a different method (integrating a circular arc, not his
  own Newton's method)
\end{itemize}

\subsection{Inverse Trigonometric
Functions}\label{inverse-trigonometric-functions}

\begin{itemize}
\tightlist
\item
  New strategy: we want f(something simple) = π
\item
  We already have: arcsin(1) = π/2 and arccos(0) = π/2
\item
  As integrals: π/2 = ∫₀¹ 1/√(1−t²) dt
\item
  Approximate via Riemann sums: Sₙ = (1/n) Σₖ₌₀ⁿ⁻¹ 1/√(1−(k/n)²)
\item
  Problem: each summand involves a square root --- irrational, needs its
  own approximation
\item
  Worse: the integrand blows up at t = 1 (improper integral), so
  convergence of Riemann sums is slow
\item
  Can we do better? Arctangent: arctan(1) = π/4, and 1/(1+t²) is a
  \emph{rational} function
\item
  Definition: arctan(x) = ∫₀ˣ 1/(1+t²) dt
\item
  Properties: arctan'(x) = 1/(1+x²) by FTC I, arctan(0) = 0, arctan is
  odd, arctan(x) → π/2 as x → ∞
\item
  Theorem: arctan(1) = π/4 (proof via substitution t = tan(u), or as
  inverse of tan on (−π/2, π/2))
\item
  So π/4 = ∫₀¹ 1/(1+t²) dt --- a proper integral of a rational function
\item
  Riemann sums: Sₙ = (1/n) Σₖ₌₀ⁿ⁻¹ 1/(1+(k/n)²) --- purely rational!
\item
  Compute explicit approximations: write out S₄, S₁₀, compare 4Sₙ to π
\item
  Convergence is O(1/n) --- still slow, but each step is now elementary
  arithmetic
\item
  Historical remark: the connection between arctan and π was known to
  James Gregory (1670s) and independently to Madhava of Sangamagrama
  (\textasciitilde1400)
\end{itemize}

\subsection{The Gregory-Leibniz
Series}\label{the-gregory-leibniz-series}

\begin{itemize}
\tightlist
\item
  The integrand 1/(1+t²) has a power series: 1/(1+t²) = 1 − t² + t⁴ − ⋯
  for \textbar t\textbar{} \textless{} 1 (geometric series with r = −t²)
\item
  Integrate term-by-term (Power Series III): arctan(x) = x − x³/3 + x⁵/5
  − ⋯ for \textbar x\textbar{} \textless{} 1
\item
  The series also has radius of convergence 1 (ratio test)
\item
  Abel's theorem (§6.3) extends to the boundary: the series converges at
  x = 1 (alternating series with terms → 0)
\item
  Theorem (Gregory-Leibniz): π/4 = 1 − 1/3 + 1/5 − 1/7 + ⋯
\item
  One of the most beautiful formulas in mathematics: π emerges from
  simple arithmetic of odd numbers
\item
  But terribly slow: alternating series error bound gives \textbar π/4 −
  Sₙ\textbar{} \textless{} 1/(2n+1), need \textasciitilde500 terms for 2
  decimal places
\item
  Compute partial sums explicitly and see the slow convergence
\item
  Historical remark: Madhava (\textasciitilde1400) discovered both the
  series and correction terms to accelerate it. Leibniz rediscovered the
  series in 1674.
\end{itemize}

\subsection{Machin-type Formulas}\label{machin-type-formulas}

\begin{itemize}
\tightlist
\item
  Gregory-Leibniz is slow because x = 1 is the boundary of the radius of
  convergence
\item
  Key idea: express π/4 using arctan of \emph{small} arguments, where
  the series converges fast
\item
  Proposition: π/4 = arctan(1/2) + arctan(1/3)
\item
  Proof: tangent addition formula
\item
  The series for arctan(1/2) has terms decaying like (1/2)²ᵏ⁺¹,
  arctan(1/3) like (1/3)²ᵏ⁺¹ --- much faster than 1/(2k+1)
\item
  Machin's formula (1706): π/4 = 4 arctan(1/5) − arctan(1/239)
\item
  Terms decay like (1/5)²ᵏ⁺¹ and (1/239)²ᵏ⁺¹ --- extremely fast
\item
  Example: 9 terms give 15 digits
\item
  Machin computed 100 digits by hand
\item
  Historical remark: Machin-type formulas dominated π computation for
  250 years. William Shanks used a variant to compute 707 digits by hand
  (1873), though the last \textasciitilde180 were wrong (error
  discovered in 1944 by Ferguson using a desk calculator).
\item
  ✎ Inline: verify Machin's formula using tangent addition
\end{itemize}

\subsection{The Wallis Product}\label{the-wallis-product-1}

\begin{itemize}
\tightlist
\item
  A completely different approach: π from an infinite product, not a
  series
\item
  Start from the reduction formula (§9.4 exercises): ∫₀\^{}\{π/2\}
  sinⁿ(x) dx = ((n−1)/n) ∫₀\^{}\{π/2\} sinⁿ⁻²(x) dx
\item
  Compute Iₙ = ∫₀\^{}\{π/2\} sinⁿ(x) dx for even and odd n separately
\item
  I₂ₙ = (π/2) · (2n)!/(2²ⁿ(n!)²), I₂ₙ₊₁ = 2²ⁿ(n!)²/(2n+1)!
\item
  Since sinⁿ⁺¹ ≤ sinⁿ on {[}0, π/2{]}: I₂ₙ₊₁ ≤ I₂ₙ ≤ I₂ₙ₋₁
\item
  The ratio I₂ₙ₋₁/I₂ₙ₊₁ → 1, so I₂ₙ/I₂ₙ₊₁ → 1
\item
  Take the ratio and rearrange:
\item
  Theorem (Wallis, 1655): π/2 = ∏ₙ₌₁\^{}∞ (2n)²/((2n−1)(2n+1)) =
  (2·2)/(1·3) · (4·4)/(3·5) · (6·6)/(5·7) · ⋯
\item
  Alternative derivation: if the Weierstrass product for sine is
  available (§9.4 guided exercise), set x = 1/2 in sin(πx)/(πx) = ∏(1 −
  x²/n²) and rearrange
\item
  Compute partial products for n = 1, 2, 3, 4, \ldots{} and observe
  convergence to π/2
\item
  Remark: converges slowly (like Gregory-Leibniz), but remarkable that π
  emerges from integer arithmetic alone
\item
  Historical remark: Wallis derived this in 1655 by a brilliant
  interpolation argument, before calculus was formalized. Our proof via
  integration is anachronistic but rigorous.
\item
  ✎ Inline: compute partial products for n = 1, 2, 3, 4 and compare to
  π/2
\end{itemize}

\textbf{Euler's Formula} \emph{(Placeholder)} - Theorem: Σ 1/n² = π²/6
(the Basel problem, solved by Euler in 1735) - Note: proof TBD. Can be
derived from the Weierstrass product (§9.4 guided exercise) by expanding
and matching x² coefficients, or proved independently via Fourier
analysis (Chapter \_\_). If an elementary proof is found that fits our
technology, it goes here. - Historical remark: the problem was posed by
Pietro Mengoli in 1650 and made famous by the Bernoulli family. Euler's
solution, at age 28, made him famous across Europe.

\subsection{Guided Exercises}\label{guided-exercises-23}

\subsection{Stirling's Approximation}\label{stirlings-approximation}

\emph{Goal:} Prove n! \textasciitilde{} √(2πn)(n/e)ⁿ

\subsection{Part I: The unknown
constant}\label{part-i-the-unknown-constant}

\begin{enumerate}
\def\labelenumi{(\alph{enumi})}
\item
  Compute ∫₁ⁿ log(t) dt = n log n − n + 1 using IBP (§9.2).
\item
  Since log(n!) = Σₖ₌₁ⁿ log(k), compare the sum to the integral. Show
  Σₖ₌₁ⁿ log(k) ≥ ∫₁ⁿ log(t) dt (since log is increasing, log(k) ≥ ∫ₖ₋₁ᵏ
  log(t) dt).
\item
  Get an upper bound using concavity of log: log(k) ≤
  ∫ₖ₋₁/₂\^{}\{k+1/2\} log(t) dt (midpoint rule overshoots for concave
  functions).
\item
  Define dₙ = log(n!) − (n + 1/2)log(n) + n.~Show (dₙ) is monotone and
  bounded.
\item
  Conclude: dₙ → some constant γ, so n! \textasciitilde{} eᵞ · √n ·
  (n/e)ⁿ. Call C = eᵞ --- we know C \textgreater{} 0 but not its value.
\end{enumerate}

\subsection{Part II: Wallis identifies
C}\label{part-ii-wallis-identifies-c}

\begin{enumerate}
\def\labelenumi{(\alph{enumi})}
\setcounter{enumi}{5}
\item
  Rewrite the Wallis product in factorial form: π/2 = lim (2ⁿn!)⁴ /
  ((2n)!² · (2n+1)).
\item
  Substitute n! \textasciitilde{} C√n(n/e)ⁿ and (2n)! \textasciitilde{}
  C√(2n)(2n/e)²ⁿ into the expression from (f).
\item
  Show everything cancels except C: the limit equals π/2 only if C² =
  2π, so C = √(2π).
\item
  Conclude: n! \textasciitilde{} √(2πn)(n/e)ⁿ.
\end{enumerate}

\subsection{Exercises}\label{exercises-42}

\emph{Newton's method} - Carry out 4 iterations of Newton's method on
cos(x) = 0, starting from x₀ = 1. Use Taylor polynomials of degree 10
for sin and cos to evaluate each step. How many digits of π/2 do you
get? - ★ Implement the Newton's method computation (with Taylor
polynomial evaluation) and compute π to 15 digits. How many iterations
are needed? How does this compare to the number of Machin series terms
needed for the same accuracy?

\emph{Riemann sums} - Compute 4·S₁₀ and 4·S₁₀₀ for the Riemann sum Sₙ =
(1/n) Σₖ₌₀ⁿ⁻¹ 1/(1+(k/n)²). How many digits of π are correct? - The
trapezoidal rule gives better approximations than left Riemann sums.
Compute the trapezoidal approximation with n = 10 for ∫₀¹ 1/(1+t²) dt
and compare.

\emph{Arctangent series} - Verify: arctan(1/2) + arctan(1/3) = π/4 using
the tangent addition formula - Compute π using arctan(1/2) + arctan(1/3)
with 5 terms of each series. How many digits are correct? - Verify
Machin's formula: 4 arctan(1/5) − arctan(1/239) = π/4 - ★ Compute π to
20 digits using Machin's formula. How many terms of each series are
needed?

\emph{Wallis product} - Prove the Wallis product directly from the
reduction formula without using the Weierstrass product - Compute the
partial products P₁, P₂, \ldots, P₁₀ and compare to π/2. What is the
error at P₁₀?

\emph{Beyond our tools} - Ramanujan's series (1914): 1/π = (√8/9801)
Σₙ₌₀\^{}∞ (4n)!(1103+26390n)/((n!)⁴ · 396⁴ⁿ). Each term gives
\textasciitilde8 new digits. Compute the first 2 terms and verify you
get \textasciitilde16 digits of 1/π. (No proof --- the proof requires
modular forms.)

\subsection{Dependencies}\label{dependencies-42}

\textbf{Requires}: §7.5 (Newton's method), §8.3 (arc length, π from
geometry), §9.1 (FTC), §9.2 (Power Series III, substitution, IBP), §9.4
(sin, cos, tan, arctan, reduction formula, Weierstrass product if
available)

\textbf{Used in}: None (terminal section)

\chapter{Extending the Integral}\label{extending-the-integral}

\section{Integrability}\label{integrability}

\textbf{Beyond Continuous Functions}

\begin{itemize}
\tightlist
\item
  Our axiomatic integral was built for continuous functions
\item
  But the construction works more broadly
\item
  Upper and lower sums make sense for any bounded function
\item
  \textbf{Definition}: \(f\) is \emph{Darboux integrable} if upper
  integral equals lower integral
\item
  What functions does this capture?
\end{itemize}

\textbf{The Class of Integrable Functions}

\textbf{Many functions are integrable}: - Continuous functions (by
construction) - Piecewise continuous functions (split at
discontinuities) - Monotone functions (upper/lower sums telescope) - The
ruler function (discontinuous on dense set, still integrable) - General
principle: discontinuities must be ``small''

\textbf{It's a vector space}: - If \(f, g\) integrable and
\(\alpha, \beta \in \mathbb{R}\), then \(\alpha f + \beta g\) integrable
- \(\int(\alpha f + \beta g) = \alpha \int f + \beta \int g\) -
Linearity inherited from the construction - We can do algebra with
integrable functions

\textbf{The Boundary}

\begin{itemize}
\tightlist
\item
  \(\chi_{\mathbb{Q}}\): discontinuous everywhere, not integrable
\item
  Upper integral \(= 1\), lower integral \(= 0\)
\item
  Not just pathological---it's the boundary of our class
\end{itemize}

\textbf{Power Series III Was Special}

\begin{itemize}
\tightlist
\item
  Callback to Chapter 9: term-by-term integration worked
\item
  \(\int \left( \sum a_n x^n \right) = \sum \frac{a_n x^{n+1}}{n+1}\)
\item
  Used for \(\ln(1+x)\), \(\arctan\), computing \(\pi\)
\item
  You might expect limits generally behave this well
\end{itemize}

\textbf{The Convergence Problem}

\begin{itemize}
\tightlist
\item
  \textbf{Failure Mode 1}: \(f_n = n \cdot \chi_{[0,1/n]} \to 0\) but
  \(\int f_n = 1\)

  \begin{itemize}
  \tightlist
  \item
    Limit integrable, integrals don't converge correctly
  \end{itemize}
\item
  \textbf{Failure Mode 2}:
  \(f_n = \chi_{\{q_1,\ldots,q_n\}} \nearrow \chi_{\mathbb{Q}}\)

  \begin{itemize}
  \tightlist
  \item
    Each \(f_n\) integrable with \(\int f_n = 0\)
  \item
    Limit not integrable
  \end{itemize}
\end{itemize}

\textbf{The class of Darboux integrable functions is not closed under
limits.}

\textbf{Why This Matters}

\begin{itemize}
\tightlist
\item
  Fourier series: need \(\int\) and \(\sum\) to interchange
\item
  Parameter integrals: need \(\int\) and \(\frac{d}{dt}\) to interchange
\item
  Power series were the exception; general limits fail
\end{itemize}

\textbf{The Key Observation}

\begin{itemize}
\tightlist
\item
  In Failure Mode 2: \(\int f_n = 0 \to 0\), obvious candidate is \(0\)
\item
  Problem isn't wrong answer---it's no answer when there's an obvious
  candidate
\item
  Could we extend by \emph{defining} \(\int \chi_{\mathbb{Q}} = 0\)?
\end{itemize}

\textbf{What We Want}

\begin{itemize}
\tightlist
\item
  \textbf{Preserve linearity}: we have a vector space, want to keep it
\item
  \textbf{Add MCT}: increasing limits with bounded integrals should work
\item
  The remarkable fact: demanding both forces a unique extension
\end{itemize}

\section{\texorpdfstring{\(\bigstar\) The Daniell
Extension}{\textbackslash bigstar The Daniell Extension}}\label{bigstar-the-daniell-extension}

\textbf{The Plan}

\begin{itemize}
\tightlist
\item
  Two stages: first close under monotone limits, then close under
  subtraction
\item
  Must verify the extension still satisfies integral properties
\item
  Must verify MCT survives the second extension
\end{itemize}

\textbf{Dini's Theorem}

\begin{itemize}
\tightlist
\item
  Lemma: If \(f_n \in C[a,b]\) and \(f_n \searrow 0\) pointwise, then
  \(I(f_n) \to 0\)
\item
  Key technical tool bridging pointwise convergence to integral
  convergence
\item
  Proof uses compactness of \([a,b]\)
\end{itemize}

\textbf{Upper Functions}

\begin{itemize}
\tightlist
\item
  Definition: \(f \in L^\uparrow\) if \(f_n \in C[a,b]\) with
  \(f_n \nearrow f\) pointwise
\item
  Examples: continuous functions, \(\chi_{(a,b]}\)
\item
  Define \(I^+(f) = \lim I(f_n)\)
\end{itemize}

\textbf{Well-Definedness of \(I^+\)}

\begin{itemize}
\tightlist
\item
  Must show: different approximating sequences give same answer
\item
  Key Lemma: If \(f_n \nearrow f\) and \(g \in C[a,b]\) with
  \(g \leq f\), then \(I(g) \leq \lim I(f_n)\)
\item
  Proof uses Dini's theorem on the gap
  \(h_n = g - (f_n \wedge g) \searrow 0\)
\item
  Corollary: \(I^+\) well-defined
\end{itemize}

\textbf{Properties of \(I^+\)}

\begin{itemize}
\tightlist
\item
  Extension: agrees with \(I\) on \(C[a,b]\)
\item
  Monotonicity: \(f \leq g \Rightarrow I^+(f) \leq I^+(g)\)
\item
  Additive: \(I^+(f + g) = I^+(f) + I^+(g)\)
\item
  Positive homogeneity: \(c \geq 0 \Rightarrow I^+(cf) = c \, I^+(f)\)
\item
  MCT for \(L^\uparrow\): \(f_n \in L^\uparrow\), \(f_n \nearrow f\),
  \(\sup I^+(f_n) < \infty \Rightarrow f \in L^\uparrow\) and
  \(I^+(f_n) \to I^+(f)\)
\end{itemize}

\textbf{The Linearity Problem}

\begin{itemize}
\tightlist
\item
  \(L^\uparrow\) is a cone, not a vector space
\item
  Closed under addition and positive scalars
\item
  Not closed under subtraction: if \(f_n \nearrow f\), then
  \(-f_n \searrow -f\) (wrong direction)
\item
  We need subtraction for basic analysis
\end{itemize}

\textbf{Integrable Functions}

\begin{itemize}
\tightlist
\item
  Definition: \(f \in L^1\) if \(f = g - h\) for some
  \(g, h \in L^\uparrow\) with \(I^+(g), I^+(h) < \infty\)
\item
  Define \(I(f) = I^+(g) - I^+(h)\)
\item
  Must verify well-defined: if \(g_1 - h_1 = g_2 - h_2\), then
  \(I^+(g_1) - I^+(h_1) = I^+(g_2) - I^+(h_2)\)
\item
  Proof: \(g_1 + h_2 = g_2 + h_1\), both in \(L^\uparrow\), apply
  additivity
\end{itemize}

\textbf{Properties of \(I\) on \(L^1\)}

\begin{itemize}
\tightlist
\item
  \(L^1\) is a vector space
\item
  \(I\) is linear
\item
  Extension, monotonicity, interval additivity inherited
\end{itemize}

\textbf{MCT Survives}

\begin{itemize}
\tightlist
\item
  Key Lemma: If \(f \in L^1\) and \(f \geq 0\), then
  \(f \in L^\uparrow\)
\item
  Theorem (MCT for \(L^1\)): \(f_n \in L^1\), \(f_n \nearrow f\),
  \(\sup I(f_n) < \infty \Rightarrow f \in L^1\) and \(I(f_n) \to I(f)\)
\item
  Proof: WLOG \(f_n \geq 0\) (subtract \(f_1\)), apply key lemma, use
  MCT for \(L^\uparrow\)
\end{itemize}

\textbf{What We've Built}

\begin{itemize}
\tightlist
\item
  An integral on \(L^1 \supset C[a,b]\)
\item
  Satisfies original properties plus linearity plus MCT
\item
  Next: harvest the consequences
\end{itemize}

\section{Convergence Theorems}\label{convergence-theorems}

\textbf{Monotone Convergence Theorem}

\begin{itemize}
\tightlist
\item
  Theorem (MCT): \(f_n \in L^1\), \(f_n \nearrow f\),
  \(\sup I(f_n) < \infty \Rightarrow f \in L^1\) and \(I(f_n) \to I(f)\)
\item
  Already proved in 10.2; state cleanly here as the first payoff
\end{itemize}

\textbf{Dominated Convergence Theorem}

\begin{itemize}
\tightlist
\item
  Theorem (DCT): \(f_n \in L^1\), \(f_n \to f\) pointwise,
  \(|f_n| \leq g\) for some \(g \in L^1 \Rightarrow f \in L^1\) and
  \(I(f_n) \to I(f)\)
\item
  The workhorse theorem---this is what you actually use
\item
  Proof via MCT applied to \(\inf\)/\(\sup\) sequences
\end{itemize}

\textbf{Revisiting the Failures}

\begin{itemize}
\tightlist
\item
  Spike functions \(f_n = n \cdot \chi_{[0,1/n]}\): no dominating
  function (would need \(g(0) = \infty\))
\item
  DCT diagnoses exactly why interchange fails
\item
  Not a pathology---a precise answer
\end{itemize}

\textbf{Application: Differentiating Under the Integral}

\begin{itemize}
\tightlist
\item
  Setup: \(F(t) = \int_a^b f(x,t) \, dx\)
\item
  Theorem: If \(\frac{\partial f}{\partial t}\) exists and
  \(\left| \frac{\partial f}{\partial t} \right| \leq g \in L^1\), then
  \(F'(t) = \int_a^b \frac{\partial f}{\partial t} \, dx\)
\item
  Proof: Apply DCT to difference quotients
\end{itemize}

\textbf{Application: Integrating Series Term-by-Term}

\begin{itemize}
\tightlist
\item
  Setup: \(\sum f_n\) where each \(f_n \in L^1\)
\item
  Theorem: If \(\sum \int |f_n| < \infty\), then
  \(\int \sum f_n = \sum \int f_n\)
\item
  Proof: Partial sums dominated by \(\sum |f_n|\), apply DCT
\end{itemize}

\textbf{The Fundamental Theorem of Calculus}

\textbf{Part 1 extends nicely}: - Theorem: If \(f \in L^1\), then
\(F(x) = \int_a^x f\) is differentiable a.e. with \(F'(x) = f(x)\) a.e.
- We can integrate any \(L^1\) function and recover it by
differentiation

\textbf{Part 2 is subtle}: - Just having \(F'\) exist and be integrable
isn't enough - Counterexample: The Cantor function (devil's staircase) -
Continuous, increasing, \(F(0) = 0\), \(F(1) = 1\) - \(F'(x) = 0\) a.e.
(flat on complement of Cantor set) - But
\(\int_0^1 F' = 0 \neq 1 = F(1) - F(0)\) - The function ``gains height''
on the Cantor set where it has no derivative

\textbf{Absolute continuity}: - Definition: \(F\) is absolutely
continuous if for every \(\varepsilon > 0\) there exists \(\delta > 0\)
such that for any finite collection of disjoint intervals \((a_i, b_i)\)
with \(\sum(b_i - a_i) < \delta\), we have
\(\sum |F(b_i) - F(a_i)| < \varepsilon\) - Stronger than uniform
continuity (handles multiple intervals simultaneously) - The Cantor
function is uniformly continuous but not absolutely continuous

\textbf{Part 2 for Lebesgue}: - Theorem: If \(F\) absolutely continuous,
then \(F'\) exists a.e., \(F' \in L^1\), and
\(\int_a^b F' = F(b) - F(a)\)

\textbf{The beautiful characterization}: -
\(L^1 = \{F' : F \text{ absolutely continuous}\}\) - Lebesgue integrable
functions are exactly the a.e. derivatives of absolutely continuous
functions - Parallels: Darboux integrable \(\leftrightarrow\)
derivatives of Lipschitz functions

\textbf{Completeness}

\begin{itemize}
\tightlist
\item
  Theorem: If \(f_n \in L^1\) and
  \(\sum \int |f_{n+1} - f_n| < \infty\), then \(f_n\) converges to some
  \(f \in L^1\)
\item
  Equivalently: Cauchy sequences in the \(\int |\cdot|\) sense converge
\item
  False for Darboux integrable functions
\item
  ``Completing the Integral''---we've filled the holes
\end{itemize}

\textbf{Series Closure}

\begin{itemize}
\tightlist
\item
  Theorem: If \(f_n \in L^1\) and \(\sum \int |f_n| < \infty\), then
  \(\sum f_n \in L^1\)
\item
  False for Darboux: \(f_n = \chi_{\{q_n\}}\), each integrable, sum is
  \(\chi_{\mathbb{Q}}\)
\item
  \(L^1\) closed under absolutely convergent series
\end{itemize}

\section{Measure Theory}\label{measure-theory}

\textbf{Null Sets}

\begin{itemize}
\tightlist
\item
  Definition: \(E\) is null if for every \(\varepsilon > 0\), \(E\)
  covered by countably many intervals of total length \(< \varepsilon\)
\item
  Equivalently: \(\chi_E \in L^1\) and \(I(\chi_E) = 0\)
\item
  Examples: finite sets, countable sets, Cantor set
\end{itemize}

\textbf{Properties of Null Sets}

\begin{itemize}
\tightlist
\item
  Subset of null is null
\item
  Countable union of null sets is null
\item
  Proof: cover \(n\)th set with intervals of total length
  \(\varepsilon / 2^n\)
\end{itemize}

\textbf{Almost Everywhere}

\begin{itemize}
\tightlist
\item
  Definition: Property holds a.e. if it fails only on a null set
\item
  ``\(f = g\) a.e.'' means \(\{x : f(x) \neq g(x)\}\) is null
\item
  ``\(f_n \to f\) a.e.'' means \(\{x : f_n(x) \not\to f(x)\}\) is null
\end{itemize}

\textbf{The A.E. Modification Theorem}

\begin{itemize}
\tightlist
\item
  Theorem: If \(f = g\) a.e. and \(f \in L^1\), then \(g \in L^1\) and
  \(I(f) = I(g)\)
\item
  Striking: modify on infinitely many points, integral unchanged
\item
  False for Darboux: \(f = 0\), \(g = \chi_{\mathbb{Q}}\), differ on
  null set, but \(\chi_{\mathbb{Q}}\) not Darboux integrable
\end{itemize}

\textbf{Measurable Sets}

\begin{itemize}
\tightlist
\item
  Definition: \(E\) is measurable if \(\chi_E \in L^1\)
\item
  Examples: intervals, open sets, closed sets, countable
  unions/intersections
\item
  Non-measurable sets exist (requires axiom of choice)
\end{itemize}

\textbf{Lebesgue Measure}

\begin{itemize}
\tightlist
\item
  Definition: \(\mu(E) = I(\chi_E)\) for measurable \(E\)
\item
  Properties: \(\mu(E) \geq 0\), \(\mu(\varnothing) = 0\),
  \(\mu([a,b]) = b - a\)
\item
  Countable additivity:
  \(\mu\left( \bigsqcup E_n \right) = \sum \mu(E_n)\)
\item
  Proof via DCT
\end{itemize}

\textbf{The Integral-First Philosophy}

\begin{itemize}
\tightlist
\item
  Traditional: measure → measurable functions → integral
\item
  Our approach: integral → measure as byproduct
\item
  For analysis, the integral is primary; measure serves it
\end{itemize}

\textbf{EXERCISES} EXTENDING THE AXIOMS TO MEET OTHER MEASURES:
STIELTJES INTEGRATION ETC

\section{Why Lebesgue?}\label{why-lebesgue}

\textbf{Uniqueness}

\begin{itemize}
\tightlist
\item
  Theorem: Any extension of \(I\) from \(C[a,b]\) satisfying linearity
  and MCT must agree with ours
\item
  Proof sketch: MCT forces values on \(L^\uparrow\), linearity forces
  values on \(L^1\)
\item
  There was no choice---MCT determined everything
\end{itemize}

\textbf{Maximality}

\begin{itemize}
\tightlist
\item
  Theorem: \(L^1\) is the largest class where \(f\) integrable
  \(\Rightarrow |f|\) integrable
\item
  ``Absolute integrability'' property
\item
  Extensions beyond (like HK) lose this
\end{itemize}

\textbf{Connection to Derivatives}

\begin{itemize}
\tightlist
\item
  Theorem: \(f \in L^1\) iff \(f = F'\) a.e. for some absolutely
  continuous \(F\)
\item
  FTC for Lebesgue: cleaner than Darboux version
\item
  Preview of absolute continuity
\end{itemize}

\textbf{Completeness Revisited}

\begin{itemize}
\tightlist
\item
  \(L^1\) with \(d(f,g) = \int |f - g|\) is complete
\item
  In fact: \(L^1\) is the completion of \(C[a,b]\) under this metric
\item
  Every element of \(L^1\) is a limit of continuous functions
\item
  We added exactly the missing limits---nothing more
\end{itemize}

\textbf{A General Machine}

\begin{itemize}
\tightlist
\item
  The construction used the length axiom \(\mu([c,d]) = d - c\) only as
  input
\item
  Positivity, linearity, Daniell's condition did the work
\item
  Replace with \(\mu((c,d]) = g(d) - g(c)\) for increasing
  right-continuous \(g\)
\item
  Same machine produces a new integral
\item
  In Chapter N: this captures all finite Borel measures on
  \([a,b]\)---the Riesz representation theorem
\end{itemize}

\textbf{Beyond Lebesgue: The Henstock-Kurzweil Integral}

\begin{itemize}
\tightlist
\item
  Lebesgue isn't the only extension
\item
  HK: every derivative is integrable, full FTC
\item
  Example: \(F(x) = x^2 \sin(1/x^2)\), differentiable everywhere, \(F'\)
  not Lebesgue integrable
\item
  Tradeoff: HK has better FTC, but no \(L^p\) theory
\item
  Different integrals answer different questions
\end{itemize}

\textbf{Closing Perspective}

\begin{itemize}
\tightlist
\item
  Darboux: simplest integral for area under curves
\item
  Lebesgue: largest integral with good limit behavior
\item
  HK: largest integral with full FTC
\item
  For functional analysis: Lebesgue is essential
\end{itemize}

\part{Function Space}

\chapter{Metric Spaces}\label{metric-spaces}

We abstract the key ideas from Parts I and II. What made analysis on
\(\mathbb{R}\) work wasn't the specific numbers --- it was having a
notion of distance where Cauchy sequences converge. The same theorems
hold in any complete metric space. This opens the door to function
spaces (Ch 12), Fourier analysis (Ch 13), and differential equations (Ch
14).

\section{Abstracting Distance}\label{abstracting-distance}

\textbf{Definition of metric space}

\begin{itemize}
\tightlist
\item
  A set \(X\) with distance function \(d: X \times X \to [0, \infty)\)
\item
  Axioms: \(d(x,y) = 0 \iff x = y\); \(d(x,y) = d(y,x)\);
  \(d(x,z) \leq d(x,y) + d(y,z)\)
\item
  We've been doing this on \(\mathbb{R}\) all along
\end{itemize}

\textbf{Examples --- Euclidean spaces}

\begin{itemize}
\tightlist
\item
  \(\mathbb{R}\) with \(d(x,y) = |x - y|\)
\item
  \(\mathbb{R}^n\) with Euclidean distance
  \(d(x,y) = \sqrt{\sum (x_i - y_i)^2}\)
\item
  \(\mathbb{C}\) with \(d(z,w) = |z - w|\)
\end{itemize}

\textbf{Examples --- Sequence spaces}

\begin{itemize}
\tightlist
\item
  \(\ell^1\): sequences with \(\sum |a_n| < \infty\), distance
  \(d(a,b) = \sum |a_n - b_n|\)
\item
  \(\ell^2\): sequences with \(\sum |a_n|^2 < \infty\), distance
  \(d(a,b) = \sqrt{\sum |a_n - b_n|^2}\)
\item
  \(\ell^\infty\): bounded sequences, distance
  \(d(a,b) = \sup_n |a_n - b_n|\)
\item
  These are infinite-dimensional analogues of \(\mathbb{R}^n\)
\end{itemize}

\textbf{Examples --- Function spaces}

\begin{itemize}
\tightlist
\item
  \(C[a,b]\) with sup norm:
  \(d(f,g) = \sup_{x \in [a,b]} |f(x) - g(x)|\)
\item
  \(C[a,b]\) with \(L^1\) norm: \(d(f,g) = \int_a^b |f(x) - g(x)|\, dx\)
\item
  Same set, different metrics! Different notions of ``close''
\end{itemize}

\textbf{Examples --- Other}

\begin{itemize}
\tightlist
\item
  Discrete metric: \(d(x,y) = 1\) if \(x \neq y\), else \(0\). Every
  subset is open!
\item
  Hausdorff metric on closed bounded subsets:
  \(d_H(A,B) = \max\{\sup_{a \in A} d(a,B), \sup_{b \in B} d(b,A)\}\)
\end{itemize}

\textbf{Balls and topology}

\begin{itemize}
\tightlist
\item
  Open ball: \(B_r(x) = \{y : d(x,y) < r\}\)
\item
  Closed ball: \(\bar{B}_r(x) = \{y : d(x,y) \leq r\}\)
\item
  Open set: every point has a ball around it contained in the set
\item
  Closed set: complement is open (equivalently: contains all its limit
  points)
\end{itemize}

\textbf{Convergence}

\begin{itemize}
\tightlist
\item
  \(x_n \to x\) means \(d(x_n, x) \to 0\)
\item
  Same definition as \(\mathbb{R}\), works in any metric space
\end{itemize}

\textbf{Continuity}

\begin{itemize}
\tightlist
\item
  \(f: X \to Y\) continuous at \(x\) if \(x_n \to x\) implies
  \(f(x_n) \to f(x)\)
\item
  Equivalently: preimage of open is open
\item
  Equivalently: \(\varepsilon\)-\(\delta\) definition
\end{itemize}

\textbf{Cauchy sequences and completeness}

\begin{itemize}
\tightlist
\item
  Cauchy: for all \(\varepsilon > 0\), exists \(N\) such that
  \(m, n > N\) implies \(d(x_m, x_n) < \varepsilon\)
\item
  Complete: every Cauchy sequence converges
\item
  \(\mathbb{R}\), \(\mathbb{R}^n\), \(\mathbb{C}\): complete
\item
  \(\ell^1\), \(\ell^2\), \(\ell^\infty\): complete
\item
  \(C[a,b]\) with sup norm: complete
\item
  \(C[a,b]\) with \(L^1\) norm: NOT complete! (Cauchy sequences can
  converge to discontinuous functions)
\item
  \(\mathbb{Q}\): NOT complete
\end{itemize}

\textbf{Sequential compactness}

\begin{itemize}
\tightlist
\item
  Definition: every sequence has a convergent subsequence
\item
  In \(\mathbb{R}^n\): compact iff closed and bounded (Heine-Borel)
\item
  Foreshadow: this will NOT hold in infinite dimensions
\end{itemize}

\textbf{Connectedness}

\begin{itemize}
\tightlist
\item
  Definition: cannot be written as union of two disjoint nonempty open
  sets
\item
  Path-connected implies connected
\item
  \(\mathbb{R}\), \(\mathbb{R}^n\), \(\mathbb{C}\): connected
\item
  \(\mathbb{Q}\): totally disconnected
\end{itemize}

\textbf{Inner products and metrics (brief paragraph)}

\begin{itemize}
\tightlist
\item
  An inner product \(\langle \cdot, \cdot \rangle\) induces a norm
  \(\|x\| = \sqrt{\langle x, x \rangle}\)
\item
  A norm induces a metric \(d(x,y) = \|x - y\|\)
\item
  Examples: Euclidean inner product on \(\mathbb{R}^n\), \(L^2\) inner
  product on functions
\item
  Not every metric comes from a norm (discrete metric doesn't)
\end{itemize}

\emph{{[}Note: End of section practice --- verifying examples are
metrics, checking completeness of specific spaces, showing specific sets
are open/closed{]}}

\section{Reproving Theorems}\label{reproving-theorems}

\textbf{Theme:} Familiar theorems from Part I were really about metric
space structure, not specifically about \(\mathbb{R}\).

\textbf{Compactness: continuous image}

\begin{itemize}
\tightlist
\item
  Theorem: If \(K\) is compact and \(f: K \to Y\) is continuous, then
  \(f(K)\) is compact
\item
  Proof: Same as Part I --- extract subsequence, use continuity
\end{itemize}

\textbf{Compactness: maximum and minimum}

\begin{itemize}
\tightlist
\item
  Theorem: If \(K\) is compact and \(f: K \to \mathbb{R}\) is
  continuous, then \(f\) achieves its max and min
\item
  Proof: \(f(K)\) is compact subset of \(\mathbb{R}\), hence closed and
  bounded, hence contains its sup
\end{itemize}

\textbf{Compactness: uniform continuity}

\begin{itemize}
\tightlist
\item
  Theorem: If \(K\) is compact and \(f: K \to Y\) is continuous, then
  \(f\) is uniformly continuous
\item
  Proof: Same as Part I
\end{itemize}

\textbf{Heine-Borel in \(\mathbb{R}^n\)}

\begin{itemize}
\tightlist
\item
  Theorem: A subset of \(\mathbb{R}^n\) is compact iff it is closed and
  bounded
\item
  Proof: Bounded means contained in a cube; cube is compact
  (Bolzano-Weierstrass iterated); closed subset of compact is compact
\end{itemize}

\textbf{Connectedness: continuous image}

\begin{itemize}
\tightlist
\item
  Theorem: If \(X\) is connected and \(f: X \to Y\) is continuous, then
  \(f(X)\) is connected
\item
  Proof: If \(f(X) = U \cup V\) separated, then
  \(X = f^{-1}(U) \cup f^{-1}(V)\) separated
\end{itemize}

\textbf{Connectedness implies IVT}

\begin{itemize}
\tightlist
\item
  Corollary: If \(f: [a,b] \to \mathbb{R}\) is continuous and
  \(f(a) < c < f(b)\), then \(f(x) = c\) for some \(x\)
\item
  Proof: \(f([a,b])\) is connected subset of \(\mathbb{R}\), hence an
  interval
\item
  The intermediate value theorem IS ``continuous image of connected is
  connected''
\end{itemize}

\textbf{Infinite-dimensional contrast: closed ball is NOT compact}

\begin{itemize}
\tightlist
\item
  Theorem: The closed unit ball in \(C[a,b]\) (sup norm) is not compact
\item
  Proof: The sequence \(f_n(x) = x^n\) on \([0,1]\) has no convergent
  subsequence (converges pointwise to discontinuous function, but
  convergence in sup norm implies uniform convergence implies continuous
  limit)
\item
  Or: construct sequence with \(\|f_n - f_m\| = 1\) for all \(n \neq m\)
\item
  This is why Arzelà-Ascoli (Ch 12) needs extra conditions
\end{itemize}

\textbf{Contraction mapping theorem}

\begin{itemize}
\tightlist
\item
  Definition: \(T: X \to X\) is a contraction if
  \(d(Tx, Ty) \leq c \cdot d(x,y)\) for some \(c < 1\)
\item
  Theorem: If \(X\) is complete and \(T\) is a contraction, then \(T\)
  has a unique fixed point
\item
  Proof: (Brief --- same as Part I, cite earlier) Start anywhere,
  iterate, Cauchy sequence, completeness gives limit, limit is fixed
  point, uniqueness by contraction property
\end{itemize}

\textbf{Application: Cantor set as fixed point}

\begin{itemize}
\tightlist
\item
  Consider \(T(A) = \frac{1}{3}A \cup (\frac{2}{3} + \frac{1}{3}A)\) on
  closed subsets of \([0,1]\) with Hausdorff metric
\item
  Claim: \(T\) is a contraction (with constant \(1/3\))
\item
  Hausdorff metric on closed subsets is complete (state; details to
  guided exercises)
\item
  By contraction mapping: unique fixed point \(C\) with \(T(C) = C\)
\item
  This fixed point IS the Cantor set!
\item
  A fractal defined as a fixed point
\end{itemize}

\emph{{[}Note: Further exploration at chapter end --- more fractals
(Sierpinski triangle, Koch curve) as fixed points of iterated function
systems{]}}

\emph{{[}Note: Guided exercises --- prove Hausdorff metric is complete;
verify contraction constant for other IFS examples{]}}

\section{Extending Series}\label{extending-series}

\textbf{Theme:} Our theory of series works wherever we have completeness
and compatible multiplication.

\textbf{Absolute convergence in complete metric spaces}

\begin{itemize}
\tightlist
\item
  Recall: In \(\mathbb{R}\), absolute convergence implies convergence
\item
  The proof used: completeness (Cauchy criterion) and triangle
  inequality
\item
  Same proof works in any complete normed space!
\item
  If \(\sum \|a_n\| < \infty\), then \(\sum a_n\) converges
\end{itemize}

\textbf{Power series in complete normed rings}

\begin{itemize}
\tightlist
\item
  Setup: \((R, \|\cdot\|)\) is complete, has multiplication, and
  \(\|xy\| \leq \|x\|\|y\|\) (submultiplicative)
\item
  Examples: \(\mathbb{R}\), \(\mathbb{C}\), \(M_n(\mathbb{R})\),
  \(M_n(\mathbb{C})\)
\item
  Claim: If \(\sum |a_n| \|x\|^n < \infty\), then \(\sum a_n x^n\)
  converges
\item
  In particular: \(e^x = \sum \frac{x^n}{n!}\) converges for all \(x\)
  (since \(\sum \frac{\|x\|^n}{n!} = e^{\|x\|} < \infty\))
\item
  Same for \(\sin x\), \(\cos x\), etc.
\end{itemize}

\textbf{The complex numbers}

\begin{itemize}
\tightlist
\item
  \(\mathbb{C} = \mathbb{R}^2\) with multiplication
  \((a,b)(c,d) = (ac - bd, ad + bc)\)
\item
  Write \((0,1) = i\); then \(i^2 = -1\)
\item
  \(|z| = \sqrt{a^2 + b^2}\) is a norm with \(|zw| = |z||w|\)
  (multiplicative, even better than submultiplicative!)
\item
  \(\mathbb{C}\) is complete (because \(\mathbb{R}^2\) is)
\end{itemize}

\textbf{Euler's formula}

\begin{itemize}
\tightlist
\item
  \(e^{i\theta} = \sum \frac{(i\theta)^n}{n!}\) converges for all
  \(\theta \in \mathbb{R}\)
\item
  Separate real and imaginary parts:

  \begin{itemize}
  \tightlist
  \item
    Real:
    \(1 - \frac{\theta^2}{2!} + \frac{\theta^4}{4!} - \cdots = \cos\theta\)
  \item
    Imaginary:
    \(\theta - \frac{\theta^3}{3!} + \frac{\theta^5}{5!} - \cdots = \sin\theta\)
  \end{itemize}
\item
  Therefore: \(e^{i\theta} = \cos\theta + i\sin\theta\)
\item
  This is WHERE complex exponentials come from
\item
  Corollary: \(|e^{i\theta}| = 1\) (lives on unit circle)
\item
  Corollary: \(e^{i\pi} + 1 = 0\) (Euler's identity)
\end{itemize}

\textbf{Matrix exponential}

\begin{itemize}
\tightlist
\item
  \(M_n(\mathbb{R})\): \(n \times n\) real matrices with operator norm
  \(\|A\| = \sup_{\|x\|=1} \|Ax\|\)
\item
  Key property: \(\|AB\| \leq \|A\| \|B\|\) (submultiplicative)
\item
  \(M_n(\mathbb{R})\) is complete (it's \(\mathbb{R}^{n^2}\)
  topologically)
\item
  Definition: \(e^A = \sum_{n=0}^\infty \frac{A^n}{n!}\)
\item
  Convergence: \(\|A^n/n!\| \leq \|A\|^n/n!\), so
  \(\sum \|A^n/n!\| \leq e^{\|A\|} < \infty\)
\end{itemize}

\textbf{Properties of matrix exponential}

\begin{itemize}
\tightlist
\item
  \(e^0 = I\)
\item
  \(\frac{d}{dt} e^{tA} = A e^{tA}\)
\item
  If \(AB = BA\), then \(e^{A+B} = e^A e^B\)
\item
  Warning: \(e^{A+B} \neq e^A e^B\) in general when \(AB \neq BA\)!
\end{itemize}

\textbf{Computing \(e^A\) via eigenvalues}

\begin{itemize}
\tightlist
\item
  If \(A = PDP^{-1}\) with \(D\) diagonal, then \(e^A = P e^D P^{-1}\)
\item
  \(e^D\) is diagonal with \((e^D)_{ii} = e^{D_{ii}}\)
\item
  Real eigenvalues: exponential growth/decay
\item
  Complex eigenvalues: oscillation!
\end{itemize}

\textbf{Complex eigenvalues and rotation}

\begin{itemize}
\tightlist
\item
  Let \(A = \begin{pmatrix} 0 & -\omega \\ \omega & 0 \end{pmatrix}\)
\item
  Eigenvalues: \(\pm i\omega\)
\item
  Compute:
  \(e^{tA} = \begin{pmatrix} \cos(\omega t) & -\sin(\omega t) \\ \sin(\omega t) & \cos(\omega t) \end{pmatrix}\)
\item
  This IS rotation by angle \(\omega t\)!
\item
  Complex eigenvalues explain oscillation in differential equations (Ch
  14)
\end{itemize}

\emph{{[}Note: Practice --- compute \(e^A\) for specific matrices;
verify properties{]}}

\emph{{[}Note: Guided exercises --- prove \(e^{A+B} = e^A e^B\) when
\(AB = BA\); explore what goes wrong when \(AB \neq BA\){]}}

\section{\texorpdfstring{\(\bigstar\) Baire
Category}{\textbackslash bigstar Baire Category}}\label{bigstar-baire-category}

\textbf{Theme:} A genuinely new theorem that we couldn't even state
before. Completeness implies ``largeness.''

\textbf{Nowhere dense sets}

\begin{itemize}
\tightlist
\item
  Definition: \(A\) is nowhere dense if its closure has empty interior
\item
  Equivalently: \(A\) contains no ball
\item
  Examples: \(\mathbb{Z}\) in \(\mathbb{R}\); any finite set; the Cantor
  set
\end{itemize}

\textbf{Meager and comeager sets}

\begin{itemize}
\tightlist
\item
  Meager (first category): countable union of nowhere dense sets
\item
  Comeager (residual): complement of a meager set
\item
  Comeager = countable intersection of open dense sets
\end{itemize}

\textbf{Baire Category Theorem}

\begin{itemize}
\tightlist
\item
  Theorem: A complete metric space is not meager
\item
  Equivalently: The intersection of countably many open dense sets is
  dense
\item
  Proof: Given open dense sets \(U_1, U_2, \ldots\) and a ball \(B_0\),
  construct nested balls \(B_1 \supset B_2 \supset \cdots\) with
  \(B_n \subset U_n\), radii \(\to 0\). Completeness gives point in all
  \(U_n\).
\end{itemize}

\textbf{The notion of ``generic''}

\begin{itemize}
\tightlist
\item
  A property is generic if it holds on a comeager set
\item
  In complete spaces: generic properties are ``most'' points
\item
  Not the same as ``full measure'' --- different notion of large
\end{itemize}

\textbf{Application: \(\mathbb{R}\) is uncountable (new proof!)}

\begin{itemize}
\tightlist
\item
  If \(\mathbb{R} = \{x_1, x_2, x_3, \ldots\}\), then
  \(\mathbb{R} = \bigcup_n \{x_n\}\)
\item
  Each singleton is closed with empty interior (nowhere dense)
\item
  So \(\mathbb{R}\) would be meager
\item
  But \(\mathbb{R}\) is complete, so Baire says it's not meager
\item
  Contradiction!
\item
  Compare to Ch 1: different proof, same conclusion
\end{itemize}

\textbf{Application: No function is continuous exactly at
\(\mathbb{Q}\)}

\begin{itemize}
\tightlist
\item
  For any \(f: \mathbb{R} \to \mathbb{R}\), let \(C_f\) = set of
  continuity points
\item
  Fact: \(C_f\) is always a \(G_\delta\) set (countable intersection of
  open sets)
\item
  Claim: \(\mathbb{Q}\) is NOT \(G_\delta\)
\item
  Proof: If \(\mathbb{Q} = \bigcap_n U_n\) with \(U_n\) open, each
  \(U_n\) contains \(\mathbb{Q}\) hence is dense. By Baire,
  \(\bigcap U_n\) is dense. But \(\mathbb{Q}\) is countable, hence not
  dense in any interval\ldots{} wait, \(\mathbb{Q}\) IS dense. Try
  again.
\item
  Better proof: \(\mathbb{Q}\) is meager (countable union of
  singletons). If \(\mathbb{Q}\) were \(G_\delta\), then
  \(\mathbb{R} \setminus \mathbb{Q}\) would be \(F_\sigma\), i.e.,
  countable union of closed sets. The irrationals would be meager. Then
  \(\mathbb{R} = \mathbb{Q} \cup (\mathbb{R} \setminus \mathbb{Q})\)
  would be meager. Contradiction!
\item
  Therefore no function has \(\mathbb{Q}\) as its continuity set
\end{itemize}

\textbf{But: there IS a function continuous exactly at the irrationals}

\begin{itemize}
\tightlist
\item
  Thomae's ruler function: \(f(p/q) = 1/q\) (in lowest terms),
  \(f(\text{irrational}) = 0\)
\item
  Continuous at every irrational, discontinuous at every rational
\item
  The asymmetry is explained by Baire!
\end{itemize}

\textbf{Foreshadow: Nowhere differentiable is generic}

\begin{itemize}
\tightlist
\item
  In Ch 12 we'll prove: the set of continuous functions differentiable
  at even one point is meager in \(C[a,b]\)
\item
  ``Most'' continuous functions are nowhere differentiable
\item
  Weierstrass's monster isn't pathological --- it's typical!
\end{itemize}

\emph{{[}Note: Guided exercises --- prove \(C_f\) is \(G_\delta\); show
other sets are/aren't \(G_\delta\); explore Thomae function
properties{]}}

\section{\texorpdfstring{\(\bigstar\) Completing
Spaces}{\textbackslash bigstar Completing Spaces}}\label{bigstar-completing-spaces}

\textbf{Theme:} We can BUILD complete spaces. This finally delivers on
the promise from Chapter 1.

\textbf{The construction}

\begin{itemize}
\tightlist
\item
  Start with \(\mathbb{Q}\), a metric space that isn't complete
\item
  Consider all Cauchy sequences in \(\mathbb{Q}\)
\item
  Define equivalence: \((a_n) \sim (b_n)\) if \(|a_n - b_n| \to 0\)
\item
  Let \(\tilde{\mathbb{Q}}\) = equivalence classes of Cauchy sequences
\end{itemize}

\textbf{Making it a metric space}

\begin{itemize}
\tightlist
\item
  For classes \([a_n]\), \([b_n]\): define
  \(d([a_n], [b_n]) = \lim |a_n - b_n|\)
\item
  Check: limit exists (Cauchy sequences of rationals have rational
  limits\ldots{} no wait, that's the problem)
\item
  Fix: \(|a_n - b_n|\) is Cauchy in \(\mathbb{R}\)\ldots{} but we're
  constructing \(\mathbb{R}\)!
\item
  Actual fix: \(|a_n - b_n|\) is Cauchy in \(\mathbb{Q}\) when both are
  Cauchy, and Cauchy sequences of rationals that would converge to the
  same real have \(|a_n - b_n| \to 0\). The limit exists as an
  equivalence class.
\item
  (This is subtle --- details in guided exercises)
\end{itemize}

\textbf{Making it a field}

\begin{itemize}
\tightlist
\item
  Addition: \([a_n] + [b_n] = [a_n + b_n]\)
\item
  Multiplication: \([a_n] \cdot [b_n] = [a_n \cdot b_n]\)
\item
  Check: well-defined, satisfies field axioms
\end{itemize}

\textbf{Embedding \(\mathbb{Q}\)}

\begin{itemize}
\tightlist
\item
  Map \(q \mapsto [(q, q, q, \ldots)]\) (constant sequence)
\item
  This embeds \(\mathbb{Q}\) in \(\tilde{\mathbb{Q}}\)
\item
  The image is dense
\end{itemize}

\textbf{Completeness}

\begin{itemize}
\tightlist
\item
  Theorem: \(\tilde{\mathbb{Q}}\) is complete
\item
  Proof: A Cauchy sequence of equivalence classes can be represented by
  a ``diagonal'' Cauchy sequence in \(\mathbb{Q}\)
\item
  (Details technical but not deep)
\end{itemize}

\textbf{This IS \(\mathbb{R}\)}

\begin{itemize}
\tightlist
\item
  \(\tilde{\mathbb{Q}}\) is a complete ordered field containing
  \(\mathbb{Q}\) densely
\item
  By uniqueness of complete ordered field:
  \(\tilde{\mathbb{Q}} \cong \mathbb{R}\)
\item
  We have BUILT the real numbers from the rationals!
\end{itemize}

\textbf{Other completions exist!}

\begin{itemize}
\tightlist
\item
  The completion depends on the metric
\item
  The \(p\)-adic metric on \(\mathbb{Q}\): \(|x|_p = p^{-n}\) where
  \(n\) is the largest power of \(p\) dividing \(x\)
\item
  ``Close'' means difference is divisible by high power of \(p\)
\item
  Completion: \(p\)-adic numbers \(\mathbb{Q}_p\)
\item
  Completely different from \(\mathbb{R}\)! (Totally disconnected, every
  triangle is isoceles, etc.)
\item
  The metric determines the world
\end{itemize}

\textbf{Same construction elsewhere}

\begin{itemize}
\tightlist
\item
  \(C[a,b]\) with \(L^1\) norm is not complete
\item
  Its completion is \(L^1[a,b]\) (Lebesgue integrable functions)
\item
  Connects to Chapter 10
\end{itemize}

\emph{{[}Note: This entire section could be guided exercises for
advanced students{]}}

\chapter{Function Spaces}\label{function-spaces}

\begin{itemize}
\tightlist
\item
  Chapter 1: We asked ``what is \(\mathbb{R}\)?'' Answer: complete
  ordered field
\item
  Now: What are spaces of functions? What structure do they have?
\item
  Same questions: completeness, convergence, optimization, density
\item
  Multiple families: \(C^k\) (differentiability), \(L^p\)
  (integrability)
\item
  Different norms measure different things
\end{itemize}

\section{\texorpdfstring{The Spaces \(C^k\), \(L^1\),
\(L^2\)}{The Spaces C\^{}k, L\^{}1, L\^{}2}}\label{the-spaces-ck-l1-l2}

\textbf{\(C[a,b]\): Continuous Functions}

\begin{itemize}
\tightlist
\item
  Definition: continuous functions on \([a,b]\)
\item
  Sup norm: \(\|f\|_\infty = \sup_{x \in [a,b]} |f(x)|\)
\item
  Measures: worst-case error
\end{itemize}

\textbf{\(C^k[a,b]\): Differentiable Functions}

\begin{itemize}
\tightlist
\item
  Definition: \(k\) times continuously differentiable
\item
  Norm: \(\|f\|_{C^k} = \sum_{j=0}^k \|f^{(j)}\|_\infty\)
\item
  Measures: control of function and its derivatives
\item
  \(C^\infty[a,b]\): smooth functions (infinitely differentiable)
\item
  Inclusions:
  \(C^\infty \subset \cdots \subset C^2 \subset C^1 \subset C^0 = C\)
\end{itemize}

\textbf{\(L^1[a,b]\): Integrable Functions}

\begin{itemize}
\tightlist
\item
  Definition: functions with \(\int_a^b |f| < \infty\)
\item
  Norm: \(\|f\|_1 = \int_a^b |f|\)
\item
  Measures: total size
\item
  Technical point: identify functions equal almost everywhere
\end{itemize}

\textbf{\(L^2[a,b]\): Square-Integrable Functions}

\begin{itemize}
\tightlist
\item
  Definition: functions with \(\int_a^b |f|^2 < \infty\)
\item
  Norm: \(\|f\|_2 = \left(\int_a^b |f|^2\right)^{1/2}\)
\item
  \textbf{Inner product}: \(\langle f, g \rangle = \int_a^b f \bar{g}\)
\item
  Cauchy-Schwarz: \(|\langle f, g \rangle| \leq \|f\|_2 \|g\|_2\)
\item
  Measures: energy / mean-square
\end{itemize}

\textbf{Inclusions on Bounded Intervals}

\begin{itemize}
\tightlist
\item
  \(C[a,b] \subset L^2[a,b] \subset L^1[a,b]\)
\item
  \(\|f\|_1 \leq \sqrt{b-a} \|f\|_2 \leq (b-a) \|f\|_\infty\)
\end{itemize}

\textbf{Guided Exercises}

\begin{itemize}
\tightlist
\item
  General \(L^p\): \(\|f\|_p = \left(\int |f|^p\right)^{1/p}\)
\item
  Hölder's inequality: \(\|fg\|_1 \leq \|f\|_p \|g\|_q\) where
  \(\frac{1}{p} + \frac{1}{q} = 1\)
\end{itemize}

\section{Completeness \& Convergence}\label{completeness-convergence}

\textbf{Completeness of \(C[a,b]\)}

\begin{itemize}
\tightlist
\item
  Theorem: Cauchy sequence in sup norm \(\Rightarrow\) converges in sup
  norm
\item
  Proof: pointwise limit exists; uniform convergence \(\Rightarrow\)
  limit is continuous
\item
  \textbf{Message}: Uniform limit of continuous functions is continuous
\end{itemize}

\textbf{Completeness of \(C^k[a,b]\)}

\begin{itemize}
\tightlist
\item
  Theorem: \(C^k\) is complete in \(C^k\) norm
\item
  Need uniform convergence of all derivatives up to order \(k\)
\end{itemize}

\textbf{Completeness of \(L^1\) and \(L^2\)}

\begin{itemize}
\tightlist
\item
  Theorem (Riesz-Fischer): \(L^1\) and \(L^2\) are complete
\item
  Proof uses DCT (Ch 10)
\item
  Key step: Cauchy sequence has subsequence converging a.e.; DCT gives
  \(L^p\) convergence
\end{itemize}

\textbf{Comparing Notions of Convergence}

\begin{itemize}
\tightlist
\item
  Uniform convergence: \(\|f_n - f\|_\infty \to 0\)
\item
  \(L^p\) convergence: \(\|f_n - f\|_p \to 0\)
\item
  Pointwise convergence: \(f_n(x) \to f(x)\) for each \(x\)
\item
  \textbf{Examples where they differ}:

  \begin{itemize}
  \tightlist
  \item
    Pointwise but not uniform: \(f_n(x) = x^n\) on \([0,1]\)
  \item
    \(L^2\) but not pointwise: sliding bumps
  \item
    Pointwise but not \(L^1\): tall thin spikes
  \end{itemize}
\end{itemize}

\textbf{Interchanging Limits}

\begin{itemize}
\tightlist
\item
  \textbf{With integration}: Uniform convergence \(\Rightarrow\)
  \(\int f_n \to \int f\)
\item
  \textbf{With integration}: \(L^1\) convergence \(\Rightarrow\)
  \(\int f_n \to \int f\)
\item
  \textbf{With differentiation}: Need uniform convergence of derivatives
\item
  \textbf{Callback to Ch 10}: DCT as the tool for \(L^p\)
\end{itemize}

\section{\texorpdfstring{\(\bigstar\)
Optimization}{\textbackslash bigstar Optimization}}\label{bigstar-optimization}

\textbf{The Problem}

\begin{itemize}
\tightlist
\item
  In calculus: find \(x\) minimizing \(f(x)\)
\item
  Here: find function \(y\) minimizing a functional \(F[y]\)
\end{itemize}

\textbf{Functionals}

\begin{itemize}
\tightlist
\item
  Definition: \(F: C^1[a,b] \to \mathbb{R}\)
\item
  Example: \(F[y] = \int_a^b L(x, y(x), y'(x))\, dx\)
\item
  This is ``analysis level 2'' optimization
\end{itemize}

\textbf{The First Variation}

\begin{itemize}
\tightlist
\item
  Analogue of derivative: perturb \(y \to y + \epsilon \eta\)
\item
  First variation:
  \(\delta F[y; \eta] = \frac{d}{d\epsilon}\Big|_{\epsilon=0} F[y + \epsilon\eta]\)
\item
  At an extremum: \(\delta F[y; \eta] = 0\) for all admissible \(\eta\)
\end{itemize}

\textbf{The Euler-Lagrange Equation}

\begin{itemize}
\tightlist
\item
  For \(F[y] = \int_a^b L(x, y, y')\, dx\) with fixed endpoints
  \(y(a) = A\), \(y(b) = B\):
\item
  \textbf{Theorem}: If \(y\) is an extremum, then
  \[\frac{\partial L}{\partial y} - \frac{d}{dx}\frac{\partial L}{\partial y'} = 0\]
\item
  Derivation via integration by parts
\item
  E-L is a second-order ODE; solutions form 2-parameter family
\item
  Boundary conditions select from this family
\end{itemize}

\textbf{Applications}

\begin{itemize}
\tightlist
\item
  \textbf{Shortest path}: \(L = \sqrt{1 + (y')^2}\) gives \(y'' = 0\)
  --- straight lines
\item
  \textbf{Brachistochrone}: fastest descent gives cycloid
\item
  \textbf{Catenary}: hanging chain minimizes potential energy
\item
  \textbf{Geodesics}: shortest paths on surfaces
\end{itemize}

\textbf{Constrained Optimization}

\begin{itemize}
\tightlist
\item
  Analogue of Lagrange multipliers
\item
  Problem: optimize \(F[y]\) subject to \(G[y] = c\)
\item
  Method: optimize \(F[y] - \lambda G[y]\)
\item
  \textbf{Isoperimetric problem}: maximize area subject to fixed
  perimeter
\item
  Setup here; proof via Fourier in Ch 13
\end{itemize}

\textbf{Existence and Uniqueness}

\begin{itemize}
\tightlist
\item
  Unlike calculus: existence of minimizers is subtle
\item
  \(F\) may be bounded below but not achieve minimum
\item
  Convexity of \(L\) in \(y'\) helps guarantee existence
\item
  Beyond our scope, but important to note
\end{itemize}

\textbf{Guided Exercises}

\begin{itemize}
\tightlist
\item
  Second variation and sufficient conditions
\item
  Multiple dependent variables
\item
  Higher derivatives in \(L\)
\end{itemize}

\section{Regularization}\label{regularization}

\textbf{The Technique}

\begin{itemize}
\tightlist
\item
  To analyze rough functions: approximate by smooth, prove for smooth,
  take limits
\item
  \textbf{Tool}: Convolution with a mollifier
\end{itemize}

\textbf{Convolution}

\begin{itemize}
\tightlist
\item
  Definition: \((f * g)(x) = \int f(x-t)g(t)\, dt\)
\item
  Properties: commutative, associative, linear in each argument
\item
  \textbf{Young's inequality}: \(\|f * g\|_p \leq \|f\|_1 \|g\|_p\)
\end{itemize}

\textbf{Smoothness Inheritance}

\begin{itemize}
\tightlist
\item
  \textbf{Theorem}: If \(\phi \in C^\infty\) and \(f \in L^1\), then
  \(f * \phi \in C^\infty\)
\item
  Proof: Differentiate under the integral; derivatives fall on \(\phi\)
\item
  \textbf{Message}: Convolution with smooth function produces smooth
  function
\end{itemize}

\textbf{Approximate Identities}

\begin{itemize}
\tightlist
\item
  Definition: Family \(\phi_\epsilon\) with

  \begin{enumerate}
  \def\labelenumi{\arabic{enumi}.}
  \tightlist
  \item
    \(\int \phi_\epsilon = 1\)
  \item
    \(\phi_\epsilon \geq 0\)
  \item
    \(\phi_\epsilon\) concentrates at 0 as \(\epsilon \to 0\)
  \end{enumerate}
\item
  \textbf{Examples}: Rescaled box functions, Gaussians
  \(\frac{1}{\sqrt{2\pi\epsilon}}e^{-x^2/2\epsilon}\)
\end{itemize}

\textbf{The Approximation Theorem}

\begin{itemize}
\tightlist
\item
  \textbf{Theorem}: \(f * \phi_\epsilon \to f\) as \(\epsilon \to 0\)

  \begin{itemize}
  \tightlist
  \item
    Uniformly for \(f \in C[a,b]\)
  \item
    In \(L^p\) norm for \(f \in L^p\)
  \end{itemize}
\item
  Proof:
  \(f * \phi_\epsilon - f = \int [f(x-t) - f(x)]\phi_\epsilon(t)\, dt\);
  use concentration
\end{itemize}

\textbf{The Density Payoffs}

\begin{itemize}
\tightlist
\item
  \textbf{Theorem}: \(C^\infty\) is dense in \(C[a,b]\) (sup norm)
\item
  \textbf{Theorem}: \(C_c^\infty\) is dense in \(L^p\)
\item
  Proof: Convolve with smooth mollifier; apply approximation theorem
\end{itemize}

\textbf{Looking Ahead}

\begin{itemize}
\tightlist
\item
  Fejér kernel is an approximate identity (Ch 13)
\item
  Heat kernel \(K_t\) is an approximate identity as \(t \to 0\) (Ch 15)
\end{itemize}

\section{Density Theorems}\label{density-theorems}

\textbf{The Question}

\begin{itemize}
\tightlist
\item
  \(C^\infty\) is dense (12.4). Can even less suffice?
\item
  Are polynomials dense? Trig polynomials?
\end{itemize}

\textbf{Weierstrass Approximation Theorem}

\begin{itemize}
\tightlist
\item
  \textbf{Theorem}: Polynomials are dense in \(C[a,b]\)
\item
  Every continuous function is a uniform limit of polynomials
\item
  Remarkable: polynomials are so rigid (analytic!), yet they approximate
  everything
\end{itemize}

\textbf{Proof via Bernstein Polynomials}

\begin{itemize}
\tightlist
\item
  \(B_n(f)(x) = \sum_{k=0}^n f(k/n) \binom{n}{k} x^k (1-x)^{n-k}\)
\item
  Probabilistic interpretation: expected value of \(f\) at binomial
  random variable
\item
  Theorem: \(B_n(f) \to f\) uniformly for continuous \(f\)
\item
  Proof uses law of large numbers flavor
\end{itemize}

\textbf{Stone-Weierstrass: The General Principle}

\begin{itemize}
\tightlist
\item
  \textbf{Theorem}: Let \(X\) be compact, \(A \subset C(X)\) a
  subalgebra. If:

  \begin{enumerate}
  \def\labelenumi{\arabic{enumi}.}
  \tightlist
  \item
    \(A\) separates points
  \item
    \(A\) contains constants
  \end{enumerate}

  Then \(A\) is dense in \(C(X)\).
\end{itemize}

\textbf{Corollaries}

\begin{itemize}
\tightlist
\item
  \textbf{Polynomials dense in \(C[a,b]\)}: \(f(x) = x\) separates
  points ✓
\item
  \textbf{Trig polynomials dense in \(C(\mathbb{T})\)}: \(e^{ix}\)
  separates points ✓
\end{itemize}

\textbf{The Narrative}

\begin{itemize}
\tightlist
\item
  12.4 showed: smooth functions approximate continuous
\item
  12.5 shows: even polynomials suffice
\item
  \textbf{Looking ahead}: Fourier gives EXPLICIT approximations (Ch 13)
\end{itemize}

\chapter{Fourier Series}\label{fourier-series}

Different problems call for different representations. We develop
orthogonal expansions in \(L^2\), starting with concrete polynomial
approximation, then Fourier series as the main example, and finally
explaining why Fourier's complex exponentials are ubiquitous: they
diagonalize the derivative.

\begin{itemize}
\tightlist
\item
  Fourier (1807): claims any function can be expanded in sines
\item
  Lagrange, Laplace skeptical: step functions as sums of smooth waves?
\item
  Fourier does it anyway, solves heat equation
\item
  Takes a century to make rigorous
\item
  We now have the tools (Ch 12): density, completeness, \(L^2\)
\end{itemize}

\section{Best Approximation}\label{best-approximation}

\textbf{Best Approximation in \(L^2\)}

\begin{itemize}
\tightlist
\item
  Given \(f \in L^2\) and subspace \(V\), find \(g \in V\) minimizing
  \(\|f - g\|_2\)
\item
  Theorem: unique best approximation exists; it's the orthogonal
  projection
\item
  Characterization: \(f - g \perp V\)
\end{itemize}

\textbf{The Projection Formula}

\begin{itemize}
\tightlist
\item
  If \(\{\phi_1, \ldots, \phi_n\}\) orthonormal:
  \[g = \sum_{k=1}^n \langle f, \phi_k \rangle \phi_k\]
\item
  Coefficients: \(c_k = \langle f, \phi_k \rangle\)
\item
  This is linear algebra in function spaces
\end{itemize}

\textbf{Concrete Example: Best Cubic Approximation to \(e^x\) on
\([-1,1]\)}

\begin{itemize}
\tightlist
\item
  Goal: find cubic \(p(x)\) minimizing
  \(\int_{-1}^1 |e^x - p(x)|^2\, dx\)
\item
  Step 1: Start with \(\{1, x, x^2, x^3\}\)
\item
  Step 2: Apply Gram-Schmidt to get orthonormal
  \(\{\phi_0, \phi_1, \phi_2, \phi_3\}\)

  \begin{itemize}
  \tightlist
  \item
    \(\phi_0 = \frac{1}{\sqrt{2}}\)
  \item
    \(\phi_1 = \sqrt{\frac{3}{2}} x\)
  \item
    \(\phi_2, \phi_3\): messier but computable
  \end{itemize}
\item
  Step 3: Compute \(c_k = \int_{-1}^1 e^x \phi_k(x)\, dx\)
\item
  Step 4: Best cubic is \(\sum_{k=0}^3 c_k \phi_k\)
\item
  Explicit answer (can verify numerically)
\end{itemize}

\textbf{From Finite to Infinite}

\begin{itemize}
\tightlist
\item
  What if we don't stop at degree 3?
\item
  Orthonormal sequence \(\{\phi_1, \phi_2, \ldots\}\)
\item
  Partial sums:
  \(S_N f = \sum_{k=1}^N \langle f, \phi_k \rangle \phi_k\)
\end{itemize}

\textbf{Bessel's Inequality}

\begin{itemize}
\tightlist
\item
  \(\sum_{k=1}^\infty |\langle f, \phi_k \rangle|^2 \leq \|f\|_2^2\)
\item
  Proof: \(\|f - S_N f\|^2 \geq 0\), expand and rearrange
\item
  Coefficients can't be too big
\end{itemize}

\textbf{Parseval's Identity}

\begin{itemize}
\tightlist
\item
  If \(\{\phi_k\}\) is complete (spans \(L^2\)):
  \[\sum_{k=1}^\infty |\langle f, \phi_k \rangle|^2 = \|f\|_2^2\]
\item
  Bessel becomes equality
\item
  \(L^2\) convergence: \(\|f - S_N f\|_2 \to 0\)
\end{itemize}

\textbf{Density Implies Completeness}

\begin{itemize}
\tightlist
\item
  If finite linear combinations of \(\{\phi_k\}\) are dense in \(L^2\),
  Parseval holds
\item
  Connects to Stone-Weierstrass (Ch 12)
\end{itemize}

\section{Fourier Series}\label{fourier-series-1}

\textbf{The Observation}

\begin{itemize}
\tightlist
\item
  In 13.1, we had to orthonormalize polynomials via Gram-Schmidt
\item
  Sines \(\{\sin(nx)\}\) on \([0, \pi]\) are ALREADY orthogonal!
\item
  \(\int_0^\pi \sin(mx)\sin(nx)\, dx = \frac{\pi}{2}\delta_{mn}\)
\item
  No Gram-Schmidt needed --- why? (Brief: eigenfunctions of
  \(-\frac{d^2}{dx^2}\); full story in 13.4)
\end{itemize}

\textbf{The Fourier Basis}

\begin{itemize}
\tightlist
\item
  Orthonormal:
  \(\left\{\sqrt{\frac{2}{\pi}}\sin(nx)\right\}_{n=1}^\infty\) on
  \([0, \pi]\)
\item
  Density: trig polynomials dense in \(C([0,\pi])\) by Stone-Weierstrass
  (Ch 12)
\item
  Therefore complete in \(L^2\)
\end{itemize}

\textbf{Fourier Coefficients}

\begin{itemize}
\tightlist
\item
  \(b_n = \frac{2}{\pi}\int_0^\pi f(x)\sin(nx)\, dx\)
\item
  Parseval: \(\sum_{n=1}^\infty b_n^2 = \frac{2}{\pi}\|f\|_2^2\)
\item
  \(L^2\) convergence:
  \(\left\|f - \sum_{n=1}^N b_n\sin(nx)\right\|_2 \to 0\)
\end{itemize}

\textbf{Interpretation: Frequency Decomposition}

\begin{itemize}
\tightlist
\item
  \(b_n\) measures ``how much frequency \(n\)'' is in \(f\)
\item
  Low \(n\): slow variation (smooth features)
\item
  High \(n\): rapid oscillation (fine details, discontinuities)
\item
  Parseval: total energy = sum of energies in each frequency
\end{itemize}

\textbf{Application: Filtering}

\begin{itemize}
\tightlist
\item
  Low-pass filter: keep only \(b_1, \ldots, b_N\), discard higher
  frequencies
\item
  Effect: smoothing, removes rapid oscillations
\item
  High-pass filter: discard low frequencies
\item
  Effect: edge detection, keeps rapid changes
\item
  Parseval controls energy lost/kept
\end{itemize}

\textbf{Application: The Basel Problem}

\begin{itemize}
\tightlist
\item
  Expand \(f(x) = x\) on \([0, \pi]\)
\item
  Compute: \(b_n = \frac{2(-1)^{n+1}}{n}\)
\item
  Compute: \(\|f\|_2^2 = \int_0^\pi x^2\, dx = \frac{\pi^3}{3}\)
\item
  Parseval:
  \(\sum_{n=1}^\infty \frac{4}{n^2} = \frac{2}{\pi} \cdot \frac{\pi^3}{3}\)
\item
  Therefore: \(\sum_{n=1}^\infty \frac{1}{n^2} = \frac{\pi^2}{6}\)
\end{itemize}

\textbf{Application: Isoperimetric Inequality}

\begin{itemize}
\tightlist
\item
  Setup from Ch 12: maximize area subject to fixed perimeter
\item
  Parametrize closed curve: \((x(t), y(t))\) periodic
\item
  Length: \(L = \int_0^{2\pi} \sqrt{x'^2 + y'^2}\, dt\)
\item
  Signed area: \(A = \frac{1}{2}\int_0^{2\pi}(xy' - yx')\, dt\)

  \begin{itemize}
  \tightlist
  \item
    (Equals geometric area for simple curves; proof via Green's theorem
    in multivariable analysis)
  \end{itemize}
\item
  Expand \(x, y\) in Fourier series
\item
  Apply Parseval to express \(L^2\) and \(A\) in terms of coefficients
\item
  Derive: \(4\pi A \leq L^2\)
\item
  Equality iff circle
\end{itemize}

\textbf{Full Fourier Series on \([-\pi, \pi]\)} (brief)

\begin{itemize}
\tightlist
\item
  Need sines AND cosines for non-odd functions
\item
  Basis: \(\{1, \cos(nx), \sin(nx)\}_{n \geq 1}\)
\item
  Coefficients: \(a_0, a_n, b_n\)
\item
  Same theory applies
\end{itemize}

\section{\texorpdfstring{\(\bigstar\) Other
Bases}{\textbackslash bigstar Other Bases}}\label{bigstar-other-bases}

\textbf{Revisiting the Orthonormal Polynomials}

\begin{itemize}
\tightlist
\item
  In 13.1, we orthonormalized \(\{1, x, x^2, \ldots\}\) on \([-1,1]\)
\item
  These are (scalar multiples of) the Legendre polynomials \(P_n(x)\)
\item
  \(P_0 = 1\), \(P_1 = x\), \(P_2 = \frac{1}{2}(3x^2 - 1)\),
  \(P_3 = \frac{1}{2}(5x^3 - 3x)\), \ldots{}
\item
  Orthogonality:
  \(\int_{-1}^1 P_m(x) P_n(x)\, dx = \frac{2}{2n+1}\delta_{mn}\)
\item
  Density: polynomials dense by Weierstrass
\end{itemize}

\textbf{Why ``Legendre''?}

\begin{itemize}
\tightlist
\item
  Not just Gram-Schmidt output --- they're eigenfunctions of an
  operator:
  \[-\frac{d}{dx}\left[(1-x^2)\frac{d}{dx}\right] P_n = n(n+1) P_n\]
\item
  That's why they're orthogonal (self-adjoint operator)
\item
  More on this pattern in 13.4
\end{itemize}

\textbf{Application: Optimal Numerical Integration}

\begin{itemize}
\tightlist
\item
  Problem: approximate \(\int_{-1}^1 f(x)\, dx\) using \(n\) sample
  points
\item
  Quadrature rule: \(\int_{-1}^1 f \approx \sum_{i=1}^n w_i f(x_i)\)
\item
  Question: which nodes \(x_i\) are best?
\end{itemize}

\textbf{Gauss-Legendre Quadrature}

\begin{itemize}
\tightlist
\item
  Theorem: If nodes are roots of \(P_n\), quadrature is exact for all
  polynomials of degree \(\leq 2n-1\)
\item
  Remarkable: \(n\) points give \(2n-1\) degrees of exactness (optimal)
\item
  Why roots of \(P_n\)? Orthogonality kills error terms
\end{itemize}

\textbf{Why This Helps for Non-Polynomials}

\begin{itemize}
\tightlist
\item
  Let \(Q(f) = \sum w_i f(x_i)\)
\item
  For any polynomial \(p\) of degree \(\leq 2n-1\):
  \(\int f - Q(f) = \int(f-p) - Q(f-p)\)
\item
  Therefore: \(\left|\int f - Q(f)\right| \leq 4\|f - p\|_\infty\)
\item
  Taking infimum: error \(\leq 4 E_{2n-1}(f)\) where \(E_k(f)\) is best
  polynomial approximation error
\item
  For smooth \(f\): \(E_k(f) \to 0\)
\item
  For analytic \(f\): \(E_k(f) \to 0\) exponentially
\item
  Gauss-Legendre inherits optimal convergence from approximation theory
\end{itemize}

\textbf{Physical Interpretation} (remark)

\begin{itemize}
\tightlist
\item
  Legendre coefficients are multipole moments
\item
  \(c_0\) = monopole (average), \(c_1\) = dipole, \(c_2\) = quadrupole,
  \ldots{}
\item
  Generating function:
  \(\frac{1}{\sqrt{1-2xt+t^2}} = \sum_{n=0}^\infty P_n(x) t^n\)
\item
  This IS the potential of a point charge --- Legendre appears naturally
  in physics
\end{itemize}

\textbf{Haar Wavelets on \([0,1]\)} (brief; details in exercises)

\begin{itemize}
\tightlist
\item
  A very different basis: piecewise constant functions
\item
  Level 0: \(\phi = 1\)
\item
  Level 1: \(\psi = +1\) on \([0, 1/2)\), \(-1\) on \([1/2, 1)\)
\item
  Higher levels: rescaled, translated copies
\item
  Orthogonality: by construction (disjoint supports or cancellation)
\item
  Density: elementary! Step functions on dyadic intervals are dense
\item
  No Stone-Weierstrass needed
\end{itemize}

\textbf{When Haar Wins}

\begin{itemize}
\tightlist
\item
  Step function \(f = \mathbf{1}_{[1/2,1]}\):

  \begin{itemize}
  \tightlist
  \item
    In Haar: \(f = \frac{1}{2}\phi - \frac{1}{2}\psi\). Two terms.
    Exact.
  \item
    In Fourier: infinitely many terms, Gibbs phenomenon
  \end{itemize}
\item
  Functions with jumps prefer Haar
\item
  Smooth periodic functions prefer Fourier
\end{itemize}

\textbf{The Message}

\begin{itemize}
\tightlist
\item
  Different bases for different problems
\item
  Fourier: smooth, periodic, differential equations
\item
  Legendre: polynomial approximation, numerical integration
\item
  Haar: piecewise constant, jumps, localized features
\item
  Same Parseval, same projection formula, different strengths
\end{itemize}

\section{Why Complex Exponentials?}\label{why-complex-exponentials}

\textbf{Fourier's Original Problem: The Heat Equation}

\begin{itemize}
\tightlist
\item
  \(u_t = u_{xx}\) on \([0, \pi]\), boundary conditions
  \(u(0,t) = u(\pi,t) = 0\)
\item
  Separation of variables: try \(u(x,t) = X(x)T(t)\)
\item
  Get \(\frac{T'}{T} = \frac{X''}{X} = -\lambda\) (both sides constant)
\item
  \(X'' = -\lambda X\) with \(X(0) = X(\pi) = 0\)
\item
  Solutions: \(X_n = \sin(nx)\), eigenvalues \(\lambda_n = n^2\)
\item
  Each mode: \(T_n' = -n^2 T_n\) gives \(T_n(t) = e^{-n^2 t}\)
\item
  General solution:
  \(u(x,t) = \sum_{n=1}^\infty b_n \sin(nx) e^{-n^2 t}\)
\end{itemize}

\textbf{The Power of Eigenfunctions}

\begin{itemize}
\tightlist
\item
  Each mode evolves independently
\item
  Infinitely many coupled PDEs → infinitely many decoupled ODEs
\item
  Initial condition determines \(b_n\); then each \(b_n(t)\) evolves on
  its own
\end{itemize}

\textbf{Sines Diagonalize \(-\frac{d^2}{dx^2}\)}

\begin{itemize}
\tightlist
\item
  \(-\frac{d^2}{dx^2}\sin(nx) = n^2\sin(nx)\)
\item
  Eigenfunctions with eigenvalue \(n^2\)
\item
  Heat equation: \(u_t = u_{xx}\) becomes \(b_n' = -n^2 b_n\)
\item
  This works for ANY even-degree constant-coefficient operator
\item
  \(P\left(-\frac{d^2}{dx^2}\right)\sin(nx) = P(n^2)\sin(nx)\)
\end{itemize}

\textbf{What About \(\frac{d}{dx}\)?}

\begin{itemize}
\tightlist
\item
  \(\frac{d}{dx}\sin(nx) = n\cos(nx)\) --- NOT a scalar multiple of
  \(\sin(nx)\)!
\item
  \(\frac{d}{dx}\cos(nx) = -n\sin(nx)\) --- also not an eigenfunction!
\item
  Sines and cosines MIX under first derivative
\end{itemize}

\textbf{The Matrix Picture}

\begin{itemize}
\tightlist
\item
  Restrict to \(V_n = \text{span}\{\cos(nx), \sin(nx)\}\)
\item
  In basis \((\cos(nx), \sin(nx))\):
  \[\frac{d}{dx}\Big|_{V_n} \sim \begin{pmatrix} 0 & -n \\ n & 0 \end{pmatrix}\]
\item
  This is a rotation matrix (times \(n\))
\item
  Eigenvalues over \(\mathbb{R}\): none! Characteristic polynomial
  \(\lambda^2 + n^2 = 0\)
\end{itemize}

\textbf{Just Like Linear Algebra}

\begin{itemize}
\tightlist
\item
  Rotation matrices don't diagonalize over \(\mathbb{R}\)
\item
  But over \(\mathbb{C}\): eigenvalues \(\pm in\), eigenvectors involve
  \(i\)
\item
  Same situation here
\end{itemize}

\textbf{The Resolution: Work Over \(\mathbb{C}\)}

\begin{itemize}
\tightlist
\item
  Eigenvalues of \(\begin{pmatrix} 0 & -n \\ n & 0 \end{pmatrix}\):
  \(\lambda = \pm in\)
\item
  Eigenvectors: \(\begin{pmatrix} 1 \\ -i \end{pmatrix}\) and
  \(\begin{pmatrix} 1 \\ i \end{pmatrix}\)
\item
  In function terms: \(\cos(nx) \pm i\sin(nx) = e^{\pm inx}\)
\end{itemize}

\textbf{Complex Exponentials ARE Eigenfunctions of \(\frac{d}{dx}\)}

\begin{itemize}
\tightlist
\item
  \(\frac{d}{dx}e^{inx} = in \cdot e^{inx}\)
\item
  Each \(e^{inx}\) is an eigenfunction with eigenvalue \(in\)
\item
  No mixing!
\end{itemize}

\textbf{Complex Fourier Series}

\begin{itemize}
\tightlist
\item
  Basis: \(\{e^{inx}\}_{n \in \mathbb{Z}}\) on \([-\pi, \pi]\)
\item
  Orthogonality:
  \(\frac{1}{2\pi}\int_{-\pi}^{\pi} e^{imx}\overline{e^{inx}}\, dx = \delta_{mn}\)
\item
  Coefficients:
  \(c_n = \frac{1}{2\pi}\int_{-\pi}^{\pi} f(x)e^{-inx}\, dx\)
\item
  Parseval:
  \(\sum_{n \in \mathbb{Z}} |c_n|^2 = \frac{1}{2\pi}\|f\|_2^2\)
\end{itemize}

\textbf{The Punchline}

\begin{itemize}
\tightlist
\item
  Complex exponentials diagonalize \(\frac{d}{dx}\)
\item
  Therefore they diagonalize ALL of \(\mathbb{C}[\frac{d}{dx}]\)
\item
  For any polynomial \(P\):
  \(P\left(\frac{d}{dx}\right)e^{inx} = P(in) \cdot e^{inx}\)
\item
  ANY constant-coefficient differential operator!
\end{itemize}

\textbf{Example: First-Order ODE}

\begin{itemize}
\tightlist
\item
  Solve \(y' + y = f\) with \(f\) periodic
\item
  Expand: \(f = \sum c_n^{(f)} e^{inx}\), try
  \(y = \sum c_n^{(y)} e^{inx}\)
\item
  Substitute: \((in + 1)c_n^{(y)} = c_n^{(f)}\)
\item
  Solution: \(c_n^{(y)} = \frac{c_n^{(f)}}{1 + in}\)
\item
  Each coefficient computed independently --- no coupling!
\end{itemize}

\textbf{Contrast: Same ODE in Sine-Cosine Basis}

\begin{itemize}
\tightlist
\item
  Expand \(f = \sum(a_n \cos(nx) + b_n \sin(nx))\)
\item
  The equation \(y' + y = f\) gives coupled \(2 \times 2\) systems for
  each \(n\)
\item
  Solvable, but messier
\end{itemize}

\textbf{Translation for Real Functions}

\begin{itemize}
\tightlist
\item
  Real \(f\) means \(c_{-n} = \overline{c_n}\)
\item
  Relationship: \(a_n = c_n + c_{-n}\), \(b_n = i(c_n - c_{-n})\)
\item
  Complex form is not ``just cleaner notation''
\item
  It's the natural eigenbasis for differentiation
\end{itemize}

\textbf{Why Fourier is Ubiquitous}

\begin{itemize}
\tightlist
\item
  Fourier discovered the right basis for working with \(\frac{d}{dx}\)
\item
  Any problem with constant-coefficient differential operators
  simplifies in this basis
\item
  Heat equation, wave equation, signal processing, quantum
  mechanics\ldots{}
\item
  Not because sines are special --- because complex exponentials
  diagonalize differentiation
\end{itemize}

\textbf{Looking Ahead}

\begin{itemize}
\tightlist
\item
  On \(\mathbb{R}\) (non-periodic): eigenfunctions \(e^{i\xi x}\) for
  continuous \(\xi\)
\item
  But \(e^{i\xi x} \notin L^2(\mathbb{R})\)!
\item
  Need distributions (Ch 15) to make sense of this
\item
  Leads to the Fourier transform
\end{itemize}

\section{Pointwise Convergence}\label{pointwise-convergence}

\textbf{The Question}

\begin{itemize}
\tightlist
\item
  We have \(L^2\) convergence: \(\|f - S_N f\|_2 \to 0\)
\item
  Does \(S_N f(x) \to f(x)\) for each \(x\)?
\end{itemize}

\textbf{The Answer: Not Always}

\begin{itemize}
\tightlist
\item
  \(L^2\) convergence does NOT imply pointwise convergence
\item
  There exist continuous \(f\) whose Fourier series diverges at a point!
\item
  This is subtle and took decades to understand
\end{itemize}

\textbf{Fourier Partial Sums as Convolution}

\begin{itemize}
\tightlist
\item
  \(S_N f(x) = \frac{1}{\pi}\int_0^\pi f(t) D_N(x-t)\, dt = (f * D_N)(x)\)
\item
  Dirichlet kernel:
  \(D_N(t) = \frac{\sin((N+\frac{1}{2})t)}{\sin(\frac{t}{2})}\)
\end{itemize}

\textbf{Why Dirichlet Fails}

\begin{itemize}
\tightlist
\item
  Recall approximate identities (Ch 12): \(\phi_\epsilon \geq 0\),
  \(\int \phi_\epsilon = 1\), concentrates at 0
\item
  Dirichlet kernel: \(\int D_N = \pi\) ✓, concentrates ✓, but
  \(D_N \not\geq 0\) ✗
\item
  In fact: \(\int |D_N| \to \infty\) as \(N \to \infty\)
\item
  NOT an approximate identity
\end{itemize}

\textbf{The Fejér Kernel}

\begin{itemize}
\tightlist
\item
  Cesàro means:
  \(\sigma_N f = \frac{1}{N+1}(S_0 f + S_1 f + \cdots + S_N f)\)
\item
  This is also a convolution: \(\sigma_N f = f * F_N\)
\item
  Fejér kernel:
  \(F_N(t) = \frac{1}{N+1}\left(\frac{\sin((N+1)t/2)}{\sin(t/2)}\right)^2\)
\item
  Key property: \(F_N \geq 0\) (it's a square!)
\item
  Also: \(\int F_N = \pi\), and \(F_N\) concentrates at 0
\item
  Fejér kernel IS an approximate identity
\end{itemize}

\textbf{Fejér's Theorem}

\begin{itemize}
\tightlist
\item
  Theorem: For continuous \(f\), \(\sigma_N f \to f\) uniformly
\item
  Proof: Apply approximate identity theorem (Ch 12)
\item
  Cesàro averaging ``smooths out'' the divergence of partial sums
\end{itemize}

\textbf{Corollary: Another Proof of Weierstrass}

\begin{itemize}
\tightlist
\item
  \(\sigma_N f\) is a trig polynomial
\item
  \(\sigma_N f \to f\) uniformly
\item
  Therefore trig polynomials are dense in \(C(\mathbb{T})\)
\end{itemize}

\textbf{The Moral}

\begin{itemize}
\tightlist
\item
  Ordinary convergence can fail
\item
  Cesàro summation recovers convergence
\item
  Averaging is powerful
\end{itemize}

\chapter{Differential Equations}\label{differential-equations}

We apply the tools of analysis to differential equations. Completeness
gives existence via Picard iteration. Matrix exponentials solve linear
systems. Fourier methods solve PDEs. And numerical methods require
analysis to prove they work.

\textbf{Historical Prelude}

\begin{itemize}
\tightlist
\item
  Newton: calculus born from differential equations (motion, gravity)
\item
  18th-19th century: solving specific equations, cataloguing tricks
\item
  Cauchy, Lipschitz, Picard: existence and uniqueness theory
\item
  Poincaré: qualitative theory when explicit solutions fail
\item
  Today: numerical methods dominate practice, but analysis proves they
  work
\end{itemize}

\section{Existence and Uniqueness}\label{existence-and-uniqueness}

\textbf{The Problem}

\begin{itemize}
\tightlist
\item
  ODE: \(y' = f(t, y)\), initial condition \(y(t_0) = y_0\)
\item
  Does a solution exist? Is it unique?
\item
  Not obvious! \(y' = y^2\), \(y(0) = 1\) has solution
  \(y = \frac{1}{1-t}\), blows up at \(t = 1\)
\end{itemize}

\textbf{Reduction to Integral Equation}

\begin{itemize}
\tightlist
\item
  \(y' = f(t, y)\) with \(y(t_0) = y_0\)
\item
  Integrate: \(y(t) = y_0 + \int_{t_0}^t f(s, y(s))\, ds\)
\item
  Fixed point problem: \(y = Ty\) where \(T\) is the integral operator
\end{itemize}

\textbf{The Picard Operator}

\begin{itemize}
\tightlist
\item
  Define
  \(T: C([t_0 - \delta, t_0 + \delta]) \to C([t_0 - \delta, t_0 + \delta])\)
\item
  \((Ty)(t) = y_0 + \int_{t_0}^t f(s, y(s))\, ds\)
\item
  Fixed point of \(T\) = solution of ODE
\end{itemize}

\textbf{The Lipschitz Condition}

\begin{itemize}
\tightlist
\item
  \(f\) is Lipschitz in \(y\):
  \(|f(t, y_1) - f(t, y_2)| \leq L|y_1 - y_2|\)
\item
  Key estimate:
  \(\|Ty_1 - Ty_2\|_\infty \leq L\delta \|y_1 - y_2\|_\infty\)
\item
  For \(\delta < 1/L\): \(T\) is a contraction!
\end{itemize}

\textbf{Picard-Lindelöf Theorem}

\begin{itemize}
\tightlist
\item
  Theorem: If \(f\) is continuous and Lipschitz in \(y\), then the IVP
  has a unique local solution
\item
  Proof: Banach fixed point theorem (Ch 11)
\item
  The payoff of abstract completeness!
\end{itemize}

\textbf{Local vs Global}

\begin{itemize}
\tightlist
\item
  Theorem gives solution on \([t_0 - \delta, t_0 + \delta]\)
\item
  Can extend until solution blows up or leaves domain
\item
  Example: \(y' = y^2\) blows up in finite time
\item
  Example: \(y' = y\) extends globally
\end{itemize}

\textbf{Gronwall's Inequality}

\begin{itemize}
\tightlist
\item
  Lemma: If \(u(t) \leq \alpha + \int_0^t \beta u(s)\, ds\) then
  \(u(t) \leq \alpha e^{\beta t}\)
\item
  Proof: Let \(U(t) = \int_0^t u(s)\, ds\), then
  \(U' \leq \alpha + \beta U\), multiply by \(e^{-\beta t}\)
\item
  The fundamental tool for controlling growth
\end{itemize}

\textbf{Continuous Dependence on Initial Conditions}

\begin{itemize}
\tightlist
\item
  Theorem: \(|y(t; y_0) - y(t; z_0)| \leq |y_0 - z_0| e^{Lt}\)
\item
  Proof: Apply Gronwall to \(u(t) = |y(t) - z(t)|\)
\item
  Message: Solutions depend continuously on initial data (short time)
\item
  Warning: Exponential growth means small errors can amplify
\end{itemize}

\textbf{Continuous Dependence on Parameters}

\begin{itemize}
\tightlist
\item
  If \(f = f(t, y; \mu)\) depends on parameter \(\mu\)
\item
  Solution \(y(t; \mu)\) depends continuously on \(\mu\)
\item
  Same proof technique
\end{itemize}

\textbf{Systems of ODEs}

\begin{itemize}
\tightlist
\item
  \(\mathbf{y}' = \mathbf{f}(t, \mathbf{y})\) where
  \(\mathbf{y} \in \mathbb{R}^n\)
\item
  Same theory applies with \(|\cdot|\) replaced by vector norm
\item
  Higher-order equations reduce to first-order systems:

  \begin{itemize}
  \tightlist
  \item
    \(y'' = g(t, y, y')\) becomes \(y_1' = y_2\),
    \(y_2' = g(t, y_1, y_2)\)
  \end{itemize}
\end{itemize}

\section{Linear ODEs}\label{linear-odes}

\textbf{Higher-Order Linear Equations}

\begin{itemize}
\tightlist
\item
  \(y^{(n)} + a_{n-1}y^{(n-1)} + \cdots + a_0 y = 0\)
\item
  Convert to system via companion matrix
\item
  Eigenvalues of companion matrix = roots of characteristic polynomial
\item
  Standard theory: \(n\) linearly independent solutions, general
  solution is linear combination
\end{itemize}

\textbf{The Problem}

\begin{itemize}
\tightlist
\item
  \(\mathbf{y}' = A\mathbf{y}\) where \(A\) is \(n \times n\) constant
  matrix
\item
  Initial condition \(\mathbf{y}(0) = \mathbf{y}_0\)
\end{itemize}

\textbf{The Matrix Exponential (Recall from Ch 11)}

\begin{itemize}
\tightlist
\item
  \(e^{A} = \sum_{k=0}^\infty \frac{A^k}{k!}\)
\item
  Converges for all \(A\) (comparison with \(e^{\|A\|}\))
\item
  Key property: \(\frac{d}{dt} e^{At} = A e^{At}\)
\end{itemize}

\textbf{Solution of \(\mathbf{y}' = A\mathbf{y}\)}

\begin{itemize}
\tightlist
\item
  Theorem: \(\mathbf{y}(t) = e^{At}\mathbf{y}_0\)
\item
  Proof: Verify by differentiation
\item
  Existence, uniqueness, and explicit formula all at once
\end{itemize}

\textbf{Computing \(e^{At}\): Diagonalizable Case}

\begin{itemize}
\tightlist
\item
  If \(A = PDP^{-1}\) with \(D\) diagonal, then
  \(e^{At} = P e^{Dt} P^{-1}\)
\item
  \(e^{Dt} = \text{diag}(e^{\lambda_1 t}, \ldots, e^{\lambda_n t})\)
\item
  Each eigenvalue contributes a mode
\end{itemize}

\textbf{Real vs Complex Eigenvalues}

\begin{itemize}
\tightlist
\item
  Real \(\lambda\): mode \(e^{\lambda t}\) (growth or decay)
\item
  Complex \(\lambda = a + bi\): mode \(e^{at}(\cos(bt) + i\sin(bt))\)
\item
  For real \(A\): complex eigenvalues come in conjugate pairs
\item
  Combined contribution: \(e^{at}(c_1\cos(bt) + c_2\sin(bt))\) ---
  oscillation!
\end{itemize}

\textbf{Example: Harmonic Oscillator}

\begin{itemize}
\tightlist
\item
  \(y'' + \omega^2 y = 0\) becomes
  \(\mathbf{y}' = \begin{pmatrix} 0 & 1 \\ -\omega^2 & 0 \end{pmatrix} \mathbf{y}\)
\item
  Eigenvalues: \(\pm i\omega\)
\item
  Solution: \(y(t) = c_1 \cos(\omega t) + c_2 \sin(\omega t)\)
\item
  Complex eigenvalues explain oscillation
\end{itemize}

\textbf{Non-Diagonalizable Case} (brief)

\begin{itemize}
\tightlist
\item
  Jordan normal form: \(A = PJP^{-1}\)
\item
  \(e^{Jt}\) involves polynomial times exponential:
  \(t^k e^{\lambda t}\)
\item
  Generalized eigenvectors
\item
  State without developing fully
\end{itemize}

\textbf{Inhomogeneous Systems: Variation of Parameters}

\begin{itemize}
\tightlist
\item
  Problem: \(\mathbf{y}' = A\mathbf{y} + \mathbf{g}(t)\)
\item
  Guess: \(\mathbf{y}(t) = e^{At}\mathbf{c}(t)\) (vary the ``constant'')
\item
  Substitute: \(\mathbf{c}'(t) = e^{-At}\mathbf{g}(t)\)
\item
  Solution:
  \(\mathbf{y}(t) = e^{At}\mathbf{y}_0 + \int_0^t e^{A(t-s)}\mathbf{g}(s)\, ds\)
\item
  The integral is a convolution!
\end{itemize}

\section{Eigenfunction Methods}\label{eigenfunction-methods}

LETS DO SOME ODE STUFF FIRST!

\textbf{The Heat Equation}

\begin{itemize}
\tightlist
\item
  \(u_t = u_{xx}\) on \([0, \pi]\), boundary conditions
  \(u(0, t) = u(\pi, t) = 0\)
\item
  Initial condition: \(u(x, 0) = f(x)\)
\item
  Physically: temperature in a rod with endpoints held at zero
\end{itemize}

\textbf{Separation of Variables}

\begin{itemize}
\tightlist
\item
  Guess: \(u(x, t) = X(x)T(t)\)
\item
  Substitute: \(\frac{T'}{T} = \frac{X''}{X} = -\lambda\) (both sides
  constant)
\item
  Two ODEs: \(X'' = -\lambda X\) and \(T' = -\lambda T\)
\end{itemize}

\textbf{The Eigenvalue Problem}

\begin{itemize}
\tightlist
\item
  \(X'' = -\lambda X\) with \(X(0) = X(\pi) = 0\)
\item
  Nontrivial solutions only for \(\lambda = n^2\),
  \(n = 1, 2, 3, \ldots\)
\item
  Eigenfunctions: \(X_n(x) = \sin(nx)\)
\item
  This is a Sturm-Liouville problem!
\end{itemize}

\textbf{The Time Evolution}

\begin{itemize}
\tightlist
\item
  For each \(n\): \(T_n' = -n^2 T_n\)
\item
  Solution: \(T_n(t) = b_n e^{-n^2 t}\)
\end{itemize}

\textbf{The General Solution}

\begin{itemize}
\tightlist
\item
  Superposition: \(u(x, t) = \sum_{n=1}^\infty b_n \sin(nx) e^{-n^2 t}\)
\item
  Coefficients from initial condition:
  \(b_n = \frac{2}{\pi}\int_0^\pi f(x)\sin(nx)\, dx\)
\item
  This is the Fourier series of \(f\)!
\end{itemize}

\textbf{Smoothing Property}

\begin{itemize}
\tightlist
\item
  Each mode decays like \(e^{-n^2 t}\)
\item
  Higher frequencies decay faster
\item
  For any \(t > 0\): solution is \(C^\infty\) even if \(f\) is rough
\item
  Heat equation ``smooths out'' irregularities
\end{itemize}

\textbf{The Wave Equation}

\begin{itemize}
\tightlist
\item
  \(u_{tt} = c^2 u_{xx}\) on \([0, \pi]\), boundary conditions
  \(u(0, t) = u(\pi, t) = 0\)
\item
  Initial conditions: \(u(x, 0) = f(x)\), \(u_t(x, 0) = g(x)\)
\item
  Physically: vibrating string with fixed endpoints
\end{itemize}

\textbf{Separation of Variables}

\begin{itemize}
\tightlist
\item
  Same eigenvalue problem for \(X\): \(X_n = \sin(nx)\),
  \(\lambda_n = n^2\)
\item
  Time equation: \(T'' = -c^2 n^2 T\)
\item
  Solution: \(T_n(t) = a_n \cos(cnt) + b_n \sin(cnt)\)
\end{itemize}

\textbf{The General Solution}

\begin{itemize}
\tightlist
\item
  \(u(x, t) = \sum_{n=1}^\infty (a_n \cos(cnt) + b_n \sin(cnt))\sin(nx)\)
\item
  Standing waves: each mode oscillates in place
\item
  Coefficients from initial conditions
\end{itemize}

\textbf{D'Alembert's Solution} (brief)

\begin{itemize}
\tightlist
\item
  Alternative: \(u(x, t) = F(x + ct) + G(x - ct)\)
\item
  Traveling waves moving left and right
\item
  Equivalent to Fourier form but different perspective
\end{itemize}

\textbf{Energy Conservation}

\begin{itemize}
\tightlist
\item
  Define \(E(t) = \frac{1}{2}\int_0^\pi (u_t^2 + c^2 u_x^2)\, dx\)
\item
  Theorem: \(E'(t) = 0\)
\item
  Proof: Differentiate, use PDE and integration by parts
\item
  Contrast with heat equation: wave conserves energy, heat dissipates
\end{itemize}

\textbf{Why Does This Work? The General Principle}

Sines appeared because they are eigenfunctions of \(-\frac{d^2}{dx^2}\)
with Dirichlet boundary conditions.

The key property of eigenfunctions: applying \(L\) just multiplies by
the eigenvalue. So the PDE decouples into independent ODEs for each
mode.

This is a general principle: given a self-adjoint operator \(L\): -
Eigenvalues are real - Eigenfunctions for distinct eigenvalues are
orthogonal - Under suitable conditions, eigenfunctions form a complete
basis

To solve \(u_t = Lu\) or \(u_{tt} = Lu\): 1. Find eigenfunctions of
\(L\) 2. Expand in that basis 3. Solve decoupled ODEs for coefficients

\textbf{Sturm-Liouville Problems}

General form:
\(-\frac{d}{dx}\left[p(x)\frac{du}{dx}\right] + q(x)u = \lambda w(x)u\)
on \([a, b]\)

With boundary conditions (Dirichlet, Neumann, or mixed).

Under appropriate conditions on \(p, q, w\): - Eigenvalues \(\lambda_n\)
form a sequence \(\lambda_1 < \lambda_2 < \cdots \to \infty\) -
Eigenfunctions \(\phi_n\) are orthogonal with respect to weight \(w\) -
Eigenfunctions are complete in \(L^2([a, b], w)\)

\textbf{Other Operators, Other Bases}

\begin{longtable}[]{@{}
  >{\raggedright\arraybackslash}p{(\linewidth - 4\tabcolsep) * \real{0.2571}}
  >{\raggedright\arraybackslash}p{(\linewidth - 4\tabcolsep) * \real{0.2857}}
  >{\raggedright\arraybackslash}p{(\linewidth - 4\tabcolsep) * \real{0.4571}}@{}}
\toprule\noalign{}
\begin{minipage}[b]{\linewidth}\raggedright
Problem
\end{minipage} & \begin{minipage}[b]{\linewidth}\raggedright
Operator
\end{minipage} & \begin{minipage}[b]{\linewidth}\raggedright
Eigenfunctions
\end{minipage} \\
\midrule\noalign{}
\endhead
\bottomrule\noalign{}
\endlastfoot
Heat/wave on interval & \(-\frac{d^2}{dx^2}\) (Dirichlet) &
\(\sin(nx)\) \\
Heat/wave on circle & \(-\frac{d^2}{dx^2}\) (periodic) & \(e^{inx}\) \\
Spherical symmetry & Legendre operator & \(P_n(x)\) \\
Cylindrical symmetry & Bessel operator & \(J_n(x)\) \\
Quantum harmonic oscillator & \(-\frac{d^2}{dx^2} + x^2\) & Hermite
functions \\
\end{longtable}

The technique is always the same: expand in eigenfunctions of your
operator.

\section{Numerical Methods}\label{numerical-methods}

\textbf{The Problem}

\begin{itemize}
\tightlist
\item
  Most ODEs can't be solved explicitly
\item
  Need numerical approximation
\item
  But does the approximation converge to the true solution?
\item
  Analysis gives the answer
\end{itemize}

\textbf{Picard Iteration as Computation}

\begin{itemize}
\tightlist
\item
  Recall: \(y_{n+1}(t) = y_0 + \int_0^t f(s, y_n(s))\, ds\)
\item
  This isn't just for existence proofs --- it computes!
\end{itemize}

\textbf{Example: \(y' = y\), \(y(0) = 1\)}

\begin{itemize}
\tightlist
\item
  \(y_0(t) = 1\)
\item
  \(y_1(t) = 1 + \int_0^t 1\, ds = 1 + t\)
\item
  \(y_2(t) = 1 + \int_0^t (1 + s)\, ds = 1 + t + \frac{t^2}{2}\)
\item
  \(y_n(t) = \sum_{k=0}^n \frac{t^k}{k!}\)
\item
  Limit: \(e^t\)
\item
  We derived the Taylor series!
\end{itemize}

\textbf{Example: \(y'' = -y\), \(y(0) = 0\), \(y'(0) = 1\)}

\begin{itemize}
\tightlist
\item
  System: \(y_1' = y_2\), \(y_2' = -y_1\), with
  \((y_1(0), y_2(0)) = (0, 1)\)
\item
  Iterate:

  \begin{itemize}
  \tightlist
  \item
    \((y_1^{(0)}, y_2^{(0)}) = (0, 1)\)
  \item
    \(y_1^{(1)} = t\), \(y_2^{(1)} = 1\)
  \item
    \(y_1^{(2)} = t\), \(y_2^{(2)} = 1 - \frac{t^2}{2}\)
  \item
    \(y_1^{(3)} = t - \frac{t^3}{6}\), \(y_2^{(3)} = 1 - \frac{t^2}{2}\)
  \item
    \(y_1^{(4)} = t - \frac{t^3}{6}\),
    \(y_2^{(4)} = 1 - \frac{t^2}{2} + \frac{t^4}{24}\)
  \end{itemize}
\item
  Limits: \(\sin(t)\) and \(\cos(t)\)
\item
  Taylor series emerges term by term!
\end{itemize}

\textbf{Euler's Method}

\begin{itemize}
\tightlist
\item
  Idea: Follow the tangent line
\item
  \(y_{n+1} = y_n + h f(t_n, y_n)\)
\item
  Step from \(t_n\) to \(t_{n+1} = t_n + h\)
\item
  Simple, explicit, easy to implement
\end{itemize}

\textbf{Local Truncation Error}

\begin{itemize}
\tightlist
\item
  True solution: \(y(t + h) = y(t) + hy'(t) + \frac{h^2}{2}y''(\xi)\)
\item
  Euler gives: \(y(t) + hf(t, y(t)) = y(t) + hy'(t)\)
\item
  Error per step: \(O(h^2)\)
\end{itemize}

\textbf{Global Error}

\begin{itemize}
\tightlist
\item
  \(N = T/h\) steps to reach time \(T\)
\item
  Naive bound: \(N \times O(h^2) = O(h)\)
\item
  But errors propagate and compound!
\item
  Need Gronwall to control accumulation
\end{itemize}

\textbf{Convergence Theorem for Euler's Method}

\begin{itemize}
\tightlist
\item
  Theorem: If \(f\) is Lipschitz with constant \(L\), and \(y\) has
  bounded second derivative, then
  \[\max_{0 \leq n \leq N} |y(t_n) - y_n| \leq Ch\] where \(C\) depends
  on \(L\), \(T\), and \(\|y''\|_\infty\)
\item
  Proof outline:

  \begin{itemize}
  \tightlist
  \item
    Let \(e_n = y(t_n) - y_n\)
  \item
    Local error: \(|e_{n+1}| \leq |e_n| + hL|e_n| + Ch^2\)
  \item
    Rearrange: \(|e_{n+1}| \leq (1 + hL)|e_n| + Ch^2\)
  \item
    Iterate: \(|e_n| \leq \frac{C h}{L}((1 + hL)^n - 1)\)
  \item
    Since \((1 + hL)^n \leq e^{nLh} = e^{LT}\): bound is \(O(h)\)
  \end{itemize}
\item
  This is Gronwall in discrete form!
\end{itemize}

\textbf{Backward Euler (Implicit Method)}

\begin{itemize}
\tightlist
\item
  \(y_{n+1} = y_n + hf(t_{n+1}, y_{n+1})\)
\item
  Note: \(y_{n+1}\) appears on both sides!
\item
  Must solve an equation at each step
\end{itemize}

\textbf{Why Implicit?}

\begin{itemize}
\tightlist
\item
  Explicit Euler can be unstable for ``stiff'' equations
\item
  Example: \(y' = -100y\) with \(h = 0.1\)

  \begin{itemize}
  \tightlist
  \item
    True solution decays
  \item
    Euler: \(y_{n+1} = (1 - 100 \cdot 0.1)y_n = -9y_n\) --- oscillates
    and grows!
  \end{itemize}
\item
  Backward Euler: \(y_{n+1} = y_n - 100h \cdot y_{n+1}\)

  \begin{itemize}
  \tightlist
  \item
    Solve: \(y_{n+1} = \frac{y_n}{1 + 100h}\) --- stable!
  \end{itemize}
\end{itemize}

\textbf{Solving the Implicit Equation}

\begin{itemize}
\tightlist
\item
  Need to solve \(y_{n+1} = y_n + hf(t_{n+1}, y_{n+1})\)
\item
  Define \(T(z) = y_n + hf(t_{n+1}, z)\)
\item
  Fixed point: \(y_{n+1} = T(y_{n+1})\)
\item
  For small \(h\):
  \(|T(z_1) - T(z_2)| = h|f(t_{n+1}, z_1) - f(t_{n+1}, z_2)| \leq hL|z_1 - z_2|\)
\item
  If \(hL < 1\): \(T\) is a contraction!
\item
  Banach fixed point theorem solves the implicit equation!
\end{itemize}

\textbf{Higher-Order Methods} (stated)

\begin{itemize}
\tightlist
\item
  Runge-Kutta methods: match more Taylor coefficients
\item
  RK2 (midpoint): \(O(h^2)\) global error
\item
  RK4 (classical): \(O(h^4)\) global error
\item
  Same convergence theory applies (more tedious to prove)
\item
  Higher order = fewer steps for same accuracy
\end{itemize}

\textbf{The Message}

\begin{itemize}
\tightlist
\item
  Picard: existence proof that also computes Taylor series
\item
  Euler: simple, explicit, \(O(h)\), uses Gronwall for convergence
\item
  Backward Euler: implicit, stable, uses contraction mapping at each
  step
\item
  Analysis tools (completeness, Gronwall, contraction) make numerical
  methods rigorous
\end{itemize}

\chapter{Beyond Functions}\label{beyond-functions}

We complete the space of functions to allow differentiation everywhere
and limits of approximate identities. Distributions extend
integration-against-a-function to general linear functionals. This makes
differentiation total, gives meaning to \(\delta\), and unlocks the
Fourier transform on \(\mathbb{R}\).

\textbf{The recurring theme: sequences that ``should'' converge}

\begin{itemize}
\tightlist
\item
  Approximate identities: \(\phi_\epsilon \to ???\) as
  \(\epsilon \to 0\)

  \begin{itemize}
  \tightlist
  \item
    Concentrated, positive, integral 1
  \item
    Not Cauchy in any \(L^p\) --- they blow up!
  \item
    But \(\int \phi_\epsilon f \to f(0)\) for every nice \(f\)
  \end{itemize}
\item
  Heat kernel: \(K_t \to ???\) as \(t \to 0\)

  \begin{itemize}
  \tightlist
  \item
    Same story
  \item
    The ``limit'' should be initial data for heat equation
  \end{itemize}
\item
  Derivatives of rough functions:

  \begin{itemize}
  \tightlist
  \item
    Step function \(H\): what is \(H'\)?
  \item
    Approximate: \(H_\epsilon\) smooth, \(H_\epsilon \to H\)
  \item
    \(H_\epsilon'\) is concentrated spike at origin
  \item
    \(H_\epsilon' \to ???\)
  \end{itemize}
\item
  Eigenfunctions that aren't functions (Ch 13):

  \begin{itemize}
  \tightlist
  \item
    \(e^{i\xi x}\) should be eigenfunction of \(\frac{d}{dx}\)
  \item
    But \(e^{i\xi x} \notin L^2(\mathbb{R})\)
  \end{itemize}
\item
  Wave equation kernel (Ch 14):

  \begin{itemize}
  \tightlist
  \item
    \(W_t = \frac{1}{2}(\delta_{ct} + \delta_{-ct})\)
  \item
    Can't even state without \(\delta\)
  \end{itemize}
\end{itemize}

\textbf{Historical note:} - Heaviside (1890s): operational calculus,
treats \(\delta\) formally - Dirac (1930s): uses \(\delta\) freely in
quantum mechanics - Mathematicians: ``That's not rigorous!'' - Schwartz
(1950s): theory of distributions makes it rigorous, wins Fields Medal

\textbf{The theme: completion}

\begin{longtable}[]{@{}
  >{\raggedright\arraybackslash}p{(\linewidth - 6\tabcolsep) * \real{0.1667}}
  >{\raggedright\arraybackslash}p{(\linewidth - 6\tabcolsep) * \real{0.1667}}
  >{\raggedright\arraybackslash}p{(\linewidth - 6\tabcolsep) * \real{0.3810}}
  >{\raggedright\arraybackslash}p{(\linewidth - 6\tabcolsep) * \real{0.2857}}@{}}
\toprule\noalign{}
\begin{minipage}[b]{\linewidth}\raggedright
Level
\end{minipage} & \begin{minipage}[b]{\linewidth}\raggedright
Space
\end{minipage} & \begin{minipage}[b]{\linewidth}\raggedright
What's missing
\end{minipage} & \begin{minipage}[b]{\linewidth}\raggedright
Completion
\end{minipage} \\
\midrule\noalign{}
\endhead
\bottomrule\noalign{}
\endlastfoot
Ch 1 & \(\mathbb{Q}\) & Limits of Cauchy sequences & \(\mathbb{R}\) \\
Ch 10 & Riemann integrable & Limits of monotone sequences & Lebesgue
integrable \\
Ch 15 & Smooth functions & Derivatives of rough functions; limits of
approximate identities & Distributions \\
\end{longtable}

``Reaching for infinity'' --- we keep completing to allow more limiting
operations.

\section{Reaching for Infinity}\label{reaching-for-infinity}

NOTE TO SELF

\begin{itemize}
\tightlist
\item
  perhaps we do the periodic case first? (or in parallel)?
\item
  some things are easier (no worries about compact support)
\item
  duality is \emph{even cleaner} for fourier
\item
  is our original kernel for the fourier series proof a periodic one?
\end{itemize}

\textbf{What approximate identities do}

\begin{itemize}
\tightlist
\item
  Recall (Ch 12): \(\phi_\epsilon\) with \(\int \phi_\epsilon = 1\),
  \(\phi_\epsilon \geq 0\), concentrating at 0
\item
  For each \(\epsilon\): \(f \mapsto \int \phi_\epsilon(x) f(x)\, dx\)
  is a linear functional
\item
  As \(\epsilon \to 0\): converges to \(f \mapsto f(0)\)
\item
  The limit is still a linear functional\ldots{} but not integration
  against any function
\end{itemize}

\textbf{The key observation}

\begin{itemize}
\tightlist
\item
  Integration against a function: \(f \mapsto \int g(x) f(x)\, dx\)
\item
  This is a linear functional on test functions
\item
  Not every linear functional arises this way
\item
  But limits of such functionals should still be ``something''
\end{itemize}

\textbf{Test functions}

\begin{itemize}
\tightlist
\item
  Definition: \(\mathcal{D} = C_c^\infty(\mathbb{R})\) --- smooth,
  compactly supported
\item
  Why smooth? Can integrate by parts freely
\item
  Why compact support? No convergence issues at infinity
\item
  Examples: bump functions
\item
  \(\mathcal{D}\) is a vector space
\end{itemize}

\textbf{Convergence in \(\mathcal{D}\)}

\begin{itemize}
\tightlist
\item
  Definition: \(\phi_n \to \phi\) in \(\mathcal{D}\) means:

  \begin{enumerate}
  \def\labelenumi{\arabic{enumi}.}
  \tightlist
  \item
    All \(\phi_n\) supported in common compact set \(K\)
  \item
    \(\phi_n^{(k)} \to \phi^{(k)}\) uniformly for all \(k \geq 0\)
  \end{enumerate}
\item
  This is strong convergence --- controls all derivatives
\item
  Remark: can build from seminorms
  \(p_{K,k}(f) = \sup_{x \in K} |f^{(k)}(x)|\)
\end{itemize}

\textbf{Distributions: the definition}

\begin{itemize}
\tightlist
\item
  A distribution is a continuous linear functional
  \(T: \mathcal{D} \to \mathbb{R}\)
\item
  Linear: \(T(a\phi + b\psi) = aT(\phi) + bT(\psi)\)
\item
  Continuous: \(\phi_n \to \phi\) in \(\mathcal{D}\) implies
  \(T(\phi_n) \to T(\phi)\)
\item
  Notation: write \(\langle T, \phi \rangle\) instead of \(T(\phi)\)
\item
  Space of distributions: \(\mathcal{D}'\)
\end{itemize}

\textbf{Regular distributions}

\begin{itemize}
\tightlist
\item
  For \(g \in L^1_{loc}\): define \(T_g\) by
  \(\langle T_g, \phi \rangle = \int g(x)\phi(x)\, dx\)
\item
  Check: linear ✓
\item
  Check: continuous (DCT) ✓
\item
  So \(T_g\) is a distribution
\item
  Notation: write \(\langle g, \phi \rangle = \int g\phi\) --- extending
  integral notation
\item
  Functions ``are'' distributions (via this embedding)
\end{itemize}

\textbf{The delta distribution}

\begin{itemize}
\tightlist
\item
  Definition: \(\langle \delta, \phi \rangle = \phi(0)\)
\item
  Linear?
  \(\langle \delta, a\phi + b\psi \rangle = a\phi(0) + b\psi(0)\) ✓
\item
  Continuous? If \(\phi_n \to \phi\) uniformly, then
  \(\phi_n(0) \to \phi(0)\) ✓
\item
  Is \(\delta = T_g\) for some function \(g\)?
\item
  \textbf{Claim: No.}
\item
  Proof: Would need \(\int g\phi = \phi(0)\) for all \(\phi\)
\item
  Take \(\phi\) with \(\phi(0) = 0\) but \(\phi \neq 0\): need
  \(\int g\phi = 0\) for all such \(\phi\)
\item
  This forces \(g = 0\) a.e., but then \(\int g\phi = 0 \neq \phi(0)\)
  for \(\phi\) with \(\phi(0) \neq 0\)
\item
  Contradiction. \(\delta\) is NOT a function.
\end{itemize}

\textbf{Delta as limit of approximate identities}

\begin{itemize}
\tightlist
\item
  \(\langle \phi_\epsilon, \psi \rangle = \int \phi_\epsilon(x) \psi(x)\, dx \to \psi(0) = \langle \delta, \psi \rangle\)
\item
  The approximate identities converge to \(\delta\) in the sense of
  distributions
\item
  This is the rigorous meaning of ``\(\phi_\epsilon \to \delta\)''
\item
  We have COMPLETED the space to include this limit
\end{itemize}

\textbf{Convergence of distributions}

\begin{itemize}
\tightlist
\item
  Definition: \(T_n \to T\) in \(\mathcal{D}'\) means
  \(\langle T_n, \phi \rangle \to \langle T, \phi \rangle\) for all
  \(\phi \in \mathcal{D}\)
\item
  This is weak-* convergence (pointwise on test functions)
\item
  Distributions form a complete space in this sense
\end{itemize}

\textbf{Other examples}

\begin{itemize}
\tightlist
\item
  \(\delta_a\): \(\langle \delta_a, \phi \rangle = \phi(a)\) (delta at
  point \(a\))
\item
  Principal value:
  \(\langle \text{p.v.}\frac{1}{x}, \phi \rangle = \lim_{\epsilon \to 0} \int_{|x| > \epsilon} \frac{\phi(x)}{x}\, dx\)
\item
  Functions with growth: \(e^{x}\) defines distribution via
  \(\langle e^x, \phi \rangle = \int e^x \phi\) (finite since \(\phi\)
  has compact support)
\end{itemize}

\textbf{Operations on distributions}

\begin{itemize}
\tightlist
\item
  Addition:
  \(\langle S + T, \phi \rangle = \langle S, \phi \rangle + \langle T, \phi \rangle\)
\item
  Scalar multiplication:
  \(\langle cT, \phi \rangle = c\langle T, \phi \rangle\)
\item
  Multiplication by smooth function:
  \(\langle gT, \phi \rangle = \langle T, g\phi \rangle\)
\item
  Translation:
  \(\langle T_a, \phi \rangle = \langle T, \phi(\cdot + a) \rangle\)
\item
  Dilation:
  \(\langle T(cx), \phi \rangle = \frac{1}{|c|}\langle T, \phi(\cdot/c) \rangle\)
\end{itemize}

\textbf{Extended Guided Exercise: Riesz Representation and Integration}

This exercise develops a parallel perspective: distributions as measures
we integrate against.

\emph{Part 1: Riesz Representation Theorem} - State: Every continuous
linear functional on \(C[a,b]\) has the form \(f \mapsto \int f\, d\mu\)
for a unique signed measure \(\mu\) - The functional \(f \mapsto f(c)\)
corresponds to the point mass \(\delta_c\) - So \(\delta\) really IS
something we can ``integrate against''

\emph{Part 2: Riemann-Stieltjes Connection} - Recall (or introduce):
\(\int f\, dg = \lim \sum f(x_i)(g(x_{i+1}) - g(x_i))\) - If \(g = H\)
(Heaviside), then \(\int f\, dH = f(0)\) - The ``jump'' in \(H\)
captures the point mass

\emph{Part 3: Revisiting Our Axioms} - Ch 10: We axiomatized integration
against Lebesgue measure - Could generalize: integration against
arbitrary measures - Point masses, weighted sums of point masses,
continuous measures\ldots{} - Distributions extend this further:
\(\delta'\) is ``integration against'' something that isn't even a
measure

\emph{Part 4: The Hierarchy}
\[\text{Functions} \subset \text{Measures} \subset \text{Distributions}\]
- Functions: \(\langle f, \phi \rangle = \int f\phi\, dx\) - Measures:
\(\langle \mu, \phi \rangle = \int \phi\, d\mu\) - Distributions:
\(\langle T, \phi \rangle = T(\phi)\) - Each level allows more
``integration against'' operations - Distributions complete the picture:
closed under differentiation

\section{Completing Differentiation}\label{completing-differentiation}

\textbf{Motivation: what should \(\delta'\) be?}

\begin{itemize}
\tightlist
\item
  \(\delta = \lim \phi_\epsilon\) where \(\phi_\epsilon\) are
  approximate identities
\item
  Each \(\phi_\epsilon\) is smooth, so \(\phi_\epsilon'\) exists
\item
  Should have \(\delta' = \lim \phi_\epsilon'\)
\item
  What does \(\phi_\epsilon'\) do to a test function? Integration by
  parts:
  \[\int \phi_\epsilon'(x) \psi(x)\, dx = -\int \phi_\epsilon(x) \psi'(x)\, dx\]
\item
  Taking limit:
  \(\langle \delta', \psi \rangle = -\langle \delta, \psi' \rangle = -\psi'(0)\)
\end{itemize}

\textbf{Definition of distributional derivative}

\[\langle T', \phi \rangle = -\langle T, \phi' \rangle\]

\begin{itemize}
\tightlist
\item
  Motivation: integration by parts, boundary terms vanish (compact
  support)
\item
  If \(T = T_f\) for differentiable \(f\): recovers ordinary derivative
\item
  Works for ANY distribution --- no smoothness required
\end{itemize}

\textbf{Theorem: Every distribution is infinitely differentiable}

\begin{itemize}
\tightlist
\item
  Define \(T^{(n)}\) by:
  \(\langle T^{(n)}, \phi \rangle = (-1)^n \langle T, \phi^{(n)} \rangle\)
\item
  Well-defined: test functions are infinitely differentiable
\item
  We have COMPLETED differentiation: every object now has all
  derivatives
\end{itemize}

\textbf{Example: Heaviside function}

\begin{itemize}
\tightlist
\item
  \(H(x) = \begin{cases} 0 & x < 0 \\ 1 & x \geq 0 \end{cases}\)
\item
  Compute:
  \(\langle H', \phi \rangle = -\langle H, \phi' \rangle = -\int_0^\infty \phi'(x)\, dx = \phi(0) - \phi(\infty) = \phi(0)\)
\item
  Therefore: \(H' = \delta\)
\item
  The derivative of a step is a spike!
\end{itemize}

\textbf{Example: Absolute value}

\begin{itemize}
\tightlist
\item
  \(\langle |x|', \phi \rangle = -\int |x| \phi'(x)\, dx\)
\item
  Split:
  \(= -\int_{-\infty}^0 (-x)\phi'(x)\, dx - \int_0^\infty x\phi'(x)\, dx\)
\item
  Integrate by parts each piece:

  \begin{itemize}
  \tightlist
  \item
    Left:
    \(= -[(-x)\phi]_{-\infty}^0 - \int_{-\infty}^0 \phi\, dx = -\int_{-\infty}^0 \phi\)
  \item
    Right:
    \(= -[x\phi]_0^\infty + \int_0^\infty \phi\, dx = \int_0^\infty \phi\)
  \end{itemize}
\item
  Sum:
  \(\int_0^\infty \phi - \int_{-\infty}^0 \phi = \int \text{sgn}(x)\phi(x)\, dx\)
\item
  Therefore: \((|x|)' = \text{sgn}(x)\)
\item
  Differentiate again: \((\text{sgn})' = 2\delta\)
\end{itemize}

\textbf{General jump formula}

\begin{itemize}
\tightlist
\item
  If \(f\) is piecewise smooth with jump of size \(a\) at \(x = c\):
\item
  \(f'_{\text{dist}} = f'_{\text{classical}} + a\delta_c\)
\item
  Jumps contribute delta functions to the derivative
\end{itemize}

\textbf{The completion perspective}

\begin{itemize}
\tightlist
\item
  Smooth functions: all derivatives exist (classically)
\item
  Continuous functions: first derivative may not exist classically
\item
  \(L^1_{loc}\) functions: derivatives may not exist classically
\item
  Distributions: ALL have ALL derivatives
\item
  The space is now closed under differentiation
\end{itemize}

\textbf{Theorem: Differentiation is continuous}

\begin{itemize}
\tightlist
\item
  If \(T_n \to T\) in \(\mathcal{D}'\), then \(T_n' \to T'\) in
  \(\mathcal{D}'\)
\item
  Proof:
  \(\langle T_n', \phi \rangle = -\langle T_n, \phi' \rangle \to -\langle T, \phi' \rangle = \langle T', \phi \rangle\)
\item
  Limits and derivatives commute in distribution space
\item
  This FAILS for ordinary functions
\end{itemize}

\textbf{Convolution with test functions}

\begin{itemize}
\tightlist
\item
  For \(T \in \mathcal{D}'\), \(\phi \in \mathcal{D}\): define
  \((T * \phi)(x) = \langle T, \phi(x - \cdot) \rangle\)
\item
  This is a smooth function!
\item
  \((T * \phi)' = T' * \phi = T * \phi'\)
\item
  Convolution with smooth function regularizes
\end{itemize}

\section{The Fourier Transform}\label{the-fourier-transform}

\textbf{Classical definition: \(L^1\) functions}

\begin{itemize}
\tightlist
\item
  For \(f \in L^1(\mathbb{R})\):
  \(\hat{f}(\xi) = \int_{-\infty}^\infty f(x) e^{-i\xi x}\, dx\)
\item
  Well-defined: \(|f(x)e^{-i\xi x}| = |f(x)|\), which is integrable
\item
  \(\hat{f}\) is bounded and continuous
\end{itemize}

\textbf{Basic properties}

\begin{itemize}
\tightlist
\item
  Linear: \(\widehat{af + bg} = a\hat{f} + b\hat{g}\)
\item
  Translation: \(\widehat{f(x-a)}(\xi) = e^{-ia\xi}\hat{f}(\xi)\)
\item
  Modulation: \(\widehat{e^{iax}f(x)}(\xi) = \hat{f}(\xi - a)\)
\item
  Dilation: \(\widehat{f(cx)}(\xi) = \frac{1}{|c|}\hat{f}(\xi/c)\)
\item
  Conjugation: \(\widehat{\bar{f}}(\xi) = \overline{\hat{f}(-\xi)}\)
\end{itemize}

\textbf{The key property: differentiation ↔ multiplication}

\begin{itemize}
\tightlist
\item
  \(\widehat{f'}(\xi) = i\xi \hat{f}(\xi)\) (if \(f, f' \in L^1\))
\item
  Proof: integrate by parts, boundary terms vanish
\item
  \(\widehat{(-ix)f(x)}(\xi) = \hat{f}'(\xi)\)
\item
  Derivatives become multiplication, multiplication becomes derivatives
\item
  This is WHY Fourier transform is powerful for differential equations
\end{itemize}

\textbf{Convolution theorem}

\begin{itemize}
\tightlist
\item
  For \(f, g \in L^1\): \(\widehat{f * g} = \hat{f} \cdot \hat{g}\)
\item
  Convolution ↔ multiplication
\item
  Proof: Fubini
\end{itemize}

\textbf{Examples}

\emph{Gaussian:} - \(f(x) = e^{-x^2/2}\) -
\(\hat{f}(\xi) = \sqrt{2\pi} e^{-\xi^2/2}\) - Gaussian transforms to
Gaussian! - Proof: \(f\) satisfies \(f' = -xf\); apply Fourier, get
\(i\xi\hat{f} = -i\hat{f}'\), same ODE

\emph{Rectangle:} - \(f = \mathbf{1}_{[-a,a]}\) -
\(\hat{f}(\xi) = \frac{2\sin(a\xi)}{\xi}\) - Sharp cutoff → slow decay
(sinc)

\emph{Heat kernel:} - \(K_t(x) = \frac{1}{\sqrt{4\pi t}} e^{-x^2/4t}\) -
\(\hat{K}_t(\xi) = e^{-t\xi^2}\) - Simple Gaussian in Fourier space

\textbf{Inversion formula}

\begin{itemize}
\tightlist
\item
  Under suitable conditions:
  \(f(x) = \frac{1}{2\pi}\int_{-\infty}^\infty \hat{f}(\xi) e^{i\xi x}\, d\xi\)
\item
  Fourier transform is (almost) self-inverse
\item
  Precise statement requires care about convergence
\end{itemize}

\textbf{Plancherel theorem}

\begin{itemize}
\tightlist
\item
  For \(f \in L^1 \cap L^2\): \(\|\hat{f}\|_2 = \sqrt{2\pi}\|f\|_2\)
\item
  Extends by density to all \(f \in L^2\)
\item
  Fourier transform is isometry \(L^2 \to L^2\) (up to constant)
\item
  Proof: compute
  \(\|\hat{f}\|_2^2 = \int \hat{f} \overline{\hat{f}} = \int \hat{f} \widehat{\bar{f}(-\cdot)} = 2\pi (f * \bar{f}(-\cdot))(0) = 2\pi \|f\|_2^2\)
\end{itemize}

\textbf{Connection to Fourier series (Ch 13)}

\begin{itemize}
\tightlist
\item
  Fourier series: periodic functions on circle, frequencies
  \(n \in \mathbb{Z}\) (discrete)
\item
  Fourier transform: functions on \(\mathbb{R}\), frequencies
  \(\xi \in \mathbb{R}\) (continuous)
\item
  Series: \(f = \sum c_n e^{inx}\)
\item
  Transform: \(f = \frac{1}{2\pi}\int \hat{f}(\xi) e^{i\xi x}\, d\xi\)
\item
  Discrete spectrum → continuous spectrum
\end{itemize}

\section{Transforming Distributions}\label{transforming-distributions}

\textbf{The problem}

\begin{itemize}
\tightlist
\item
  Want to extend Fourier transform to distributions
\item
  First attempt:
  \(\langle \hat{T}, \phi \rangle = \langle T, \hat{\phi} \rangle\)?
\item
  Issue: if \(\phi \in \mathcal{D}\) (compact support), is
  \(\hat{\phi} \in \mathcal{D}\)?
\item
  \textbf{No!} Fourier transform of compactly supported function is
  analytic, hence NOT compactly supported (unless zero)
\end{itemize}

\textbf{Schwartz space (the fix)}

\begin{itemize}
\tightlist
\item
  Definition: \(f \in \mathcal{S}\) if \(f \in C^\infty\) and for all
  \(n, k\): \[|x|^n |f^{(k)}(x)| \to 0 \text{ as } |x| \to \infty\]
\item
  ``Rapidly decreasing'' --- decays faster than any polynomial
\item
  Examples: \(e^{-x^2}\), \(e^{-x^2}P(x)\) for any polynomial \(P\)
\item
  \(\mathcal{D} \subset \mathcal{S}\) (compact support implies rapid
  decay)
\end{itemize}

\textbf{Key property}

\begin{itemize}
\tightlist
\item
  Theorem: \(\mathcal{F}: \mathcal{S} \to \mathcal{S}\) is a bijection
\item
  Fourier transform preserves Schwartz class
\item
  Proof idea: \(\widehat{x^n f} = i^n \hat{f}^{(n)}\) and
  \(\widehat{f^{(k)}} = (i\xi)^k \hat{f}\); decay ↔ smoothness
\end{itemize}

\textbf{Tempered distributions}

\begin{itemize}
\tightlist
\item
  Definition: \(\mathcal{S}^\prime\) is continuous linear functionals on
  \(\mathcal{S}\)
\item
  Larger test space → smaller distribution space
\item
  \(\mathcal{S}' \subset \mathcal{D}'\)
\item
  Examples of tempered distributions:

  \begin{itemize}
  \tightlist
  \item
    All \(L^p\) functions
  \item
    All polynomials
  \item
    \(\delta\) and its derivatives
  \item
    Everything we care about!
  \end{itemize}
\end{itemize}

\textbf{Fourier transform of tempered distributions}

\begin{itemize}
\tightlist
\item
  Definition:
  \(\langle \hat{T}, \phi \rangle = \langle T, \hat{\phi} \rangle\) for
  \(T \in \mathcal{S}'\), \(\phi \in \mathcal{S}\)
\item
  Well-defined since \(\hat{\phi} \in \mathcal{S}\)
\item
  Extends the classical Fourier transform
\end{itemize}

\textbf{The fundamental examples}

\emph{Fourier transform of \(\delta\):} -
\(\langle \hat{\delta}, \phi \rangle = \langle \delta, \hat{\phi} \rangle = \hat{\phi}(0) = \int \phi(x)\, dx = \langle 1, \phi \rangle\)
- Therefore: \(\hat{\delta} = 1\) - \textbf{The delta function
transforms to the constant function!} - Infinitely concentrated in space
→ spread over all frequencies

\emph{Fourier transform of \(1\):}

\begin{itemize}
\tightlist
\item
  \(\langle \hat{1}, \phi \rangle = \langle 1, \hat{\phi} \rangle = \int \hat{\phi}(\xi)\, d\xi = 2\pi \phi(0) = \langle 2\pi\delta, \phi \rangle\)
\item
  Therefore: \(\hat{1} = 2\pi\delta\)
\item
  \textbf{The constant function transforms to (scaled) delta!}
\item
  Spread over all space → concentrated at frequency zero
\end{itemize}

\emph{Fourier transform of \(e^{iax}\):}

\begin{itemize}
\tightlist
\item
  \(\widehat{e^{iax}} = 2\pi\delta_a\)
\item
  Pure frequency \(a\) ↔ delta at \(\xi = a\)
\end{itemize}

\emph{Fourier transform of \(\delta'\):}

\begin{itemize}
\tightlist
\item
  \(\widehat{\delta'} = i\xi \cdot \hat{\delta} = i\xi\)
\item
  More generally: \(\widehat{\delta^{(n)}} = (i\xi)^n\)
\end{itemize}

\textbf{Properties preserved}

\begin{itemize}
\tightlist
\item
  \(\widehat{T'} = i\xi \hat{T}\)
\item
  \(\widehat{(-ix)T} = \hat{T}'\)
\item
  \(\widehat{T * S} = \hat{T} \cdot \hat{S}\) (when defined)
\item
  Differentiation ↔ multiplication extends to distributions
\end{itemize}

\textbf{The Dirac comb}

\begin{itemize}
\tightlist
\item
  Definition: \(\text{III}(x) = \sum_{n \in \mathbb{Z}} \delta_n\)
\item
  A distribution (infinite sum of deltas)
\end{itemize}

\textbf{Poisson summation formula}

\begin{itemize}
\tightlist
\item
  Theorem: \(\widehat{\text{III}} = 2\pi \cdot \text{III}\)
\item
  The Dirac comb is (essentially) its own Fourier transform!
\item
  Proof sketch:

  \begin{itemize}
  \tightlist
  \item
    For \(f \in \mathcal{S}\): \(\sum_n f(n) = \sum_n \hat{f}(2\pi n)\)
    (classical Poisson)
  \item
    Rewrite as statement about distributions
  \end{itemize}
\item
  This is the bridge between Fourier series and Fourier transform
\end{itemize}

\section{Fundamental Solutions and
PDEs}\label{fundamental-solutions-and-pdes}

\textbf{The strategy}

\begin{itemize}
\tightlist
\item
  Want to solve \(Lu = f\) for differential operator \(L\)
\item
  Idea: find \(G\) with \(LG = \delta\) (fundamental solution)
\item
  Then \(u = G * f\) solves \(Lu = f\)
\item
  Why? \(L(G * f) = (LG) * f = \delta * f = f\)
\end{itemize}

\textbf{Finding fundamental solutions via Fourier}

\begin{itemize}
\tightlist
\item
  \(LG = \delta\) where \(L = P(\frac{d}{dx})\)
\item
  Take Fourier: \(P(i\xi)\hat{G}(\xi) = 1\)
\item
  Solve: \(\hat{G}(\xi) = \frac{1}{P(i\xi)}\)
\item
  Invert to find \(G\)
\item
  Poles of \(\frac{1}{P(i\xi)}\) require care (principal value, etc.)
\end{itemize}

\textbf{Heat equation on \(\mathbb{R}\)}

\emph{Setup:} - \(u_t = u_{xx}\) for \(x \in \mathbb{R}\), \(t > 0\) -
Initial condition: \(u(x, 0) = f(x)\)

\emph{Fourier method:} - Take Fourier transform in \(x\):
\(\hat{u}_t = (i\xi)^2 \hat{u} = -\xi^2 \hat{u}\) - This is an ODE in
\(t\)! - Solution: \(\hat{u}(\xi, t) = \hat{f}(\xi) e^{-\xi^2 t}\) - By
convolution theorem: \(\hat{u} = \hat{f} \cdot \hat{K}_t\) where
\(\hat{K}_t(\xi) = e^{-\xi^2 t}\) - Therefore: \(u = f * K_t\)

\emph{The heat kernel:} - \(\hat{K}_t(\xi) = e^{-\xi^2 t}\) - Invert:
\(K_t(x) = \frac{1}{\sqrt{4\pi t}} e^{-x^2/4t}\) - This is the
fundamental solution

\emph{Heat kernel as delta limit:} - As \(t \to 0\):
\(\hat{K}_t(\xi) = e^{-\xi^2 t} \to 1 = \hat{\delta}(\xi)\) - Therefore:
\(K_t \to \delta\) in \(\mathcal{S}'\) - Direct proof: \(K_t\) is
approximate identity (Ch 12), so
\(\langle K_t, \phi \rangle \to \phi(0)\) - Initial condition:
\(u(x,0) = (K_0 * f)(x) = (\delta * f)(x) = f(x)\) ✓

\emph{Smoothing property:} - \(K_t \in C^\infty\) for \(t > 0\) -
\(u = K_t * f\) is \(C^\infty\) for \(t > 0\), even if \(f\) is rough -
Heat flow smooths instantly

\textbf{Wave equation on \(\mathbb{R}\)}

\emph{Setup:} - \(u_{tt} = c^2 u_{xx}\) for \(x \in \mathbb{R}\),
\(t > 0\) - Initial conditions: \(u(x,0) = f(x)\), \(u_t(x,0) = g(x)\)

\emph{Fourier method:} - Take Fourier in \(x\):
\(\hat{u}_{tt} = -c^2\xi^2 \hat{u}\) - ODE:
\(\hat{u}_{tt} + c^2\xi^2 \hat{u} = 0\) - General solution:
\(\hat{u} = A(\xi)\cos(c\xi t) + B(\xi)\sin(c\xi t)\) - Initial
conditions: \(A = \hat{f}\), \(c\xi B = \hat{g}\)

\emph{D'Alembert's formula:} - Invert to get:
\(u(x,t) = \frac{1}{2}(f(x+ct) + f(x-ct)) + \frac{1}{2c}\int_{x-ct}^{x+ct} g(y)\, dy\)
- Two waves traveling in opposite directions

\emph{Fundamental solution:} - Response to \(\delta\) initial position,
zero initial velocity: -
\(W_t(x) = \frac{1}{2}(\delta(x-ct) + \delta(x+ct))\) - Two delta
functions propagating outward at speed \(c\) - Distributions essential
to even write this!

\emph{Finite vs infinite speed:} - Heat: \(K_t(x) > 0\) for all \(x\),
all \(t > 0\) --- infinite propagation speed - Wave: \(W_t\) supported
on \(\{x : |x| = ct\}\) --- finite propagation speed \(c\) - Information
travels at most at speed \(c\) for wave equation

\textbf{General constant-coefficient operators}

\begin{itemize}
\tightlist
\item
  \(P(\partial_t, \partial_x) u = f\)
\item
  Fourier transforms to algebraic/ODE:
  \(P(\partial_t, i\xi)\hat{u} = \hat{f}\)
\item
  Solve in Fourier space, invert
\item
  This is the general method
\end{itemize}

\section{Applications}\label{applications}

\textbf{Application 1: The Sampling Theorem}

\emph{Bandlimited functions:} - Definition: \(f\) is bandlimited to
\([-\Omega, \Omega]\) if \(\hat{f}(\xi) = 0\) for \(|\xi| > \Omega\) -
Only frequencies up to \(\Omega\) present - Example: audio filtered to
remove high frequencies

\emph{The Nyquist-Shannon theorem:} - Theorem: If \(f\) is bandlimited
to \([-\Omega, \Omega]\), then \(f\) is completely determined by its
samples \(f(n\pi/\Omega)\) for \(n \in \mathbb{Z}\) - Reconstruction
formula:
\[f(x) = \sum_{n=-\infty}^\infty f\left(\frac{n\pi}{\Omega}\right) \text{sinc}\left(\frac{\Omega x}{\pi} - n\right)\]
- where \(\text{sinc}(x) = \frac{\sin(\pi x)}{\pi x}\)

\emph{Proof using distributions:} - Sampling = multiplying by Dirac
comb: \(f_{\text{sampled}} = f \cdot \text{III}_{\pi/\Omega}\) -
Fourier: multiplication → convolution -
\(\hat{f}_{\text{sampled}} = \frac{1}{2\pi}\hat{f} * \widehat{\text{III}_{\pi/\Omega}}\)
- Poisson:
\(\widehat{\text{III}_{\pi/\Omega}} = \frac{2\Omega}{\pi} \text{III}_{2\Omega}\)
- Result: \(\hat{f}\) gets periodically replicated with period
\(2\Omega\) - If \(\hat{f}\) supported in \([-\Omega, \Omega]\):
replicas don't overlap! - Can recover \(\hat{f}\) by multiplying by
rectangle \(\mathbf{1}_{[-\Omega,\Omega]}\) - Inverse Fourier gives
reconstruction formula

\emph{The message:} - Sample rate must exceed \(2\Omega\) (Nyquist rate)
to avoid aliasing - Distributions (Dirac comb) essential to the proof -
Foundation of digital audio, images, communications

\begin{center}\rule{0.5\linewidth}{0.5pt}\end{center}

\textbf{Application 2: The Central Limit Theorem}

\emph{Setup (no probability needed):} - Let \(f \geq 0\) with
\(\int f = 1\) (normalized) - Assume \(\int xf(x)\, dx = 0\) (centered)
- Assume \(\int x^2 f(x)\, dx = 1\) (unit variance) - Consider
\(n\)-fold convolution: \(f_n = f * f * \cdots * f\) (\(n\) times) -
Rescale: \(g_n(x) = \sqrt{n} f_n(\sqrt{n} x)\)

\emph{The theorem:} - \(g_n \to \frac{1}{\sqrt{2\pi}} e^{-x^2/2}\) (the
Gaussian) - Convergence in \(L^1\) and pointwise

\emph{Proof via Fourier:} - \(\hat{f}_n(\xi) = \hat{f}(\xi)^n\)
(convolution → multiplication) - Rescaling:
\(\hat{g}_n(\xi) = \hat{f}(\xi/\sqrt{n})^n\) - Expand \(\hat{f}\) near
\(\xi = 0\): - \(\hat{f}(0) = \int f = 1\) -
\(\hat{f}'(0) = -i\int xf = 0\) (centered) -
\(\hat{f}''(0) = -\int x^2 f = -1\) (unit variance) - So
\(\hat{f}(\xi) = 1 - \frac{\xi^2}{2} + O(\xi^4)\) - Therefore:
\(\hat{g}_n(\xi) = \left(1 - \frac{\xi^2}{2n} + O(n^{-2})\right)^n \to e^{-\xi^2/2}\)
- This is the Fourier transform of \(\frac{1}{\sqrt{2\pi}}e^{-x^2/2}\)!

\emph{The message:} - The Gaussian is the universal attractor under
convolution - No matter what bump you start with, repeated convolution →
Gaussian - This is WHY normal distributions appear everywhere in nature
- Pure Fourier analysis proof --- no probability theory needed

\begin{center}\rule{0.5\linewidth}{0.5pt}\end{center}

\textbf{Application 3: The Uncertainty Principle}

\emph{Statement:} \[\sigma_x \sigma_\xi \geq \frac{1}{2}\]

where \(\sigma_x^2 = \int x^2 |f(x)|^2\, dx\) (spread in position) and
\(\sigma_\xi^2 = \int \xi^2 |\hat{f}(\xi)|^2\, d\xi\) (spread in
frequency), assuming \(\|f\|_2 = 1\).

\emph{Meaning:} You cannot simultaneously localize a function in both
position and frequency.

\emph{Proof sketch:} - Use \(\widehat{f'}(\xi) = i\xi\hat{f}(\xi)\) -
Apply Cauchy-Schwarz to \(\int |xf| \cdot |f'|\, dx\) - Integration by
parts gives the bound

\emph{Equality case:} - Achieved when \(f(x) = ce^{-ax^2}\) (Gaussian) -
Gaussian is the ``most concentrated'' function possible given the
constraint

\emph{Implications beyond quantum mechanics:}

In signal processing: - Sharp cutoff in time → slow decay in frequency
(ringing) - Sharp cutoff in frequency → slow decay in time (sinc
oscillations) - Can't have both compact support in time AND bandlimited

In audio: - A very short sound pulse must contain many frequencies - A
pure tone must extend over long time - Musical notes are a tradeoff:
finite duration = spread of frequencies

In communication: - Bandwidth and time-duration of signals are
constrained - Shorter pulses require more bandwidth - Fundamental limit
on data transmission

\emph{The message:} - Information has fundamental limits - Not just
quantum mechanics --- pure Fourier analysis - Gaussian achieves the
optimal tradeoff

\begin{center}\rule{0.5\linewidth}{0.5pt}\end{center}

\textbf{Application 4: Classical Integrals}

Fourier transforms can compute integrals that are difficult or
impossible by other means.

\emph{The Dirichlet integral:}
\[\int_0^\infty \frac{\sin x}{x}\, dx = \frac{\pi}{2}\]

A subtlety: This is NOT Lebesgue integrable!
\(\int_0^\infty \left|\frac{\sin x}{x}\right|\, dx = \infty\)

Only converges as improper integral (conditionally convergent).

\emph{Fourier proof:} -
\(\widehat{\mathbf{1}_{[-1,1]}}(\xi) = \int_{-1}^1 e^{-i\xi x}\, dx = \frac{2\sin\xi}{\xi}\)
- Fourier inversion at \(x = 0\):
\(\mathbf{1}_{[-1,1]}(0) = 1 = \frac{1}{2\pi}\int_{-\infty}^\infty \frac{2\sin\xi}{\xi}\, d\xi\)
- Therefore: \(\int_{-\infty}^\infty \frac{\sin\xi}{\xi}\, d\xi = \pi\)
- By symmetry: \(\int_0^\infty \frac{\sin x}{x}\, dx = \frac{\pi}{2}\)

\emph{The Gaussian integral:}
\[\int_{-\infty}^\infty e^{-x^2}\, dx = \sqrt{\pi}\]

\emph{Fourier proof:} - Let \(f(x) = e^{-x^2/2}\). Then \(f' = -xf\). -
Take Fourier: \(i\xi\hat{f} = -i\hat{f}'\), so
\(\hat{f}' = -\xi\hat{f}\) - Same ODE! So
\(\hat{f}(\xi) = c \cdot e^{-\xi^2/2}\) for some \(c\) - Plancherel:
\(\|\hat{f}\|_2^2 = 2\pi\|f\|_2^2\), but \(\hat{f} = cf\) - So
\(|c|^2 = 2\pi\), hence \(c = \sqrt{2\pi}\) (positive since
\(\hat{f}(0) > 0\)) - Evaluate at \(\xi = 0\):
\(\int e^{-x^2/2}\, dx = \sqrt{2\pi}\) - Rescale:
\(\int e^{-x^2}\, dx = \sqrt{\pi}\)

\emph{The Borwein integrals --- a surprise:}

\[\int_0^\infty \frac{\sin x}{x}\, dx = \frac{\pi}{2}\]
\[\int_0^\infty \frac{\sin x}{x} \cdot \frac{\sin(x/3)}{x/3}\, dx = \frac{\pi}{2}\]
\[\int_0^\infty \frac{\sin x}{x} \cdot \frac{\sin(x/3)}{x/3} \cdot \frac{\sin(x/5)}{x/5}\, dx = \frac{\pi}{2}\]
\[\vdots\]
\[\int_0^\infty \prod_{k=0}^{6} \frac{\sin(x/(2k+1))}{x/(2k+1)}\, dx = \frac{\pi}{2}\]

Pattern holds! But then:
\[\int_0^\infty \prod_{k=0}^{7} \frac{\sin(x/(2k+1))}{x/(2k+1)}\, dx = \frac{\pi}{2} - \frac{\pi}{467807924713440738696537864469}\]

Slightly LESS than \(\frac{\pi}{2}\)!

\emph{Fourier explanation:} - \(\frac{\sin(ax)}{ax}\) is (up to scaling)
the Fourier transform of \(\mathbf{1}_{[-a,a]}\) - Product in frequency
domain = convolution in time domain - Convolving rectangles of widths
\(2, \frac{2}{3}, \frac{2}{5}, \ldots\) - The integral equals (up to
scaling) the value of this convolution at 0 - Key: convolution of
rectangles is a ``trapezoidal'' shape - At origin, stays at maximum
height as long as total width \(< 2\) -
\(1 + \frac{1}{3} + \frac{1}{5} + \frac{1}{7} + \frac{1}{9} + \frac{1}{11} + \frac{1}{13} = \frac{4436}{3003} < 2\)
✓ -
\(1 + \frac{1}{3} + \frac{1}{5} + \cdots + \frac{1}{15} = \frac{88069}{45045} > 2\)
✗ - When sum exceeds 2, the convolved shape at 0 drops below maximum -
The tiny deviation from \(\frac{\pi}{2}\) is exactly computable!

\emph{The message:} - Fourier transforms mysterious integrals into
geometry of rectangles - A pattern that holds for 7 terms and fails on
the 8th --- Fourier explains why - Beautiful interplay between product
(frequency) and convolution (time)

\bookmarksetup{startatroot}

\chapter*{Coda}\label{coda}
\addcontentsline{toc}{chapter}{Coda}

\markboth{Coda}{Coda}

Analysis Text, Draft III; concluding remarks.


\backmatter


\end{document}
